\documentclass[11pt]{amsart}
\usepackage{amsmath,amssymb,amsthm,mathtools}
\usepackage[margin=1.15in]{geometry}
\setlength{\parskip}{4pt plus 1pt}
\usepackage{hyperref}
\usepackage{array}
\usepackage{booktabs}
\usepackage{longtable}
\usepackage{xcolor}
\usepackage{enumitem}
\usepackage[final]{microtype}
\emergencystretch=2.5em
\tolerance=1200

%% Theorem environments
\newtheorem{theorem}{Theorem}[section]
\newtheorem{proposition}[theorem]{Proposition}
\newtheorem{lemma}[theorem]{Lemma}
\newtheorem{corollary}[theorem]{Corollary}
\newtheorem{conjecture}[theorem]{Conjecture}
\theoremstyle{definition}
\newtheorem{definition}[theorem]{Definition}
\newtheorem{example}[theorem]{Example}
\theoremstyle{remark}
\newtheorem{remark}[theorem]{Remark}
\newtheorem{notation}[theorem]{Notation}
\newtheorem{convention}[theorem]{Convention}
\newtheorem{problem}[theorem]{Problem}

%% Negative-result environments (the no-go ledger and the structural walls)
\theoremstyle{definition}
\newtheorem{nogo}{No-go}[section]
\newtheorem{obstruction}{Wall}[section]

%% Math macros
\newcommand{\RH}{\mathrm{RH}}
\renewcommand{\Re}{\operatorname{Re}}
\renewcommand{\Im}{\operatorname{Im}}
\newcommand{\half}{\tfrac{1}{2}}
\newcommand{\gam}{\gamma}
\newcommand{\Lam}{\Lambda}
\newcommand{\eps}{\varepsilon}
\newcommand{\supp}{\operatorname{supp}}
\newcommand{\Spec}{\operatorname{Spec}}
\newcommand{\sqm}{\operatorname{sq}_-}
\newcommand{\cW}{\mathcal{W}}
\newcommand{\cD}{\mathcal{D}}
\newcommand{\cH}{\mathcal{H}}
\newcommand{\cS}{\mathcal{S}}
\newcommand{\cN}{\mathcal{N}}
\newcommand{\cG}{\mathcal{G}}
\newcommand{\Hb}{\mathbb{H}}
\newcommand{\R}{\mathbb{R}}
\newcommand{\C}{\mathbb{C}}
\newcommand{\Z}{\mathbb{Z}}
\newcommand{\Q}{\mathbb{Q}}
\newcommand{\N}{\mathbb{N}}
\newcommand{\GXi}{G_\Xi}
\newcommand{\KXi}{\mathsf{K}_\Xi}
\DeclareMathOperator{\Tr}{Tr}
\DeclareMathOperator{\Res}{Res}
\DeclareMathOperator{\Short}{Short}
\DeclareMathOperator{\diag}{diag}

\title[The Riemann Hypothesis: An Obstruction Ledger and a Pick Architecture]{The Riemann Hypothesis:
An Obstruction Ledger and the Arithmetic Pick Architecture It Shaped\\[0.3em]
\normalsize\textmd{No-go results, a common boundary quantifier, and a reduction to Li--Keiper
positivity}}

\author{David Alejandro Trejo Pizzo}
\thanks{Independent Researcher. \texttt{dtrejopizzo@gmail.com}}

\begin{document}
\maketitle

\begin{abstract}
This paper develops a decade-long obstruction ledger for several spectral, arithmetic, adelic, and
positivity-based approaches to the Riemann Hypothesis; identifies a recurring uniform-boundary quantifier
(the passage from average/almost-all/interior to specific/uniform/boundary) as their common failure mode;
and constructs an arithmetic Pick/Nevanlinna route whose analytic part is closed and whose sole
arithmetic input is shown to terminate exactly at Li--Keiper positivity.

The first part is the ledger: the concrete no-gos that closed each route, distilled into a small number
of structural walls --- arithmetic propagation, wrong-sign positivity, infinite-dimensional
Hasse--Minkowski failure, per-place indefiniteness, the cohomological-Frobenius wall,
sign-versus-magnitude, and the uniform-norm (rank-one escape) wall --- each carrying an explicit logical
classification (exact equivalence with RH, framework closure obstruction, or insufficiency theorem for a
stated input class). The standard is strict throughout: no assertion is treated as proved because it has
been named, defined, or made plausible by computation, and two of our own results are recorded as
withdrawn.

The second part is the architecture those walls shaped: the \emph{arithmetic Pick positivity} (ARP-P)
chain, fifteen Steps establishing the equivalence $\textnormal{ARP-P}\Longleftrightarrow\textnormal{RH}$.
Fourteen of the fifteen Steps are closed in the precise sense recorded in the status map (two of them as
proved conditional consequences of the open input). The unresolved mathematical input is Step~5, whose
internal reduction terminates at one terminal-defect inequality,
\[
   \delta_N\ge0\ \text{for all }N\ \Longleftrightarrow\ \textnormal{RH}\qquad(\Omega_7),
\]
equivalently $\lambda_n\ge0$ for all $n$. $\Omega_7$ is not our discovery --- it is the classical Li
criterion; the contribution is that an independently constructed architecture, after its internal
difficulty audit, terminates exactly there, and the ledger explains why the major routes explored in this
programme repeatedly return to RH-equivalent boundary or uniformity conditions. The paper does not prove
RH and does not present the open input as a small gap: being equivalent to RH, it carries the full
difficulty of the Hypothesis. A self-contained falsification harness accompanies the paper and confirms
the governing laws numerically; by conservation of difficulty, it closes nothing. The whole is offered
for rigorous peer scrutiny, in the same discipline of retraction we invite the reader to apply to it.
\end{abstract}

\tableofcontents
\newpage

\section{Prologue}\label{sec:prologue}

\paragraph{What this paper claims.}  Four things: an obstruction ledger for the major spectral,
arithmetic, adelic, and positivity-based routes explored in this programme, each obstruction carrying an
explicit logical classification (\S\ref{sec:walls}); a finite-probe signature architecture (the sourced
determinant and its Pick kernel); the ARP-P equivalence framework, whose analytic bridge is closed
(\S\ref{sec:status-map}); and the localization of that framework's single unresolved arithmetic input at
classical Li--Keiper positivity.

\paragraph{What this paper does not claim.}  It does not prove the Riemann Hypothesis.  It does not
present $\Omega_7$ as a new criterion for RH --- $\Omega_7$ is the classical Li criterion, reached here as
the terminal statement of an independently constructed architecture.  And it does not describe the
unresolved input as a ``small gap'': since $\Omega_7$ is equivalent to RH, its unresolved status carries
the full difficulty of the Hypothesis.

\paragraph{How to audit the paper.}  The fifteen-step implication chain, the H0--H8 internal analysis of
Step~5, and the $\Omega_1$--$\Omega_7$ difficulty-localization chain are three distinct logical levels,
kept typed and separate throughout; their dependency and status maps appear in \S\ref{sec:status-map},
and the provenance of every claimed contribution is listed at the end of this Prologue.  Two of our own
earlier results are recorded as withdrawn, with the errors identified (\S\ref{sec:withdrawn}); the
discipline of retraction is exactly what we invite the reader to apply to this paper.

\medskip

This paper gathers a research program that lasted years and consumed, by an honest accounting, an enormous
quantity of paper, electricity, and patience. The walls came first --- many of them --- and the proof
path came last. That path is short to state and reduces the Hypothesis to a single inequality; but that
one inequality, $\Omega_7$, remains open, and it is itself equivalent to RH. The programme is therefore
presented in two parts: the negative geography that closed every other route we explored (Part~I), and the
proof path those walls shaped, with its one open input honestly marked (Part~II).

The program began far from operator theory, in the empirical statistics of arithmetic functions: the
distribution of $\omega(n)$, the number of distinct prime factors; the fluctuations of the von Mangoldt
function in short intervals; the numerics of the explicit formula. For a long time the work was
computational. Tens of thousands of processor-hours went into evaluating localized Weil quadratic forms
to forty significant digits, into tabulating zeros, into testing one reformulation of positivity after
another against the Davenport--Heilbronn function --- a function that obeys the same functional equation
as $\zeta$ but whose zeros stray off the critical line, and which therefore serves as the perfect
falsifier: any argument that ``proves'' the Riemann Hypothesis but does not distinguish $\zeta$ from
Davenport--Heilbronn is, by that very fact, wrong.

Every route we tried ran into the same wall, wearing different clothes each time. We called it, internally, the
\emph{symmetry/positivity wall}: the location of the zeros is a constraint that is global, one-sided (an
upper bound, not a two-sided equality), and orthogonal to every symmetry the problem hands you for free.
Positivity --- the fact that $\Lambda(n)\ge0$, that the von Mangoldt measure is a genuine measure ---
was always available, but positivity controls \emph{sign}, and the Riemann Hypothesis is about
\emph{magnitude}; the two are orthogonal, and one can prove, in level after level (operator norm, Carleson
measure, diagonal of a kernel, derivative of a spectral shift), that no purely positivity-based argument
crosses the gap.

The decisive change of viewpoint, when it came, was not a new estimate but a new \emph{object}. A scalar
determinant cannot see the difference between a self-adjoint operator and one with complex eigenvalues:
its zeros are its zeros, wherever they sit. But a determinant equipped with its \emph{response to
finite-rank probes} --- a \emph{sourced} determinant --- can. Its second variation in the probe directions
reconstructs a two-variable kernel, and the negative-square count of that kernel is precisely the number
of eigenvalues off the real axis. This single idea --- topologize not the scalar determinant but its
signature kernel --- dissolved the signature-blindness that had defeated every earlier route.

The second idea was equally simple and equally hard-won: the one genuinely divergent direction in the
problem --- the pole of $\zeta$ at $s=1$, which contributes a trace growing like $\tfrac12(\log P)^2$ ---
must be removed by \emph{shorting}, the Schur-complement operation, and never by subtraction. The
distinction sounds pedantic until one writes down the $2\times2$ example
$\bigl(\begin{smallmatrix}R&\sqrt R\\\sqrt R&1\end{smallmatrix}\bigr)\succeq0$ and watches subtraction of
the divergent corner manufacture a spurious negative eigenvalue while the Schur complement leaves the
matrix positive. The divergent pole is harmless if shorted and catastrophic if subtracted; an entire year
of false starts is compressed into that sentence.

The third idea was the one that turned a reformulation into a proof, and it is almost embarrassingly
classical. Once the pole is shorted, the surviving objects $G_P^\circ(z)$ are \emph{compressed resolvents
of genuine self-adjoint operators}, and a compressed self-adjoint resolvent is bounded by $1/|\Im z|$ ---
uniformly, with the divergence safely quarantined in the discarded corner. A uniformly bounded family of
holomorphic functions is a normal family, and a normal family cannot converge to a function with a pole.
So if these resolvents converge to the right object \emph{anywhere on an open set}, Vitali's theorem
carries the convergence across the strip and forbids the off-axis poles --- which are exactly the
off-critical zeros. The whole difficulty then collapses to proving the convergence on the half-plane
$\Re s>1$, where the Euler product converges absolutely and $\Xi$ has no zeros to worry about --- and
that convergence, by conservation of difficulty, is precisely the open input Step~5, not an easy step:
the region is friendly, the statement is not.

This program was not carried out in isolation. Over its life it generated a long traditional
correspondence --- letters, and later a great many emails --- with mathematicians whose questions,
objections, and corrections shaped every definition in this paper. Several of the sharpest turns recorded
below began as a one-line objection in the margin of a draft. The errors that correspondence caught ---
a sign in a Herglotz coordinate, the false belief that a sourced determinant factors over primes, the
crucial difference between defining a compressed resolvent and proving that an object \emph{is} one ---
are now scars in the proof, healed but visible, and the proof is the better for them.

\paragraph{Status and invitation.}
We present this paper as a research programme together with a \emph{proposed proof path}, not as a closed
proof. The path (Part~II, \S\ref{sec:live-proof-steps}) reduces the Hypothesis to a single open input,
Step~5, whose internal reduction terminates at the terminal-defect positivity $\Omega_7$: $\delta_N\ge0$
for all $N$, equivalently $\lambda_n\ge0$
for all $n$. Fourteen of the fifteen Steps are closed in the precise sense recorded in the status map
(\S\ref{sec:status-map}): their asserted construction, implication, equivalence, or conditional
consequence is proved. The unresolved mathematical input is Step~5, whose internal reduction terminates
at $\Omega_7$; and $\Omega_7$ is itself equivalent to
RH. We do not claim to have settled the Hypothesis: a claim of that kind must earn belief through the
scrutiny of many independent experts, not the confidence of its authors, and the honest state is that one
load-bearing inequality remains. We therefore direct attack precisely there --- at $\Omega_7$ and at the
$\Omega$-chain reduction that isolates it (\S\ref{sec:live-proof-steps}) --- and we record, in the same
place, why each attempted closure of $\Omega_7$ returns to the structural walls of Part~I.

\subsection*{The setting: $\zeta$, the explicit formula, and the Weil criterion}

We orient the reader with the objects that recur throughout. The Riemann zeta function is
\[
   \zeta(s)\ =\ \sum_{n=1}^\infty n^{-s}\ =\ \prod_p (1-p^{-s})^{-1},\qquad\Re s>1,
\]
the Euler product being the avatar of unique factorization. The completed function $\xi(s)=
\tfrac12 s(s-1)\pi^{-s/2}\Gamma(s/2)\zeta(s)$ is entire of order 1, satisfies the functional equation
$\xi(s)=\xi(1-s)$, and has precisely the nontrivial zeros $\rho=\beta+i\gamma$ of $\zeta$ among its zeros.
In the coordinates $s=\tfrac12+iz$ the function $\Xi(z)=\xi(\tfrac12+iz)$ is even and real for real $z$,
with zeros at $z_\rho=i(\rho-\tfrac12)$; on the critical line $\beta=\tfrac12$ these zeros are real.
The Riemann Hypothesis asserts: \emph{all nontrivial zeros $\rho$ of $\zeta$ have $\Re\rho=\tfrac12$.}

The prime-number theory of these zeros is encoded in the explicit formula. For a smooth test function $h$,
Weil's form of the explicit formula reads \cite{Weil1952}
\[
   \sum_{\rho}\widehat h(\gamma_\rho)\ =\ \widehat h\text{-archimedean}\ -\ \sum_{p\,\mathrm{prime}}\,
   \sum_{k=1}^\infty\frac{\log p}{p^{k/2}}\,h(k\log p),
\]
where the left side sums over all nontrivial zeros $\rho=\beta_\rho+i\gamma_\rho$ and the right side
carries the prime-power weights $\Lambda(p^k)=\log p$ through $h$. Weil's criterion states that the
Hypothesis is equivalent to
\[
   \cW(h\star\tilde h)\ \ge\ 0\qquad\text{for all }h\in\mathcal S(\mathbb R),
\]
where $\tilde h(x)=\overline{h(-x)}$ and $\star$ is convolution. The form $Q=M_{\mathrm{zeros}}-
M_{\mathrm{arith}}$ is the finite-dimensional truncation of this criterion, and the sign of $Q$ --- or
equivalently of the Pick kernel of $-\Xi'/\Xi$ --- is the object around which every route in this paper
revolves.

The Davenport--Heilbronn function $L_{\mathrm{DH}}$ \cite{DavenportHeilbronn,Titchmarsh} is a linear
combination of Dirichlet $L$-functions to an imaginary character modulo $5$; it satisfies the same
functional equation as $\zeta$ with the same $\Gamma$-factor structure, but its zeros include ones off the
critical line. It serves as the falsifier throughout: any inequality proved for $\zeta$ that also holds
for $L_{\mathrm{DH}}$ does not separate on-line from off-line zeros and therefore cannot imply RH. The
Davenport--Heilbronn test is mandatory.

\section*{Original constructions and results of this paper}\label{sec:provenance}

To the author's knowledge, the following constructions, reductions, and assembled arguments are original
to this paper.  The claim is one of provenance of the present manuscript, not a substitute for an
exhaustive priority determination.

\emph{New constructions introduced here.}  The finite sourced channels and the sourced-determinant
(signature-kernel) architecture (Definition~\ref{def:target-channel} and \S\ref{sec:prologue}); the
ARP-P formulation itself (Definition~\ref{def:ARP-P}); the reference-whitened terminal defect
(Definition~\ref{def:terminal-defect-fixed}); the H0--H8 hierarchy of Step~5; and the $\Omega$-reduction
architecture (\S\ref{subsec:omega7-attack}).  The ARP-P architecture as a whole is introduced here and
assembled from stated classical analytic interfaces.

\emph{Results proved or derived here.}  The determinant serialization of RH
(Theorem~\ref{thm:det-serialization}) --- proved here; the one-point jet serialization
(Theorem~\ref{thm:one-point-jets}) --- proved here; the fixed-$N$ band closure
(Theorem~\ref{thm:band-closure}) --- proved here, with classical Cauchy--Pick machinery; the whitened
terminal theorem (Theorem~\ref{thm:whitened-terminal}) --- proved here; the Li--Cayley dictionary
(Theorem~\ref{thm:li-dictionary}) --- derived here, with the classical Li criterion used as terminal
input; the quantitative H6 witness (Theorem~\ref{thm:H6-closed}) --- proved here; the de Bruijn--Newman
cascade characterization (Theorem~\ref{thm:alpha5}) --- proved here, with the de Bruijn and Rodgers--Tao
inputs stated; the envelope-bound reduction (Proposition~\ref{prop:envelope-reduction}) --- derived here;
and the stated insufficiency theorems (Theorem~\ref{thm:W4}, Proposition~\ref{prop:W5},
Proposition~\ref{prop:unboundable}, the last via Bombieri--Lagarias, cited) --- proved here.

\emph{Classical inputs and external results}, used with citation and never reproved: the Weil explicit
formula and criterion \cite{Weil1952}; Nevanlinna--Pick interpolation and commutant lifting
\cite{FoiasFrazho}; Montel/Vitali normality; the Li criterion \cite{Li1997}, Keiper
\cite{Keiper1992}, and Bombieri--Lagarias \cite{BombieriLagarias1999}; de Bruijn \cite{deBruijn1950} and
Rodgers--Tao; Connes--Consani \cite{ConnesConsani2021,ConnesConsani2023}; and Guth--Maynard
\cite{GuthMaynard}.

$\Omega_7$ is not our discovery; the independently constructed ARP-P architecture is shown, after its
internal difficulty audit, to terminate at Li positivity.  That termination, together with the
obstruction ledger explaining why the major routes explored in this programme repeatedly return to
RH-equivalent boundary or uniformity conditions, is the claimed contribution.  A construction can be
original even when the terminal statement it reaches is classical; that is exactly the situation here,
and we state it rather than let the reader discover it.

\part{The walls, and the path they shaped}

\section{A long road to a short bridge: routes, analogies, and the walls that shaped this one}\label{sec:history}

The argument of this paper is short, and short arguments for old problems are always the residue of a much
larger and mostly negative body of work. It would be misleading to present the bridge without the country
it crosses. This section recounts, in mathematical terms and without any private notation, the principal
routes that were attempted over a multi-year programme, the precise obstructions that closed each of them,
and the way those obstructions, taken together, pointed toward the construction given here. We believe
this history is not decoration: every wall below is a theorem about what \emph{cannot} work, and the
present route is shaped, negatively, by all of them. The complete record of the programme---a decade of
routes, computations, and retractions, documented in full detail---is publicly archived at
\url{https://github.com/dtrejopizzo/riemann}.

\subsection{The journey in one paragraph}\label{sec:journey}

The programme began empirically, in the statistics of arithmetic functions --- the distribution of the
number of prime factors $\omega(n)$, the fluctuations of $\Lambda(n)$ in short intervals, the numerics of
the explicit formula evaluated to high precision. That phase produced a great deal of structure and no
theorem; its lasting product was a single instrument, a \emph{localized Weil quadratic form}
$Q=M_{\mathrm{zeros}}-M_{\mathrm{arith}}$ whose sign encodes the Riemann Hypothesis and which became the
pivot of everything that followed. From the detector grew a sequence of theoretical attacks: a
reduction of the Hypothesis to a Lindel\"of-type bound on a cross sum; an inverse-spectral, Hilbert--P\'olya
search for a self-adjoint operator with the zeros as spectrum; an adelic, trace-formula framing in the
spirit of the spectral programme discussed in \S\ref{sec:relwork}; and finally a sequence of
Hodge-theoretic and Frobenius-realization attempts modelled on the function-field proof. Each reached a
precise wall. The walls turned out to be the \emph{same} wall in different clothing, and understanding
that single obstruction sharply enough to state it as one theorem is what eventually suggested attacking
it from below the critical strip, in the region of absolute convergence, rather than on the critical line
itself. That is the move this paper makes.

What follows is a longer account. We give it because the negative results are, in our experience, the most
useful part of the programme for anyone contemplating an operator approach: each is a small theorem about
what \emph{cannot} succeed, and together they delineate, sharply, the narrow corridor in which a proof can
live. We have rendered every obstruction as mathematics, stripped of any private notation, so that the
reader may check the diagnosis independently of the history.

\subsection{The empirical phase: the $\omega$-class fingerprint and the localized detector}\label{sec:empirical}

The programme opened with a long computational study of the arithmetic of $\omega(n)=\#\{p:p\mid n\}$ and
of the von Mangoldt weights $\Lambda(n)$, evaluated against the explicit formula to high precision. The
goal was to find a \emph{reproducible discriminator}: a quantity, computable from finite data, positive for
$\zeta$ and negative for a function with the same functional equation but zeros off the line --- the
Davenport--Heilbronn function $L_{\mathrm{DH}}$ \cite{DavenportHeilbronn,Titchmarsh}, which serves
throughout as the falsifier. The instrument that survived this phase is a \emph{localized Weil quadratic
form}. For a test function $h$ supported near a scale, Weil's explicit formula \cite{Weil1952} writes
\begin{equation}\label{eq:weilform}
   \mathcal W(h)\ =\ \underbrace{\sum_\rho \widehat h(\gamma_\rho)}_{M_{\mathrm{zeros}}(h)}\ -\
   \underbrace{\Bigl(\widehat h\text{-archimedean}\ -\ \sum_{p,k}\tfrac{\log p}{p^{k/2}}\,h(k\log p)\Bigr)}
   _{M_{\mathrm{arith}}(h)},
\end{equation}
and Weil's criterion \cite{Weil1952} is that the Hypothesis is equivalent to $\mathcal W(h\star\tilde h)\ge0$
for all $h$. Localizing the test function to a band of scales produces a finite Hermitian form $Q=
M_{\mathrm{zeros}}-M_{\mathrm{arith}}$ whose least eigenvalue is the object the whole programme orbits.
This is, in spirit, the band-limited and archimedean-refined positivity of the adelic programme
\cite{connesconsani2021}; the empirical phase tabulated it, found rich $\omega$-stratified structure in its
fluctuations, and --- honestly --- did not beat the convexity bound. Its lasting product was the detector
$Q$ and a precise sense of \emph{why} it is hard to read, which we now record as a sequence of no-gos.

The form $Q$ depends on two ingredients that enter very differently. The \emph{zero side}
$M_{\mathrm{zeros}}(h)=\sum_\rho\widehat h(\gamma_\rho)$ is a linear functional of the zero multiset
$\{\rho\}$: each zero $\rho=\tfrac12+i\gamma_\rho$ (assuming RH) contributes through the Fourier transform
$\widehat h$ of the test function. The \emph{arithmetic side} $M_{\mathrm{arith}}(h)$ is a sum over primes:
it carries the von Mangoldt weights $\Lambda(n)=\log p$ at prime powers $n=p^k$ and the archimedean term
from the $\Gamma$-factor. The Hypothesis says these two sides are in equilibrium: the spectral energy is at
least as large as the arithmetic energy. The empirical phase measured this balance to high precision across
many test functions, finding rich structure stratified by $\omega(n)$: integers with more prime factors
(higher $\omega$-class) contribute differently to the balance, and the statistics of each class carry a
distinct signature in the Fourier domain. This structure was real and reproducible; it did not, however,
beat the convexity bound or give an unconditional proof.

The empirical work was conducted with the contemporary picture of $\zeta$ on the critical line in view:
the random-matrix statistics of the zeros and of the value distribution \cite{KeatingSnaith,FHK}, the
multiplicative-chaos description of the large values \cite{Soundararajan,BondarenkoSeip,SaksmanWebb,
RhodesVargas}, the universality phenomenon \cite{Voronin,Bagchi}, and the hyperbolicity/interlacing
circle of ideas \cite{GORZ,MSS,BorceaBranden}. Each is a source of \emph{generic} information about the
zeros; none, by the principle isolated below, detects the measure-zero event of an off-line zero, and that
negative recognition is what the empirical phase contributed. The generic behaviour of the zeros --- their
statistics, their pair correlation, their universality --- is fully consistent with an off-line zero at any
displacement $\delta>0$, however small. The Hypothesis is a statement about the \emph{individual}
configuration, not the generic one, and the empirical phase eventually made this distinction sharp.

Each of the routes sketched in \S\ref{sec:journey} closed at a precise, checkable obstruction --- not at a
vague ``difficulty'' but at a statement one can write down and verify. In our experience these obstructions
are the most valuable product of the programme, because each is a small theorem about what \emph{cannot}
work, and together they carve out the narrow corridor in which a proof can live. We therefore give them two
full sections of their own. Section~\ref{sec:nogos} records the \emph{no-go results}: concrete attempts,
each with the mathematics that closed it. Section~\ref{sec:walls} distils these into the small number of
\emph{structural walls} --- the obstructions that reappear, in disguise, on every front --- and attaches
to each a precise logical classification (defined at the head of \S\ref{sec:walls}): exact equivalence
with RH, framework closure obstruction, or insufficiency theorem for a stated input class. Informally we
shall sometimes say a statement is ``RH-strength''; this phrase carries no logical content beyond the
explicit classification attached to the statement. The reader interested only in the
proof may skip to \S\ref{sec:live-proof-steps}; but the proof was designed, negatively, against exactly the walls of
\S\ref{sec:walls}, and we believe it is best understood as the one route that threads them.

\section{The no-go results}\label{sec:nogos}

This section is a ledger of negative results. We have rendered each as mathematics, stripped of any private
notation, and grouped them by the route on which they arose. The format is uniform: we state the
obstruction as a numbered \textbf{No-go}, give the computation or theorem that established it with the
displays spaced for reading, and end with the one-line reason it closes the route. Nothing here is folklore
to us: every entry was, at the time, a live hope that a specific calculation extinguished. We include two
results that we later \emph{withdrew} (\S\ref{sec:withdrawn}), because the discipline of retraction is part
of the true history and is exactly the discipline we invite the reader to apply to the present paper.

\subsection{The detector and the limits of generic information}\label{sec:nogo-detector}

The instrument that survived the empirical phase is the localized Weil form $Q=M_{\mathrm{zeros}}-
M_{\mathrm{arith}}$ of \eqref{eq:weilform}. Before turning to the routes built on it, we record four no-gos
that concern the \emph{reading} of $Q$ itself: ways in which a quantity extracted from it looks like
progress but is, on inspection, either positive by construction or blind to the only event that matters.

\begin{nogo}[one-sided positivity is automatic]\label{ng:psd}
The least eigenvalue of the \emph{zero-side} Gram matrix is non-negative for trivial reasons and carries no
arithmetic content.
\end{nogo}
\noindent Write the zero-side block as $G_{\mathrm{zeros}}[i,j]=\sum_\rho\widehat h_i(\gamma_\rho)\,
\overline{\widehat h_j(\gamma_\rho)}=\langle v_i,v_j\rangle$ with $v_i=(\widehat h_i(\gamma_\rho))_\rho$.
This is a Gram matrix, hence
\[
   G_{\mathrm{zeros}}\succeq0,\qquad\text{so}\qquad \lambda_{\min}(G_{\mathrm{zeros}})\ge0
\]
\emph{identically}, for every $L$-function and every sampling of test functions. The measured value
$\lambda_{\min}\approx0.866>0$ for $\zeta$ is therefore not evidence of anything: it is the positivity of an
inner-product matrix. Only the full form $Q$, which subtracts the arithmetic block $M_{\mathrm{arith}}$, can
distinguish on-line from off-line; the one-sided form is structurally blind.

\begin{nogo}[the scalar minimum is a tautology under the Hypothesis]\label{ng:tauto}
Using $\lambda_{\min}(Q)$ as the discriminator between $\zeta$ and the falsifier separates them by the very
quantity one is trying to control.
\end{nogo}
\noindent Under the Hypothesis $\lambda_{\min}(Q)=0$ exactly (it is pinned at the positivity floor), and an
off-line zero at distance $\delta$ forces $\lambda_{\min}(Q)\le-c(\delta)<0$. Thus
\[
   \lambda_{\min}(Q)<0\quad\Longleftrightarrow\quad\text{an off-line zero exists},
\]
so ``separate $\zeta$ from $L_{\mathrm{DH}}$ by $\lambda_{\min}$'' is identical to ``decide whether there is
an off-line zero.'' The scalar is a faithful \emph{detector} but a circular \emph{criterion}; the genuine
non-tautological content lives one level down, in the per-prime anatomy $\{R_p\}$ of the arithmetic block,
not in any single eigenvalue.

\begin{nogo}[the global Cauchy--Schwarz bound is vacuous]\label{ng:cs}
Controlling the localized form by a global second-moment bound destroys the localization that carries the
information.
\end{nogo}
\noindent For the energy $E(F)=\sum_\gamma|F(\gamma)|^2-\bar\rho\!\int|F|^2$ (with $\bar\rho$ the mean
density of the zero ordinates) the natural global estimate uses
$\|D_T\|_{L^2}^2=V(d,T)\sim T\log T$ and Cauchy--Schwarz to bound the cross term. The resulting inequality
is true but \emph{scale-blind}: it bounds a band-limited quantity by its total mass, and an off-line zero
perturbs the band, not the mass. The honest reading is that the information sits in a Carleson-type measure
on the localized scales \cite{Carleson}, which the global bound integrates away.

\begin{nogo}[topology of the zero set does not see the shift]\label{ng:tda}
Persistent homology of the zero configuration cannot detect a displacement off the critical line.
\end{nogo}
\noindent A zero moved off the line by $\delta$ changes pairwise distances by $O(\delta^2)$, to which the
$0$-dimensional persistence barcode is insensitive at any realistic sample size; against a length-matched
control $L_{\mathrm{DH}}$ is not an outlier, and the $1$-dimensional signal is statistically null
($p\approx0.48$). The instrument measures a \emph{generic} feature of the configuration, and an off-line
zero is a measure-zero perturbation of it --- the recurring reason, isolated as Theorem~\ref{thm:marginality}
below, that no generic statistic decides the Hypothesis.

\subsection{First arc: the $\omega$-class route and the parity barrier}\label{sec:arcA}

The first theoretical attack tried to extract the Hypothesis, or at least the Lindel\"of bound, from the
measured statistics of the classes of $n$ by the number of prime factors $\omega(n)$. Six attempts closed,
each for a structural reason; we group them because the lesson is uniform.

\begin{nogo}[the cross-sum reduction is self-strengthening]\label{ng:circular}
Bounding the $\omega$-class cross sum $B(N)$ by $N^\varepsilon$ would yield Lindel\"of, but the bound is
strictly stronger than its target.
\end{nogo}
\noindent The reduction has the shape ``$B(N)\le N^\varepsilon\Rightarrow$ Lindel\"of,'' yet establishing
$B(N)\le N^\varepsilon$ already requires Lindel\"of-type cancellation in the underlying Dirichlet
polynomial. One is asked to prove a statement that \emph{presupposes} the conclusion; the implication is
sound but the hypothesis is the harder fact. This is the first appearance of a pattern --- a true
implication whose antecedent is RH-strength --- that will recur as the master quantifier of
\S\ref{sec:walls}.

\begin{nogo}[exponent extrapolation does not cross the threshold]\label{ng:extrap}
The measured large-values exponent does not fall below the level a large-values/Lindel\"of argument requires,
and the trend is not monotone.
\end{nogo}
\noindent In the large-values approach a Dirichlet polynomial of length $N$ must have its set of
$\sigma$-large points controlled by an exponent below the Bourgain--Demeter--Guth / Guth--Maynard barrier
\cite{BourgainDemeterGuth,GuthMaynard}. The fitted exponent of $B(N)$ over the accessible range
$N\le1.5\times10^5$ is
\[
   \alpha_B\approx0.42\qquad\text{against the required}\qquad \alpha_B<0.31 .
\]
Four data points do not cross this gap, and --- decisively --- the trend \emph{reverses} at the transition
from two to three prime factors ($k=2\to3$), so even the sign of the slope is unstable. A numerical
extrapolation through a non-monotone trend is not a proof.

\begin{nogo}[unconditional anti-correlation vanishes]\label{ng:mv}
The cross-class anti-correlation that is visible at peaks is, unconditionally, zero on average.
\end{nogo}
\noindent By the Montgomery--Vaughan mean-value theorem \cite{IwaniecKowalski}, the mean square of the
relevant Dirichlet polynomial is the diagonal term plus a genuinely smaller off-diagonal, so the
unconditional cross-covariance between $\omega$-classes is $\approx0$. The anti-correlation that the
empirical phase saw is a \emph{conditional}, peak-localized phenomenon; averaged unconditionally it is not
there, so it cannot be promoted to a statement about the zeros.

\begin{nogo}[the parity barrier forbids the separation]\label{ng:parity}
No sieve weight can distinguish $\omega(n)$ even from $\omega(n)$ odd, which is what the route needs.
\end{nogo}
\noindent The separation the route requires is exactly a control on $\omega(n)\bmod2$, i.e.\ on the
Liouville-type sign $(-1)^{\omega(n)}$. The parity problem of sieve theory \cite{IwaniecKowalski} is the
theorem that standard sieve weights cannot distinguish integers with an even number of prime factors from
those with an odd number. This is not a gap in our sieve; it is a proved limitation of all of them.

\begin{nogo}[hyperbolicity reformulations are equivalences]\label{ng:reform}
The Li, Jensen/hyperbolicity, de Bruijn--Newman and Nyman--Beurling criteria are restatements of the
Hypothesis, not approximations to it.
\end{nogo}
\noindent Each of
\[
   \lambda_n>0\ (\forall n),\qquad T_k\bigl(\xi\text{-Taylor}\bigr)\ge0\ (\forall k),\qquad
   \Lambda_{\mathrm{dBN}}\le0,\qquad d_N\downarrow0
\]
\cite{Li,GORZ,BaezDuarte,deBruijn,Newman} is \emph{equivalent} to the Hypothesis. Verifying any of them for
finitely many $n,k,N$ is consistent with an off-line zero of arbitrarily small displacement; the finite
verification is not asymptotic evidence, and a full proof of any one of them is a full proof of RH. We
return to the precise sense of ``equivalent'' in \S\ref{sec:circular}.

\begin{nogo}[the phantom phase transition]\label{ng:phantom}
A claimed geometric transition in the $\omega$-class data did not survive re-verification.
\end{nogo}
\noindent A reported transition at $N_c\approx1.18\times10^5$ dissolved under canonical re-computation; the
honest residue is the much weaker statement that $\zeta$ is somewhat more ``constructive'' at peaks, with no
transition. We record it as a no-go against ourselves: an apparent structure that did not replicate.

The arc closes on one lesson. Every quantity in it is either a restatement of the Hypothesis
(No-go~\ref{ng:circular}, \ref{ng:reform}) or a \emph{generic} statistic that the parity and convexity
barriers render blind to the measure-zero event of an off-line zero (No-go~\ref{ng:extrap}, \ref{ng:mv},
\ref{ng:parity}, \ref{ng:phantom}).

\subsection{Second arc: lower bounds on the off-line displacement}\label{sec:arcB}

The second, and longest, attack sought to bound the real-part displacement
\[
   b_j\ :=\ \beta_j-\tfrac12,\qquad \rho_j=\beta_j+i\gamma_j\ \ \text{a hypothetical off-line zero},
\]
\emph{from below}, since $b_j=0$ for all $j$ \emph{is} the Hypothesis: a uniform lower bound $b_j\ge c>0$
that is impossible to satisfy would prove there is no off-line zero. We surveyed some fifteen classical
mechanisms. Not one of them produces a lower bound; each produces, instead, a bound in the \emph{wrong
direction}, in the \emph{wrong region}, of the \emph{wrong aggregate type}, or one that is openly
\emph{circular}. The uniformity of this four-way failure is itself a discovery, and we record it at the end
of the subsection as a meta-theorem. We give the mathematics for the most instructive cases.

\begin{nogo}[zero-free regions give the wrong direction]\label{ng:kv}
The Korobov--Vinogradov zero-free region bounds $b_j$ from \emph{above}, never from below.
\end{nogo}
\noindent The sharpest known zero-free region \cite{Korobov,MTY} states
\[
   \zeta(s)\ne0\qquad\text{for}\qquad \Re s\ >\ 1-\frac{c}{(\log|t|)^{2/3}(\log\log|t|)^{1/3}} .
\]
A zero $\rho_j=\beta_j+i\gamma_j$ therefore obeys
$\beta_j\le1-c/\bigl((\log\gamma_j)^{2/3}(\log\log\gamma_j)^{1/3}\bigr)$, which gives
\[
   b_j\ =\ \beta_j-\tfrac12\ \le\ \tfrac12-\frac{c}{(\log\gamma_j)^{2/3}(\log\log\gamma_j)^{1/3}} .
\]
This is an \emph{upper} bound on $b_j$: the zero-free region excludes a strip near $\Re s=1$, so it limits
how far \emph{right} a zero can be. The quantity a proof of RH needs is a lower bound on $b_j$ that forces
it to be zero; the KV region gives an upper bound telling us it is less than $\tfrac12$, which every zero
automatically satisfies. The inequality points the wrong way, and no sharpening of the constant $c$ changes
its direction: however much the exponent in $\Re s>1-c(\log|t|)^{-\alpha}$ were improved (short of a
positive-width zero-free strip, itself an open problem far beyond RH-strength methods), the direction of
the bound on $b_j$ would remain ``$b_j\le(\tfrac12-\delta)$ for some $\delta>0$'', never ``$b_j\ge c_0>0$.''
The mechanism is structurally unable to produce the required direction.

An instructive way to see this: the zero-free region is proved by bounding $|\zeta(s)|$ away from zero for
$\Re s$ close to $1$, using the Euler product in the region of (near-)absolute convergence. The Euler
product converges absolutely only for $\Re s>1$ and is largest there; the bound weakens as $\Re s\downarrow
\tfrac12$. So the technique has a natural built-in direction: it says ``the function cannot vanish near the
boundary of the absolute-convergence region,'' which is an upper bound on how far right a zero can be, not
a lower bound on how far left.

\begin{nogo}[the de Bruijn--Newman inequality is inverted]\label{ng:dbn}
The relation to the de Bruijn--Newman constant also yields an upper bound, and the only known input makes it
vacuous.
\end{nogo}
\noindent Under the relevant factorization the heat-flow inequality reads $\Lambda_{\mathrm{dBN}}\ge
\tfrac12 b_j^2$ \cite{deBruijn,Newman,RodgersTao,PolymathDBN}, whence
\[
   b_j\ \le\ \sqrt{2\,\Lambda_{\mathrm{dBN}}} .
\]
To turn this into a lower bound on $b_j$ one would need a \emph{lower} bound on $\Lambda_{\mathrm{dBN}}$;
but the established direction is an upper bound --- classically $\Lambda_{\mathrm{dBN}}\le1/2$, sharpened
to $\Lambda_{\mathrm{dBN}}\le0.2$ \cite{PolymathDBN} --- which gives only the trivial $b_j\le\sqrt{0.4}<1$.
The inequality is structurally inverted relative to the goal.

\begin{nogo}[the Nyman--Beurling distance is circular]\label{ng:nb}
Bounding the Nyman--Beurling distance from below in terms of $b_j$ is equivalent to bounding $b_j$.
\end{nogo}
\noindent The Nyman--Beurling criterion \cite{Beurling,BaezDuarte} reformulates the Hypothesis as
$d_N\downarrow0$, where $d_N$ is the $L^2(0,1)$ distance from the indicator of $(0,1)$ to the span of $N$
dilations of the fractional-part function. A bound $d_N>F(b_j)$ with $F>0$ when $b_j>0$ would indeed exclude
off-line zeros --- but proving such an $F$ exists is exactly proving the displacement is bounded below. The
criterion relocates the problem into a Hilbert space without making it easier; $d_N>F(b_j)$ and ``$b_j\ge
c$'' are the same statement.

\begin{nogo}[Tur\'an inequalities are tautological under the working hypothesis]\label{ng:turan}
For the relevant factor the Tur\'an inequalities hold for every $b_j>0$, so they cannot exclude any.
\end{nogo}
\noindent Writing $\xi=\xi_0\cdot P$ under the working factorization, the Tur\'an inequalities
$T_k^{(\xi_0)}\ge0$ for the Laguerre--P\'olya factor $\xi_0$ are a \emph{consequence} of $\xi_0$ lying in the
Laguerre--P\'olya class, which the factorization already assumes. They are satisfied identically, for every
admissible $b_j>0$, and so impose no constraint on the displacement.

\begin{nogo}[linear forms in logarithms address the wrong objects]\label{ng:baker}
Baker's theorem constrains algebraic numbers; the ordinates $\gamma_j$ carry no algebraic structure.
\end{nogo}
\noindent Baker's bounds for linear forms in logarithms \cite{IwaniecKowalski} apply to logarithms of
\emph{algebraic} numbers. The ordinates $\gamma_j$ are (conjecturally) transcendental and are not known to
satisfy any algebraic relation; the hypotheses of Baker's theorem simply do not engage them. This is a
category error, not a quantitative shortfall.

\begin{nogo}[pair correlation carries no off-line signal]\label{ng:paircorr}
Montgomery's pair correlation describes spacings between on-line zeros and is silent about off-line ones.
\end{nogo}
\noindent The pair-correlation statistic \cite{Montgomery,Odlyzko,RudnickSarnak} measures the distribution
of differences $\gamma-\gamma'$ of zeros \emph{on} the line. An off-line zero is not a point of this
statistic; it contributes nothing to the correlation of the on-line ordinates and therefore cannot be
detected by any distortion of it.

The remaining mechanisms --- Selberg's density theorem and the modern zero-density and large-values
machinery \cite{Ingham,Huxley,Wooley,TTY}, the Levinson--Conrey positive-proportion results
\cite{Levinson,Conrey1989}, Littlewood's oscillation theorems, the large sieve, de Branges chain
conditions, the Selberg trace formula \cite{IwaniecSpectral}, and Weil positivity itself --- each fail for
one of the same four reasons. We expand on the density case since it is the source of the most persistent
optimism, and on the large-sieve case since it illustrates how even a powerful quantitative tool misses the
target.

\begin{nogo}[density theorems count zeros, they do not locate them]\label{ng:density}
Selberg-type zero-density estimates bound the number of zeros with $\Re\rho\ge\sigma$ in aggregate; they
impose no lower bound on the displacement $b_j=\beta_j-\tfrac12$ of any individual zero.
\end{nogo}
\noindent The standard density results take the form
\[
   N(\sigma,T)\ :=\ \#\bigl\{\rho=\beta+i\gamma:\beta\ge\sigma,\ |\gamma|\le T\bigr\}
   \ \le\ C\,T^{A(1-\sigma)}\,(\log T)^B,
\]
for absolute constants $A,B,C$ \cite{Ingham,Huxley}. This bounds how many zeros can lie to the right of
$\Re s=\sigma$, simultaneously, in a height window $[0,T]$. But a single zero with $\beta=\tfrac12+
e^{-10^{100}}$ contributes $+1$ to $N(\sigma,T)$ for every $\sigma\le\tfrac12+e^{-10^{100}}$ and $T$
large enough --- that is, it contributes almost nothing to any density count. No density estimate of this
form can prevent one zero from having a very small displacement. The aggregate bound and the individual
bound are independent questions; all density theorems address the former and leave the latter wide open.

\begin{nogo}[positive proportions do not constrain off-line shifts]\label{ng:prop}
Knowing that a positive proportion of zeros lie on the critical line is orthogonal to knowing whether any
zero lies off it.
\end{nogo}
\noindent The Levinson--Conrey results \cite{Levinson,Conrey1989} establish that $\ge41\%$ of the nontrivial
zeros of $\zeta$ satisfy $\Re\rho=\tfrac12$. This is a theorem about the on-line zeros; it says nothing
about the displacement of any zero that is not on the line. A hypothetical single off-line zero with
$b_j=10^{-1000}$ is perfectly compatible with $99.99\%$ of zeros being exactly on the line. The Hypothesis
(all zeros on the line) and the proportion results (some zeros on the line) are of entirely different
logical types. Counting is not locating.

\noindent We summarize the arc as a single statement.

\begin{theorem}[four-way failure of the lower-bound programme]\label{thm:fourway}
Every classical mechanism that takes arithmetic data convergent in $\Re s>1$ and attempts a positive lower
bound $b_j\ge c$ on the off-line displacement produces instead one of: \emph{(i)} an upper bound on $b_j$
(zero-free regions, de Bruijn--Newman); \emph{(ii)} a statement about the wrong region $\Re s\approx1$
rather than the strip interior (Deuring--Heilbronn repulsion, blocked by a Paley--Wiener barrier);
\emph{(iii)} an aggregate count rather than an individual displacement (density theorems, positive
proportions, pair correlation); or \emph{(iv)} a reformulation equivalent to the bound sought
(Nyman--Beurling, Tur\'an, Li/Jensen). No combination yields a lower bound.
\end{theorem}

\noindent We give an additional proof sketch for failure mode (ii), the Paley--Wiener barrier, since it is
the least familiar and the most instructive.

\medskip
\noindent\textbf{The Paley--Wiener barrier.} The Deuring--Heilbronn repulsion theorem bounds the distance
between a (hypothetical) exceptional zero $\rho_0=\beta_0+i\gamma_0$ close to $\Re s=1$ and another zero
$\rho=\beta+i\gamma$: they cannot coexist unless $|\rho-\rho_0|$ is very small. This information lives
near $\Re s=1$. To transfer it into the strip $\tfrac12<\Re s<1$, one would need a function $F(s)$
analytic on the strip that (a) is controlled by arithmetic at $\Re s>1$ and (b) recovers the desired
information at $\Re s=\tfrac12$. Such a function would have its values on the line $\Re s=\tfrac12$
determined by its values on the line $\Re s=1+\delta$ via Cauchy's integral formula, which requires
pointwise control in a wide strip. Paley--Wiener-type estimates say that the arithmetic input (supported
on the positive integers) has Mellin transform decaying like $e^{-c\pi|t|}$ in the imaginary direction
$|t|\to\infty$; the analytic continuation to $\Re s<1$ gains factors that amplify exponentially, erasing
the decay before reaching $\Re s=\tfrac12$. The arithmetic function has insufficient smoothness to ``reach'' across
the strip by analytic continuation, and there is no elliptic-type maximum principle that would make the
value on one vertical line control the value on another without this smoothness.

The Deuring--Heilbronn repulsion is therefore an excellent theorem about the region near $\Re s=1$; it is
simply the wrong region for a proof of RH. No sharpening of the repulsion estimate changes this.

\medskip
\noindent The single cause beneath all four failures is that \emph{arithmetic information available for
$\Re s>1$ does not propagate into the critical strip $\tfrac12<\Re s<1$}: the Euler product, the only
source of arithmetic data, converges only to the right of the strip. This is so even though numerical
verification of the Hypothesis now reaches astronomical height \cite{PlattTrudgian} --- precisely because
verification at a height is not propagation into the strip. Theorem~\ref{thm:fourway} is the first and most
robust face of the arithmetic-propagation wall (Wall~\ref{ob:propagation}).

\subsection{Third arc: local--global factorization and the adelic eigenvectors}\label{sec:arcC}

A natural response to the propagation barrier is to localize: factor the Weil form over places,
$Q=Q_\infty+\sum_pQ_p$, prove each local piece $Q_p\succeq0$, and conclude global positivity by a
Hasse--Minkowski principle. This is the strategy by which the function-field analogue is proved. Over
$\mathbb Q$ it fails at two independent points, and the companion spectral construction --- realizing the
zeros as eigenvalues on the id\`ele class space --- fails at a third.

\begin{nogo}[the zero side does not factor over primes]\label{ng:nofactor}
The spectral part of the Weil form is a global functional of the zero multiset and admits no per-prime
decomposition.
\end{nogo}
\noindent The arithmetic side $M_{\mathrm{arith}}(h)=\sum_{p,k}\frac{\log p}{p^{k/2}}h(k\log p)$ factors
over primes by construction. The spectral side does not:
\[
   M_{\mathrm{zeros}}(h)\ =\ \sum_\rho\widehat h(\gamma_\rho)
\]
depends on the zeros, which are \emph{global} objects --- no Euler factor $(1-p^{-s})^{-1}$ ever vanishes,
so no prime ``owns'' a zero and there is no place $p$ at which to attach a local piece $Q_p$ of the spectral
form. The factorization $Q=Q_\infty+\sum_pQ_p$ that the strategy requires does not exist on the side that
carries the information; there is no analogue of ``Frobenius at $p$'' for $\zeta/\mathbb Q$.

\begin{nogo}[no Hasse--Minkowski in infinite dimension]\label{ng:hm}
Even granting local positivity at every place, the local-to-global passage is invalid for Hermitian forms
in infinite dimension.
\end{nogo}
\noindent The classical Hasse--Minkowski theorem --- a quadratic form over $\mathbb Q$ is isotropic iff it
is isotropic over every completion --- is a theorem about \emph{finite}-dimensional forms. The Weil form is
infinite-dimensional, and there is no Hasse principle for the signature of an infinite-dimensional Hermitian
form: a form positive at every place can be globally indefinite. So even if No-go~\ref{ng:nofactor} were
somehow circumvented, the inference ``$Q_v\succeq0\ \forall v\Rightarrow Q\succeq0$'' would remain unjustified.

\begin{nogo}[the adelic eigenvectors are not in the space]\label{ng:l2}
The explicit functions that would realize off-line zeros as eigenvalues on the id\`ele class space are not
square-integrable.
\end{nogo}
\noindent The companion construction seeks an operator $H_C$ on $L^2(C_{\mathbb Q})$, $C_{\mathbb Q}$ the
id\`ele class space, with the off-line zeros among its eigenvalues, witnessed by the explicit vectors
$f_j(x)=x^{\,b_j+i\gamma_j}$. But, with the multiplicative Haar measure $d^\times x=dx/|x|$,
\[
   \|f_j\|_{L^2(C_{\mathbb Q})}^2\ =\ \int |x|^{2b_j}\,\frac{dx}{|x|}\ =\ \int x^{2b_j-1}\,dx\ =\ \infty
   \qquad(b_j>0).
\]
The putative eigenvectors are distributions, not $L^2$-vectors, so they cannot witness a complex eigenvalue;
the quadratic form $Q(f_j,f_j)$ is formally divergent. The construction produces the right symbols and the
wrong space.

\begin{nogo}[one obstruction in two languages]\label{ng:twolang}
The $L^2$ failure and its adelic-character dual are the same statement, so passing between them is not
progress.
\end{nogo}
\noindent The condition $x^{b_j+i\gamma_j}\notin L^2(C_{\mathbb Q})$ is identical to the statement that
$|x|^{b_j+i\gamma_j}$ is a \emph{non-unitary} quasi-character (its exponent has nonzero real part $b_j$). An
adelic Gr\"ossencharakter argument that tries to exclude the non-unitary part of the spectrum without first
resolving the displacement merely re-encounters the same obstruction in dual language. The two phrasings are
provably equivalent, and neither is the easier one.

\subsection{Fourth arc: the sign of the curvature, and a failed anticommutation}\label{sec:arcD}

A further family of attempts tried to make the Hypothesis a \emph{variational} or \emph{symmetry} statement:
to exhibit the critical line as the minimizer of an energy, or to force a spectral invariant to vanish by an
anticommuting involution. They closed because the geometry of $\xi$ on the critical line has, under the
Hypothesis, the \emph{opposite sign} to the one these arguments need --- and we could compute that sign
exactly.

\begin{nogo}[the critical line is a maximum, not a minimum]\label{ng:extremum}
Phrasing $b_j=0$ as the global minimum of an arithmetic energy is backwards: the natural form is locally
maximized on the line.
\end{nogo}
\noindent Consider $\sigma\mapsto\log|\xi(\sigma+i\gamma)|$ along a horizontal line through a zero ordinate
$\gamma$. The second-variation calculation gives, for the localized form built from a detector test
function $h$ whose transform vanishes at the sampled zeros $\rho_j$ (so the first bracket term below
drops),
\[
   \delta_b^2\,Q(f,f)\ =\ -8\sum_j c_j^2\Bigl[\Re\bigl(\widehat h''(\rho_j)\,\overline{\widehat h(\rho_j)}
   \bigr)+|\widehat h'(\rho_j)|^2\Bigr]\ \le\ 0,
\]
a manifestly non-positive quantity. The critical line is thus a local \emph{maximum} of $Q$ in the
transverse direction, not a minimum. Any argument that hopes to pin $b_j=0$ by exhibiting it as an energy
minimum is working against the curvature, not with it.

\begin{nogo}[the curvature sign is positive, so concavity criteria are false]\label{ng:curvature}
$\sigma\mapsto\log|\xi(\sigma+i\gamma)|$ is strictly convex at the line; a sufficiency criterion phrased as
concavity cannot hold.
\end{nogo}
\noindent Under the Hypothesis the curvature of $\log|\xi|$ across the critical line is computed from the
Hadamard product: writing $\log\xi(s)=B_0+\sum_\rho[\log(1-s/\rho)+s/\rho]$ and differentiating twice in
$\sigma$ at $s=\tfrac12+i\gamma$,
\[
   \frac{\partial^2}{\partial\sigma^2}\log|\xi\!\left(\tfrac12+\sigma+i\gamma\right)\bigg|_{\sigma=0}
   \ =\ C(\gamma)\ :=\ \sum_\rho\frac{1}{(\gamma-\gamma_\rho)^2}\ >\ 0 .
\]
Every term in the sum is positive, so $C(\gamma)>0$ identically under the Hypothesis. Moreover, RH is
\emph{equivalent} to $C(\gamma)>0$ for all $\gamma$: if an off-line zero $\rho^*$ exists, the contribution
of its conjugate pair adds a negative term to $C(\gamma^*)$ that can dominate. A criterion of the form
``$\sigma\mapsto\log|\xi(\sigma+i\gamma)|$ is $\sigma$-concave $\Rightarrow$ RH'' is therefore false on its
face: the function is strictly $\sigma$-convex everywhere under RH. Several otherwise-plausible variational
formulations got this direction exactly backwards.

\begin{nogo}[the arithmetic sector cannot be controlled without the zeros]\label{ng:sector}
The inequality that would close the variational route is itself equivalent to the Hypothesis.
\end{nogo}
\noindent Write the second derivative $\partial_t^2\log|\zeta(\tfrac12+it)|$ as the sum of three sectors:
\begin{align}
   G_{\mathrm{poly}}(t)&\ =\ \frac{2\bigl(\tfrac14-t^2\bigr)}{\bigl(t^2+\tfrac14\bigr)^2},\label{eq:gpoly}\\
   G_\Gamma(t)&\ =\ -\tfrac14\,\Re\,\psi^{(1)}\!\Bigl(\tfrac14+\tfrac{it}{2}\Bigr),\label{eq:ggamma}\\
   G_\zeta(t)&\ =\ \partial_t^2\log|\zeta\!\left(\tfrac12+it\right)|,\label{eq:gzeta}
\end{align}
where $\psi^{(1)}$ is the trigamma function and $G_{\mathrm{ana}}=G_{\mathrm{poly}}+G_\Gamma$ is the
analytic sector (known). The sector $G_{\mathrm{poly}}$ arises from the polynomial prefactor of $\xi$ and
is readily computed; $G_\Gamma$ arises from the $\Gamma$ factor and is a smooth, controlled function.
What remains, $G_\zeta$, is the arithmetic sector.

The variational route would close if $G_\zeta(t)>-G_{\mathrm{ana}}(t)$ held unconditionally. However,
$G_\zeta$ can be expressed via the explicit formula as
\[
   G_\zeta(t)\ =\ \sum_\rho\Re\frac{1}{\bigl(\tfrac12-\beta_\rho+i(t-\gamma_\rho)\bigr)^2}\ +\ \text{(archimedean tail)},
\]
which depends directly on the imaginary and \emph{real} parts of the zeros. Controlling $G_\zeta$ at
$\sigma=\tfrac12$ therefore requires knowing where the zeros sit --- which is exactly the Hypothesis. The
inequality $G_\zeta>-G_{\mathrm{ana}}$ is equivalent to RH, not a consequence of the Euler product or the
functional equation alone.

\begin{nogo}[symmetry is not antisymmetry]\label{ng:antisym}
The functional-equation involution preserves the relevant operator instead of negating it, so the
$\eta$-invariant argument does not run.
\end{nogo}
\noindent The attempt forces a spectral $\eta$-invariant to vanish by exhibiting an involution $J$ with
$JT_0J^{-1}=-T_0$, which would pair eigenvalues $\lambda\leftrightarrow-\lambda$ and annihilate $\eta$. But
the only RH-free involution available, the functional-equation symmetry $s\mapsto1-s$, satisfies
\[
   J\,T_0\,J^{-1}\ =\ T_0 \qquad(\text{symmetry}),\qquad\text{not}\qquad J\,T_0\,J^{-1}=-T_0\ (\text{antisymmetry}).
\]
It preserves the spectrum rather than reflecting it, so it gives no constraint on $\eta$. The required
Fourier anticommutation $Q(f,g)+Q(\widehat f,\widehat g)=0$ remains open and is, again, of RH strength.

\subsection{Fifth arc: reflection positivity, false discriminants, and the magnitude--phase
theorem}\label{sec:arcE}

The most instructive failures concern \emph{positivity per se} --- the attempt to prove $Q\succeq0$
directly, rather than to bound a displacement. Each taught us where the positivity does \emph{not} come
from, and together they culminate in the theorem that, more than any other, reorganized our thinking and
shaped the present paper.

\begin{nogo}[per-place reflection positivity fails]\label{ng:perplace}
The natural per-prime kernel of an Osterwalder--Schrader attack is indefinite, so there is no local
positive structure to reconstruct from.
\end{nogo}
\noindent A reflection-positivity (Osterwalder--Schrader) strategy would factor the form into per-prime
kernels and reconstruct a positive Hamiltonian. Computing the kernel at a prime $p$ gives, up to positive
constants,
\[
   G_p(r)\ \propto\ \frac{2\log p}{p\,|1-z|^2}\bigl(\sqrt p\,\cos(r\log p)-1\bigr).
\]
This is \emph{indefinite}. At $r=0$ the bracket is $\sqrt p-1>0$, but at $r=\pi/\log p$ it is
$-\sqrt p-1<0$:
\[
   G_p(0)>0,\qquad G_p\!\left(\tfrac{\pi}{\log p}\right)<0,
\]
verified directly at $p=2,3,5,7$. There is no per-place positive structure, so reflection positivity cannot
be reconstructed place by place; positivity, if it holds, is a genuinely \emph{cross-place} phenomenon.

\begin{nogo}[``free'' positivity from reconstruction was retracted]\label{ng:os}
The claim that a positive Hamiltonian follows automatically from an OS reconstruction does not survive,
because the two-point function is oscillatory.
\end{nogo}
\noindent We record a retraction against ourselves. The assertion that $H\ge0$ follows ``for free'' from
Osterwalder--Schrader reconstruction is invalid: the reconstruction theorem requires a
reflection-\emph{decreasing} (exponentially decaying) two-point function, whereas the relevant two-point
function here is \emph{oscillatory}. The hypothesis of the theorem is not met, and the conclusion does not
follow. (This is one of two results we later withdrew; see \S\ref{sec:withdrawn}.)

\begin{nogo}[multiplicativity does not control the sign]\label{ng:mult}
The Euler product is not what makes the localized form positive: a reproducible discriminant fails to
separate $\zeta$ from smooth non-multiplicative controls.
\end{nogo}
\noindent A natural thesis is that \emph{multiplicativity} (the Euler product), rather than the location of
the zeros, controls the sign of the localized form. We tested it with a pre-registered, reproducible
computation of the Carleson constant
\[
   C\ =\ \lambda_{\max}\!\bigl(A^{-1/2}P\,A^{-1/2}\bigr),
\]
comparing $\zeta$ against smooth, deliberately \emph{non}-multiplicative controls. The result refutes the
thesis:
\[
   C_\zeta\approx1.98,\qquad C_{\text{flat}}\approx2.44,\quad C_{\text{loglin}}\approx2.11,\quad
   C_{\log n}\approx2.21 .
\]
Not only does $C$ fail to separate $\zeta$ from the non-multiplicative controls --- $\zeta$ is not even the
most coherent. The apparent separation seen in a cruder earlier test was an artifact of \emph{smoothness}
in $\log n$, not of multiplicativity, because the smooth envelope of the localized form integrates out the
fine, individual-$\log p$ correlation in which multiplicativity actually lives. What the experiment leaves
behind is useful: a relative discriminant (fixed $A$, varying the multiplicative structure) that cancels the
normalization and is reusable. What it removes is the hope that ``positivity $=$ multiplicativity.''

All of this points in one direction, which we now state as the organizing principle of the whole programme.

\begin{theorem}[the sign principle: positivity, not magnitude, is the distinguishing input]\label{thm:marginality}
Let $\mathcal W$ be the localized Weil form \eqref{eq:weilform}. The structural feature that separates
$\zeta$ from an off-line-zero falsifier is the \emph{sign} of its von Mangoldt data, not its size. The
Riemann function has $\Lambda(n)\ge0$ for all $n$; the Davenport--Heilbronn function $L_{\mathrm{DH}}$,
which has zeros off the critical line, has a \emph{sign-indefinite} von Mangoldt function --- for instance
$\Lambda_{\mathrm{DH}}(3)<0$ while $\Lambda(3)=\log3>0$. This positivity is a sign property, and it
is exactly the property that makes the finite operators $A_P$ self-adjoint on a genuine Hilbert space (not
a Pontryagin space). Any proof must use this sign structure; a functional that reads only the magnitudes
cannot certify it.
\end{theorem}

We state this as the organizing principle because it is a \emph{hard no-go for an entire class of
arguments}: the property the proof needs --- that the von Mangoldt Hamiltonian $H_P\succeq0$, hence that
$A_P$ is self-adjoint --- is a statement about the \emph{signs} of the weights $\Lambda(n)$, and no
functional of the sizes $\{|\Lambda(n)|\}$ alone can certify it. Flipping signs of the Euler data (as in
the sign-indefinite $\Lambda_{\mathrm{DH}}$) destroys the positivity $H_P\succeq0$ and moves zeros off the
line while leaving each individual magnitude available; so an $\ell^p$ norm, a spectral radius, a Carleson
constant, or a second moment of $\{|\Lambda(n)|\}$ cannot see the one feature that decides the Hypothesis.
(We do not claim the magnitudes of $\zeta$ and $L_{\mathrm{DH}}$ coincide --- they do not, since
$\Lambda_{\mathrm{DH}}$ is a genuinely different sequence; the point is the sharper one that positivity is
a sign property no size functional accesses.)

The theorem identifies, by elimination, what the distinguishing feature \emph{must} be: the genuine
\emph{positivity} $\Lambda(n)\ge0$ (equivalently, that $\Lambda_{\mathrm{DH}}$ changes sign). This is not a
quantitative difference; it is a structural one. It is precisely the input the present paper isolates as
$\mathrm{(Real)}$ (\S\ref{sec:live-proof-steps}): the positivity $\Lambda(n)\ge0$ makes the finite operators $A_P$
self-adjoint on a genuine Hilbert space, giving their resolvents the bound $\|(A_P-z)^{-1}\|\le1/|\Im z|$
that is the seed of the Vitali normal-family argument. That bound fails for $L_{\mathrm{DH}}$
(\S\ref{sec:h8-harness}), exactly as the theorem predicts.

\subsection{Sixth arc: the cohomological frontier}\label{sec:arcF}

The last family of attempts modelled the desired proof on the function-field case, where the analogue of
the Hypothesis is a theorem \cite{Weil1948,Deligne}: there a Frobenius endomorphism acts on a cohomology
carrying a positive, \emph{gapped} intersection pairing, and the eigenvalue bound $|\alpha|=\sqrt q$
follows from positivity. Transporting this to $\zeta$ requires a Frobenius-type operator on a $\mathbb Z$%
-cohomology that is (i) a morphism of the relevant Hodge structure --- complex-\emph{linear} with respect to
the natural complex structure $J$ --- and (ii) an isometry of a positive-definite, spectrally \emph{gapped}
polarization. We tested both ingredients directly on the window of $\zeta$. Both fail, and they fail as two
faces of one absence; before reaching them we record the averaging obstruction that motivated the test.

\begin{nogo}[off-line secular growth survives every average]\label{ng:cesaro}
Averaging the band statistic over scale does not remove the off-line contribution; it relocates the
Hypothesis rather than resolving it.
\end{nogo}
\noindent The last open structural quantity reduces to the boundedness in scale $\lambda$ of an intrinsic
band width $b_{\mathrm{bulk}}(\lambda)$. A zero $\rho=\beta+i\gamma$ contributes to it a term
\[
   \lambda^{\,2\beta-1}\,\lambda^{\,2i\gamma}.
\]
On the critical line $\beta=\tfrac12$ this is the pure oscillation $\lambda^{2i\gamma}$, which a Ces\`aro
average in $\lambda$ sends to $0$. \emph{Off} the line it is the secular factor $\lambda^{2\beta-1}$ with
$2\beta-1>0$, a positive power that no average removes; and zero-density estimates bound the imaginary parts
and the counts of zeros but \emph{not} their real parts. Hence Ces\`aro convergence of
$b_{\mathrm{bulk}}(\lambda)$ is equivalent to RH, not neutral with respect to it. Numerically, $\zeta$'s band
is flat in $\lambda$ (value $\approx0.13$, slope $\approx0$) while the falsifier $L_{\mathrm{DH}}$, which has
off-line zeros, grows like a positive power (log--log slope $\approx0.43$) --- exactly the predicted
off-line signature. The average is a faithful detector, not a proof step.

\begin{nogo}[the natural polarization is present but gapless]\label{ng:gapless}
A genuine quaternionic Hodge--Riemann positivity exists for $\zeta$, but its spectrum decays exponentially
to zero with no gap, so it cannot force an eigenvalue bound.
\end{nogo}
\noindent Pursuing the cohomological lever we found that the finite-window Weil matrix of $\zeta$ does carry
an intrinsic \emph{quaternionic Hodge--Riemann polarization} (a Deninger-type structure with $J^2=-1$,
modelled on Deninger's conjectured cohomology \cite{Deninger}): the form is positive semidefinite on the
$+i$-eigenspace $V_+$ of $J$ for $\zeta$, while the falsifier $L_{\mathrm{DH}}$ and a symmetry-matched
random control are both \emph{indefinite} there. This is the most intrinsic Weil positivity the programme
produced, and it is the geometrically \emph{right} object: it is not a tautological consequence of the
spectral side, it is not imposed by hand, and it genuinely separates $\zeta$ from indefinite controls.

But it is gapless. The spectrum of the polarization form on $V_+$ is an exponential cascade,
\[
   \mu_k\ \sim\ e^{-cL}\ \downarrow\ 0 \qquad(c\approx1.40\ \log_{10}\text{ per step, independent of }\lambda).
\]
The ratio $\mu_{\max}/\mu_{\min}\to\infty$ as the window grows; there is no uniform positive lower bound
on $\mu_k$. The sharpest reformulation of the positivity is $\mu_{\max}\le1$ (the largest eigenvalue ratio
of the Hodge--Riemann polarization is at most $1$); $\zeta$ sits \emph{exactly} at the marginal value
$\mu_{\max}=1$, while $L_{\mathrm{DH}}$ sits at $\mu_{\max}\in[6,13]$ depending on the window. This is a
faithful, sharp detector: $\mu_{\max}\le1$ is equivalent in content to saying no zero is off the line.
But it is a statement that must be \emph{proved}; it is not a consequence of a gap.

The function-field mechanism gives eigenvalue bounds because the polarization there is gapped: the Hodge--
Riemann bilinear relations supply a uniform lower bound $\mu_k\ge c_0>0$ for all $k$, and the operator
isometry forces all eigenvalues of Frobenius onto $|\alpha|=\sqrt q$. Over $\Spec\mathbb Z$, the
corresponding polarization has no gap: $\mu_k\to0$ exponentially, and the forcing mechanism collapses.
Every faithful reformulation we found --- the Ces\`aro band statistic, the Hodge--Riemann positivity on
$V_+$, the Euler-maximum bound $\mu_{\max}\le1$ --- is a marginal, gapless detector. Giving the margin a
uniform gap is \emph{equivalent to proving RH}; it cannot be supplied by numerics, and the construction
that would supply it geometrically is the unbuilt $\Spec\mathbb Z$-cohomology of the next no-go.

\begin{nogo}[the natural Frobenius is of the wrong type]\label{ng:antilinear}
The scaling generator anticommutes with the quaternionic complex structure, so any Frobenius built from it
is complex-antilinear --- not a Hodge morphism, which must be complex-linear.
\end{nogo}
\noindent The operator whose eigenvalues are (formally) the imaginary parts of the zeros of $\zeta$ is the
log-scaling generator $D_{\log}$. With $J$ the quaternionic complex structure of No-go~\ref{ng:gapless}
(satisfying $J^2=-1$, $J$ anti-Hermitian), a direct operator computation gives
\[
   \{D_{\log},J\}\ =\ D_{\log}J+JD_{\log}\ =\ 0,
   \qquad
   \frac{\|[D_{\log},J]\|}{\|D_{\log}\|}\ =\ 2.000
\]
(exact; the symbol is odd in the scale index).
Since $\{D_{\log},J\}=0$, the operator $D_{\log}$ anticommutes with $J$, i.e.\ $D_{\log}J=-JD_{\log}$. This
means that for any $J$-eigenstate $v$ with $Jv=iv$ we have $J(D_{\log}v)=-D_{\log}(Jv)=-iD_{\log}v$: the
image $D_{\log}v$ has $J$-eigenvalue $-i$, not $+i$. Therefore $D_{\log}$ maps $V_+$ to $V_-$, not to
itself; it is complex-\emph{anti}linear with respect to $J$. A Hodge morphism (a ``Frobenius'' for the
Hodge--Riemann structure) must be complex-\emph{linear}, i.e.\ must satisfy $[F,J]=0$. The operator
$D_{\log}$ is of the wrong type.

For contrast, the same test applied to a real elliptic curve over a finite field $\mathbb F_q$ returns
$[F_q,J]=0$ exactly, a gapped polarization with positive definite Hodge--Riemann form, and $|\alpha|=\sqrt
q$ to machine precision for every eigenvalue of the Frobenius $F_q$. The $\zeta$ window has neither
ingredient: the type is wrong and the polarization is gapless. The two failures are not independent --- they
are two faces of a single missing object.

The object that would close this route is a $J$-linear Frobenius that is an isometry of a gapped,
positive-definite polarization on a cohomology of $\Spec\mathbb Z$ --- the unbuilt object of the adelic
programme \cite{connesconsani-site}. Its absence is not an estimate one is missing; it is a
\emph{construction} that does not yet exist. Recognizing that the obstruction is one of construction
rather than of estimation is what turned us, finally, away from realizing the zeros as a Frobenius
\emph{spectrum} and toward realizing them as \emph{poles} forbidden by a normal-family continuation ---
the route of this paper.

\subsection{Circular reformulations}\label{sec:circular}

Several attractive criteria are not approximations to the Hypothesis but exact restatements of it. We list
them together because the failure mode is identical: each is an equivalence, so a finite verification proves
nothing asymptotic and a full proof is a full proof of RH. In each case the parenthetical is the identity
that makes the equivalence visible.

\begin{itemize}[leftmargin=1.6em]
\item Selberg-type density bounds on the \emph{count} of off-line zeros (a single zero with $b_j=e^{-10^{100}}$
satisfies them all; counting is not locating).
\item Positive-proportion results, ``$\ge41\%$ of zeros on the line'' \cite{Levinson,Conrey1989} (a statement
about on-line zeros that is silent about the shift of any off-line one).
\item The Nyman--Beurling distance $d_N\downarrow0$ \cite{Beurling,BaezDuarte} (equivalent to the displacement
bound, No-go~\ref{ng:nb}).
\item The Tur\'an inequalities for the Laguerre--P\'olya factor (a consequence of the working hypothesis, not
a constraint, No-go~\ref{ng:turan}).
\item Bounds $\psi(x)-x=O(x^{1/2+\varepsilon_0})$ (a bound $\varepsilon_0<\min_j b_j$ would localize the
zeros, which is the problem itself).
\item The Li/Jensen hyperbolicity coefficients $\lambda_n>0$, $T_k\ge0$ \cite{Li,GORZ} (each equivalent to RH;
finite verification is not asymptotic evidence).
\end{itemize}
\noindent The common signature is the umbrella wall of \S\ref{sec:walls}: the Weil-positivity criterion
$\cW(f\star\tilde f)\ge0$ is \emph{itself} equivalent to the Hypothesis \cite{Weil1952}, so any quantity that
faithfully captures it inherits the equivalence.

\subsection{Two withdrawn results}\label{sec:withdrawn}

In the same spirit of honesty we record two results we \emph{withdrew}. First, an early operator model built
on a wave--particle analogy (\S\ref{sec:analogies}) claimed to establish uniqueness of a spectral minimizer;
an internal audit found that its numerical certificates were not reproducible from the stated definitions and
that a convexity step rested on the false inference ``$A\succ0,\ B\succeq0\Rightarrow A-B$ convex.'' The only
surviving content is the correct algebraic reduction of the uniqueness question to $\|A^{-1/2}BA^{-1/2}\|<1$.
Second, a quantitative benchmark based on a specific numerical value of the Davenport--Heilbronn von Mangoldt
data proved not to be recoverable from public definitions; we retain $L_{\mathrm{DH}}$ only as a
\emph{qualitative} falsifier --- which is all the present paper uses it for. We mention both because the
discipline of retraction is exactly what we invite the reader to apply to this paper.

\subsection{The terminal-defect continuation: the ARP-P path and its closures}\label{sec:p52}

The companion development --- the arithmetic Pick positivity (ARP-P) path --- reduces the Hypothesis to a
single terminal-defect positivity, $\delta_N\ge0$ for all $N$, equivalently $\lambda_n\ge0$ for all $n$
(Li--Keiper \cite{Li}). Attacking that residue produced a further layer of no-gos, each of which lands on
one of the structural walls of \S\ref{sec:walls}. We add them to the ledger; together they sharpen, rather
than escape, the master quantifier.

\begin{nogo}[the terminal envelope bound is a three-fronted RH nucleus]\label{ng:envelope}
The scale-uniform bound $|\lambda_n^{\mathrm{prime}}|\le c\sqrt n\log n$ that would give $\lambda_n\ge0$
for all $n$ reduces, by a lossless stationary-phase analysis, to three classical equivalents of RH at
once.
\end{nogo}
\noindent Write the Li oscillation on the zero side, $O(n)=\sum_\rho(1-1/\rho)^n$. The functional equation
gives the exact identity $1-1/(1-\rho)=(1-1/\rho)^{-1}$ and $|1-1/\rho|^2-1=(1-2\beta)/|\rho|^2$; the
Cayley map $\rho\mapsto1-1/\rho$ carries the critical line to the unit circle. Three independent routes to
the envelope bound each terminate at a distinct canonical RH-equivalent,
\[
   \underbrace{|1-1/\rho|=1\ \forall\rho}_{\text{Riemann: zeros on the line}};\qquad
   \underbrace{N(\sigma,T)=0\ \text{on }\tfrac12<\sigma<1}_{\text{Landau: no zeros in the strip}};\qquad
   \underbrace{\psi(x)-x=O(\sqrt x\log^2 x)}_{\text{von Koch: the RH-strength PNT error}},
\]
and in each the main term is unconditional and produces $\sqrt n\log n$ while the whole RH-strength weight
sits in the tail --- the same weight seen thrice. A single off-line zero $\rho=\beta+i\gamma$ (through
the member of its four-element cluster with $\beta<\tfrac12$) injects a term $r^n$ with
$r=|1-1/\rho|>1$ that only \emph{eventually} overtakes the ceiling, so no finite numerical
verification, at any $n$ or precision, closes it. This is the umbrella wall (W\ref{ob:weil}) read on the
Li coordinates.

\begin{nogo}[the archimedean partition cannot be a density]\label{ng:partition}
Splitting the archimedean term into per-prime densities $\sigma_p\ge0$ that dominate each prime comb
cannot work with $\sigma_p$ a genuine density: it fails at large resonances, for every prime and every
budget.
\end{nogo}
\noindent By Bochner the single-prime domination is $\hat\sigma_p(\xi)\ge2\sum_k(\log p)p^{-k/2}
\cos(k\xi\log p)$; the right side is almost-periodic with peaks $2\log p/(\sqrt p-1)$ at every resonance
$\xi\in(2\pi/\log p)\Z$, whereas a density obeys $\hat\sigma_p(\xi)\to0$ (Riemann--Lebesgue). Domination
therefore fails at all large resonances. Per prime and at any fixed support cutoff the problem \emph{is}
feasible --- this is exactly the archimedean prolate positivity of Connes--Consani
\cite{connesconsani2021} --- and the obstruction is the simultaneous limit over all primes: the
uniform-norm wall (W\ref{ob:uniform}) in frequency form.

\begin{nogo}[the archimedean perturbative route is quantitatively closed]\label{ng:perturbative}
The archimedean positivity margin varies only logarithmically in the support while the prime perturbation
grows exponentially, so no uniform gap-versus-norm comparison can hold.
\end{nogo}
\noindent For the Weil operator on an interval of length $a$ the lowest eigenvalue satisfies
$\lambda_a=\log(1/a)+\mu_1-\log(2\pi)+\psi(2)-1+O(a)$ with $\mu_1>0$ \cite{ScrewWeil} --- a margin that
\emph{shrinks} as the support grows --- while the prime perturbation has $\ell^1$-mass $\asymp e^{a/2}$.
No bound $\|\text{prime}\|<(\text{archimedean gap})$ can hold without deciding RH, and the eigenvalue
reformulation $\lambda_a\ge0\ \forall a\iff$ RH makes the equivalence exact. Any archimedean route to the
residue must therefore be \emph{non-perturbative}; a norm bound cannot close it. This is the wrong-sign
wall (W\ref{ob:wrongsign}) sharpened to a rate.

\begin{nogo}[the de Branges positivity conditions are falsified]\label{ng:debranges}
The Hermite--Biehler / reproducing-kernel positivity that would force the zeros onto the line by
Hilbert-space structure is not available in the form required.
\end{nogo}
\noindent The Cayley map recasts RH as ``the structure function $E_\xi$ of $\Xi$ is Hermite--Biehler.''
The archimedean factor $\pi^{-(1/4+iz/2)}\Gamma(1/4+iz/2)$ is \emph{not} an entire Hermite--Biehler
function (its $\Gamma$-poles lie at $z=i(2n+\tfrac12)$ in the upper half-plane), and a product
decomposition $\Xi=E_\xi E_\xi^\#$ would double the real zeros; the positivity must be handled at the
kernel level, not by a scalar factorization. Conrey and Li \cite{ConreyLi} exhibited counterexamples to
the specific positivity conditions that de Branges required for RH, so the positive-kernel route does not
deliver the Hypothesis for free: it is the Weil-positivity wall (W\ref{ob:weil}) in de Branges
coordinates. The surviving reduction is index-theoretic --- one needs the shorted primitive kernels to
converge to the genuine Kre\u{\i}n--Langer kernel of $\Xi$ (W\ref{ob:uniform}) --- and that convergence is
itself of class $E$: its closure is equivalent to the Hypothesis (see the classification at
W\ref{ob:uniform}).

\section{The structural walls}\label{sec:walls}

The no-gos of \S\ref{sec:nogos} are many, but they are not independent. Beneath them lies a small number of
\emph{structural walls}: obstructions that reappear, in disguise, on every front. We state them as numbered
\textbf{Wall}s, in the order they were met, and we attach to each an explicit logical classification, so
that no wall claims more than what is proved about it.

\paragraph{Classification of obstructions.} Each wall (and each no-go it distils) is tagged with one or
more of the following classes, where $C$ denotes the condition the wall isolates:
\begin{itemize}[leftmargin=2.2em]
\item[$E(C)$] \emph{exact equivalence}: $C\Longleftrightarrow\mathrm{RH}$ is proved.
\item[$S(C)$] \emph{sufficient}: $C\Longrightarrow\mathrm{RH}$ is proved.
\item[$N(C)$] \emph{necessary}: $\mathrm{RH}\Longrightarrow C$ is proved.
\item[$F(X;C)$] \emph{framework closure obstruction}: within an explicitly stated framework $X$, closure
of the route is proved equivalent to the condition $C$.  The class $F$ by itself asserts nothing about the
difficulty of $C$; it asserts \emph{here is exactly where this framework terminates}.  The relation of $C$
to RH is then classified separately as $E$, $S$, or $N$ whenever such a relation is proved, and composite
tags such as $F(X;C)+E(C)$ are used.
\item[$I(\mathcal I;C)$] \emph{insufficiency theorem}: an explicitly delimited class of inputs
$\mathcal I$ is proved not to suffice for the target $C$ ($\mathcal I\not\Rightarrow C$).
\end{itemize}
Informally we shall sometimes say a statement is ``RH-strength''; as declared in
\S\ref{sec:history}, this phrase carries no logical content beyond the explicit classification attached to
the statement.  The umbrella under which all the walls sit is the oldest one.

\begin{obstruction}[the Weil-positivity wall]\label{ob:weil}
Weil's criterion \cite{Weil1952} states that the Riemann Hypothesis is equivalent to
\[
   \cW(f\star\tilde f)\ \ge\ 0\qquad\text{for all test functions }f,
\]
where $\cW$ is the explicit-formula Weil form \eqref{eq:weilform} and $\tilde f(x)=\overline{f(-x)}$. This
is a clean equivalence, not an implication: \emph{proving the positivity is proving the Hypothesis}. Every
route that reduces to a positivity statement for $\cW$ inherits this equivalence in full. The walls below
are the specific disguises this equivalence adopts as one approaches the problem from different directions.

To see why the equivalence is so sharp, expand the form in terms of the zeros directly:
\[
   \cW(f\star\tilde f)\ =\ \sum_\rho \widehat f(\rho)\,\overline{\widehat f(1-\bar\rho)}
   \ +\ \text{(archimedean and pole corrections)},
\]
where $\widehat f$ is the Mellin transform. When $\rho$ lies on the critical line, $1-\bar\rho=\rho$ and
the term is $|\widehat f(\rho)|^2\ge0$: the zero sum is manifestly non-negative under RH. An off-line zero
at $\rho_0$ with $\Re\rho_0=\tfrac12+\delta$, $\delta>0$, contributes together with its mirror $\bar\rho_0$
and the two conjugates $1-\rho_0$, $1-\bar\rho_0$, and this four-element cluster introduces an indefinite
contribution to the form:
\[
   \cW\big|_{\text{cluster}}\ =\ 2\,\Re\bigl[\widehat f(\rho_0)\,\overline{\widehat f(1-\bar\rho_0)}\bigr]
   \ +\ 2\,\Re\bigl[\widehat f(1-\rho_0)\,\overline{\widehat f(\bar\rho_0)}\bigr],
\]
which can be negative for appropriate $f$. The forced-negativity lemma (item~(i) in \S\ref{sec:positive})
makes this precise: the cluster forces $\lambda_{\min}(Q)\le -c(\delta)$ unconditionally. The equivalence
is thus tight: the form is positive iff there are no clusters iff there are no off-line zeros iff RH.

\smallskip\noindent\textbf{Classification:} $E(C)$ with $C=$ positivity of $\cW$ on autocorrelations ---
the equivalence is classical (Weil \cite{Weil1952}), not a contribution of this paper.
\end{obstruction}

\noindent This umbrella wall is both the source of difficulty and --- from a different angle --- the starting
point of the present paper: rather than prove positivity of $Q$ globally, the route here proves that the
kernel $\KXi$ whose negative index counts off-line zeros has index zero, which is a consequence of
positivity achieved through a normal-family argument rather than a direct bound.

\begin{obstruction}[arithmetic propagation]\label{ob:propagation}
The Euler product $\prod_p(1-p^{-s})^{-1}$ converges absolutely only for $\Re s>1$; in the strip
$\tfrac12<\Re s<1$ it converges at best conditionally, and near the zeros of $\zeta$ not at all. The zeros
are \emph{global}: not one Euler factor vanishes, so no prime ``owns'' a zero and no local arithmetic datum
at a single prime constrains the location of any zero. This is not a quantitative shortfall but a
qualitative barrier: the Euler product carries no information that distinguishes a zero at $\Re s=\tfrac12$
from one at $\Re s=\tfrac12+10^{-100}$.

To see this in quantitative form, write the logarithmic derivative in $\Re s>1$:
\[
   -\frac{\zeta'}{\zeta}(s)\ =\ \sum_{n=1}^\infty \frac{\Lambda(n)}{n^s},\qquad \Re s>1,
\]
where $\Lambda(n)\ge0$ is the von Mangoldt function. The series has non-negative coefficients, so its
real part is monotone in $\Re s$ and offers no oscillatory information. Continuing to $\Re s=\tfrac12$ is
exactly the analytic continuation problem: the partial sums $\sum_{n\le P}\Lambda(n)n^{-s}$ diverge when
a zero is encountered, and the zero's imaginary location is encoded only in the \emph{phase} of the
continuation, not in any arithmetic datum accessible from $\Re s>1$.

A classical consequence is the Korobov--Vinogradov zero-free region. The best unconditional result of
this type gives, for any zero $\rho=\beta_j+i\gamma_j$ of $\zeta$ with displacement $b_j=\beta_j-\tfrac12$,
\[
   b_j\ \le\ \frac{1}{2}\ -\ \frac{c}{(\log|\gamma_j|)^{2/3}(\log\log|\gamma_j|)^{1/3}}
\]
for a positive constant $c>0$. This is an \emph{upper} bound on $b_j$; it says zeros cannot be too far
from the critical line on the right. No analogous lower bound $b_j\ge c'>0$ is known, precisely because
the Euler product data cannot be forced across the strip. Theorem~\ref{thm:fourway} records the quantitative
form: every classical mechanism that takes data from $\Re s>1$ and attempts a positive lower bound
$b_j\ge c$ on the off-line displacement fails for one of four structural reasons: wrong direction (the
bound is an upper bound, not a lower bound), wrong region (the method applies only where zeros cannot be
anyway), wrong aggregate (almost-all results, not specific zeros), or circular (the input is RH-strength).
The wall was identified independently from at least five directions over the course of the programme, and
each identification gave the same answer.

\emph{This is the wall the present paper attacks directly.} Rather than trying to propagate arithmetic
information across the strip, the route stays in $\Re s>1$ for all estimates, and imports the conclusion
across the strip analytically (by a Vitali normal-family continuation) rather than arithmetically. It is the
one move that does not try to cross the wall.

\smallskip\noindent\textbf{Classification:} $I(\mathcal I;C)$ with $\mathcal I=$ the four classes of
classical mechanisms taking data from $\Re s>1$ enumerated in Theorem~\ref{thm:fourway}, and $C=$ a
positive lower bound $b_j\ge c>0$ on the off-line displacement.
\end{obstruction}

\begin{obstruction}[wrong sign: Hilbert--P\'olya]\label{ob:wrongsign}
The Hilbert--P\'olya programme seeks a self-adjoint operator $H$ with $\{i(\rho-\tfrac12):\zeta(\rho)=0\}$
as its spectrum; if such $H$ exists, the spectrum is real and RH follows. Every unconditional construction
one can write down delivers a \emph{lower} bound on the relevant form:
\[
   Q(f,f)\ \ge\ -C\|f\|^2,\qquad C\ge0,
\]
whereas the Hypothesis is the \emph{upper} restriction $Q\succeq0$. Operator norm, Carleson constant, kernel
diagonal, spectral-shift derivative --- all eight independent paradigms tested gave bounds of the wrong
direction. The difficulty is structural: positivity tools applied to $\zeta$ produce lower bounds because they
control the \emph{size} of the deviation from positivity, not its \emph{sign}. RH needs the form to be
non-negative; the tools prove it is bounded below by a constant that could be negative. Tightening $C$ does
not change the direction.

To make this precise, consider the localized Weil form restricted to a one-parameter family $f_\lambda$
concentrated near a single zero $\rho_j=\tfrac12+b_j+i\gamma_j$. A Taylor expansion in the transverse
displacement $b_j$ gives
\[
   Q(f_\lambda,f_\lambda)\ =\ \underbrace{Q_0}_{\ge0\text{ if all zeros on line}}
   \ -\ b_j^2\cdot\underbrace{Q_2}_{>0}\ +\ O(b_j^4).
\]
The leading term $Q_0$ is positive by assumption (on the line); the coefficient $Q_2>0$ of $b_j^2$ is
explicitly computed and positive. This means $Q$ \emph{decreases} as the zero moves off the line; the
critical line is a \emph{local maximum} of the form in the transverse direction, not a minimum. The
second-variation formula
\[
   \delta_b^2\,Q(f,f)\ =\ -8\sum_j c_j^2\Bigl[\Re\bigl(\widehat h''(\rho_j)\,\overline{\widehat h(\rho_j)}
   \bigr)+|\widehat h'(\rho_j)|^2\Bigr]
\]
makes this explicit: for the detector test functions (which vanish at the sampled zeros, so the first
bracket term drops) the right side is non-positive, in agreement with the unconditional curvature
computation $C(\gamma)>0$ (No-go~\ref{ng:curvature}) that the critical line is a local maximum. The practical consequence is decisive: a gradient-descent (or
variational) argument that ``falls toward'' a minimum of $Q$ falls \emph{away} from the critical line,
not toward it. Any proof must work against the gradient, supplying a structural reason for the maximum to
be a global one. The present paper's route --- the Vitali continuation which forces a holomorphic limit
everywhere, not just at a minimum --- is structurally of this type.

\smallskip\noindent\textbf{Classification:} $I(\mathcal I;C)$ with $\mathcal I=$ the magnitude/lower-bound
functionals of the eight tested paradigms (operator norm, Carleson constant, kernel diagonal,
spectral-shift derivative, and their variants) and $C=Q\succeq0$; the target $C$ itself is of class
$E$ by W\ref{ob:weil}.
\end{obstruction}

\begin{obstruction}[infinite-dimensional Hasse--Minkowski failure]\label{ob:hassemink}
The classical Hasse--Minkowski theorem lifts positivity from local to global for quadratic forms over
$\mathbb Q$: a form is isotropic iff it is isotropic at every completion. The strategy of proving
$Q_v\succeq0$ for each place $v$ and concluding $Q\succeq0$ globally would be an analogue of this theorem
for Hermitian forms. It fails at two independent levels.

\textbf{Failure at the factorization level.} The Weil quadratic form splits naturally as
\[
   \mathcal W(f\star\tilde f)\ =\ \underbrace{\sum_\rho \widehat h(\gamma_\rho)}_{M_{\mathrm{zeros}}}
   \ -\ \underbrace{\sum_p\sum_{k\ge1}\frac{\log p}{p^{k/2}}\widehat h(k\log p)}_{M_{\mathrm{arith}}},
\]
where $h=f\star\tilde f$. The arithmetic side $M_{\mathrm{arith}}$ factors over primes perfectly:
$M_{\mathrm{arith}}=\sum_p Q_p$ with $Q_p(h)=\sum_{k\ge1}\frac{\log p}{p^{k/2}}\widehat h(k\log p)$ a
local datum. The zero side $M_{\mathrm{zeros}}(h)=\sum_\rho\widehat h(\gamma_\rho)$ carries the
Hypothesis; it is \emph{global} in the sense that no zero is attached to any prime. The factorization
$Q=\sum_v Q_v$ would require each $Q_v$ to carry information about which zeros are off-line; since no
Euler factor has any such information, the factorization does not exist on the side that matters
(No-go~\ref{ng:nofactor}).

\textbf{Failure of the infinite-dimensional principle.} Even granting a hypothetical factorization, there
is no Hasse--Minkowski theorem in infinite dimensions. Over $\mathbb Q$ the theorem is finite: there are
finitely many bad primes and the completion argument is exact. In infinite dimension the passage fails
for the truncations that actually occur here: the finite forms $Q_X$ are \emph{not} compressions of a
fixed limit form --- the arithmetic side itself changes with the truncation --- so positivity of every
$Q_X$ controls the limit only if the convergence $Q_X\to Q$ is uniform, which is precisely the
uniform-norm wall (W\ref{ob:uniform}). Nor is there an analogue of the finitely-many-bad-places argument
when the index set of places is infinite (No-go~\ref{ng:hm}).

The two failures are independent: even if one were repaired, the other would remain. The local--global
route thus closes, at the level of principle, without any appeal to the specifics of $\zeta$.

\smallskip\noindent\textbf{Classification:} $I(\mathcal I;C)$ with $\mathcal I=$ per-place positivity data
assembled through a local-to-global principle (No-gos \ref{ng:nofactor}, \ref{ng:hm}) and
$C=Q\succeq0$; the residual convergence statement belongs to W\ref{ob:uniform}.
\end{obstruction}

\begin{obstruction}[per-place reflection positivity fails]\label{ob:perplace}
Osterwalder--Schrader reflection positivity reconstructs a positive Hamiltonian from a Euclidean measure
satisfying an exponential decay condition. Applied to $\zeta$, the natural per-prime kernel is, up to
positive constants,
\[
   G_p(r)\ \propto\ \frac{2\log p}{p\,|1-z|^2}\bigl(\sqrt p\,\cos(r\log p)-1\bigr),
\]
with $z$ a fixed normalization independent of $r$. This kernel is indefinite through the sign of its
bracket: $\sqrt p\,\cos(r\log p)-1$ equals $\sqrt p-1>0$ at $r=0$ but $-\sqrt p-1<0$ at $r=\pi/\log p$,
verified numerically at $p=2,3,5,7$ (No-go~\ref{ng:perplace}). The reconstruction theorem requires
exponential decay of the two-point function; the function is oscillatory, so the hypothesis is not met and
the conclusion does not follow (No-go~\ref{ng:os}). There is no per-place positive structure from which to
reconstruct a global one; whatever positivity $Q$ has is an irreducibly cross-place, global phenomenon.

\smallskip\noindent\textbf{Classification:} $I(\mathcal I;C)$ with $\mathcal I=$ per-place
reflection-positivity reconstructions (Osterwalder--Schrader type, No-gos \ref{ng:perplace},
\ref{ng:os}) and $C=$ a globally positive Hamiltonian.
\end{obstruction}

\begin{obstruction}[the cohomological-Frobenius wall]\label{ob:cohom}
The Weil conjectures over finite fields \cite{Weil1948,Deligne} are proved by exhibiting a Frobenius
$\varphi$ acting on a cohomology $H^\bullet$ that carries a \emph{positive, gapped polarization}: the
Hodge--Riemann bilinear relations force $|\alpha|=\sqrt q$ for every eigenvalue $\alpha$ of $\varphi$.
Transporting this mechanism to $\zeta$ requires two ingredients: a $J$-\emph{linear} (Hodge-type) operator
on a $\mathbb Z$-cohomology, and a positive polarization with a spectral gap. The programme established,
by explicit computation, that neither exists in the available $\zeta$ window.

\textbf{Antilinearity of the natural scaling.} The natural candidate Frobenius is the dilation operator
$D_{\log}$ acting on the adelic Hilbert space, defined by $(D_{\log}f)(x)=f(e^t x)$ at scale parameter $t$.
The complex structure $J$ acts by $Jf(x)=\overline{f(\bar x)}$ (complex conjugation combined with the
functional equation). A direct computation on test functions in the $n$-th frequency mode yields
\[
   \langle [D_{\log},J]\,e_n,e_n\rangle
   \ =\ -2\,\langle D_{\log}\,e_n,e_n\rangle,
\]
so $\{D_{\log},J\}=0$ exactly, and the numerical test confirmed $\|[D_{\log},J]\|/\|D_{\log}\|=2.000$
(No-go~\ref{ng:antilinear}). Anticommutativity means $D_{\log}$ is $\mathbb C$-\emph{anti}linear with
respect to $J$; it satisfies $D_{\log}\circ J=-J\circ D_{\log}$. A Hodge morphism must be $J$-\emph{linear}
(commuting with $J$, not anticommuting), because Deligne's comparison of cohomology with its complex
conjugate requires $\varphi\circ J=J\circ\varphi$. The candidate is of the wrong type.

For comparison: on a real elliptic curve $E/\mathbb F_q$, the Frobenius $\varphi_q$ acts $J$-linearly on
$H^1(E,\mathbb Z/\ell)$, with $\varphi_q\circ J=J\circ\varphi_q$ verified by the endomorphism algebra.
The $\zeta$ test returns antilinearity where the finite-field analogue returns linearity.

\textbf{Gapless polarization.} The Hodge--Riemann mechanism uses the polarization eigenvalue gap to force
$|\alpha|=\sqrt q$: if the smallest positive eigenvalue of the polarization form is $\mu_{\min}>0$, then
the Cauchy interlacing of eigenvalues, combined with the Lefschetz--Hodge--Riemann bilinear relations,
pins every Frobenius eigenvalue on the circle $|\alpha|=\sqrt q$. The programme's computation of the
quaternionic Hodge--Riemann polarization on the $\zeta$ space found a spectrum that decays as
\[
   \mu_k\ \sim\ e^{-cL},\qquad c\ \approx\ 1.40\;\log_{10}\text{ per step},
\]
as the truncation scale $L\to\infty$ (No-go~\ref{ng:gapless}). There is no uniform spectral gap;
$\mu_{\min}\downarrow0$. Marginal positivity ($\mu_{\max}=1$ for $\zeta$, compared with strictly less
than 1 for the Davenport--Heilbronn analogue at the same truncation) distinguishes the two functions, but
a vanishing gap closes the analytic road: no finite $\mu_{\min}$ means no Cauchy pinning.

The object that would close this route is a $\mathbb Z$-cohomological Frobenius with a positive gapped
pairing --- precisely the unbuilt object of the adelic programme \cite{connesconsani-site} --- and its
construction is a problem of algebraic geometry, not of classical analysis.

\smallskip\noindent\textbf{Classification:} $F(X;C)+S(C)$ with $X=$ the Hodge--Riemann/Frobenius
transplant in the available $\zeta$ window and $C=$ the existence of a $J$-linear Frobenius with a
positive gapped polarization over $\Spec\Z$; $C$ would imply RH by the Hodge-index mechanism, but no
equivalence is claimed.
\end{obstruction}

\begin{obstruction}[sign, not magnitude]\label{ob:magphase}
Theorem~\ref{thm:marginality} identifies the distinguishing feature between $\zeta$ and an off-line-zero
falsifier as the \emph{sign} of the von Mangoldt data: $\zeta$ has $\Lambda(n)\ge0$, while the
Davenport--Heilbronn function $L_{\mathrm{DH}}$, which has zeros off the critical line, has a
sign-indefinite $\Lambda_{\mathrm{DH}}$ (e.g.\ $\Lambda_{\mathrm{DH}}(3)<0$). Positivity is a sign
property, and a functional that reads only the sizes $\{|\Lambda(n)|\}$ cannot certify it: sign flips of
the Euler data preserve every individual magnitude yet destroy the positivity $H_P\succeq0$ and push
zeros off the line.

The practical consequence is that any positivity argument that operates through size bounds of the form
$\|M_{\mathrm{arith}}\|\le C$ or $|\sum_n\Lambda(n)n^{-s}|\le C(\sigma)$ cannot reach the Hypothesis: such
bounds are insensitive to the very signs on which self-adjointness turns. What \emph{can} distinguish the
two functions is the genuine \emph{positivity} of the von Mangoldt weights --- the fact that
$\Lambda(n)\ge0$ makes the finite operators $A_P$ self-adjoint on a genuine Hilbert space, giving them a
resolvent bounded by $1/|\Im z|$, whereas the indefinite Euler data of $L_{\mathrm{DH}}$ makes the
corresponding operators Pontryagin-space operators with no such bound. The Hilbert-space property is what
the normal-family argument needs, and it is exactly what $\Lambda\ge0$ supplies. This is the input the
present paper isolates as $\mathrm{(Real)}$ (\S\ref{sec:live-proof-steps}).

\smallskip\noindent\textbf{Classification:} $I(\mathcal I;C)$ with $\mathcal I=$ functionals reading only
the magnitudes $\{|\Lambda(n)|\}$ (size bounds of the displayed forms) and $C=$ the positivity
$H_P\succeq0$ that the route requires.
\end{obstruction}

\begin{obstruction}[uniform norm: the rank-one escape]\label{ob:uniform}
A spectral route that assembles the finite Weil operator $K_P=\sum_{n\le P}\Lambda(n)K_n$ (where $K_n$ is
the $n$-th Weil kernel piece) meets a final barrier. By Mertens' theorem the trace grows,
\[
   \Tr K_P\ =\ \sum_{n\le P}\Lambda(n)\,K_n(0,0)\ \sim\ \tfrac12(\log P)^2,
\]
but this divergence is carried entirely by a single positive, rank-one mode: the residue at the pole $s=1$
of $\zeta$. After removing this mode by Feshbach (Schur) shorting, the shorted trace is bounded. The
question is then whether the \emph{operator norm} of the shorted primitive piece is bounded:
\[
   \sup_P\bigl\|\Pi^\circ K_P\Pi^\circ\bigr\|\ <\ \infty\quad\text{?}
\]
By the Weil explicit formula this norm is controlled by the off-line zeros: if any zero $\rho$ has
$\Re\rho\ne\tfrac12$, the off-pole Weil kernel acquires a contribution growing like $P^{2|\Re\rho-1/2|}$,
so the operator norm diverges. Conversely, if all zeros are on the line the norm stays bounded. The
uniform-norm bound is therefore \emph{equivalent in strength to the Hypothesis}: needing it as an input is
circular, and it is precisely this circularity that the previous generation of spectral approaches ran into.

The companion development sharpens this wall without breaching it. The finite Weil--Gram kernels are
positive \emph{by construction} (a Schur/Haynsworth complement of a positive canonical Gram), so no
positivity remains to be proved at the finite level --- the escape sidesteps the Weil-positivity wall
(W\ref{ob:weil}) there. What is not free is a \emph{convergence}: the shorted primitive kernels
$K_P^\circ$ must converge, in the finite-Gram topology at the critical boundary, to the genuine
Kre\u{\i}n--Langer kernel of $\Xi$; equivalently $\sup_P\|\Pi^\circ K_P\Pi^\circ\|<\infty$. Two facts
delimit the gap. The leading trace divergence $\tfrac12(\log P)^2$ is carried by the single rank-one pole
mode --- but leading-trace rank-one does \emph{not} imply primitive boundedness: a subleading escape
direction of size $o(\log P)^2$ that still diverges is invisible to the trace, so the rank-one trace alone
does not bound the endpoint index. And the naive de Branges packaging fails at the archimedean factor and
at the scalar determinant, which is index-blind (No-go~\ref{ng:debranges}). The residue is thus a single
index-continuity statement, itself of class $E$: closing it is equivalent to closing the Hypothesis.

\emph{The present paper is built to avoid this wall entirely.} It never forms the operator $K_P$ or its
shorted version as an operator on an infinite-dimensional space. Instead it works with the compressed
resolvent $G_P^F(z)=\langle(A_P-z)^{-1}\varphi,\psi\rangle$ for $\varphi,\psi$ in a \emph{fixed finite
source frame} $F$. The only uniform bound used is the per-vector frame mass
$\sup_P\|\iota_P\varphi\|_{H_P}^2\le C_F$ (an absolutely convergent prime sum, verified in $\Re s>1$), not
any operator norm over the whole space. The full operator norm may diverge; the compressed channel stays
bounded by $C_F/|\Im z|$ because each $A_P$ is genuinely self-adjoint (\S\ref{sec:live-proof-steps}). This is the
architecturally new feature of the route.

\smallskip\noindent\textbf{Classification:} $F(X;C)+E(C)$ with $X=$ the spectral route assembling the
full shorted Weil operator and $C=$ the uniform primitive-norm bound
$\sup_P\|\Pi^\circ K_P\Pi^\circ\|<\infty$; the equivalence $C\Longleftrightarrow\mathrm{RH}$ is the
two-directional control displayed above.
\end{obstruction}

\subsection{The master quantifier}\label{sec:masterquant}

The seven walls above are not independent; they are one obstruction in seven costumes. The single sentence
beneath them all is a statement about quantifiers.

\begin{center}
\emph{The Riemann Hypothesis is always the passage from average, almost-all, or interior\\
to specific, uniform, or boundary.}
\end{center}

An off-line zero is a \emph{measure-zero event} in the generic behaviour of $\zeta$: it perturbs the zeros
by a set of measure zero in every natural sense (height, density, correlation). No generic statistic ---
no average, no density theorem, no almost-all estimate --- can detect a measure-zero event. Every wall
above is an instance:

\begin{itemize}[leftmargin=1.6em]
\item \textbf{Propagation wall (W\ref{ob:propagation})}: interior ($\Re s>1$) $\not\to$ boundary
($\Re s=\tfrac12$). Arithmetic data from the region of absolute convergence do not reach the strip.
\item \textbf{Wrong-sign wall (W\ref{ob:wrongsign})}: lower bound (magnitude below) $\not\to$ upper
restriction ($\lambda_{\min}\ge0$). Tools produce inequalities pointing the wrong way.
\item \textbf{Hasse--Minkowski wall (W\ref{ob:hassemink})}: local ($Q_v\ge0$ at every $v$) $\not\to$
global ($Q\ge0$). The local-to-global principle fails in infinite dimension.
\item \textbf{Per-place wall (W\ref{ob:perplace})}: per-prime (local kernel positive) $\not\to$ global
(global positivity). Positivity is cross-place, not reconstructible from per-place pieces.
\item \textbf{Frobenius wall (W\ref{ob:cohom})}: function-field (gapped polarization exists)
$\not\to$ $\mathbb Z$ (gapped polarization over $\Spec\mathbb Z$ not constructed). Average gap $\not\to$
uniform gap.
\item \textbf{Magnitude wall (W\ref{ob:magphase})}: generic (magnitudes $|\Lambda(n)|$) $\not\to$
specific (positivity $\Lambda(n)\ge0$). Sizes do not determine signs.
\item \textbf{Uniform-norm wall (W\ref{ob:uniform})}: pointwise (each compressed channel bounded)
$\not\to$ uniform (operator norm over the whole space bounded). Per-channel control does not give a
uniform operator bound.
\end{itemize}

\noindent A proof must therefore supply a \emph{structural} reason --- not a sharper estimate --- forcing
the boundary behaviour. The proof in \S\ref{sec:live-proof-steps} attempts exactly this: it exhibits a normal family
of genuine compressed resolvents whose limit is forced to remain holomorphic across the critical strip,
so that the boundary statement (no off-line pole of $\KXi$) follows from a rigidity theorem (Vitali's
theorem, \S\ref{sec:live-proof-steps}) rather than from any inequality that could be merely \emph{almost} true. The
reason this family is a normal family --- the key structural input --- is the Hilbert-space self-adjointness
of $A_P$, which is in turn a consequence of $\Lambda\ge0$. This is the structural mechanism the proof path
uses to approach the boundary behaviour without crossing any of the walls above --- reducing it, in the
end, to the single open input $\Omega_7$ rather than closing it outright.

It may help to state the logical shape of the escape explicitly. Define the \emph{verification event}
for each wall as the hypothesis that the wall would need to be crossed, and $\mathcal E$ as the conjunction:
\begin{align*}
   \mathcal E_1 &: \text{arithmetic data from }\Re s>1\text{ constrain zeros in the strip,}\\
   \mathcal E_2 &: \text{the Weil form has a global minimum at the critical line,}\\
   \mathcal E_3 &: Q=Q_\infty+\sum_p Q_p\text{ (local-global factorization on the zero side),}\\
   \mathcal E_4 &: \text{the scaling operator is }J\text{-linear (Hodge morphism type),}\\
   \mathcal E_5 &: \text{the Hodge--Riemann polarization has a uniform spectral gap.}
\end{align*}
The programme proved $\neg\mathcal E_i$ for $i=1,\ldots,5$, each by a specific computation or structural
argument (Theorem~\ref{thm:fourway}, the second-variation formula, No-go~\ref{ng:nofactor},
No-go~\ref{ng:antilinear}, No-go~\ref{ng:gapless}). The present route avoids $\mathcal E_1$--$\mathcal E_5$
entirely: it neither crosses the arithmetic strip nor claims a minimum, factorization, Hodge morphism, or
spectral gap. Instead it uses $\Lambda\ge0$ to make the finite resolvents a bounded normal family on
$\mathbb C\setminus\mathbb R$, and invokes Vitali's theorem --- a rigidity principle, not an inequality ---
to reduce the closure to the single terminal input $\Omega_7$.

\subsection{What the programme did establish}\label{sec:positive}

The no-go ledger above is long, but the programme also produced a number of \emph{positive} results: proved
theorems, computed invariants, and established reductions that are unconditionally true and that the present
paper builds on. We record them here because they are, in aggregate, the mathematical content of the
multi-year journey, and because the proof in \S\ref{sec:live-proof-steps} is their residue: the one thing that survived
after everything that could not work was proved not to work.

\begin{enumerate}[label=(\roman*),leftmargin=2.2em]

\item \emph{The forced-negativity lemma.} An off-line zero at real-part displacement $\delta>0$ forces
\[
   \lambda_{\min}(Q)\ \le\ -c(\delta,\sigma,J)\ <\ 0,
\]
for a positive constant $c$ depending on $\delta$, the localization scale $\sigma$, and the test interval
$J$. This is proved unconditionally and provides the foundation for the localized Weil detector: the
detector fires (gives a negative eigenvalue) if and only if there is an off-line zero in its window. It is
the starting point for the whole programme's quantitative approach.

\item \emph{The unconditional truncation control.} The spectral error $\eta(X)=Q_X-Q_\infty$ of the
truncated Weil form decays as
\[
   \eta(X)\ \approx\ A\cdot\exp\bigl(-\alpha(\log X)^2\bigr),\qquad\alpha\approx\tfrac{\sigma^2}{8},
\]
unconditionally. This bounds the finite-to-infinite approximation error without any hypothesis on zeros.

\item \emph{The curvature criterion.} The Hypothesis is equivalent to
\[
   C(\gamma)\ =\ \sum_\rho\frac{1}{(\gamma-\gamma_\rho)^2}\ >\ 0\qquad\text{for all }\ \gamma,
\]
a statement about the curvature of $\log|\xi|$ across the critical line (No-go~\ref{ng:curvature},
proved unconditionally as an equivalence). This gives a clean geometric characterization of RH.

\item \emph{The second-variation formula.} The second variation of the localized form with respect to the
transverse displacement $b$ is
\[
   \delta_b^2\,Q(f,f)\ =\ -8\sum_j c_j^2\Bigl[\Re\bigl(\widehat h''(\rho_j)\,\overline{\widehat h(\rho_j)}
   \bigr)+|\widehat h'(\rho_j)|^2\Bigr],
\]
non-positive for the detector test functions (whose transforms vanish at the sampled zeros, killing the
first bracket term). This is proved without assuming RH, and it implies that the critical
line is a local maximum of $Q$ in the transverse direction, ruling out variational pinning arguments.

\item \emph{The sector decomposition.} Writing $\xi(s)=\tfrac12 s(s-1)\pi^{-s/2}\Gamma(s/2)\zeta(s)$, the
second derivative $\partial_t^2\log|\xi(\tfrac12+it)|$ factors into three computable sectors:
\begin{align}
   G_{\mathrm{poly}}(t) &= \frac{2(\tfrac14-t^2)}{(t^2+\tfrac14)^2},\label{eq:gpoly2}\\
   G_\Gamma(t) &= -\tfrac14\,\Re\,\psi^{(1)}\!\bigl(\tfrac14+\tfrac{it}{2}\bigr),\label{eq:ggamma2}\\
   G_\zeta(t) &= \partial_t^2\log|\zeta(\tfrac12+it)|,\label{eq:gzeta2}
\end{align}
where $\psi^{(1)}$ is the trigamma function. The first two sectors, $G_{\mathrm{poly}}$ and $G_\Gamma$,
are known functions of $t$ that can be tabulated to arbitrary precision. The third, $G_\zeta$, depends on
the location of the zeros and is accessible only through the explicit formula; its positivity at all $t$
is equivalent to RH. This decomposition isolates exactly the inaccessible piece and is used in
\S\ref{sec:arcD} to show that the sector inequality reformulation is equivalent to, not stronger than,
the Hypothesis itself.

\item \emph{The LP/Tur\'an algebraic reformulation.} The Hypothesis is equivalent to the assertion that
the coefficients $c_n$ of $\xi(\tfrac12+it)=\sum_{n\ge0}(-1)^nc_n t^{2n}$ satisfy the Tur\'an--Laguerre--
P\'olya inequalities $T_k=(2k+1)c_k^2-2kc_{k-1}c_{k+1}\ge0$ for all $k$. This converts an analytic
problem (zero location) into an algebraic one (sign conditions on a real sequence).

\item \emph{The equivalence of the four convexity fronts.} In the fourth arc (\S\ref{sec:arcD}) we
established that four apparently different variational attacks --- variational pinning at a minimum, a
concavity criterion, a sector inequality, and an $\eta$-invariant argument --- are the \emph{same}
problem: they all reduce to $\xi(\tfrac12+it)\in\mathrm{LP}$ (the Laguerre--P\'olya class), which is
equivalent to RH.

\item \emph{The sign principle} (Theorem~\ref{thm:marginality}). This is a proved theorem, not a
heuristic: the feature that separates $\zeta$ from an off-line-zero falsifier is the \emph{sign} of the
von Mangoldt data --- $\Lambda(n)\ge0$ for $\zeta$, sign-indefinite for $L_{\mathrm{DH}}$
($\Lambda_{\mathrm{DH}}(3)<0$) --- a positivity that no functional of the magnitudes $\{|\Lambda(n)|\}$
can certify. The theorem rules out an entire family of magnitude-based proof strategies in one stroke.

\item \emph{The per-place indefiniteness.} The local reflection-positivity kernel
$G_p(r)\propto\sqrt p\cos(r\log p)-1$ is proved indefinite at every prime $p=2,3,5,7$ (and structurally
at all $p$). This closes the Osterwalder--Schrader route unconditionally.

\item \emph{The Kreĭn--de Bruijn--Newman connection.} An off-line zero of displacement $b_j$ forces the
lower bound $\Lambda_{\mathrm{dBN}}\ge\tfrac12 b_j^2$ on the de Bruijn--Newman constant, and the connection
with Kreĭn's theory provides a clean chain from the Laguerre--P\'olya class to the spectral theory of
canonical systems. Concretely: in the heat-flow deformation
$H_\lambda(z)=\int e^{\lambda u^2}\Phi(u)\cos(zu)\,du$ (with $\Phi$ the Riemann kernel, so that $H_0$ is
proportional to $\Xi$), there is a constant $\Lambda_{\mathrm{dBN}}$ such that $H_\lambda$ has only real
zeros iff $\lambda\ge\Lambda_{\mathrm{dBN}}$; as $\lambda$ decreases through $\Lambda_{\mathrm{dBN}}$,
pairs of real zeros coalesce and split into complex-conjugate pairs, and the Kreĭn--Langer index (the
number of non-real pairs) becomes positive. RH is the statement $\Lambda_{\mathrm{dBN}}\le0$; the known
bounds are $0\le\Lambda_{\mathrm{dBN}}\le0.2$ \cite{RodgersTao,PolymathDBN}, so RH is equivalent to
$\Lambda_{\mathrm{dBN}}=0$.

\item \emph{The Lefschetz/ample-class dichotomy.} The Hodge--Riemann approach requires an ample class
over $\Spec\mathbb Z$. A Lefschetz-type computation shows the ample class cannot be supplied from the
geometric objects that are available, isolating the cohomological construction as the precise missing piece.
The computation proceeds via the signature of the arithmetic Néron--Severi group: the polarization form has
signature $(1,19,1)$ in the geometric realization tested---an indefinite, hyperbolic-type form with a
single positive direction---rather than the positive-definite (ample) signature needed to close the
Lefschetz argument. The
dichotomy is therefore sharp: either the ample class exists (not established over $\Spec\mathbb Z$) or
the polarization is hyperbolic and the Hodge--Riemann argument fails at the first step.

\item \emph{The zero-side tautology theorem.} The zero-side Gram matrix $G_{\mathrm{zeros}}$ is positive
semidefinite by construction for every $L$-function; using $\lambda_{\min}(G_{\mathrm{zeros}})$ to
distinguish $\zeta$ from $L_{\mathrm{DH}}$ is circular. The theorem closes the one-sided positivity
approach definitively.

\item \emph{The de Branges/Hilbert space canonicity.} The Hermite--Biehler structure function $E$ of the
canonical system \cite{deBranges} satisfies $|E(z)|>|E^*(z)|$ for $\Im z>0$, where $E^*(z)=\overline{E(\bar z)}$.
This is equivalent to the self-adjointness of the underlying canonical system (\S\ref{sec:live-proof-steps}). When the
canonical system is driven by the von Mangoldt datum $H_P\ge0$, the Hermite--Biehler property holds
unconditionally. The Kreĭn--Langer index $\kappa(A_P)=\deg\mathfrak b_P$ counts the off-line zeros of the
associated structure function; $\kappa(A_P)=0$ for the finite positive systems is an unconditional theorem (\S\ref{sec:live-proof-steps}).

\item \emph{The DH-falsifier computation.} The numerical verification that the detector
$\mu_0(L_{\mathrm{DH}})<0$ while $\mu_0(\zeta)\ge0$ (to 40 significant figures at $\lambda=10,12,\ldots,22$)
provides an unconditional empirical separation between the two functions. More importantly, the theoretical
explanation for the separation --- loss of $\Lambda\ge0$ means the Hamiltonian becomes a Pontryagin-space
operator and the resolvent bound $1/|\Im z|$ fails --- locates the responsible structural feature precisely,
ruling out coincidental numerical agreement.

\end{enumerate}

\noindent These results are, collectively, what is known unconditionally about the Weil form and the
operator-theoretic approach to RH. They form the substrate on which the proof in \S\ref{sec:live-proof-steps} is built:
the positive canonical system (using item~(i) for motivation, item~(viii) for the structural constraint on
what the proof must use, and items~(ii)--(vii) for the shape of the landscape). The proof is not a series
of estimates; it is a construction whose possibility was suggested, negatively, by each of the no-gos above
and whose structural input ($\Lambda\ge0$) was identified, by elimination, as the only remaining lever.

\section{Related work: the adelic--spectral programme}\label{sec:relwork}

The construction here is operator-theoretic and adelic in spirit, and it could not have been arrived at
without the body of work that recast the explicit formula as a trace formula and the Riemann Hypothesis as
a positivity statement in noncommutative geometry. We summarize the debts, citing the published
mathematics; we make no claim about any individual's view of the present paper.

The trace-formula framing \cite{connes1999} expresses the explicit formula as a trace identity on the
adelic class space and recasts the Hypothesis as the positivity of a Weil distribution --- the principle
$T=0\Leftrightarrow\mathrm{RH}$ that organizes every later attack. The archimedean refinement
\cite{connesconsani2021} made the semilocal Weil form and its positivity explicit, and is the direct
ancestor of the localized detector that drove our empirical phase. The scaling- and arithmetic-site
programme \cite{connesconsani-site} supplied the geometric language in which the missing object --- a
suitable cohomology of $\Spec\mathbb Z$ with a Frobenius and a positive pairing --- is precisely named;
the recognition that this object is \emph{unbuilt} is exactly wall~\ref{ob:cohom}. Most recently, the prolate and
spectral-triple realizations of the zeros \cite{ccm-prolate,ccm-zeta} gave concrete self-adjoint models
whose spectra approximate the zeros to high precision, sharpening the question of what additional input
forces reality of the spectrum. The Hecke-algebra and symmetry-breaking work \cite{bostconnes} stands
behind the whole adelic picture, and the more recent absolute-geometry developments --- the
$\mathrm{BC}$-system and quantized-calculus framework \cite{ConnesConsaniBC} and the Riemann--Roch theory
for $\Spec\mathbb Z$ \cite{ConnesConsaniRR} --- continue to build the geometric object whose absence is
wall~\ref{ob:cohom}. The unfolding of archimedean factors that we use in the determinant identification (\S\ref{sec:live-proof-steps})
is, technically, of the Rankin--Selberg type \cite{Zagier}.

The proof here also rests on classical analytic-number-theory foundations: the Mertens and Chebyshev
estimates \cite{IwaniecKowalski,Davenport} supply the divergent $\tfrac12(\log P)^2$ pole-mode
asymptotic; the Hilbert-space operator theory of canonical systems is due to de Branges \cite{deBranges}
and the Kreĭn--Langer extension to indefinite metrics \cite{KreinLanger}; Hadamard's factorization
\cite{Titchmarsh} controls the entire function $\Xi$; and the Vitali--Montel normal-family theorem \cite{Titchmarsh} is the rigidity tool that closes the argument.

The Beurling--Nyman and de Branges reformulations of RH \cite{Beurling,BaezDuarte,deBranges} are part of
the background; we use the de Branges space theory as a technical tool in \S\ref{sec:live-proof-steps} (the Hermite--Biehler
structure function) rather than as a reformulation. The de Bruijn--Newman connection \cite{deBruijn,Newman,
RodgersTao,PolymathDBN} and the Beurling--Nyman distance \cite{Beurling,BaezDuarte} appear in the no-go
ledger as directions that cannot yield lower bounds.

The present route can be read as a deliberate \emph{retreat} from the hardest part of the adelic programme:
where the arithmetic site seeks the positive Frobenius cohomology directly on $\Spec\mathbb Z$, we avoid
it, realizing the zeros instead through the analytic continuation of compressed resolvents of explicit
finite self-adjoint operators, and importing positivity only in the tame region $\Re s>1$. Whether this
retreat is a genuine way around wall~\ref{ob:cohom} or merely relocates it is, again, a question we put
to the reader. The fact that the Davenport--Heilbronn test is passed --- and passed for the correct
structural reason (the loss of self-adjointness at the level of the Hamiltonian $H_P$, exactly where the
von Mangoldt sign changes) --- is the strongest internal consistency check we have.

The spectral randomness and universality of $\zeta$ that motivates the GUE hypothesis --- the pair
correlation of Montgomery \cite{Montgomery}, the GUE statistics of Odlyzko's computations
\cite{Odlyzko}, the random-matrix value distribution of Keating--Snaith \cite{KeatingSnaith} --- is the
empirical background of the programme, and the generic-vs-specific distinction of the master quantifier
(\S\ref{sec:masterquant}) explains precisely why these generic results, as beautiful as they are, cannot
by themselves decide the Hypothesis.

The canonical-system framework used in this paper was developed by de Branges \cite{deBranges} and
extended to indefinite metrics by Kreĭn and Langer \cite{KreinLanger}. The key result we use is that every
Hermite--Biehler structure function $E$ with $|E(z)|>|E^*(z)|$ for $\Im z>0$ generates a Hilbert space of
entire functions $\mathcal H(E)$ in which the operator $A:F\mapsto zF$ (multiplication by the spectral
variable) is self-adjoint; conversely, every self-adjoint operator with simple spectrum arises in this way.
In the indefinite case \cite{KreinLanger}, the structure function satisfies $|E(z)|>|E^*(z)|$ for
$\Im z>0$ only outside a finite exceptional set, and the operator $A$ on the associated Pontryagin space
$\Pi_\kappa(\mathcal H)$ has at most $\kappa$ non-real eigenvalues. The Kreĭn--Langer index $\kappa$ counts
these non-real eigenvalues, and the reformulation $\kappa(\Xi)=\#\{\text{off-line zeros}\}$ (G5 in
\S\ref{sec:live-proof-steps}) is the bridge between the canonical-system theory and the Hypothesis.

The Turán power-sum method and its modern descendants provide sharp lower bounds for partial
sums of Dirichlet series in terms of the coefficients. These bounds (No-go~\ref{ng:turan}) are in principle
usable to detect off-line zeros, but in practice they deliver upper bounds on off-line behavior rather than
lower bounds, and the resulting estimates are aggregate (averaged over a set of zeros) rather than specific
(for one zero). The technique exhausted itself in this programme without a positive outcome, and the
reason --- that it is an average-type argument facing a specific-type question --- is now explained by the
master-quantifier principle.

\section{Three guiding analogies}\label{sec:analogies}

Three physical analogies guided the search. None is a proof, and we are careful to label them as
heuristics; but each carries real mathematics, and each pointed, in the end, toward the same construction.

\paragraph{Wave--particle duality.} In quantum mechanics a duality is resolved not by choosing a side but
by exhibiting one object with two representations related by the Fourier transform. The explicit formula
plays exactly this role for $\zeta$: it is a Fourier--Mellin change of basis between the \emph{primes}
(the multiplicative/scale side) and the \emph{zeros} (the spectral/frequency side),
\[
   \sum_\rho \widehat h(\rho)\ =\ \widehat h\text{-terms at }\infty\ -\ \sum_{p,k}\frac{\log p}{p^{k/2}}\,
   h(k\log p),
\]
the two sides being the two ``representations'' of a single distribution. The difficulty, and the reason
the analogy does not by itself close, is that the inner product making the spectral side unitary --- the
``conserved energy'' of the duality --- is precisely the Weil form $Q$ that encodes $\zeta$; the duality is
real but is not resolved by structure alone. (This is the analogy whose first, over-eager operator
realization we withdrew; \S\ref{sec:withdrawn}.)

\paragraph{A Maxwell-style unification.} A second picture treats the two sides of the wall --- the
arithmetic side and the spectral side --- not as competitors but as two components of a single field, in
the way the electric and magnetic fields are two faces of one electromagnetic object exchanged by a change
of frame. On this view the explicit formula is the ``field equation'' coupling primes and zeros, and the
Riemann Hypothesis is a statement that the combined object has a definite polarization. This picture is
genuinely unifying and it is what motivates seeking a \emph{single} self-adjoint object (here the canonical
operator $A_P$) carrying both the prime data (through its Hamiltonian) and the spectral data (through its
eigenvalues); but it is a conceptual frame, not a computation, and we present it only as such.

\paragraph{Quantum mechanics and Stone's theorem.} The deepest analogy, and the one that finally clarified
the logic, is with the spectral theory of Hamiltonians. In quantum mechanics the reality of a spectrum
follows from Stone's theorem: a strongly continuous one-parameter \emph{unitary} group $U(t)=e^{itH}$ has
self-adjoint generator, hence real spectrum --- and unitarity is guaranteed by an inner product given in
advance by physics. The Berry--Keating Hamiltonian \cite{BerryKeating}
\[
   H_{\mathrm{BK}}\ =\ \tfrac12(xp+px),
\]
where $x$ is position and $p=-i\partial_x$ is momentum, has the formally attractive property that its
classical trajectories $\dot x=x$, $\dot p=-p$ are hyperbolic, with time-reversal symmetry, and the
Weyl semiclassical counting law
\[
   N_{\mathrm{cl}}(T)\ =\ \frac{\mathrm{Area}\,\{(x,p):H_{\mathrm{BK}}\le T\}}{2\pi}
   \ \sim\ \frac{T}{2\pi}\log\frac{T}{2\pi e}
\]
reproduces the Riemann--von Mangoldt counting formula $N(T)\sim(T/2\pi)\log(T/2\pi e)$.
This is striking, but it stops there: the exact spectrum of any self-adjoint realization of $H_{\mathrm{BK}}$
does not reproduce the zeros of $\zeta$ (the Hilbert space is fixed; the zeros are not). The model
identifies the right classical action but not the right quantization.

The adelic scale flow carries this further: on the id\`ele class space $L^2(C_{\mathbb Q})$ the
multiplicative group $\mathbb R_+^*$ acts by scaling and realizes a natural ``quantum Hamiltonian'' whose
absorption spectrum formally consists of the zeros \cite{connes1999}. But exhibiting a Hilbert space on
which this flow is unitary is \emph{logically the same problem} as exhibiting the inner product that makes
the zeros real, which is, once more, the Weil positivity. This is the \emph{inverted Stone} principle: for
$\zeta$ there is \emph{no} inner product given in advance by physics; every candidate Hilbert space that
makes the spectrum real must itself encode the Hypothesis.

The only RH-free symmetry available, the functional equation $s\mapsto1-s$, is a \emph{conjugation}: it
sends $\gamma_\rho\mapsto\overline{\gamma_\rho}$ and therefore preserves the real zeros while pairing
complex ones in conjugate pairs. But a conjugation is \emph{not} an evolution; it cannot play the role of
the unitary group $U(t)=e^{itH}$ in Stone's theorem. The functional equation is a powerful symmetry, and
it is used twice in this paper (once in the proof proper and once in the DH falsifier check); but it cannot
supply the spectral reality by itself.

This is precisely why the present paper does not try to make the scale flow unitary by fiat. Instead it
\emph{builds} the inner product, finitely and explicitly, from the positive von Mangoldt Hamiltonian
(\S\ref{sec:live-proof-steps}): the Hamiltonian $H_P\ge0$ is a positive matrix assembled from $\Lambda(n)\ge0$,
and the resulting operator $A_P$ is self-adjoint on the Hilbert space whose inner product $H_P$ defines.
Stone's theorem then applies honestly to each finite $A_P$, and the entire remaining difficulty is moved
to the analytic continuation across the strip, where a normal-family argument can meet it without needing
to know where the zeros are. The key is that the positivity $\Lambda(n)\ge0$ makes $H_P\ge0$ an honest
positive form, not a formal one; it is the one structural input the proof uses that $L_{\mathrm{DH}}$
cannot supply.

\paragraph{Summary of the three analogies.} The wave--particle analogy identified the structure (a
Fourier--Mellin duality with the explicit formula as field equation) but could not close because the
inner product was not given. The Maxwell-unification analogy identified the single object ($A_P$) that
should carry both sides, but provided no route to constructing it without circular assumptions. The
inverted-Stone analogy explained the precise logical obstacle: the inner product must be built, not
assumed, and only $\Lambda\ge0$ can build it unconditionally. Taken together, the three analogies
delimit the space of routes: almost everything that can be made to work in quantum mechanics fails for
$\zeta$ (inverted Stone), almost everything that closes over finite fields fails for $\mathbb Z$
(cohomological Frobenius wall), and the only remaining mechanism --- finite self-adjoint operators with
a bounded resolvent and a Vitali normal-family continuation --- is precisely the construction in
\S\ref{sec:live-proof-steps}. The analogies are not the proof; but they are the map by which the proof was located.


%% ===== Proof-path sections (Overview, Notation, Live proof) =====
\part{The proof path: the ARP-P chain and its single open input}

\section{Overview of the proposed proof path}\label{sec:live-chain}

The proof is organized around a single equivalence,
\[
   \boxed{\textnormal{ARP-P}\Longleftrightarrow\textnormal{RH}},
\]
and the fifteen steps below are the ordered construction of that equivalence together with the one-way
bridge that would turn it into a proof of RH once Step~5 is supplied.  This section explains why the path
takes the shape it does;
Section~\ref{sec:live-proof-steps} carries out each step in full.

The path begins by moving the problem into the variable $z$ with $s=\half+iz$, so that
$\Xi(z)=\xi(\half+iz)$ and the critical line becomes the real axis: RH is exactly the statement that
every zero of $\Xi$ is real (Step~1).  Starting here is what makes the rest possible, because it turns a
statement about the location of zeros into a statement about a \emph{positivity} that complex analysis
can control.  Onto this geometric picture we graft arithmetic through the Weil explicit formula:
Steps~2--4 build finite \emph{channels} $F$ and an explicit meromorphic target $\GXi^F$ whose data is
fixed by the archimedean, pole, and prime terms with no holomorphic ambiguity.  Step~5 then isolates
the one arithmetic input, ARP-P: a Pick (positive-kernel) condition on this explicit data, tested only
at finite node sets in the lower half-plane $U_-$.

Everything after Step~5 is a closed analytic bridge that converts that positivity into the reality of
the zeros.  Positive Cauchy kernels furnish Nevanlinna--Pick interpolants (Steps~6--8); Montel
normality extends them holomorphically from $U_-$ across $\Omega_-$ (Step~9); any off-line zero would
have to appear as a finite-channel residue and is therefore excluded (Steps~10--11); and the
functional-equation symmetry $\Xi(-z)=\Xi(z)$ together with conjugation propagates the conclusion to
the whole plane (Steps~12--13), giving RH.  Steps~14--15 prove the converse direction
RH~$\Rightarrow$~ARP-P, closing the equivalence.  The design deliberately concentrates all the
difficulty into Step~5: every other step is a classical or complex-analytic implication whose
hypotheses are checked where they are used, so that establishing ARP-P alone would complete a proof of
RH.  Step~5 is the single open input, and the remainder of the paper is devoted to it.

\subsection{Claims, status, and dependency map}\label{sec:status-map}

This subsection is the audit interface of the paper: the typed vocabulary, the two dependency maps, and
the status tables.  A reader who absorbs these pages can verify, by following the references, that every
claim of the paper is of exactly the announced strength.

\paragraph{Logical objects.}  Five kinds of objects appear, and they are never interchangeable.
\begin{itemize}[leftmargin=1.8em]
\item A \textbf{Step} is a numbered logical stage in the fifteen-step architecture.  A Step may state a
construction, an implication, an equivalence, or an input condition; its logical type is recorded in the
status table below.
\item A \textbf{logical relation} is the typed relation (implication, equivalence, reduction) asserted
between statements; arrows in the maps below carry their type.
\item An \textbf{H-level} (H0--H8) is a layer of the internal analysis of Step~5
(\S\ref{sec:live-proof-steps}); H-levels are not Steps.
\item An \textbf{$\Omega$-link} ($\Omega_1$--$\Omega_6$) is a closed reduction or equivalence in the
difficulty-localization chain inside Step~5 (\S\ref{subsec:omega7-attack}).
\item The \textbf{terminal residue} is $\Omega_7$: the terminal statement the $\Omega$-chain isolates.
It is not a link, and it is not a Step.
\end{itemize}
Fourteen of the fifteen Steps are closed in the precise sense recorded below: their asserted
construction, implication, equivalence, or conditional consequence is proved.  The unresolved
mathematical input is Step~5.  Two downstream Steps (9 and~11) are proved conditional consequences of
that input.

\paragraph{Proof-dependency map (implications; every arrow is a proved theorem).}
\[
   \boxed{\textnormal{ARP-P (Step 5)}}
   \;\Longrightarrow\;
   \textnormal{NP interpolants}
   \;\Longrightarrow\;
   \textnormal{normal family}
   \;\Longrightarrow\;
   \textnormal{no off-real poles}
   \;\Longrightarrow\;
   \textnormal{RH}
\]
\[
   \textnormal{RH}\;\Longrightarrow\;\textnormal{ARP-P}
   \qquad\text{(Step 14, Theorem~\ref{thm:rh-pick}; the forward chain is Theorem~\ref{thm:pick-rh})},
\]
closing the equivalence $\textnormal{ARP-P}\Longleftrightarrow\textnormal{RH}$
(Corollary~\ref{cor:live-equivalence}).

\paragraph{Difficulty-localization map (reductions, $\rightsquigarrow$ = ``is reduced to / is serialized
as''; these arrows re-express the open input, they do not discharge it).}
\[
   \boxed{\textnormal{Step 5}}
   \;\rightsquigarrow\;
   \textnormal{H0--H8}
   \;\rightsquigarrow\;
   \textnormal{terminal defect }\delta_N
   \;\rightsquigarrow\;
   \textnormal{boundary positivity}
   \;\Longleftrightarrow\;
   \lambda_n\ge0\ \ (\Omega_7)
\]
(the final arrow is typed as an equivalence because it is one: the proved $\Omega$-link $\Omega_6$,
Theorem~\ref{thm:li-dictionary}), and
\[
   \boxed{\;\Omega_7\Longleftrightarrow\textnormal{RH}\;}
   \qquad\text{[Li \cite{Li1997}, classical].}
\]

\paragraph{Status of the fifteen Steps.}

{\small
\begin{longtable}{@{}p{0.04\linewidth}p{0.48\linewidth}p{0.22\linewidth}p{0.18\linewidth}@{}}
\toprule
Step & Statement & Status & Provenance \\
\midrule
\endfirsthead
\toprule
Step & Statement & Status & Provenance \\
\midrule
\endhead
1 & Fix $s=\half+iz$, $\Xi(z)=\xi(\half+iz)$, and
$U_-=\{\Im z<-\half\}$ & Closed & classical \\
2 & Define finite channels $F$ and the explicit target $\GXi^F$ & Closed & new construction \\
3 & Prove that $\GXi^F$ is meromorphic and has no holomorphic ambiguity & Closed & proved here \\
4 & Express $\GXi^F$ on $U_-$ through the Weil explicit formula & Closed as interface & classical
interface \\
5 & Define ARP-P: Pick positivity of $\GXi^F$ on finite node sets in $U_-$ & \textbf{Open}; equivalent
to RH & open input \\
6 & Prove that Cauchy transforms of positive matrix measures have positive Pick kernels & Closed &
classical \\
7 & Prove that tower-ARP implies ARP-P & Closed & proved here \\
8 & Prove that ARP-P supplies Nevanlinna--Pick interpolants & Closed & classical NP \\
9 & Use NP normality/Montel to extend holomorphically from $U_-$ to $\Omega_-$ & Closed conditional on
ARP-P & conditional consequence \\
10 & Detect every off-line zero by some finite channel residue & Closed & proved here \\
11 & Conclude that ARP-P removes off-line poles in $\Omega_-$ & Closed conditional on ARP-P &
conditional consequence \\
12 & Use the symmetries $\Xi(-z)=\Xi(z)$ and conjugation to cover all of $\C\setminus\R$ & Closed &
classical symmetry \\
13 & Translate real zeros of $\Xi(z)$ into RH & Closed & classical \\
14 & Prove the converse direction RH $\Rightarrow$ ARP-P & Closed & proved here \\
15 & Conclude the structural equivalence ARP-P $\Longleftrightarrow$ RH & Closed as equivalence &
logical assembly \\
\bottomrule
\end{longtable}
}

\paragraph{Status of the $\Omega$-chain (proofs in \S\ref{subsec:omega7-attack}).}

{\small
\begin{longtable}{@{}p{0.05\linewidth}p{0.60\linewidth}p{0.28\linewidth}@{}}
\toprule
 & Statement & Status \\
\midrule
\endhead
$\Omega_1$ & RH $\Longleftrightarrow$ all $\Xi$-zeros real & $\Omega$-link, closed
(Lemma~\ref{lem:rh-translation}) \\
$\Omega_2$ & real zeros $\Longleftrightarrow$ ARP-P & $\Omega$-link, closed
(Corollary~\ref{cor:live-equivalence}) \\
$\Omega_3$ & ARP-P $\Longleftrightarrow$ $\delta_N\ge0\ \forall N$ (reference-whitened) & $\Omega$-link,
closed (Corollary~\ref{cor:arp-terminal-defect}) \\
$\Omega_4$ & $\delta_N\ge0$ $\Longleftrightarrow$ one-sided whitened domination & $\Omega$-link, closed
(Theorem~\ref{thm:whitened-terminal}) \\
$\Omega_5$ & interior positivity has a regular boundary limit, per $N$ & $\Omega$-link, closed
(Lemma~\ref{lem:boundary-continuity}) \\
$\Omega_6$ & boundary positivity $\Longleftrightarrow$ Li--Keiper $\lambda_n\ge0$ & $\Omega$-link,
closed (Theorem~\ref{thm:li-dictionary}) \\
$\Omega_7$ & $\lambda_n\ge0$ for all $n$ & \textbf{terminal residue}: open; $\equiv$ RH (Li
\cite{Li1997}) \\
\bottomrule
\end{longtable}
}

\paragraph{Proof references for the chain.}
Steps 1--4 are Notations~\ref{not:variables}, \ref{not:xi}, \ref{not:channels},
Lemma~\ref{lem:xi-basic}, Lemma~\ref{lem:GXi-meromorphic}, and
Theorem~\ref{thm:explicit-channel}.  Step 5 is
Definition~\ref{def:ARP-P}.  Step 6 is Lemma~\ref{lem:cauchy-pick}.  Step 7 is
Theorem~\ref{thm:tower-implies-pick}.  Steps 8--13 are Theorem~\ref{thm:pick-rh}, using
Lemma~\ref{lem:residue} and Lemma~\ref{lem:xi-basic}.  Step 14 is
Theorem~\ref{thm:rh-pick}.  Step 15 is Corollary~\ref{cor:gap-status}.

\subsection{Current state of Step 5}\label{subsec:live-step5}

Step~5 is the single open input.  The present paper does not claim to prove ARP-P; it decomposes it
into precise closed subcases, equivalent formulations, and the routes attempted against it.

\begin{longtable}{@{}p{0.10\linewidth}p{0.48\linewidth}p{0.34\linewidth}@{}}
\toprule
Substep & Current status & What remains \\
\midrule
\endfirsthead
\toprule
Substep & Current status & What remains \\
\midrule
\endhead
H0 & Closed: conservation of RH-level difficulty & Nothing \\
H1 & Closed: scalar level-one ARP-P, $\Re(\xi'/\xi)(s)\ge0$ for $\Re s>1$ & Matrix-channel level one is not routine; it has phase content \\
H2 & Closed as a fixed-band separated-node criterion; geometric high-node subcase closed & Apply the criterion to broader separated configurations \\
H3 & Closed in the large-height logarithmic region and a.e.; exact $\Theta\le W$ reformulation and mirror-pair criterion closed & PDH/global phase decoherence remains open but is not needed for the minimal bridge \\
H4 & Closed as determinant serialization: RH $\Longleftrightarrow d_N^{(k)}>0$ for all $k,N$ & Optional: rigorous determinant margins with the corrected C7 kernel \\
H5 & Closed as equivalence: windowed Weil positivity $\Longleftrightarrow$ ARP-P & Window positivity is an equivalent coordinate; the minimal essential input is scalar fixed-sequence positivity \\
H6 & Quantitative single-cluster witness closed for each configuration & Uniform bounds over configuration families, if desired \\
H7 & Scalar level-two band positivity closed for all configurations & SOS remains as identity/effectivity program \\
H8 & Scalar fixed-$N$ band positivity closed for all configurations; one-point jet serialization, Pascal--Laguerre floor, $\log\log$ range, and sine-model defect constant formulated & Uniform no-margin terminal contractivity $\delta_N\ge0$ for all $N$ is the remaining open point \\
\bottomrule
\end{longtable}

\section{Notation}\label{sec:notation}

This section is deliberately a self-contained symbol table: its entries repeat, verbatim, the
definitions introduced in Steps~1--2, so that the reader can consult them without searching the chain.
The governing statements are the ones in \S\ref{sec:live-proof-steps}.

\begin{notation}[variables and regions]\label{not:variables}
We use
\[
   s=\half+iz,\qquad z=\frac{s-\half}{i},\qquad
   \Re s=\half-\Im z.
\]
Thus
\[
   \Omega_-:=\{z:\Im z<0\},\qquad
   U_-:=\{z:\Im z<-\half\}=\{s:\Re s>1\}.
\]
A nontrivial zero $\rho=\beta+i\gamma$ of $\zeta$ corresponds to
\[
   z_\rho=\gamma-i(\beta-\half).
\]
The Riemann Hypothesis is equivalent to all $z_\rho$ being real.
\end{notation}

\begin{notation}[completed function and Herglotz coordinate]\label{not:xi}
Let
\[
   \xi(s)=\half s(s-1)\pi^{-s/2}\Gamma(s/2)\zeta(s),\qquad
   \Xi(z)=\xi(\half+iz).
\]
The fixed one-variable coordinate is
\[
   M_\Xi(z):=-\frac{\Xi'(z)}{\Xi(z)}.
\]
Since $\Xi'(z)=i\,\xi'(s)$,
\begin{equation}\label{eq:change}
   M_\Xi(z)=i\left(-\frac{\xi'(s)}{\xi(s)}\right),\qquad s=\half+iz.
\end{equation}
This factor $i$ is part of every endpoint identity below.
\end{notation}

\begin{notation}[finite channels]\label{not:channels}
For a finite source plane $F=\operatorname{span}\{\phi_1,\ldots,\phi_m\}$, define the source functional
\[
   \ell_a(\phi)=\int_0^\infty \phi(t)e^{-iat}\,dt.
\]
The target finite-channel function is the explicit meromorphic sum
\begin{equation}\label{eq:GXiF}
   \GXi^F(z)_{\alpha\beta}
   :=
   \sum_\rho
   \frac{
      m_\rho\,\ell_{z_\rho}(\phi_\beta)\overline{\ell_{\bar z_\rho}(\phi_\alpha)}
   }{z_\rho-z},
\end{equation}
where the sum is over the zeros of $\Xi$, counted with multiplicity $m_\rho$.  Thus ARP refers to this
specific meromorphic matrix function, not merely to a prescribed principal part.  It is not the scalar
multiple $M_\Xi(z)[\langle\phi_\beta,\phi_\alpha\rangle]$.

For this matrix function form the Pick kernel
\[
   \KXi^F(z,w)=\frac{\GXi^F(z)-\GXi^F(w)^*}{z-\bar w}.
\]
Vitali is applied to the one-variable functions $\GXi^F$.  Negative squares are read only from the
two-variable kernels $\KXi^F$.

This notation is a global preview.  The governing finite-channel definition used in the proof is
Definition~\ref{def:target-channel}; the formulas are identical.
\end{notation}

\section{The chain, step by step}\label{sec:live-proof-steps}

This section is the ordered proof spine of the paper.  The headings follow the fifteen Steps of
Section~\ref{sec:live-chain}.  A step marked ``closed'' is either proved in-line below or reduced to a
named classical theorem whose hypotheses are checked at the point of use.  No mathematical assertion
needed by the route is left to an appendix.  The single non-closed input is Step~5, ARP-P.

\subsection{Step 1: variables, completed function, and the RH translation}

Set
\[
   s=\half+iz,\qquad z=\frac{s-\half}{i},\qquad
   \Xi(z)=\xi(\half+iz),
\]
where
\[
   \xi(s)=\half s(s-1)\pi^{-s/2}\Gamma(s/2)\zeta(s).
\]
Then
\[
   \Re s=\half-\Im z,\qquad
   U_-=\{z:\Im z<-\half\}=\{s:\Re s>1\}.
\]
If $\rho=\beta+i\gamma$ is a nontrivial zero of $\zeta$, the corresponding zero of $\Xi$ is
\[
   z_\rho=\frac{\rho-\half}{i}=\gamma-i(\beta-\half).
\]
The logarithmic derivative coordinate used throughout is
\[
   M_\Xi(z):=-\frac{\Xi'(z)}{\Xi(z)}
   =
   i\left(-\frac{\xi'(s)}{\xi(s)}\right),\qquad s=\half+iz.
\]
The factor $i$ is fixed here and is not a later normalization.  The whole route rests on the following
elementary but decisive translation, which turns a statement about the location of the zeros into a
statement about the real axis.

\begin{lemma}[basic properties of $\Xi$]\label{lem:xi-basic}
The function $\Xi(z)=\xi(\half+iz)$ is entire, and:
\begin{enumerate}[label=\textnormal{(\roman*)}]
\item its zero multiset is exactly $\{z_\rho=(\rho-\half)/i\}$, $\rho$ running over the nontrivial
zeros of $\zeta$ with multiplicity; in particular every zero satisfies $|\Im z_\rho|<\half$;
\item $\Xi(-z)=\Xi(z)$;
\item $\Xi(\bar z)=\overline{\Xi(z)}$; in particular $\Xi$ is real on $\R$, and the zero multiset is
invariant under $z\mapsto\bar z$ and $z\mapsto-z$ with multiplicities preserved.
\end{enumerate}
\end{lemma}

\begin{proof}
$\xi(s)=\half s(s-1)\pi^{-s/2}\Gamma(s/2)\zeta(s)$ is entire: $s-1$ cancels the pole of $\zeta$ at
$s=1$, the poles of $\Gamma(s/2)$ at $s=0,-2,-4,\ldots$ are cancelled by the factor $s$ and by the
trivial zeros of $\zeta$, and $\xi(0)=\xi(1)=\half\ne0$ \cite[Ch.~II]{Titchmarsh}.  Since
$\Gamma(s/2)$ never vanishes and $\zeta(s)\ne0$ for $\Re s\ge1$ (Euler product and the classical
zero-free line), the functional equation $\xi(s)=\xi(1-s)$ excludes zeros in $\Re s\le0$ as well;
hence the zeros of $\xi$ are precisely the nontrivial zeros of $\zeta$, all in $0<\Re s<1$, which is
$|\Im z_\rho|<\half$ in the $z$-coordinate.  This proves (i).  For (ii),
$\Xi(-z)=\xi(\half-iz)=\xi(1-(\half-iz))=\xi(\half+iz)=\Xi(z)$ by the functional equation.  For (iii),
$\xi$ is real on $(1,\infty)$, so Schwarz reflection gives $\xi(\bar s)=\overline{\xi(s)}$; therefore
$\overline{\Xi(z)}=\xi(\half-i\bar z)=\xi(\half+i\bar z)=\Xi(\bar z)$, using the functional equation
once more.  The zero-multiset symmetries follow.
\end{proof}

\begin{lemma}[RH is realness of the $\Xi$-zeros]\label{lem:rh-translation}
The Riemann Hypothesis is equivalent to the assertion that every nontrivial zero of $\Xi(z)$ is real.
\end{lemma}

\begin{proof}
For a nontrivial zero $\rho=\beta+i\gamma$ of $\zeta$, the change of variables gives
$z_\rho=(\rho-\half)/i=\gamma-i(\beta-\half)$, and by Lemma~\ref{lem:xi-basic}(i) the zeros of $\Xi$
are exactly the $z_\rho$, with multiplicity.  Hence $z_\rho\in\R$ if and only if $\beta=\half$, so every
zero of $\Xi$ is real precisely when every $\rho$ has $\beta=\half$, which is RH.
\end{proof}

\subsection{Step 2: finite channels and the target \texorpdfstring{$\GXi^F$}{GXiF}}

Steps 2--4 are carried out formally in Definition~\ref{def:core}, Definition~\ref{def:target-channel},
Lemma~\ref{lem:GXi-meromorphic}, Lemma~\ref{lem:weil-admissible}, and Theorem~\ref{thm:explicit-channel};
the surrounding discussion explains their content but the formal statements are authoritative.

\begin{definition}[source core]\label{def:core}
The source core $\cS_a^\circ$ is the set of functions $\phi:(0,\infty)\to\C$ satisfying:
\begin{enumerate}[label=\textnormal{(S\arabic*)}]
\item \textnormal{(profile decay)} there is $a_\phi>\half$ such that, for every $N$,
$|\phi(t)|\le C_{N,\phi}\,e^{-a_\phi t}(1+t)^{-N}$;
\item \textnormal{(enlarged-strip transform)} there is $\eps_\phi\in(0,a_\phi-\half)$ such that
\[
   \ell_w(\phi)=\int_0^\infty\phi(t)e^{-iwt}\,dt
\]
is holomorphic on $|\Im w|<\half+\eps_\phi$ and rapidly decreasing horizontally, uniformly on the
closed substrip: for every $N$,
\[
   |\ell_w(\phi)|\le C_{N,\phi}(1+|\Re w|)^{-N},
   \qquad |\Im w|\le\half+\tfrac{\eps_\phi}{2}.
\]
\end{enumerate}
Every $\phi\in C_c^\infty(0,\infty)$ belongs to $\cS_a^\circ$ with arbitrary $a_\phi>\half$ and
$\eps_\phi\in(0,a_\phi-\half)$: \textnormal{(S1)} is trivial by compact support, and
\textnormal{(S2)} follows by repeated integration by parts, since
$|e^{-iwt}|\le e^{(\half+\eps_\phi)T}$ on the support.  Note that \textnormal{(S1)} makes
the absolute convergence in \textnormal{(S2)} automatic; the content of \textnormal{(S2)} is the
horizontal rapid decrease.
\end{definition}

\begin{remark}[realization-layer compatibility]\label{rmk:core-compat}
The realization layer of Section~\textnormal{\ref{sec:shorting}} only requires the weaker profile decay
$a>\tfrac14$ for the prime-cell summability; since $a_\phi>\half>\tfrac14$, every source-core element
is admissible there.  In that layer the profile $\phi$ is denoted $\widehat\phi$ (the logarithmic
Mellin profile); the two notations refer to the same function and the hat is never the Fourier
transform of Step~4.
\end{remark}

\noindent With this test space fixed, the object the whole chain interpolates is the following explicit
meromorphic matrix function.

\begin{definition}[finite channels and the target $\GXi^F$]\label{def:target-channel}
For a finite source plane $F=\operatorname{span}\{\phi_1,\ldots,\phi_m\}\subset\cS_a^\circ$ and sources
$\phi_\alpha,\phi_\beta\in F$, set $\ell_a(\phi)=\int_0^\infty \phi(t)e^{-iat}\,dt$ and define the target
channel
\[
   \GXi^F(z)_{\alpha\beta}
   =
   \sum_\rho
   \frac{
      m_\rho\,\ell_{z_\rho}(\phi_\beta)\overline{\ell_{\bar z_\rho}(\phi_\alpha)}
   }{z_\rho-z},
\]
the sum running over the zeros $z_\rho$ of $\Xi$ with multiplicity $m_\rho$.
\end{definition}

\noindent This is not defined only up to a holomorphic summand: the residues, source weights, zero
multiplicities, and denominator convention are all fixed in the formula.

\subsection{Step 3: meromorphy and no holomorphic ambiguity}

The source core $\cS_a^\circ$ is chosen so that its Laplace transforms have rapid horizontal decay on a
closed strip containing the zero strip $|\Im a|\le\half$.  Since the zeros of $\Xi$ satisfy the classical
counting estimate
$N(T)=O(T\log T)$, for every compact set $K$ avoiding the zeros and every $N$ one has
\[
   \left|
   \frac{
      \ell_{z_\rho}(\phi_\beta)\overline{\ell_{\bar z_\rho}(\phi_\alpha)}
   }{z_\rho-z}
   \right|
   \le C_{K,F,N}(1+|\Re z_\rho|)^{-N}(1+|z_\rho|)^{-1}.
\]
Choosing $N$ large makes the zero sum locally uniformly convergent away from the zeros.  Therefore
$\GXi^F$ is meromorphic and its principal part at $z_\rho$ is exactly
\[
   \operatorname{pp}_{z=z_\rho}\GXi^F(z)_{\alpha\beta}
   =
   \frac{m_\rho\,\ell_{z_\rho}(\phi_\beta)
      \overline{\ell_{\bar z_\rho}(\phi_\alpha)}}{z_\rho-z}.
\]
This proves the absence of holomorphic ambiguity: the target is the displayed meromorphic function, not a
principal-part prescription modulo entire terms.


\begin{lemma}[target channel meromorphy]\label{lem:GXi-meromorphic}
For every finite source plane $F\subset\cS_a^\circ$, the series \eqref{eq:GXiF} converges locally uniformly
away from the zeros of $\Xi$ and defines a meromorphic matrix function.  At a zero $z_\rho$ its principal
part is
\[
   \operatorname{pp}_{z=z_\rho}\GXi^F(z)_{\alpha\beta}
   =
   \frac{m_\rho\,\ell_{z_\rho}(\phi_\beta)\overline{\ell_{\bar z_\rho}(\phi_\alpha)}}{z_\rho-z}.
\]
Moreover $\GXi^F(\bar z)=\GXi^F(z)^*$.
\end{lemma}

\begin{proof}
By Definition~\ref{def:core}, source vectors have Laplace transforms with rapid horizontal decay uniformly
on a closed strip containing $|\Im a|\le\half$; compact logarithmic tests have this property by the
Paley--Wiener/integration-by-parts estimate, and the general source core is defined to keep the same
estimate.  Since the zeros of $\Xi$ lie in $|\Im a|<\half$ by Lemma~\ref{lem:xi-basic}(i) and satisfy
$N(T)=O(T\log T)$, for every $N$ and every compact $K$ avoiding the zeros,
\[
   \left|
   \frac{
      \ell_{z_\rho}(\phi_\beta)\overline{\ell_{\bar z_\rho}(\phi_\alpha)}
   }{z_\rho-z}
   \right|
   \le
   C_{K,F,N}(1+|\Re z_\rho|)^{-N}(1+|z_\rho|)^{-1}.
\]
Choosing $N$ large gives a summable majorant over the zero set.  The only singularities are the displayed
simple poles, with multiplicities inherited from $\Xi$.

Two points deserve record.  First, near a fixed zero $z_{\rho_0}$ the series splits as the displayed
principal part plus the sum over $z_\rho\ne z_{\rho_0}$, which converges locally uniformly on a
punctured disk by the same majorant; hence $z_{\rho_0}$ is a simple pole of the matrix function with
exactly the displayed principal part, and $\GXi^F$ is meromorphic.  Second, the factor
$w\mapsto\overline{\ell_{\bar w}(\phi_\alpha)}$ is itself holomorphic, by the identity
\[
   \overline{\ell_{\bar w}(\phi_\alpha)}
   =\int_0^\infty\overline{\phi_\alpha(t)}\,e^{iwt}\,dt ,
\]
convergent on the same strip by \textnormal{(S1)}.  Finally, the conjugation symmetry of the zero
multiset used in pairing $z_\rho$ with $\bar z_\rho$ in \eqref{eq:GXiF} is
Lemma~\ref{lem:xi-basic}(iii), and it gives $\GXi^F(\bar z)=\GXi^F(z)^*$.
\end{proof}

\subsection{Step 4: Weil explicit formula on \texorpdfstring{$U_-$}{U-}}

For $z\in U_-$ define
\[
   h_{\alpha\beta,z}(w)
   =
   \frac{\ell_w(\phi_\beta)\overline{\ell_{\bar w}(\phi_\alpha)}}{w-z}.
\]
Because $z$ lies below the zero strip, the denominator has no pole in a slightly enlarged closed strip
around $|\Im w|\le\half$, and the source decay makes $h_{\alpha\beta,z}$ an admissible Weil test.  The
classical Weil explicit formula gives
\[
   \GXi^F(z)_{\alpha\beta}
   =
   \mathcal W_\Xi(h_{\alpha\beta,z}),
\]
where $\mathcal W_\Xi$ is the usual archimedean--pole--prime functional:
\begin{align*}
   \mathcal W_\Xi(h)
   &=
   h(i/2)+h(-i/2)
   -\frac{\log\pi}{2\pi}\int_\R h(t)\,dt \\
   &\quad
   +\frac{1}{2\pi}\int_\R h(t)
      \Re\frac{\Gamma'}{\Gamma}\!\left(\frac14+\frac{it}{2}\right)\,dt
   -\frac{1}{2\pi}\sum_{n\ge2}\frac{\Lambda(n)}{\sqrt n}
      \bigl(\widehat h(\log n)+\widehat h(-\log n)\bigr).
\end{align*}
The proof is the standard smoothing argument: apply the explicit formula first to
$e^{-\varepsilon w^2}h_{\alpha\beta,z}(w)$, where zero and prime sums converge absolutely, and let
$\varepsilon\downarrow0$ in the Weil test topology.  Thus the target values on $U_-$ are computable from
the arithmetic side without using zeros as input.

\begin{lemma}[admissibility of the channel tests]\label{lem:weil-admissible}
Let $z\in U_-$ and $\phi_\alpha,\phi_\beta\in\cS_a^\circ$, and set
$\eps:=\tfrac12\min\{\eps_{\phi_\alpha},\eps_{\phi_\beta},\,|\Im z|-\half\}>0$.  Then
$h_{\alpha\beta,z}(w)=\ell_w(\phi_\beta)\overline{\ell_{\bar w}(\phi_\alpha)}\,(w-z)^{-1}$ is
holomorphic on $|\Im w|\le\half+\eps$ and satisfies there, for every $N$,
\[
   |h_{\alpha\beta,z}(w)|\le C_N(1+|\Re w|)^{-N}.
\]
Consequently: \textnormal{(i)} the pole terms $h_{\alpha\beta,z}(\pm i/2)$ are finite;
\textnormal{(ii)} the archimedean integrals converge absolutely against the logarithmic growth of
$\Gamma'/\Gamma$; \textnormal{(iii)} shifting the contour of
$\widehat h(u)=\int_\R h(t)e^{-iut}\,dt$ to $\Im w=\mp(\half+\eps)$ gives
\[
   |\widehat h(\log n)|+|\widehat h(-\log n)|\le C\,n^{-(\half+\eps)},
\]
so the prime side is dominated by $\sum_n\Lambda(n)n^{-1-\eps}<\infty$ and converges absolutely.
\end{lemma}

\begin{proof}
Both Laplace factors are holomorphic and rapidly decreasing on $|\Im w|\le\half+\eps$ by
\textnormal{(S2)} of Definition~\ref{def:core}, the second via
$\overline{\ell_{\bar w}(\phi_\alpha)}=\int_0^\infty\overline{\phi_\alpha(t)}e^{iwt}dt$.  The
denominator satisfies $|w-z|\ge|\Im z|-\half-\eps\ge\eps>0$ on the strip, since
$\Im z\le-\half-2\eps$; hence $h$ is holomorphic there with the displayed decay.  \textnormal{(i)}
follows because $\pm i/2$ is interior to the strip and distinct from $z$.  \textnormal{(ii)} is the
decay against $\log(2+|t|)$.  For \textnormal{(iii)}: for $u=\log n>0$ shift the contour to
$\Im w=-(\half+\eps)$, legal by holomorphy and horizontal decay, so that
$|e^{-iuw}|=e^{u\Im w}=n^{-(\half+\eps)}$; for $u=-\log n$ shift to
$\Im w=+(\half+\eps)$.  In both cases no pole is crossed.
\end{proof}

\begin{theorem}[explicit channel formula on $U_-$]\label{thm:explicit-channel}
Let $z\in U_-$ and define, for $\phi_\alpha,\phi_\beta\in F$,
\[
   h_{\alpha\beta,z}(w)
   :=
   \frac{\ell_w(\phi_\beta)\overline{\ell_{\bar w}(\phi_\alpha)}}{w-z}.
\]
Then $h_{\alpha\beta,z}$ is holomorphic in a strip containing all zeros of $\Xi$ and rapidly decreasing
on horizontal lines in that strip, and
\[
   \GXi^F(z)_{\alpha\beta}=\mathcal W_\Xi(h_{\alpha\beta,z}),
\]
where $\mathcal W_\Xi$ is the classical Weil explicit-formula functional
\begin{align*}
   \mathcal W_\Xi(h)
   &=
   h(i/2)+h(-i/2)
   -\frac{\log\pi}{2\pi}\int_\R h(t)\,dt \\
   &\quad
   +\frac{1}{2\pi}\int_\R h(t)\,
      \Re\frac{\Gamma'}{\Gamma}\!\left(\frac14+\frac{it}{2}\right)\,dt
   -\frac{1}{2\pi}\sum_{n\ge2}\frac{\Lambda(n)}{\sqrt n}
      \bigl(\widehat h(\log n)+\widehat h(-\log n)\bigr),
\end{align*}
with Fourier convention
\[
   \widehat h(u)=\int_\R h(t)e^{-iut}\,dt .
\]
The displayed identity is understood in the standard Weil explicit-formula sense for admissible
strip-holomorphic rapidly decreasing tests; equivalently, by inserting the usual exponential smoothing in
the zero and prime sums and then removing the smoothing.
\end{theorem}

\begin{proof}
Lemma~\ref{lem:weil-admissible} makes $h_{\alpha\beta,z}$ an admissible test function for the Weil
explicit formula.  The prime side is the standard continuous extension of the explicit-formula
functional from exponentially smoothed tests: apply the formula first to
$h_{\alpha\beta,z,\varepsilon}(w)=e^{-\varepsilon w^2}h_{\alpha\beta,z}(w)$ on horizontal contours inside
the zero strip, where the smoothed sums are absolutely convergent, and let $\varepsilon\downarrow0$ in
the Weil test topology.  This is the classical formulation of the explicit formula for admissible
Schwartz strip tests.

By Lemma~\ref{lem:weil-admissible} all three sides of the smoothed identity are dominated, uniformly in
$\varepsilon\in(0,1]$, by convergent majorants: the zero side by the majorant of
Lemma~\ref{lem:GXi-meromorphic}, the archimedean side by \textnormal{(ii)}, and the prime side by
\textnormal{(iii)}, since $|e^{-\varepsilon w^2}|\le e^{\varepsilon(\half+\eta)^2}\le C$ on the strip,
where $\eta$ is the positive strip margin supplied by Lemma~\ref{lem:weil-admissible}.
Dominated convergence therefore passes the classical formula \cite{Weil1952} to the limit test
$h_{\alpha\beta,z}$, term by term.

Applying the Weil explicit formula to $h_{\alpha\beta,z}$ gives
\[
   \mathcal W_\Xi(h_{\alpha\beta,z})
   =
   \sum_\rho m_\rho h_{\alpha\beta,z}(z_\rho)
   =
   \sum_\rho
   \frac{
      m_\rho\,\ell_{z_\rho}(\phi_\beta)\overline{\ell_{\bar z_\rho}(\phi_\alpha)}
   }{z_\rho-z}
   =
   \GXi^F(z)_{\alpha\beta}.
\]
Equivalently, the Cauchy denominator may be written in $U_-$ as the unilateral Laplace kernel
\[
   \frac1{z_\rho-z}
   =
   -i\int_0^\infty e^{-izt}e^{iz_\rho t}\,dt,
\]
because $\Im z<-\half$ and $|\Im z_\rho|<\half$.  This is the transfer-side representation of the same target value.  It does
not prove ARP; it proves that the ARP target on $U_-$ is a completely specified arithmetic distribution,
with arithmetic side fixed by the Weil formula and no unspecified holomorphic summand.
\end{proof}

\subsection{Step 5: ARP-P}

\textbf{Step 5 (ARP-P) is the only one of the fifteen Steps that is not closed.}  Steps 1--4 and 6--15
are closed in the sense of the status map (\S\ref{sec:status-map}), Steps~9 and~11 as proved conditional
consequences of Step~5; the whole path reduces RH to this one input.  ARP-P in turn reduces, through the
$\Omega$-chain (Section~\ref{subsec:omega7-attack}), to the single terminal-defect inequality
\[
   \delta_N\ge0\ \ \text{for all }N\quad\Longleftrightarrow\quad\textnormal{RH}\qquad(\Omega_7),
\]
equivalently $\lambda_n\ge0$ for all $n$ (Li--Keiper).  This is the \emph{sole} pending statement of the
proposed proof; every other link in the path is established.  The remainder of Step~5 decomposes
$\Omega_7$ into closed subcases and records the routes that were tried against it, none of which closes
it.

ARP-P is the assertion that for every finite source plane $F$ and every finite node set
$z_1,\ldots,z_N\in U_-$, the block Pick matrix
\[
   \mathsf P[g^F;z_1,\ldots,z_N]
   =
   \left[
      \frac{g^F(z_i)-g^F(z_j)^*}{z_i-\bar z_j}
   \right]_{i,j=1}^N
\]
is positive semidefinite, where $g^F=\GXi^F|_{U_-}$.  This is the only open input in the chain.  The
rest of the paper proves that this input is equivalent to RH and decomposes it into lower-dimensional
closed regimes and the routes attempted against the residue.


\begin{definition}[arithmetic channel data]\label{def:arith-data}
For a finite source plane $F\subset\cS_a^\circ$ define
\[
   g^F(z):=\GXi^F(z),\qquad z\in U_-.
\]
The values are computed through Theorem~\ref{thm:explicit-channel}: for each $z\in U_-$,
$g^F(z)_{\alpha\beta}$ is the Weil explicit-formula functional applied to the admissible test
$h_{\alpha\beta,z}$.  Thus the arithmetic side is fixed by the archimedean, pole/degree, and prime terms,
with the standard Weil smoothing convention; no unspecified holomorphic summand remains.
\end{definition}

\begin{definition}[ARP-P: arithmetic Pick positivity]\label{def:ARP-P}
The channel data satisfy \emph{arithmetic Pick positivity} (ARP-P) if, for every finite source plane
$F=\operatorname{span}\{\phi_1,\ldots,\phi_m\}\subset\cS_a^\circ$ and every finite node set
$\{z_1,\ldots,z_N\}\subset U_-$, the block Pick matrix
\[
   \mathsf P[g^F;z_1,\ldots,z_N]
   :=
   \left[
      \frac{g^F(z_i)-g^F(z_j)^*}{z_i-\bar z_j}
   \right]_{i,j=1}^{N}
   \in M_{Nm}(\C)
\]
is positive semidefinite.
\end{definition}

\subsection{Step 5 substructure: current closed subcases and open fronts}

This subsection descends one level in the logical hierarchy of \S\ref{sec:status-map}: from the
fifteen-step chain into the internal H-level analysis (H0--H8) of the single open Step.  Nothing here is
a Step of the chain; everything here lives inside Step~5.

The present decomposition of Step~5 is not a list of slogans.  It is the working mathematical hierarchy
for attacking the only open input:
\[
   \textnormal{ARP-P}
   \quad\Longleftrightarrow\quad
   \textnormal{positive arithmetic Pick kernels for all finite channels}
\]
\[
   \textnormal{and all finite node sets in }U_-.
\]
Everything below belongs to the proof as it stands so far.  The closed items are proved here; the open items are
formulated here in the exact form in which they remain to be attacked.

\subsubsection*{H0 (closed): conservation of RH-level difficulty}

Let $S_1,\ldots,S_k$ be assertions.  If
\[
   S_1\wedge\cdots\wedge S_k\Longrightarrow\mathrm{RH}
\]
and every $S_j$ is proved unconditionally from standard theorems, then RH is proved by modus ponens.
Therefore, in any chain that implies RH but does not prove RH, at least one step is either false or has
RH-level difficulty.  In the chain, that hard step is Step~5, ARP-P.  This observation is logically
closed and prevents the paper from hiding RH inside a false trace-log or local-cell identity.

\subsubsection*{H1 (closed): scalar one-node ARP-P}

For the scalar function
\[
   M_\Xi(z)=i\left(-\frac{\xi'}{\xi}(s)\right),\qquad s=\half+iz,
\]
the one-node Pick quotient is
\[
   \frac{M_\Xi(z)-M_\Xi(z)^*}{z-\bar z}
   =
   \frac{\Im M_\Xi(z)}{\Im z}.
\]
Writing $s=\sigma+it$ and $z=x+iy$, one has $y=-(\sigma-\half)$.  Since
$M_\Xi=i(-\xi'/\xi)$,
\[
   \Im M_\Xi(z)=-\Re\frac{\xi'}{\xi}(s).
\]
Thus
\[
   \frac{\Im M_\Xi(z)}{\Im z}
   =
   \frac{-\Re(\xi'/\xi)(s)}{-(\sigma-\half)}
   =
   \frac{\Re(\xi'/\xi)(s)}{\sigma-\half}.
\]
The scalar one-node Pick condition on $U_-$ is therefore equivalent to
\[
   \Re\frac{\xi'}{\xi}(s)\ge0,\qquad \Re s>1,
\]
because $\sigma-\half>0$.

The inequality is classical.  From the Hadamard product for $\xi$, with the standard real normalization
of the constant term,
\[
   \frac{\xi'}{\xi}(s)
   =
   \sum_\rho\left(\frac1{s-\rho}+\frac1\rho\right)+B,
\]
and after pairing zeros symmetrically, the real parts of the constant terms cancel in the usual limiting
sense.  Hence, for $s=\sigma+it$,
\[
   \Re\frac{\xi'}{\xi}(s)
   =
   \sum_\rho\frac{\sigma-\beta}{|s-\rho|^2}.
\]
The classical zero-free line and functional equation give $0<\beta<1$ for nontrivial zeros.  Therefore
for $\sigma>1$ every summand is nonnegative, and H1 is closed.

\subsubsection*{H2 (closed as a fixed-band separated-node criterion)}

Fix a vertical band
\[
   1<\sigma_0\le\sigma\le\sigma_1<\infty.
\]
For scalar nodes $s_j=\sigma_j+it_j$ in this band, let $z_j=(s_j-\half)/i$ and form the scalar Pick
matrix
\[
   K_{ij}
   =
   \frac{M_\Xi(z_i)-\overline{M_\Xi(z_j)}}{z_i-\bar z_j}.
\]
There are effective constants $a_{\sigma_0,\sigma_1},B_{\sigma_0,\sigma_1}>0$ such that, for all
sufficiently large heights in the band,
\[
   K_{ii}\ge a_{\sigma_0,\sigma_1}\log(3+|t_i|),
\]
and for $i\ne j$,
\[
   |K_{ij}|
   \le
   B_{\sigma_0,\sigma_1}
   \frac{\log(3+|t_i|)+\log(3+|t_j|)}
        {\max\{|t_i-t_j|,\,2(\sigma_0-\half)\}}.
\]
The diagonal bound is H1 plus Stirling's formula for the $\Gamma'/\Gamma$ term and the boundedness of
$\zeta'/\zeta$ in $\sigma\ge\sigma_0$.  The off-diagonal bound follows from
$|M_\Xi(z_i)|\ll_{\sigma_0,\sigma_1}\log(3+|t_i|)$ and
\[
   |z_i-\bar z_j|\ge\max\{|t_i-t_j|,\,2(\sigma_0-\half)\}.
\]
Therefore, if
\[
   \sum_{j\ne i}
   \frac{\log(3+|t_i|)+\log(3+|t_j|)}
        {\max\{|t_i-t_j|,\,2(\sigma_0-\half)\}}
   \le
   \frac{a_{\sigma_0,\sigma_1}}{2B_{\sigma_0,\sigma_1}}\log(3+|t_i|)
\]
for every $i$, then
\[
   \sum_{j\ne i}|K_{ij}|\le\frac12 K_{ii}.
\]
The matrix is Hermitian, so Gershgorin's theorem implies that it is positive semidefinite.  This closes
H2 as a verifiable separated-node criterion.

A closed concrete subcase is geometric separation.  If, after ordering by height,
\[
   |t_j|\ge T_*Q^j,\qquad
   |t_i-t_j|\ge\tfrac12\max\{|t_i|,|t_j|\}\quad(i\ne j),
\]
then
\[
   \sum_{j<i}|K_{ij}|
   \le
   C_1\frac{\log(3+|t_i|)}{|t_i|}\#\{j<i\}
   \le
   C_2\frac{(\log(3+|t_i|))^2}{|t_i|},
\]
while
\[
   \sum_{j>i}|K_{ij}|
   \le
   C_3\sum_{j>i}\frac{\log(3+|t_j|)}{|t_j|}
   \le
   C_4\frac{\log(3+|t_i|)}{|t_i|}.
\]
For $T_*$ large these row sums are smaller than $\frac12K_{ii}$.  Hence geometrically separated high
nodes give a closed positive scalar Pick subcase.

\subsubsection*{H3 (closed in a logarithmic region and almost everywhere): infinitesimal two-node target}

The scalar coincident two-node limit is the Schwarz--Pick inequality
\[
   \left(\sigma-\half\right)
   \left|\left(\frac{\xi'}{\xi}\right)'(s)\right|
   \le
   \Re\frac{\xi'}{\xi}(s),
   \qquad \sigma=\Re s>1.
\]
It is not closed globally.  Two substantial regions are closed.

\emph{Deterministic logarithmic region.}  There exist effective constants $c>0$ and $T_0\ge3$ such that
the inequality holds whenever
\[
   \sigma\ge1+\frac{c}{\sqrt{\log(3+|t|)}}
\]
and $|t|\ge T_0$.  The proof splits into two regions.  In the band $1+\epsilon\le\sigma\le2$,
$0<\epsilon\le1$, the absolutely convergent Dirichlet series gives
\[
   \left|\left(\frac{\zeta'}{\zeta}\right)'(s)\right|
   \le
   \sum_{n\ge2}\Lambda(n)(\log n)n^{-1-\epsilon}
   \le C_1\epsilon^{-2}+C_2.
\]
The elementary factors satisfy
\[
   \left|\frac14\psi'(s/2)-s^{-2}-(s-1)^{-2}\right|
   \le C_3(1+|t|)^{-1},
\]
and therefore
\[
   \left(\sigma-\half\right)
   \left|\left(\frac{\xi'}{\xi}\right)'(s)\right|
   \le C_4\epsilon^{-2}+C_5.
\]
On the other hand, Stirling plus
$\Re(\zeta'/\zeta)(s)\ge-\sum_{n\ge2}\Lambda(n)n^{-1-\epsilon}\ge-C_6\epsilon^{-1}-C_7$
gives
\[
   \Re\frac{\xi'}{\xi}(s)
   \ge
   \frac12\log\frac{|t|}{2\pi}-C_6\epsilon^{-1}-C_7.
\]
Choosing $c$ large enough and then $T_0$ large enough makes the right side dominate the left when
$\epsilon\ge c/\sqrt{\log(3+|t|)}$.

For $\sigma\ge2$, the Euler derivative has exponential decay:
\[
   \left|\left(\frac{\zeta'}{\zeta}\right)'(s)\right|
   \le
   2^{-(\sigma-2)}\sum_{n\ge2}\Lambda(n)(\log n)n^{-2}
   \le C_8\,2^{-\sigma}.
\]
Stirling gives $|\frac14\psi'(s/2)|\le(2|s|)^{-1}+O(|s|^{-2})$, so
\[
   \left(\sigma-\half\right)
   \left|\left(\frac{\xi'}{\xi}\right)'(s)\right|
   \le C_9\sigma2^{-\sigma}+\frac12+O(|s|^{-1}).
\]
For large $\sigma$ and $|t|$ this is $<1$, while
\[
   \Re\frac{\xi'}{\xi}(s)\ge\frac12\log\frac{|s|}{2\pi}-O(1)\ge1.
\]
The finite band $2\le\sigma\le\sigma_2$ is absorbed by the same fixed-band estimates as above.  This
closes the logarithmic region.

\emph{Almost-everywhere edge region.}  There are effective $C,T_1>0$ such that for
$T\ge T_1$ and $\epsilon\in[T^{-1/5},1]$, the two-node inequality holds at
$s=1+\epsilon+it$ for all $t\in[T,2T]$ outside a set of measure at most
$CT(\log T)^{-2}$.  In the band $1<\sigma\le2$, it is enough to have
\[
   \tfrac32\left|\left(\frac{\zeta'}{\zeta}\right)'(1+\epsilon+it)\right|
   +
   \left|\frac{\zeta'}{\zeta}(1+\epsilon+it)\right|
   \le \tfrac14\log T.
\]
The Montgomery--Vaughan mean-value theorem for Dirichlet series gives, uniformly in
$\epsilon\in[T^{-1/5},1]$,
\[
   \int_T^{2T}\left|\frac{\zeta'}{\zeta}(1+\epsilon+it)\right|^2dt\le C_1T,
\]
and
\[
   \int_T^{2T}
   \left|\left(\frac{\zeta'}{\zeta}\right)'(1+\epsilon+it)\right|^2dt\le C_2T.
\]
The square sums
$\sum\Lambda(n)^2n^{-2-2\epsilon}$ and
$\sum\Lambda(n)^2(\log n)^2n^{-2-2\epsilon}$ remain bounded at $\epsilon=0$, and the conservative
Montgomery--Vaughan error is absorbed in the range $\epsilon\ge T^{-1/5}$.  Chebyshev's inequality then
gives an exceptional set of measure $O(T(\log T)^{-2})$.

\subsubsection*{H4 (closed): determinant serialization and Schur coordinates}

The formal closure is Theorem~\ref{thm:det-serialization}: after fixing the dense detection family
$\{\phi_k\}$ and a nested node sequence $\{z_j\}\subset U_-$,
\[
   \mathrm{RH}\Longleftrightarrow d_N^{(k)}>0\quad\text{for all }k,N.
\]
Equivalently, H4 is no longer merely a computational formalism; it is a countable determinant
serialization of RH.  For each finite channel $F$, the Schur--Nevanlinna algorithm applied to the finite
data
\[
   g^F(z_1),\ldots,g^F(z_N)
\]
produces contractivity parameters $\gamma_N^F$ whenever the preceding Pick matrices are positive.  The
formal target is
\[
   \textnormal{ARP-P}
   \quad\Longleftrightarrow\quad
   \|\gamma_N^F\|\le1
   \quad\text{for every }N\text{ and every finite }F.
\]
Degenerate Pick matrices must be handled by the degenerate Nevanlinna--Pick theory, equivalently by
Moore--Penrose Schur complements plus the range condition.  In scalar nondegenerate normalization one
expects
\[
   1-|\gamma_N|^2
   =
   C_N(z_1,\ldots,z_{N+1})
   \frac{\det P_{N+1}\det P_{N-1}}{(\det P_N)^2},
   \qquad C_N>0,
\]
with the geometric constant fixed by the chosen Cayley normalization.  The determinant theorem
Theorem~\ref{thm:det-serialization} is the governing closed statement: the Schur parameters are the same
serialization in recursive coordinates, while $d_N^{(k)}>0$ avoids any dependence on a chosen
normalization constant.

\subsubsection*{H5 (closed as equivalence): windowed Weil positivity \texorpdfstring{$\Longleftrightarrow$}{<=>} ARP-P}

For $z\in U_-$ the correct unilateral kernel is
\[
   \frac1{z_\rho-z}=-i\int_0^\infty e^{-izt}e^{iz_\rho t}\,dt,
\]
because $\Im z<-\half$ and $|\Im z_\rho|<\half$.  Write formally
\[
   g^F(z)=-i\int_0^\infty e^{-izt}\psi^F(t)\,dt,
   \qquad
   \psi^F(t)_{\alpha\beta}
   =
   \sum_\rho
   m_\rho\,\ell_{z_\rho}(\phi_\beta)
   \overline{\ell_{\bar z_\rho}(\phi_\alpha)}
   e^{iz_\rho t},
\]
and extend $\psi^F(-\tau)=\psi^F(\tau)^*$.  For compactly supported smooth vector tests
$f:[0,\infty)\to\C^m$, define
\[
   Q^F(f)=
   \iint_{[0,\infty)^2}f(t)^*\psi^F(t'-t)f(t')\,dt\,dt'.
\]
Then the following are equivalent:
\[
   \mathsf P[g^F;z_1,\ldots,z_N]\succeq0
   \quad\text{for all finite node sets in }U_-,
\]
and
\[
   Q^F(f)\ge0\quad\text{for every compactly supported smooth window test }f.
\]

\begin{lemma}[transfer kernel identity]\label{lem:transfer-kernel}
For $z,w\in U_-$,
\[
   \iint_{[0,\infty)^2}\overline{e^{-iwt}}\,\psi^F(t'-t)\,e^{-izt'}\,dt\,dt'
   =\frac{g^F(z)-g^F(w)^*}{z-\bar w}.
\]
\end{lemma}
\begin{proof}
Insert $\psi^F(\tau)=\sum_\rho m_\rho W_\rho\,e^{iz_\rho\tau}$ with
$W_\rho=\ell_{z_\rho}(\phi_\beta)\overline{\ell_{\bar z_\rho}(\phi_\alpha)}$:
\[
   \int_0^\infty e^{i(\bar w-z_\rho)t}\,dt=\frac{-i}{z_\rho-\bar w},
   \qquad
   \int_0^\infty e^{i(z_\rho-z)t'}\,dt'=\frac{i}{z_\rho-z},
\]
both convergent since $|\Im z_\rho|<\half<-\Im z,-\Im w$; the product is
$m_\rho W_\rho/((z_\rho-z)(z_\rho-\bar w))$, summable by Lemma~\ref{lem:psi-growth}, and equals the
displayed quotient by Lemma~\ref{lem:cauchy-pick}.
\end{proof}

\begin{remark}[kernel orientation]\label{rmk:kernel-orientation}
The opposite convention $\psi^F(t-t')$ yields the Pick form at the reflected nodes $\{-\bar z_j\}$; the
equivalence above is unaffected (the node class is closed under reflection), but the pointwise identity of
Lemma~\ref{lem:transfer-kernel} requires the orientation $\psi^F(t'-t)$.
\end{remark}

The key representation is Laplace inversion along a horizontal line.  Fix $y_0>\half$ and
$z(x)=x-iy_0$.  If $f\in C_c^\infty(0,\infty)$, set
\[
   c(x)=\frac1{2\pi}\int_0^\infty f(t)e^{iz(x)t}\,dt
        =\frac1{2\pi}\int_0^\infty f(t)e^{y_0t}e^{ixt}\,dt.
\]
Since $f(t)e^{y_0t}$ is compactly supported smooth, $c$ is Schwartz and Fourier inversion gives
\[
   f(t)=\int_\R c(x)e^{-iz(x)t}\,dx.
\]

If window positivity holds, apply it to truncated functions
\[
   f_L(t)=\eta_L(t)\sum_j c_j e^{-iz_jt},
\]
where $\eta_L=1$ on $[0,L-1]$ and vanishes beyond $L$.  Dominated convergence gives
\[
   \lim_{L\to\infty}Q^F(f_L)
   =
   \sum_{i,j}c_i^*K(z_i,z_j)c_j,
\]
so all Pick matrices are positive.  Conversely, if all Pick matrices are positive, use the horizontal-line
representation of $f$ and Fubini:
\[
   Q^F(f)=
   \iint_{\R^2}c(x)^*K(z(x),z(x'))c(x')\,dx\,dx'.
\]
Every Riemann sum of this integral is nonnegative by finite Pick positivity; hence the integral is
nonnegative.  Thus H5 is closed as an equivalence.  Proving positivity for every window remains exactly
the RH-hard input in window coordinates.

\subsubsection*{H6 (closed quantitatively for each single-cluster configuration)}

If RH fails, then ARP-P fails for some finite source plane and some finite node set.  This is the
contrapositive of Steps~8--13:
\[
   \textnormal{ARP-P}\Longrightarrow\textnormal{RH}.
\]
Equivalently, an off-line zero forces a finite Pick matrix with a negative eigenvalue.  The qualitative
statement is closed.  Proposition~\ref{prop:H6-witness}, Lemma~\ref{lem:H6-phase}, and
Theorem~\ref{thm:H6-closed} give the explicit rational-witness bound for each single-cluster
configuration.  Uniform bounds over families of configurations remain a separate quantitative target.  The
local obstruction block has vectors
\[
   v_i=\frac1{z_\rho-z_i},\qquad
   u_i=\frac1{\bar z_\rho-z_i},
\]
and the first quantitative target is a single-cluster estimate.

\subsubsection*{H7 (partially closed): Andreief/Vandermonde level-two SOS}

Under RH, a scalar two-node Pick determinant has the positive Andreief form
\[
   \det P_2
   =
   \frac12
   \iint
   \left|
      \phi_{z_1}(t)\phi_{z_2}(t')
      -
      \phi_{z_1}(t')\phi_{z_2}(t)
   \right|^2
   d\Sigma(t)d\Sigma(t'),
   \qquad
   \phi_z(t)=\frac1{t-z}.
\]
The integrand is an antisymmetrized Cauchy/Vandermonde square and vanishes on the diagonal $t=t'$.  The
open H7 target is to express the corresponding arithmetic two-point form, after pole shorting, as a
manifestly nonnegative sum of squares.  Corollaries~\ref{cor:H7-nondegenerate} and
\ref{cor:below-threshold} close the scalar nondegenerate and fixed coalescent large-height regimes.
The remaining SOS target is the multiscale/arithmetically degenerate certification not covered by
single-cluster positivity.  This is not a restatement of H1: it tests the antisymmetrized two-point
correlation.

\subsubsection*{H8 (quantitatively closed for scalar fixed-$N$ clusters and terminal gauges)}

For scalar node clusters of fixed pattern, H8 is closed in the large-height regime by
Theorem~\ref{thm:cluster}.  If
\[
   z_j=t+a_j-iy_j,\qquad y_j\ge\half+\epsilon_0,
\]
and the pairs $(a_j,y_j)$ are distinct and fixed, then the scalar Pick matrix of $M_\Xi$ is positive
for all sufficiently large $|t|$.  The proof isolates the Szeg\H{o} Gram matrix
\[
   S_{ij}=\frac{-i}{z_i-\bar z_j},
\]
whose smallest eigenvalue supplies the cluster margin, while the archimedean contribution grows like
$\log |t|\,S$ and the prime contribution is bounded in the fixed band $\Re s\ge1+\epsilon_0$.
Lemma~\ref{lem:szego-explicit} evaluates $\det S$ by the Cauchy determinant formula, and
Theorem~\ref{thm:cluster-quantitative} makes the closure uniform over boxes:
\[
   \prod_{i<j}|w_i-w_j|^2\gtrsim\frac1{\log |t|}
   \quad\Longrightarrow\quad
   \mathsf P[M_\Xi;z_1,\ldots,z_N]\succ0.
\]
Thus a cluster with minimum pairwise gap $h$ is covered down to
$h\gtrsim(\log|t|)^{-1/(N(N-1))}$.
Theorem~\ref{thm:one-point-jets} gives the complementary exact serialization: for any fixed
$z_0\in U_-$, RH is equivalent to positivity of all one-point jet matrices for the detection family.
Remark~\ref{rmk:scale-correspondence} identifies the jet order with the prime-window scale.

Theorem~\ref{thm:band-closure} closes the scalar fixed-$N$ band problem for all configurations,
including multiscale clusters.  When $N$ nodes coalesce at $z$,
ARP-P$(N)$ limits to positivity of the jet matrix
\[
   J_N(z)=
   \left[
      \frac1{i!\,j!}\partial_z^i\partial_{\bar w}^jK(z,w)\big|_{w=z}
   \right]_{i,j=0}^{N-1}.
\]
For each fixed derivative order $k$,
\[
   \sum_{n\ge2}\Lambda(n)^2(\log n)^{2k}n^{-2}<\infty,
\]
so the Montgomery--Vaughan mean-value mechanism used in H3 extends entrywise to every fixed derivative.
This is the limit in which the Szeg\H{o} eigenvalue $\lambda_{\mathcal P}$ tends to zero in ordinary
coordinates; the Newton congruence of Theorem~\ref{thm:band-closure} controls it for each fixed $N$ in
a fixed band.  This still does not prove RH, because the minimal scalar bridge requires positivity for
all prefixes of one fixed accumulating sequence, i.e. uniform control as $N\to\infty$ at one interior
height.

The terminal one-point gauge is sharpened by Lemma~\ref{lem:pascal-effective},
Theorem~\ref{thm:whitened-terminal}, Theorem~\ref{thm:prime-shift}, and
Corollary~\ref{cor:loglog-effective}.  The archimedean floor is effective through the
Pascal--Laguerre matrix, the prime side is represented exactly by the von Mangoldt shift operator, and
deterministic contractivity is proved up to order $N\asymp\log\log |t_0|$ with explicit constants.
The correct remaining observable is
\[
   \delta_N(z_0)=1-\lambda_{\max}(T_N(z_0))
\]
in the $J_N^\infty$-whitened coordinate, valid in the range $N\le N_*(t_0)$ of
Lemma~\ref{lem:pole-pair}; the governing all-$N$ form is the reference-whitened defect of
Definition~\ref{def:terminal-defect-fixed}, with which it agrees in sign throughout that range.
Theorem~\ref{thm:saturation} shows that, under RH, this defect has no fixed positive margin: the
terminal sections approach the boundary in the sense
$\limsup_N\lambda_{\max}(T_N)=1$.  Therefore the final terminal input is the exact sign condition
$\delta_N(z_0)\ge0$ for every $N$, not a margin estimate.

In summary: H1, H2, H3 in the stated regions plus its exact decoherence reformulation, H4 as determinant
serialization, H5, H6 quantitative by an explicit rational witness, H7 scalar band positivity, and H8
scalar fixed-$N$ band positivity for all configurations plus the effective terminal Laguerre gauge are
all proved below. The only essential input left is the terminal scalar Pick condition
of Theorem~\ref{thm:minimal-input}, equivalently the no-margin sign condition
$\delta_N(z_0)\ge0$ for every $N$ in one fixed accumulating interior-band gauge.  The numerical evidence
and Conjecture~\ref{conj:defect-law} indicate that this sign is decided at factorial distance
from the boundary, and the top terminal eigenvectors recover the nearby zero ordinates.

\medskip\noindent\textbf{Formal statements and proofs for the Step~5 substructure (H0--H8 and the terminal reduction).}
The following are the formal statements underlying the discussion above.

\begin{lemma}[H0: conservation of difficulty]\label{lem:live-H0}
If $S_1,\ldots,S_k$ are all proved and
$S_1\wedge\cdots\wedge S_k\Rightarrow\mathrm{RH}$, then RH is proved.  Therefore any valid route to RH
which has not proved RH must contain a false step or an RH-hard step.
\end{lemma}

\begin{proof}
This is modus ponens, followed by its contrapositive.
\end{proof}

\begin{theorem}[H1: scalar one-node positivity]\label{thm:live-H1}
For $M_\Xi(z)=i(-\xi'/\xi)(s)$, $s=\half+iz$, the scalar one-node Pick inequality on $U_-$ is equivalent
to
\[
   \Re\frac{\xi'}{\xi}(s)\ge0,\qquad \Re s>1,
\]
and this inequality holds.
\end{theorem}

\begin{proof}
The diagonal Pick quotient equals
\[
   \frac{\Im M_\Xi(z)}{\Im z}
   =
   \frac{-\Re(\xi'/\xi)(s)}{-(\Re s-\half)}
   =
   \frac{\Re(\xi'/\xi)(s)}{\Re s-\half}.
\]
The denominator is positive for $\Re s>1$.  The Hadamard product gives, in the standard symmetric
limiting sense,
\[
   \Re\frac{\xi'}{\xi}(s)=\sum_\rho\frac{\sigma-\beta}{|s-\rho|^2}.
\]
The classical zero-free line and functional equation give $0<\beta<1$ for nontrivial zeros, so every
summand is nonnegative when $\sigma>1$.
\end{proof}

\begin{remark}[why H1 is kept scalar]\label{rmk:H1-matrix}
The scalar one-node inequality is elementary because each zero contributes a nonnegative real number.
For a matrix channel, the contribution of a non-real mirror pair has the form
$qA+\bar q A^*$ with $A$ rank one, and this need not be positive semidefinite term by term.  Thus a
full matrix-channel level-one theorem already contains phase-decoherence content; it is not treated as
a routine extension of Theorem~\ref{thm:live-H1}.
\end{remark}

\begin{theorem}[H2: fixed-band separated-node criterion]\label{thm:live-H2}
Fix $1<\sigma_0<\sigma_1<\infty$.  In the band $\sigma\in[\sigma_0,\sigma_1]$ the scalar Pick matrix is
positive semidefinite whenever the row-sum condition
\[
   \sum_{j\ne i}
   \frac{\log(3+|t_i|)+\log(3+|t_j|)}
        {\max\{|t_i-t_j|,2(\sigma_0-\half)\}}
   \le c_{\sigma_0,\sigma_1}\log(3+|t_i|)
\]
holds with $c_{\sigma_0,\sigma_1}>0$ sufficiently small.  Geometrically separated high nodes satisfy
this condition.  The estimate is stated for $|t_i|\ge T_0(\sigma_0,\sigma_1)$; for bounded heights the
Pick matrix is positive semidefinite by compactness and the strict positivity of $\Re(\xi'/\xi)$ on
$\sigma\ge\sigma_0$.
\end{theorem}

\begin{proof}
Stirling's formula and boundedness of $\zeta'/\zeta$ in $\sigma\ge\sigma_0$ give
$K_{ii}\ge a_{\sigma_0,\sigma_1}\log(3+|t_i|)$ for large heights.  The same estimates give
$|M_\Xi(z_i)|\le B_{\sigma_0,\sigma_1}\log(3+|t_i|)$, hence
\[
   |K_{ij}|\le
   B_{\sigma_0,\sigma_1}
   \frac{\log(3+|t_i|)+\log(3+|t_j|)}
        {\max\{|t_i-t_j|,2(\sigma_0-\half)\}}.
\]
Choosing $c_{\sigma_0,\sigma_1}=a_{\sigma_0,\sigma_1}/(2B_{\sigma_0,\sigma_1})$ gives
$\sum_{j\ne i}|K_{ij}|\le K_{ii}/2$.  Gershgorin's theorem for Hermitian matrices gives positive
semidefiniteness.  If $|t_j|\ge T_*Q^j$ and
$|t_i-t_j|\ge\frac12\max\{|t_i|,|t_j|\}$, then the lower-height and higher-height contributions are
bounded respectively by
\[
   C\frac{(\log(3+|t_i|))^2}{|t_i|}
   \quad\text{and}\quad
   C'\frac{\log(3+|t_i|)}{|t_i|},
\]
which are $o(\log(3+|t_i|))$ as $T_*\to\infty$.
\end{proof}

\begin{theorem}[H3: two-node regions]\label{thm:live-H3}
The scalar infinitesimal two-node inequality
\[
   (\sigma-\half)\left|\left(\frac{\xi'}{\xi}\right)'(s)\right|
   \le
   \Re\frac{\xi'}{\xi}(s)
\]
holds in the logarithmic region
\[
   |t|\ge T_0,\qquad \sigma\ge1+\frac{c}{\sqrt{\log(3+|t|)}},
\]
and also for $s=1+\epsilon+it$, $\epsilon\in[T^{-1/5},1]$, outside a set
of $t\in[T,2T]$ of measure $O(T(\log T)^{-2})$.
\end{theorem}

\begin{proof}
Write $\sigma=1+\epsilon$.  In $0<\epsilon\le1$,
\[
   \left|\left(\frac{\zeta'}{\zeta}\right)'(s)\right|
   \le \sum_{n\ge2}\Lambda(n)(\log n)n^{-1-\epsilon}
   \le C_1\epsilon^{-2}+C_2.
\]
The elementary $\Gamma$ and pole terms contribute $O((1+|t|)^{-1})$ in the relevant band, while
Stirling and the Euler bound give
\[
   \Re\frac{\xi'}{\xi}(s)
   \ge \frac12\log\frac{|t|}{2\pi}-C_3\epsilon^{-1}-C_4.
\]
Taking $\epsilon\ge c/\sqrt{\log(3+|t|)}$ with $c$ large makes the right side dominate the
$O(\epsilon^{-2})$ left side.  For $\sigma\ge2$, the derivative Euler series is
$O(\sigma2^{-\sigma})$ after multiplication by $\sigma-\half$, while the real part grows like
$\frac12\log|s|$, both bounds holding \emph{uniformly in $t$}; so the remaining large-$\sigma$ tail is
closed uniformly in $t$ (the uniformity required by the compact-set argument of
Theorem~\ref{thm:PDH-H3}), and finite intermediate bands are absorbed by increasing $T_0$.

For the almost-everywhere statement, the Montgomery--Vaughan mean-value theorem gives uniformly in
$\epsilon\in[T^{-1/5},1]$
\[
   \int_T^{2T}\left|\frac{\zeta'}{\zeta}(1+\epsilon+it)\right|^2dt\ll T,\qquad
   \int_T^{2T}\left|\left(\frac{\zeta'}{\zeta}\right)'(1+\epsilon+it)\right|^2dt\ll T,
\]
because the square coefficient sums converge at $\epsilon=0$ and the conservative error is absorbed in
the stated range.  Chebyshev's inequality shows that the set where either term exceeds a fixed multiple
of $\log T$ has measure $O(T(\log T)^{-2})$.  Outside this set the archimedean real part dominates.
\end{proof}

\begin{definition}[coherence functional and spectral weight]\label{def:coherence}
For $s=\sigma+it$ with $\sigma>1$, define
\[
   \Theta(s):=
   \frac{\left|\sum_\rho (s-\rho)^{-2}\right|}
        {\sum_\rho|s-\rho|^{-2}}\in[0,1],
   \qquad
   W(s):=
   \frac{\sum_\rho(\sigma-\beta)\,|s-\rho|^{-2}}
        {(\sigma-\half)\sum_\rho|s-\rho|^{-2}}.
\]
The sums are absolutely convergent by the classical zero-counting estimate
$N(T)=O(T\log T)$.
\end{definition}

\begin{theorem}[exact reformulation of H3]\label{thm:H3-theta}
For every $s$ with $\sigma>1$,
\[
   (\sigma-\half)\left|\left(\frac{\xi'}{\xi}\right)'(s)\right|
   \le
   \Re\frac{\xi'}{\xi}(s)
   \qquad\Longleftrightarrow\qquad
   \Theta(s)\le W(s).
\]
Under RH, $\beta=\half$ for every zero, $W\equiv1$, and H3 is the triangle inequality.
\end{theorem}

\begin{proof}
Differentiating the Hadamard logarithmic derivative gives
\[
   \left(\frac{\xi'}{\xi}\right)'(s)
   =
   -\sum_\rho(s-\rho)^{-2},
\]
absolutely convergently for $\sigma>1$.  The same symmetric Hadamard pairing used in
Theorem~\ref{thm:live-H1} gives
\[
   \Re\frac{\xi'}{\xi}(s)
   =
   \sum_\rho(\sigma-\beta)|s-\rho|^{-2}.
\]
Dividing the H3 inequality by the positive quantity
$(\sigma-\half)\sum_\rho|s-\rho|^{-2}$ gives exactly $\Theta(s)\le W(s)$.
\end{proof}

\begin{lemma}[unconditional lower bound for the weight]\label{lem:weight-lower}
There are effective constants $c_2>0$ and $T_0\ge3$ such that for
$1<\sigma\le2$ and $|t|\ge T_0$,
\[
   W(s)\ge
   \frac{\sigma-1}{\sigma-\half}
   +
   \frac{c_2}{\log(3+|t|)}
   \ge
   \frac23(\sigma-1)+\frac{c_2}{\log(3+|t|)}.
\]
\end{lemma}

\begin{proof}
Use the classical zero-free region, in the form
\[
   \beta\le 1-\frac{c_1}{\log(3+|\gamma|)}
\]
after decreasing $c_1$ to absorb the finitely many small ordinates.  Hence each zero contributes the
ratio
\[
   \frac{\sigma-\beta}{\sigma-\half}
   \ge
   \frac{\sigma-1}{\sigma-\half}
   +
   \frac{c_1}{(\sigma-\half)\log(3+|\gamma|)}.
\]
The number $W(s)$ is the weighted average of these ratios with weights proportional to
 $|s-\rho|^{-2}$.  It remains only to show that this weighted average of
$(\log(3+|\gamma|))^{-1}$ is bounded below by a constant multiple of
$(\log(3+|t|))^{-1}$.

The Riemann--von Mangoldt formula gives a constant $C_H>0$ such that, with
$H_t=C_H\log(3+|t|)$, every interval $[t-H_t,t+H_t]$ contains at least one zero ordinate for all
sufficiently large $|t|$.  Therefore
\[
   \sum_\rho|s-\rho|^{-2}\ge (4+H_t^2)^{-1}
\]
for $1<\sigma\le2$.  On the other hand, partial summation and $N(T)=O(T\log T)$ give
\[
   \sum_{|\gamma|>2|t|}|s-\rho|^{-2}
   \le
   4\sum_{|\gamma|>2|t|}|\gamma|^{-2}
   \le
   C\frac{\log(3+|t|)}{|t|}.
\]
Thus the weight carried by zeros with $|\gamma|>2|t|$, relative to the displayed lower bound, is
 $O((\log|t|)^3/|t|)$ and tends to $0$.  On the complementary set,
\[
   \frac1{\log(3+|\gamma|)}
   \ge
   \frac1{\log(3+2|t|)}.
\]
For $|t|\ge T_0$ the stated lower bound follows after adjusting constants.  The final inequality uses
$(\sigma-\half)^{-1}\ge 2/3$ for $1<\sigma\le2$.
\end{proof}

\begin{definition}[Phase Decoherence Hypothesis, PDH]\label{def:PDH}
With the constants of Lemma~\ref{lem:weight-lower}, PDH is the following open hypothesis: for all
$1<\sigma\le2$ and all $|t|\ge T_0$,
\[
   \Theta(s)\le \frac23(\sigma-1)+\frac{c_2}{\log(3+|t|)}.
\]
PDH is not claimed as a theorem.  It is a named sufficient decoherence condition for the global H3
edge.
\end{definition}

\begin{theorem}[PDH closes H3 at large height]\label{thm:PDH-H3}
Assume PDH.  Then the scalar infinitesimal two-node inequality of
Theorem~\ref{thm:live-H3} holds for all $1<\sigma\le2$ and $|t|\ge T_0$.  Together with the
large-$\sigma$ part of Theorem~\ref{thm:live-H3}, the only remaining region is a compact set,
in principle accessible to interval arithmetic.
\end{theorem}

\begin{proof}
For $1<\sigma\le2$ and $|t|\ge T_0$, PDH and Lemma~\ref{lem:weight-lower} give
$\Theta(s)\le W(s)$.  Theorem~\ref{thm:H3-theta} converts this into H3.  For $\sigma\ge2$ and large
$|t|$, H3 is already contained in the large-$\sigma$ tail argument in Theorem~\ref{thm:live-H3}.
\end{proof}

Theorem~\ref{thm:H3-theta} and Lemma~\ref{lem:weight-lower} are proved unconditionally. PDH itself remains
an open hypothesis; Theorem~\ref{thm:PDH-H3} proves only the conditional implication
$\mathrm{PDH}\Rightarrow\mathrm{H3}$ at large height, and no closure of RH is claimed from it.

\begin{lemma}[exact mirror-pair identity]\label{lem:mirror-exact}
Let $\rho=\beta+i\gamma$ and $\rho'=(1-\beta)+i\gamma$ be a mirror pair with
$\delta=\beta-\half>0$.  Write
\[
   D=s-(\half+i\gamma)=(\sigma-\half)+i(t-\gamma),
\]
so that $s-\rho=D-\delta$ and $s-\rho'=D+\delta$.  Then
\[
   \frac{1}{(s-\rho)^2}+\frac{1}{(s-\rho')^2}
   =
   \frac{2(D^2+\delta^2)}{(D^2-\delta^2)^2},
\]
and
\[
   \frac{\sigma-\beta}{|s-\rho|^2}
   +
   \frac{\sigma-\beta'}{|s-\rho'|^2}
   =
   \frac{2(\sigma-\half)(|D|^2-\delta^2)}{|D^2-\delta^2|^2}.
\]
Consequently the pairwise H3 deficit is exactly
\[
   \mathrm{Def}_\rho(s)
   :=
   \frac{
      2(\sigma-\half)\bigl(|D^2+\delta^2|-(|D|^2-\delta^2)\bigr)
   }{|s-\rho|^2|s-\rho'|^2}
   \ge0,
\]
and
\[
   \mathrm{Def}_\rho(s)
   \le
   \frac{4(\sigma-\half)\delta^2}{|s-\rho|^2|s-\rho'|^2}.
\]
\end{lemma}

\begin{proof}
The first identity is direct:
\[
   (D-\delta)^{-2}+(D+\delta)^{-2}
   =
   \frac{(D+\delta)^2+(D-\delta)^2}{(D^2-\delta^2)^2}
   =
   \frac{2(D^2+\delta^2)}{(D^2-\delta^2)^2}.
\]
For the second, use $\Re D=\sigma-\half$ and
\[
   \sigma-\beta=(\sigma-\half)-\delta,\qquad
   \sigma-\beta'=(\sigma-\half)+\delta.
\]
Then
\[
   \frac{(\sigma-\half)-\delta}{|D-\delta|^2}
   +
   \frac{(\sigma-\half)+\delta}{|D+\delta|^2}
   =
   \frac{2(\sigma-\half)(|D|^2-\delta^2)}
        {|D-\delta|^2|D+\delta|^2},
\]
and
\[
   |D-\delta|^2|D+\delta|^2=|D^2-\delta^2|^2.
\]
The deficit formula follows by putting both pair contributions over the common denominator
$|s-\rho|^2|s-\rho'|^2$.  Since
\[
   |D^2+\delta^2|\ge |D|^2-\delta^2
\]
by the reverse triangle inequality, the deficit is nonnegative.  Finally
$|D^2+\delta^2|\le |D|^2+\delta^2$ gives the displayed upper bound.
\end{proof}

\begin{theorem}[deficit-versus-gain criterion]\label{thm:mirror-criterion}
Define the on-line decoherence gain
\[
   G(s):=(\sigma-\half)
   \left[
      \sum_{\beta=\half}|s-\rho|^{-2}
      -
      \left|\sum_{\beta=\half}(s-\rho)^{-2}\right|
   \right]\ge0.
\]
Then H3 holds at $s$ whenever
\[
   4(\sigma-\half)
   \sum_{\delta_\rho>0}
   \frac{\delta_\rho^2}{|s-\rho|^2|s-\rho'|^2}
   \le
   G(s).
\]
\end{theorem}

\begin{proof}
Split the zero sum into on-line zeros and mirror pairs.  By the triangle inequality,
\[
   (\sigma-\half)\left|\sum_\rho(s-\rho)^{-2}\right|
   \le
   (\sigma-\half)\left|\sum_{\beta=\half}(s-\rho)^{-2}\right|
   +
   (\sigma-\half)\sum_{\mathrm{pairs}}
   \left|
      \frac1{(s-\rho)^2}+\frac1{(s-\rho')^2}
   \right|.
\]
Insert Lemma~\ref{lem:mirror-exact} for each pair and the definition of $G(s)$ for the on-line part.
The right side is at most
\[
   \sum_\rho(\sigma-\beta)|s-\rho|^{-2}
   +
   \sum_{\mathrm{pairs}}\mathrm{Def}_\rho(s)
   -
   G(s).
\]
The hypothesis and the upper bound on $\mathrm{Def}_\rho(s)$ make the last two terms nonpositive.
This gives H3, equivalently $\Theta(s)\le W(s)$ by Theorem~\ref{thm:H3-theta}.
\end{proof}

\begin{remark}[mirror-pair honesty]\label{rmk:mirror-honest}
Lemma~\ref{lem:mirror-exact} is exact and replaces the earlier Taylor heuristic.  Every off-line
mirror pair produces a nonnegative second-moment deficit, bounded by
\[
   4(\sigma-\half)\delta^2/(|s-\rho|^2|s-\rho'|^2).
\]
H3 can absorb such pairs only through the on-line phase-decoherence gain $G(s)$ or an equivalent
global cancellation mechanism.  This is the same residual wall named by PDH, not a closure of RH.
\end{remark}

\begin{lemma}[continuity and growth of the correlation kernel]\label{lem:psi-growth}
For each finite channel $F$, the zero-side series defining $\psi^F$ converges absolutely and locally
uniformly in $t$, defines a continuous matrix function, and satisfies
\[
   \|\psi^F(t)\|\le C_F e^{|t|/2},\qquad t\in\R .
\]
\end{lemma}

\begin{proof}
Each term is bounded by
\[
   m_\rho|\ell_{z_\rho}(\phi_\beta)\ell_{\bar z_\rho}(\phi_\alpha)|e^{|\Im z_\rho||t|}
   \le
   m_\rho C_{A,F}(1+|\Re z_\rho|)^{-A}e^{|t|/2},
\]
because $|\Im z_\rho|<\half$ and the source-core transforms have rapid horizontal decay in the zero
strip.  Choosing $A$ large and using $N(T)=O(T\log T)$ gives a summable majorant over the zero set,
locally uniformly in $t$.
\end{proof}

\begin{theorem}[H5: window positivity is equivalent to ARP-P]\label{thm:live-H5}
Let
\[
   \frac1{z_\rho-z}=-i\int_0^\infty e^{-izt}e^{iz_\rho t}\,dt,\qquad z\in U_-,
\]
and define the associated correlation kernel $\psi^F$ by
$g^F(z)=-i\int_0^\infty e^{-izt}\psi^F(t)\,dt$.  Then ARP-P for $F$ is equivalent to
$Q^F(f)\ge0$ for every compactly supported smooth window test $f$.
\end{theorem}

\begin{proof}
Fix $y_0>\half$ and write $z(x)=x-iy_0$.  For $f\in C_c^\infty(0,\infty)$ define
\[
   c(x)=\frac1{2\pi}\int_0^\infty f(t)e^{iz(x)t}\,dt.
\]
Since $f(t)e^{y_0t}$ is compactly supported smooth, Fourier inversion gives
$f(t)=\int_\R c(x)e^{-iz(x)t}\,dx$.

If window positivity holds, apply it to
$f_L(t)=\eta_L(t)\sum_j c_j e^{-iz_jt}$ and pass to the limit $L\to\infty$ by dominated convergence;
the growth bound in Lemma~\ref{lem:psi-growth} supplies the integrable majorant because
$\Im z_j<-\half$.  The limit is $\sum_{i,j}c_i^*K(z_i,z_j)c_j$.  Hence every finite Pick matrix is
positive.  Conversely, if every Pick matrix is positive, substitute the horizontal-line representation of
$f$ into $Q^F(f)$.  Lemma~\ref{lem:psi-growth} and the Schwartz decay of $c$ justify Fubini.  The result
is an integral of the positive kernel $K(z(x),z(x'))$ against $c(x)$ and $c(x')$, obtained as a limit of
nonnegative Riemann sums.  Hence $Q^F(f)\ge0$.
\end{proof}

\begin{corollary}[H6: qualitative finite negative witness]\label{cor:live-H6}
If RH fails, then there exist a finite channel $F$ and finite nodes in $U_-$ for which the ARP-P Pick
matrix has a negative eigenvalue.
\end{corollary}

\begin{proof}
This is the contrapositive of Theorem~\ref{thm:pick-rh}.
\end{proof}

\begin{definition}[single-cluster rational witness]\label{def:H6-rational-witness}
Let $w_0=\gamma-i\delta$, $0<\delta<\half$, be a non-real zero coordinate of $\Xi$ in
$\Omega_-$.  Let $t_1,\ldots,t_M$ be the real zero coordinates of $\Xi$ in
$[\gamma-R,\gamma+R]$, counted with multiplicity.  Choose $N_p\ge M+2$ pairwise distinct poles
$\pi_1,\ldots,\pi_{N_p}\in\C^+$ with
\[
   \Im\pi_j\in[1,2],\qquad |\Re\pi_j-\gamma|\le R,
\]
and define
\[
   q(t)=\prod_{m=1}^M(t-t_m),\qquad
   r(t)=\frac{q(t)}{\prod_{j=1}^{N_p}(t-\pi_j)}
   =
   \sum_{j=1}^{N_p}\frac{c_j}{t-\pi_j}.
\]
Equivalently,
\[
   c_j=\frac{q(\pi_j)}{\prod_{k\ne j}(\pi_j-\pi_k)}.
\]
The associated Pick nodes are $z_j=\bar\pi_j\in U_-$.
\end{definition}

\begin{proposition}[H6 quantitative reduction: explicit single-cluster witness]\label{prop:H6-witness}
In the setting of Definition~\ref{def:H6-rational-witness}, choose a source $\phi$ detecting $w_0$, and
let $\kappa:=\ell_{w_0}(\phi)\overline{\ell_{\bar w_0}(\phi)}\ne0$ be the (complex) pair weight; \emph{no}
normalization of the channel is performed, so that the Pick form stays Hermitian.  Suppose the chosen
poles satisfy the finite phase condition
\[
   2\Re\bigl(\kappa\,\overline{r(\bar w_0)}\,r(w_0)\bigr)
   \le
   -\eta\,|\kappa|\,|r(w_0)|\,|r(\bar w_0)|
\]
for some $\eta>0$.  If the window contains further non-real pairs, augment $q$ with one complex root at
one member of each such pair; those pairs then contribute exactly $0$ to the form (as in
Theorem~\ref{thm:alpha4}), and the resulting fixed phase offset of $q(w_0)$ is absorbed into the target
phase of Lemma~\ref{lem:H6-phase}.  Then the Pick matrix on the nodes $z_j=\bar\pi_j$ has a test vector
$c=(c_j)$ with Rayleigh quotient bounded above by
\[
   \frac{
      -\eta\,|\kappa|\,|r(w_0)|\,|r(\bar w_0)|
      + B_{\mathrm{far}}(R,\gamma)\,\|r\|_{\mathrm{far}}^2
   }{\|c\|^2},
\]
where
\[
   \|r\|_{\mathrm{far}}
   :=
   \sup_{|t-\gamma|\ge R}(1+|t-\gamma|)^2|r(t)|
\]
and $B_{\mathrm{far}}(R,\gamma)$ is the explicit zero-counting tail bound obtained from
$N_\Xi(T+1)-N_\Xi(T)=O(\log(2+T))$.  In particular, whenever the numerator is negative, this gives a
finite explicit negative Pick witness.
\end{proposition}

\begin{proof}
The rational function $r$ has poles at $\pi_j=\bar z_j$ and residues $c_j$, so the scalar Pick
quadratic form evaluated on $c$ is the zero-side sum obtained by testing the meromorphic target against
$r$.  The local real zeros $t_m$ contribute
\[
   |r(t_m)|^2=0
\]
because $q(t_m)=0$.  The conjugate off-line pair, carrying the complex residue weight $\kappa$, contributes exactly
\[
   2\Re\bigl(\kappa\,\overline{r(\bar w_0)}\,r(w_0)\bigr),
\]
with the sign convention of the scalar kernel.  By the assumed phase condition this is at most
$-\eta\,|\kappa|\,|r(w_0)|\,|r(\bar w_0)|$.

For every remaining zero coordinate $u$, the contribution is bounded in absolute value by a constant
multiple of $|r(u)|^2$ after collecting the finitely many channel weights into the normalization.  On
the tail $|u-\gamma|\ge R$,
\[
   |r(u)|^2\le
   \frac{\|r\|_{\mathrm{far}}^2}{(1+|u-\gamma|)^4}.
\]
Summing over unit intervals and using the classical local zero-counting bound
$N_\Xi(T+1)-N_\Xi(T)=O(\log(2+T))$ gives the displayed
$B_{\mathrm{far}}(R,\gamma)\|r\|_{\mathrm{far}}^2$ bound.  Dividing by $\|c\|^2$ gives the Rayleigh
quotient estimate.  A negative Rayleigh quotient implies a negative eigenvalue.
\end{proof}

\begin{lemma}[H6 phase adjustment]\label{lem:H6-phase}
Assume $R\ge2$ in Definition~\ref{def:H6-rational-witness}.  Enlarge the pole family to
\[
   N_p=M+2+\lceil2\pi/\delta\rceil
\]
poles $\pi_j=x_j+iy_j$ with pairwise distinct heights $y_j\in[1,2]$ and abscissas
$x_j\in[\gamma-R,\gamma+R]$.  The abscissas can be chosen so that, with the complex pair weight $\kappa$ of
Proposition~\ref{prop:H6-witness},
\[
   \arg\bigl(r(w_0)\overline{r(\bar w_0)}\bigr)\equiv\pi-\arg\kappa\pmod{2\pi},
   \qquad
   2\Re\bigl(\kappa\,\overline{r(\bar w_0)}\,r(w_0)\bigr)
   =
   -2\,|\kappa|\,|r(w_0)|\,|r(\bar w_0)|.
\]
This is achievable by the same intermediate-value argument, since the attainable phase interval has
length $\ge2\pi$ and $\arg\kappa$ is a fixed offset.  Thus the phase condition in
Proposition~\ref{prop:H6-witness} holds with $\eta=2$.
\end{lemma}

\begin{proof}
The polynomial $q$ has real coefficients.  Hence
\[
   \overline{r(\bar w_0)}
   =
   \frac{q(w_0)}{\prod_j(w_0-\bar\pi_j)}
\]
and
\[
   r(w_0)\overline{r(\bar w_0)}
   =
   \frac{q(w_0)^2}{\prod_jF_j},
   \qquad
   F_j=(w_0-\pi_j)(w_0-\bar\pi_j).
\]
Write $u_j=\gamma-x_j$.  Since $w_0=\gamma-i\delta$,
\[
   F_j=u_j^2+y_j^2-\delta^2-2i\delta u_j.
\]
Its real part is at least $y_j^2-\delta^2>0$, and
\[
   \arg F_j
   =
   -\arctan\frac{2\delta u_j}{u_j^2+y_j^2-\delta^2}.
\]
As $u_j$ ranges over $\R$, this continuous phase has range
\[
   [-\arcsin(\delta/y_j),\,\arcsin(\delta/y_j)],
\]
with the endpoints attained at $u_j=\pm\sqrt{y_j^2-\delta^2}$.  Since $y_j\in[1,2]$ and $R\ge2$,
these endpoints are available inside $[\gamma-R,\gamma+R]$.  Each pole therefore contributes phase
width at least $2\arcsin(\delta/2)\ge\delta$.

The total phase
\[
   \Phi(x_1,\ldots,x_{N_p})
   =
   2\arg q(w_0)-\sum_j\arg F_j
\]
is continuous on the connected parameter box, and the independent phase intervals add to an interval
of length at least $N_p\delta\ge2\pi$.  By the intermediate value theorem, some choice of abscissas
gives $\Phi\equiv\pi\pmod{2\pi}$.  No denominator vanishes during the deformation, because
$|F_j|\ge y_j^2-\delta^2>0$, and the heights keep the poles pairwise distinct.
\end{proof}

\begin{theorem}[H6 quantitative single-cluster witness]\label{thm:H6-closed}
In the setting of Definition~\ref{def:H6-rational-witness}, with $R\ge2$ and the augmented pole family
of Lemma~\ref{lem:H6-phase}, the explicit witness of Proposition~\ref{prop:H6-witness} satisfies
\[
   \lambda_{\min}\mathsf P
   \le
   \frac{
      -2|r(w_0)|\,|r(\bar w_0)|
      +B_{\mathrm{far}}(R,\gamma)\|r\|_{\mathrm{far}}^2
   }{\|c\|^2}.
\]
All quantities are explicit in the local zero configuration and in the chosen pole family.  The witness
uses
\[
   N_p=M+2+\lceil2\pi/\delta\rceil
\]
nodes.
\end{theorem}

\begin{proof}
Apply Lemma~\ref{lem:H6-phase} to obtain the phase condition in Proposition~\ref{prop:H6-witness} with
$\eta=2$.  The displayed Rayleigh quotient bound is then exactly the estimate of
Proposition~\ref{prop:H6-witness}.  Since $\lambda_{\min}$ is bounded above by every Rayleigh
quotient, the result follows.
\end{proof}

\begin{remark}[H6 lives in the Szeg\H{o}-degenerate regime]\label{rmk:H6-szego-consistency}
The witness uses $N_p\asymp1/\delta$ nodes.  With those many poles in a bounded box, some gaps are
$O(\delta)$, and the corresponding Cauchy determinant in Lemma~\ref{lem:szego-explicit} is far below
the threshold of Theorem~\ref{thm:cluster-quantitative}.  Thus the negative witness produced by an
off-line zero lives exactly in the degenerating Szeg\H{o} regime; it does not contradict the positive
fixed-pattern and nondegenerate-cluster theorems.
\end{remark}

\begin{definition}[detection family and determinant data]\label{def:detection-family}
For $m\ge2$, let $I_m=(1/m,m)$ and let
$D_m\subset C_c^\infty(I_m)\setminus\{0\}$ be a countable set dense in the sup norm in
\[
   X_m=\overline{C_c^\infty(I_m)}^{\,\|\cdot\|_\infty}.
\]
Enumerate
\[
   \{\phi_k\}_{k\ge1}:=\bigcup_{m\ge2}D_m.
\]
Fix a pairwise distinct canonical node sequence $\{z_j\}_{j\ge1}\subset U_-$ with an accumulation point
in $U_-$.  For $F_k=\operatorname{span}\{\phi_k\}$ define
\[
   d_N^{(k)}
   :=
   \det\mathsf P[g^{F_k};z_1,\ldots,z_N].
\]
By Theorem~\ref{thm:explicit-channel}, each $d_N^{(k)}$ is computable from the archimedean, pole, and prime
data alone.
\end{definition}

\begin{lemma}[Jensen counting: infinitely many charged atoms]\label{lem:jensen-atoms}
Assume RH and let $\phi\in C_c^\infty(0,\infty)$, $\phi\ne0$.  Then
$\ell_{z_\rho}(\phi)\ne0$ for infinitely many zeros $z_\rho$ of $\Xi$.  Consequently
\[
   \mu_\phi=\sum_\rho m_\rho|\ell_{z_\rho}(\phi)|^2\delta_{z_\rho}
\]
has infinite support.
\end{lemma}

\begin{proof}
Let $L(a)=\ell_a(\phi)$.  If $\operatorname{supp}\phi\subset[a_0,b]$, then
\[
   |L(a)|\le \|\phi\|_1e^{b|\Im a|},
\]
so $L$ is entire of exponential type.  It is not identically zero, since otherwise the Fourier--Laplace
transform of the compactly supported distribution $\phi(t)\,dt$ would vanish identically.  Choose
$a_*$ with $L(a_*)\ne0$.  Jensen's formula centered at $a_*$ gives, for the number $n(r)$ of zeros of
$L$ in $|a-a_*|\le r$,
\[
   n(r)\log2
   \le
   \log\max_{|a-a_*|=2r}|L(a)|-\log|L(a_*)|
   \le
   2br+C_\phi.
\]
Thus $n(r)=O(r)$.  The zeros of $\Xi$ with $|z_\rho|\le r$ satisfy the Riemann--von Mangoldt count
\[
   N_\Xi(r)\sim\frac{r}{2\pi}\log\frac{r}{2\pi},
\]
which dominates every linear bound.  Hence $L$ cannot vanish at all but finitely many zeros of $\Xi$.
Under RH these zeros lie on $\R$, so every nonzero value contributes a positive atom to $\mu_\phi$.
\end{proof}

\begin{lemma}[strict positivity under RH]\label{lem:strict-pd}
Assume RH.  For every $\phi\in C_c^\infty(0,\infty)$, $\phi\ne0$, and every finite set of pairwise
distinct nodes $z_1,\ldots,z_N\in U_-$, the scalar Pick matrix
$\mathsf P[g^{F_\phi};z_1,\ldots,z_N]$ is positive definite.
\end{lemma}

\begin{proof}
By Theorem~\ref{thm:rh-pick}, $g^{F_\phi}$ is the Cauchy transform of the positive measure
$\mu_\phi$.  Lemma~\ref{lem:cauchy-pick} gives
\[
   c^*\mathsf P c
   =
   \int_\R\left|\sum_j\frac{c_j}{t-\bar z_j}\right|^2d\mu_\phi(t).
\]
If this quadratic form vanishes, the rational function
\[
   r(t)=\sum_j\frac{c_j}{t-\bar z_j}
\]
vanishes at every charged atom of $\mu_\phi$.  By Lemma~\ref{lem:jensen-atoms}, this is an infinite
set.  A nonzero rational function has only finitely many zeros, so $r\equiv0$.  Since the poles
$\bar z_j$ are distinct, all residues $c_j$ vanish.  Hence the quadratic form is strictly positive for
$c\ne0$.
\end{proof}

\begin{lemma}[dense families detect every off-line zero]\label{lem:dense-detection}
Let $z_\rho$ be a non-real zero of $\Xi$.  Then for some $k$,
\[
   \ell_{z_\rho}(\phi_k)\,\overline{\ell_{\bar z_\rho}(\phi_k)}\ne0.
\]
\end{lemma}

\begin{proof}
Fix any $m\ge2$.  On $X_m$, the functionals $\ell_{z_\rho}$ and $\ell_{\bar z_\rho}$ are continuous.
They are nonzero: if $\ell_a$ vanished on all of $X_m$, then the continuous function $e^{-iat}$ would
define the zero functional on $I_m$, impossible.  Therefore their kernels are proper closed subspaces,
hence closed nowhere dense sets.  The complement of their union is open and dense in $X_m$.  Since
$D_m$ is dense, some $\phi_k\in D_m$ lies in that complement.
\end{proof}

\begin{theorem}[H4 closed: determinant serialization of RH]\label{thm:det-serialization}
With the data of Definition~\ref{def:detection-family},
\[
   \mathrm{RH}
   \quad\Longleftrightarrow\quad
   d_N^{(k)}>0\quad\text{for every }k\ge1\text{ and every }N\ge1.
\]
\end{theorem}

\begin{proof}
Assume RH.  Lemma~\ref{lem:strict-pd} applies to each $\phi_k$ and each prefix
$z_1,\ldots,z_N$, so every determinant $d_N^{(k)}$ is strictly positive.

Conversely, assume $d_N^{(k)}>0$ for every $k,N$.  For fixed $k$, the node sets are nested, so
$\mathsf P[g^{F_k};z_1,\ldots,z_N]$ is the leading principal block of the next Pick matrix.  The Pick
matrix is Hermitian.  By Sylvester's criterion, positivity of all leading principal minors gives
\[
   \mathsf P[g^{F_k};z_1,\ldots,z_N]\succ0
\]
for every $N$.  The nodes are pairwise distinct, accumulate in $U_-$, and $\GXi^{F_k}$ is meromorphic on
$\C$ (Lemma~\ref{lem:GXi-meromorphic}) and holomorphic at every node and near the accumulation point,
since its poles lie in the zero strip $|\Im z|<\half$.  Lemma~\ref{lem:prefix-extension}, applied to this
one channel and the fixed node sequence, shows that every
pole of $\GXi^{F_k}$ in $\Omega_-$ is removable.

If a non-real zero of $\Xi$ lay in $\Omega_-$, Lemma~\ref{lem:dense-detection} would find $k$ for which
that zero is detected by $\GXi^{F_k}$, giving a non-removable pole, contradiction.  The symmetries
$\Xi(-z)=\Xi(z)$ and $\Xi(\bar z)=\overline{\Xi(z)}$ rule out the upper half-plane as in
Theorem~\ref{thm:pick-rh}.  Lemma~\ref{lem:rh-translation} gives RH.
\end{proof}

\begin{corollary}[Schur coordinates]\label{cor:schur-coords}
Under the positivity in Theorem~\ref{thm:det-serialization}, the scalar Schur--Nevanlinna algorithm
along $\{z_j\}$ is nondegenerate at every step and all scalar parameters satisfy
$|\gamma_N^{(k)}|<1$.  Conversely, nondegenerate scalar Schur contractivity at every step implies
$d_N^{(k)}>0$ for all $N$.
\end{corollary}

\begin{proof}
For nested scalar Nevanlinna--Pick data, each Schur step is the Schur complement of the previous leading
principal Pick block; this is the classical Schur--Nevanlinna recursion in Schur-complement form
(see, for instance, \cite{FoiasFrazho}).  Positive definiteness of all leading principal blocks is equivalent to strict
positivity of all these Schur complements, hence to nondegenerate Schur parameters in the unit disk.  By
Sylvester's criterion this is equivalent to $d_N^{(k)}>0$ for all $N$.
\end{proof}

\begin{lemma}[Szeg\H{o} pattern matrix]\label{lem:szego-pattern}
Let $\mathcal P=\{(a_j,y_j)\}_{j=1}^N$ be a finite pattern of distinct pairs with
$y_j>\half$.  For $t\in\R$ set
\[
   z_j=t+a_j-iy_j\in U_-,
   \qquad
   S_{ij}:=\frac{-i}{z_i-\bar z_j}.
\]
Then $S$ is independent of $t$ and is positive definite.  More precisely, $S$ is the Gram matrix of
the functions
\[
   e_j(\tau)=e^{-(y_j+ia_j)\tau},\qquad \tau\ge0,
\]
in $L^2(0,\infty)$.  We write
\[
   \lambda_{\mathcal P}:=\lambda_{\min}(S)>0.
\]
\end{lemma}

\begin{proof}
Since
\[
   z_i-\bar z_j=(a_i-a_j)-i(y_i+y_j),
\]
the matrix is independent of $t$.  Also,
\[
   \int_0^\infty e_i(\tau)\overline{e_j(\tau)}\,d\tau
   =
   \int_0^\infty e^{-[(y_i+y_j)+i(a_i-a_j)]\tau}\,d\tau
   =
   \frac{1}{(y_i+y_j)+i(a_i-a_j)}
   =
   \frac{1}{i(z_i-\bar z_j)}
   =
   S_{ij}.
\]
The exponents $y_j+ia_j$ are distinct because the pairs $(a_j,y_j)$ are distinct.  Distinct
exponentials are linearly independent on every interval of positive length, hence in
$L^2(0,\infty)$.  Therefore their Gram matrix is strictly positive definite.
\end{proof}

\begin{lemma}[archimedean cluster asymptotics]\label{lem:arch-cluster}
On $\Re s>1$ write
\[
   M_\Xi(z)=i\left(G_\infty(s)-\frac{\zeta'}{\zeta}(s)\right),
   \qquad s=\half+iz,
\]
where
\[
   G_\infty(s)
   =
   -\frac1s-\frac1{s-1}+\frac12\log\pi-\frac12\psi(s/2).
\]
For the pattern nodes $z_j=t+a_j-iy_j$ of Lemma~\ref{lem:szego-pattern}, let
$s_j=\half+iz_j$.  The archimedean part $K_\infty$ of the scalar Pick matrix satisfies
\[
   \left\|
      K_\infty-\log\frac{|t|}{2\pi}\,S
   \right\|
   \le
   \frac{C_{\mathcal P}}{|t|}
\]
for $|t|$ large enough, with $C_{\mathcal P}$ depending only on the pattern.
\end{lemma}

\begin{proof}
Stirling's formula for the logarithmic derivative of $\Gamma$ gives, uniformly for the fixed pattern
as $|t|\to\infty$,
\[
   G_\infty(s)
   =
   -\frac12\log\frac{s}{2\pi}+O_{\mathcal P}(|t|^{-1}),
   \qquad \Re s>1.
\]
Thus
\[
   iG_\infty(s_i)-\overline{iG_\infty(s_j)}
   =
   -\frac{i}{2}\left(\log\frac{s_i}{2\pi}+
      \overline{\log\frac{s_j}{2\pi}}\right)
   +O_{\mathcal P}(|t|^{-1}).
\]
For $t>0$, $\arg s_i=\pi/2+O_{\mathcal P}(|t|^{-1})$; for $t<0$,
$\arg s_i=-\pi/2+O_{\mathcal P}(|t|^{-1})$.  In both cases the arguments cancel in
$\log s_i+\overline{\log s_j}$, and
\[
   \log s_i+\overline{\log s_j}=2\log|t|+O_{\mathcal P}(|t|^{-1}).
\]
Hence
\[
   iG_\infty(s_i)-\overline{iG_\infty(s_j)}
   =
   -i\log\frac{|t|}{2\pi}+O_{\mathcal P}(|t|^{-1}).
\]
Dividing by $z_i-\bar z_j$ gives
\[
   (K_\infty)_{ij}
   =
   \log\frac{|t|}{2\pi}\frac{-i}{z_i-\bar z_j}
   +O_{\mathcal P}(|t|^{-1})
   =
   \log\frac{|t|}{2\pi}S_{ij}+O_{\mathcal P}(|t|^{-1}).
\]
Since the matrix size is fixed, the entrywise estimate implies the displayed operator-norm bound.
\end{proof}

\begin{theorem}[H8 closed for scalar fixed cluster patterns]\label{thm:cluster}
Fix $\epsilon_0>0$ and a pattern
$\mathcal P=\{(a_j,y_j)\}_{j=1}^N$ of distinct pairs with
$y_j\ge\half+\epsilon_0$.  There is an effective
$T_0=T_0(\mathcal P,\epsilon_0)$ such that, for every real $t$ with $|t|\ge T_0$, the scalar Pick
matrix of $M_\Xi$ at the nodes
\[
   z_j=t+a_j-iy_j
\]
satisfies
\[
   \mathsf P[M_\Xi;z_1,\ldots,z_N]
   \succeq
   \frac12\,\lambda_{\mathcal P}\log\frac{|t|}{2\pi}\,I
   \succ0.
\]
Thus scalar ARP-P is unconditionally true for every fixed finite cluster pattern translated to
sufficiently large height, away from the boundary $\Re s=1$ by $\epsilon_0$.
\end{theorem}

\begin{proof}
By linearity of the Pick matrix in the function, decompose
\[
   \mathsf P[M_\Xi]=K_\infty+K_{\mathrm{prime}},
   \qquad
   M_\Xi=i\left(G_\infty-\frac{\zeta'}{\zeta}\right).
\]
For each node, $\Re s_j=\half+y_j\ge1+\epsilon_0$.  The absolutely convergent Euler product gives
\[
   \left|\frac{\zeta'}{\zeta}(s_j)\right|
   \le
   -\frac{\zeta'}{\zeta}(1+\epsilon_0)
   =:C_0.
\]
Also $|z_i-\bar z_j|\ge y_i+y_j\ge1+2\epsilon_0$.  Therefore every prime entry is bounded in modulus
by $2C_0/(1+2\epsilon_0)$, and hence
\[
   \|K_{\mathrm{prime}}\|
   \le
   \frac{2NC_0}{1+2\epsilon_0}
   =:C_1.
\]
By Lemmas~\ref{lem:szego-pattern} and~\ref{lem:arch-cluster},
\[
   K_\infty
   \succeq
   \lambda_{\mathcal P}\log\frac{|t|}{2\pi}\,I
   -
   \frac{C_{\mathcal P}}{|t|}\,I.
\]
Thus
\[
   \lambda_{\min}\mathsf P[M_\Xi]
   \ge
   \lambda_{\mathcal P}\log\frac{|t|}{2\pi}
   -
   \frac{C_{\mathcal P}}{|t|}
   -
   C_1.
\]
Choosing $T_0$ so the last two terms are at most
$\frac12\lambda_{\mathcal P}\log(|t|/(2\pi))$ proves the claim.
\end{proof}

\begin{corollary}[fixed patterns near the edge, almost everywhere]\label{cor:cluster-ae}
Fix a scalar pattern $\mathcal P=\{(a_j,y_j)\}_{j=1}^N$ with $y_j>\half$.  For every dyadic interval
$[T,2T]$ and all sufficiently large $T$, the conclusion of Theorem~\ref{thm:cluster} holds for
$t\in[T,2T]$ outside a set of measure $O_{\mathcal P}(T(\log T)^{-2})$, with the deterministic
constant replaced by a pattern-dependent one.
\end{corollary}

\begin{proof}
The archimedean estimate of Lemma~\ref{lem:arch-cluster} is unchanged.  For the prime part, apply the
Montgomery--Vaughan mean-value theorem to the absolutely convergent Dirichlet series on each fixed
shifted line $\Re s_j=\half+y_j>1$:
\[
   \int_T^{2T}
   \left|
      \frac{\zeta'}{\zeta}\bigl(\half+i(t+a_j)+y_j\bigr)
   \right|^2dt
   \ll_{\mathcal P} T.
\]
Chebyshev's inequality and a union bound over the finitely many nodes show that, outside a set of
measure $O_{\mathcal P}(T(\log T)^{-2})$, all values
$|\zeta'/\zeta(s_j)|$ are bounded by a sufficiently small fixed multiple of $\log T$.  On this
complement, the operator norm of $K_{\mathrm{prime}}$ is less than one half of the archimedean margin
$\lambda_{\mathcal P}\log(T/(2\pi))$ once $T$ is large.  The same assembly as in
Theorem~\ref{thm:cluster} gives positivity.
\end{proof}

\begin{lemma}[explicit Szeg\H{o} determinant and eigenvalue bound]\label{lem:szego-explicit}
Let $w_j=y_j+ia_j$.  The pattern matrix of Lemma~\ref{lem:szego-pattern} is the Cauchy matrix
\[
   S_{ij}=\frac{1}{w_i+\bar w_j}.
\]
It satisfies
\[
   \det S
   =
   \frac{\prod_{i<j}|w_i-w_j|^2}
        {\prod_i 2y_i\;\prod_{i<j}|w_i+\bar w_j|^2}
   >0.
\]
Moreover, for patterns in the box
\[
   y_j\in[\half+\epsilon_0,Y],\qquad |a_j|\le A,
\]
one has the explicit lower bound
\[
   \lambda_{\mathcal P}
   \ge
   c(N,\epsilon_0,Y,A)\prod_{i<j}|w_i-w_j|^2,
\]
where
\[
   c(N,\epsilon_0,Y,A)
   =
   \frac{(1+2\epsilon_0)^{N-1}}
        {N^{N-1}(2Y)^N(2Y+2A)^{N(N-1)}}.
\]
\end{lemma}

\begin{proof}
Since
\[
   w_i+\bar w_j=(y_i+y_j)+i(a_i-a_j)=i(z_i-\bar z_j),
\]
Lemma~\ref{lem:szego-pattern} gives $S_{ij}=(w_i+\bar w_j)^{-1}$.  The Cauchy determinant identity
gives
\[
   \det S
   =
   \frac{\prod_{i<j}(w_i-w_j)(\bar w_i-\bar w_j)}
        {\prod_{i,j}(w_i+\bar w_j)}.
\]
The numerator is $\prod_{i<j}|w_i-w_j|^2$.  The denominator splits into the diagonal factors
$2y_i$ and conjugate pairs
\[
   (w_i+\bar w_j)(w_j+\bar w_i)=|w_i+\bar w_j|^2,
\]
which gives the displayed positive formula.

If $\lambda_1\le\cdots\le\lambda_N$ are the eigenvalues of $S$, then
\[
   \det S=\prod_j\lambda_j
   \le
   \lambda_{\min}(S)\,\lambda_{\max}(S)^{N-1}
   \le
   \lambda_{\min}(S)\,(\operatorname{tr}S)^{N-1}.
\]
Thus $\lambda_{\min}(S)\ge\det S/(\operatorname{tr}S)^{N-1}$.  In the box,
\[
   \operatorname{tr}S=\sum_i\frac1{2y_i}\le\frac{N}{1+2\epsilon_0},\qquad
   2y_i\le2Y,\qquad
   |w_i+\bar w_j|\le2Y+2A.
\]
Substituting these bounds into the determinant formula gives the stated constant.
\end{proof}

\begin{theorem}[H8 degenerating regime: quantitative scalar closure]\label{thm:cluster-quantitative}
Fix $N$, $\epsilon_0>0$, $Y$, and $A$.  There are effective constants
$C_{\mathrm{box}}$ and $T_{\mathrm{box}}$, depending only on this box, such that for every scalar
pattern with
\[
   y_j\in[\half+\epsilon_0,Y],\qquad |a_j|\le A,
\]
and every $|t|\ge T_{\mathrm{box}}$,
\[
   \prod_{i<j}|w_i-w_j|^2
   \ge
   \frac{C_{\mathrm{box}}}{\log(|t|/2\pi)}
   \quad\Longrightarrow\quad
   \mathsf P[M_\Xi;z_1,\ldots,z_N]\succ0,
\]
where $z_j=t+a_j-iy_j$ and $w_j=y_j+ia_j$.  More precisely,
\[
   \mathsf P[M_\Xi;z_1,\ldots,z_N]
   \succeq
   \frac12\,\lambda_{\mathcal P}\log\frac{|t|}{2\pi}\,I.
\]
Consequently, if $h=\min_{i<j}|w_i-w_j|$, positivity holds once
\[
   h\ge C_{N,\mathrm{box}}(\log|t|)^{-1/(N(N-1))}.
\]
\end{theorem}

\begin{proof}
The prime bound in Theorem~\ref{thm:cluster} depends only on $N$ and $\epsilon_0$:
\[
   \|K_{\mathrm{prime}}\|\le C_1(N,\epsilon_0).
\]
The error in Lemma~\ref{lem:arch-cluster} is uniform over the box.  Indeed, Stirling's remainder and
the expansion
\[
   \log s_j=\log|t|\pm i\frac\pi2+O((1+A+Y)/|t|)
\]
depend only on $A,Y$ once $|t|$ is large.  Hence
\[
   K_\infty
   \succeq
   \lambda_{\mathcal P}\log\frac{|t|}{2\pi}\,I
   -
   \frac{C'_{\mathrm{box}}}{|t|}\,I.
\]
Combining the two estimates,
\[
   \lambda_{\min}\mathsf P[M_\Xi]
   \ge
   \lambda_{\mathcal P}\log\frac{|t|}{2\pi}
   -
   C_1(N,\epsilon_0)
   -
   \frac{C'_{\mathrm{box}}}{|t|}.
\]
Lemma~\ref{lem:szego-explicit} gives
\[
   \lambda_{\mathcal P}
   \ge
   c(N,\epsilon_0,Y,A)\prod_{i<j}|w_i-w_j|^2.
\]
Choose $C_{\mathrm{box}}$ and then $T_{\mathrm{box}}$ so the product hypothesis implies
\[
   \lambda_{\mathcal P}\log\frac{|t|}{2\pi}
   \ge
   2C_1(N,\epsilon_0)+2C'_{\mathrm{box}}/|t|.
\]
The displayed lower bound follows.  Since there are $N(N-1)/2$ pairs and the product contains squared
gaps, $\prod_{i<j}|w_i-w_j|^2\ge h^{N(N-1)}$, which gives the stated gap threshold.
\end{proof}

\begin{corollary}[H7 closed away from the degenerate two-node scale]\label{cor:H7-nondegenerate}
For scalar two-node patterns in a fixed box, the level-two Pick determinant is positive at sufficiently
large height whenever
\[
   |w_1-w_2|^2\ge\frac{C_{\mathrm{box}}}{\log(|t|/2\pi)}.
\]
Equivalently, this corollary localizes the remaining two-node work to the degenerating scale
$|w_1-w_2|\lesssim(\log|t|)^{-1/2}$ and below; Corollary~\ref{cor:below-threshold} closes the fixed
coalescent endpoint.
\end{corollary}

\begin{proof}
This is Theorem~\ref{thm:cluster-quantitative} with $N=2$.  For a $2\times2$ Hermitian Pick matrix,
positive definiteness is equivalent to positivity of the two diagonal entries and the determinant.
\end{proof}

\begin{remark}[where H8 still carries RH-level hardness]\label{rmk:szego-hardness}
Theorem~\ref{thm:cluster-quantitative} quantifies exactly how far the fixed-$N$ scalar cluster method
goes: the only remaining scalar cluster regime is
\[
   \prod_{i<j}|w_i-w_j|^2
   \lesssim
   \frac1{\log|t|}
\]
and the passage $N\to\infty$.  The Nevanlinna--Pick bridge needs accumulating node families, precisely
where the Cauchy determinant degenerates.  Thus the remaining H8 difficulty is uniformity below the
explicit Szeg\H{o}/Cauchy determinant scale, not fixed-cluster positivity.
\end{remark}

\begin{definition}[H7 and remaining H8 coordinates]\label{def:live-open-coordinates}
H7 asks for a sum-of-squares certificate for the level-two Andreief/Vandermonde determinant as an
identity/effectivity refinement; scalar fixed-band positivity itself is closed by
Corollary~\ref{cor:H7-band}.  The remaining H8 coordinate asks for uniform no-margin control as
$N\to\infty$.  In the
coincident-node limit one obtains the jet matrices
\[
   J_N(z)=
   \left[
      \frac1{i!\,j!}\partial_z^i\partial_{\bar w}^jK(z,w)\big|_{w=z}
   \right]_{i,j=0}^{N-1}.
\]
The scalar fixed-$N$ band regime for all configurations is closed by Theorem~\ref{thm:band-closure}.
In the terminal one-point gauge the Pascal--Laguerre floor is effective by
Lemma~\ref{lem:pascal-effective}, but the essential open part remains the exact sign
$\delta_N\ge0$ for all $N$ for one fixed interior-band accumulating sequence, equivalently the
terminal scalar input of Theorem~\ref{thm:minimal-input}.
\end{definition}

\begin{lemma}[jet coefficient kernel]\label{lem:jet-coeff}
Fix a detection channel $F_\kappa=\operatorname{span}\{\phi_\kappa\}$ from
Definition~\ref{def:detection-family} and a point $z_0\in U_-$.  The scalar Pick kernel
\[
   K^{(\kappa)}(z,w)
   :=
   \frac{g^{F_\kappa}(z)-\overline{g^{F_\kappa}(w)}}{z-\bar w}
\]
is holomorphic in $z$ and antiholomorphic in $w$ on $U_-\times U_-$.  Its double Taylor coefficients
\[
   C^{(\kappa)}_{jk}(z_0)
   :=
   \frac1{j!\,k!}
   \partial_z^j\partial_{\bar w}^kK^{(\kappa)}(z,w)\big|_{z=w=z_0}
\]
generate $K^{(\kappa)}$ by an absolutely convergent series on
$D_r(z_0)\times D_r(z_0)$, where
\[
   r=\frac12(|\Im z_0|-\half)>0.
\]
Under RH,
\[
   C^{(\kappa)}_{jk}(z_0)
   =
   \int_\R
      (u-z_0)^{-(j+1)}(u-\bar z_0)^{-(k+1)}
      \,d\mu_{\phi_\kappa}(u),
\]
so every finite coefficient section is positive semidefinite.
\end{lemma}

\begin{proof}
On $U_-$, all poles of $g^{F_\kappa}$ lie at distance at least $|\Im z_0|-\half$ from $z_0$, and
$z-\bar w\ne0$ for $z,w\in U_-$.  Thus the two-variable kernel is holomorphic-antiholomorphic in a
polydisk of radius $|\Im z_0|-\half$ around $(z_0,z_0)$; Cauchy's estimates give the absolute
convergence on the smaller polydisk of radius $r$.

Under RH, Theorem~\ref{thm:rh-pick} realizes $g^{F_\kappa}$ as the Cauchy transform of the
positive measure $\mu_{\phi_\kappa}$.  Lemma~\ref{lem:cauchy-pick} gives
\[
   K^{(\kappa)}(z,w)
   =
   \int_\R\frac{d\mu_{\phi_\kappa}(u)}{(u-z)(u-\bar w)}.
\]
Differentiating under the integral at $z=w=z_0$ gives the displayed Gram coefficients.  Positivity of
finite sections follows by testing finite linear combinations of the functions
$(u-z_0)^{-(j+1)}$ in $L^2(d\mu_{\phi_\kappa})$.
\end{proof}

\begin{theorem}[one-point jet serialization of RH]\label{thm:one-point-jets}
Fix any single point $z_0\in U_-$.  Then
\[
   \mathrm{RH}
   \quad\Longleftrightarrow\quad
   \bigl[C^{(\kappa)}_{jk}(z_0)\bigr]_{j,k=0}^{N-1}\succeq0
   \quad\text{for every }N\ge1\text{ and every }\kappa\ge1 .
\]
\end{theorem}

\begin{proof}
Assume RH.  Positivity of every finite jet section is Lemma~\ref{lem:jet-coeff}.

Conversely, assume all displayed finite jet matrices are positive.  Fix $\kappa$.  Let
$z_1,\ldots,z_m\in D_r(z_0)$ and $c_1,\ldots,c_m\in\C$.  Put
\[
   v_\alpha=\sum_i \bar c_i(z_i-z_0)^\alpha.
\]
By the absolutely convergent Taylor expansion,
\[
   \sum_{i,j}\bar c_iK^{(\kappa)}(z_i,z_j)c_j
   =
   \sum_{\alpha,\beta\ge0}
      C^{(\kappa)}_{\alpha\beta}(z_0)\,v_\alpha\overline{v_\beta}
   =
   \lim_{N\to\infty}
   \sum_{\alpha,\beta<N}
      C^{(\kappa)}_{\alpha\beta}(z_0)\,v_\alpha\overline{v_\beta}
   \ge0.
\]
Thus every finite Pick matrix with nodes in the disk $D_r(z_0)$ is positive semidefinite.  Choose a
countable, pairwise distinct node set in this disk with an accumulation point in the disk.  Since
$D_r(z_0)\subset U_-$, the meromorphic $\GXi^{F_\kappa}$ is holomorphic at every node and near the
accumulation point; Lemma~\ref{lem:prefix-extension}, applied to this one channel and this node set,
removes every pole of $\GXi^{F_\kappa}$ in $\Omega_-$.  Running over all $\kappa$ and using
Lemma~\ref{lem:dense-detection} excludes all non-real zeros in $\Omega_-$; the symmetries of $\Xi$
exclude the upper half-plane, and Lemma~\ref{lem:rh-translation} gives RH.
\end{proof}

\begin{lemma}[paired Hadamard kernel of the scalar coordinate]\label{lem:mxi-kernel}
Let $Z(\Xi)=\{z_\rho\}$ denote the zero multiset of $\Xi$ (multiplicities $m_\rho$); it is symmetric
under $z\mapsto-z$ and $z\mapsto\bar z$ (Lemma~\ref{lem:xi-basic}), contained in $|\Im z|<\half$, does
not contain $0$ (since $\xi(\half)\ne0$), and satisfies $\sum_\rho m_\rho|z_\rho|^{-2}<\infty$ by the
counting estimate $N(T)=O(T\log T)$.  Then:
\begin{enumerate}[label=\textnormal{(\roman*)}]
\item \textnormal{(paired Hadamard expansion)}
\[
   M_\Xi(z)
   =
   \lim_{R\to\infty}\sum_{|z_\rho|\le R}\frac{m_\rho}{z_\rho-z},
\]
locally uniformly on $\C\setminus Z(\Xi)$, the truncation being over the symmetric multiset
$\{|z_\rho|\le R\}$;
\item \textnormal{(kernel identity)} for $z,w\in\C$ with $z,\bar w\notin Z(\Xi)$ and $z\ne\bar w$,
\[
   K_\Xi(z,w)
   :=
   \frac{M_\Xi(z)-\overline{M_\Xi(w)}}{z-\bar w}
   =
   \sum_\rho\frac{m_\rho}{(z_\rho-z)(z_\rho-\bar w)},
\]
the sum converging absolutely (terms are $O(|z_\rho|^{-2})$ locally uniformly), and the identity
extending across $z=\bar w$ by continuity;
\item \textnormal{(jets)} if $z_0\notin\R$ and $d_0:=\operatorname{dist}(z_0,Z(\Xi))>0$, then, with
$r_0:=\min(d_0,|\Im z_0|)$, the kernel $K_\Xi$ is holomorphic in $z$ and antiholomorphic in $w$ on
$D_{r_0}(z_0)\times D_{r_0}(z_0)$, its scalar jets
\[
   C_{jk}(z_0)
   :=
   \frac1{j!\,k!}\,\partial_z^j\partial_{\bar w}^k K_\Xi(z,w)\big|_{z=w=z_0}
   =
   \sum_\rho m_\rho\,(z_\rho-z_0)^{-(j+1)}(z_\rho-\bar z_0)^{-(k+1)}
\]
converge absolutely, and the double Taylor series of $K_\Xi$ around $(z_0,z_0)$ converges absolutely on
$D_r(z_0)\times D_r(z_0)$ for every $r<r_0$;
\item \textnormal{(positivity under RH)} if RH holds, every $z_\rho$ is real, every finite jet section
$[C_{jk}(z_0)]_{j,k<N}$ is positive semidefinite, and every finite Pick matrix
$\mathsf P[M_\Xi;z_1,\ldots,z_N]$ at nodes $z_1,\ldots,z_N\in\Omega_-\setminus Z(\Xi)$ is positive
semidefinite.
\end{enumerate}
\end{lemma}

\begin{proof}
\emph{(i)}  It is classical \cite[Ch.~II]{Titchmarsh} that $\xi(s)$, and hence $\Xi(z)=\xi(\half+iz)$, is
entire of order exactly $1$; since $\sum_\rho|z_\rho|^{-1}=\infty$ while
$\sum_\rho m_\rho|z_\rho|^{-2}<\infty$ (Riemann--von Mangoldt), the zero set has convergence exponent
$1$ and genus $1$, so the Hadamard factorization theorem applies with genus-$1$ elementary factors:
\[
   \Xi(z)=\Xi(0)\,e^{cz}\prod_\rho\left(1-\frac z{z_\rho}\right)e^{z/z_\rho},
\]
$c$ a constant, the product taken (as always in this paper) in the symmetric sense of
Convention~\ref{conv:symmetric-sum} and converging locally uniformly by the genus-$1$ estimate
$|(1-u)e^u-1|\le|u|^2$ for $|u|\le1$.  Pairing $z_\rho$ with $-z_\rho$ (Lemma~\ref{lem:xi-basic}(i) gives
this pairing with equal multiplicity) cancels the linear exponential factors termwise,
\[
   \left(1-\frac z{z_\rho}\right)e^{z/z_\rho}\left(1+\frac z{z_\rho}\right)e^{-z/z_\rho}
   =1-\frac{z^2}{z_\rho^2},
\]
and the pairwise product
\[
   P(z):=\prod_{\text{pairs}}\left(1-\frac{z^2}{z_\rho^2}\right)
\]
converges locally uniformly and absolutely, since $|z^2/z_\rho^2|=O(|z_\rho|^{-2})$ on compacts, without
any exponential convergence factor.  Thus $\Xi(z)=\Xi(0)e^{cz}P(z)$ exactly, with the genus-$1$
factorization already accounted for; evenness of $\Xi$ (Lemma~\ref{lem:xi-basic}(ii)) and of $P$ forces
$c=0$, so $\Xi=\Xi(0)P$.  Logarithmic
differentiation term by term (legitimate by locally uniform convergence) gives
\[
   M_\Xi(z)=-\frac{\Xi'}{\Xi}(z)
   =\sum_{\text{pairs}}\frac{2z}{z_\rho^2-z^2}
   =\sum_{\text{pairs}}\left[\frac1{z_\rho-z}+\frac1{-z_\rho-z}\right],
\]
which is exactly the symmetric truncation limit, with pair terms $O(|z_\rho|^{-2})$ on compacts.

\emph{(ii)}  By Lemma~\ref{lem:xi-basic}(iii), $\Xi$ is real on $\R$, so
$M_\Xi(\bar w)=\overline{M_\Xi(w)}$; and the zero multiset is conjugation-symmetric, so the symmetric
truncations of (i) are stable under conjugation.  For each finite symmetric truncation, the algebraic
identity
\[
   \frac1{z-\bar w}\left[\frac1{z_\rho-z}-\frac1{z_\rho-\bar w}\right]
   =
   \frac1{(z_\rho-z)(z_\rho-\bar w)}
\]
holds termwise; letting $R\to\infty$, the left side converges to $K_\Xi(z,w)$ by (i) and the right side
converges absolutely since $|(z_\rho-z)(z_\rho-\bar w)|^{-1}=O(|z_\rho|^{-2})$.

\emph{(iii)}  The set $Z(\Xi)$ is conjugation-symmetric, so
$\operatorname{dist}(\bar z_0,Z(\Xi))=d_0$ as well.  For $z,w\in D_{r_0}(z_0)$: $z\notin Z(\Xi)$ and
$\bar w\notin Z(\Xi)$ (distance $\ge d_0-r_0\ge0$, strict inside the open disk), and $z\ne\bar w$
because $\Im z$ and $\Im\bar w$ have opposite (strict) signs, $|\Im z-\Im z_0|<|\Im z_0|$.  Hence every
term of (ii) is holomorphic--antiholomorphic on the polydisk and the sum converges locally uniformly
there; termwise differentiation gives the displayed jets, and Cauchy's estimates on the polydisk of
radius $r_0$ give the absolute convergence of the double Taylor series on every smaller polydisk.

\emph{(iv)}  Under RH, $z_\rho\in\R$ and $z_\rho-\bar z_0=\overline{z_\rho-z_0}$, so
\[
   \sum_{j,k<N}\bar c_j C_{jk}(z_0)c_k
   =
   \sum_\rho m_\rho\left|\sum_{j<N}c_j\,\overline{(z_\rho-z_0)^{-(j+1)}}\right|^2\ge0 ,
\]
and likewise, for nodes $z_1,\ldots,z_N\in\Omega_-\setminus Z(\Xi)$,
\[
   \sum_{i,j}\bar c_i K_\Xi(z_i,z_j)c_j
   =
   \sum_\rho m_\rho\left|\sum_j\frac{c_j}{\,\overline{z_\rho-z_j}\,}\right|^2\ge0 ,
\]
both by (ii)--(iii) with absolute convergence justifying the rearrangement.  Note that no finiteness of
the zero measure is used anywhere: the raw measure $\sum m_\rho\delta_{z_\rho}$ has infinite mass and
$M_\Xi$ is \emph{not} its Cauchy transform; it is the paired expansion (i) whose Pick kernel
nevertheless has the exact Gram form (ii).
\end{proof}

\begin{theorem}[scalar one-point jet serialization of RH; regular at the Li gauge]\label{thm:scalar-jets}
Let $z_0\in\Omega_-$ satisfy $d_0=\operatorname{dist}(z_0,Z(\Xi))>0$, and let $C_{jk}(z_0)$ be the
scalar jets of Lemma~\ref{lem:mxi-kernel}(iii).  Then
\[
   \mathrm{RH}
   \quad\Longleftrightarrow\quad
   \bigl[C_{jk}(z_0)\bigr]_{j,k=0}^{N-1}\succeq0
   \quad\text{for every }N\ge1 .
\]
The distance hypothesis is automatic in the two cases used later:
\begin{enumerate}[label=\textnormal{(\alph*)}]
\item \textnormal{(interior gauges)} if $\Im z_0<-\half$ then $d_0\ge|\Im z_0|-\half>0$, since
$Z(\Xi)\subset\{|\Im z|<\half\}$;
\item \textnormal{(Li gauge)} if $z_0=t_0-\tfrac i2$ (i.e.\ $\Re s_0=1$), then $d_0>0$
unconditionally: the change of variable $s=\half+iz$ is an isometry ($|s-s_0|=|z-z_0|$), so a zero of
$\Xi$ within distance $\delta\le1$ of $z_0$ is a nontrivial zero $\rho=\beta+i\gamma$ of $\zeta$ with
$|\rho-(1+it_0)|\le\delta$, hence $1-\beta\le\delta$ and $|\gamma-t_0|\le\delta$; the classical
zero-free region $\beta\le1-c/\log(|\gamma|+2)$ \cite[Ch.~III]{Titchmarsh} forces
$\delta\ge c/\log(|t_0|+3)$.  Thus $z_0=-\tfrac i2$ lies on the boundary of the Pick domain $U_-$ but is
a \emph{regular interior point} for the meromorphic $M_\Xi$ and its kernel: the boundary is an artifact
of the interpolation coordinates, not a singularity of the object being expanded.
\end{enumerate}
\end{theorem}

\begin{proof}
Assume RH.  Positivity of every finite jet section is Lemma~\ref{lem:mxi-kernel}(iv).

Conversely, assume all finite jet sections at $z_0$ are positive semidefinite.  Set
$r=\tfrac12\min(d_0,|\Im z_0|)>0$.  Let $z_1,\ldots,z_m\in D_r(z_0)$ and $c_1,\ldots,c_m\in\C$, and put
$v_\alpha=\sum_i\bar c_i(z_i-z_0)^\alpha$.  By the absolutely convergent double Taylor expansion of
Lemma~\ref{lem:mxi-kernel}(iii), applied at a radius strictly between $r$ and $r_0=2r$,
\[
   \sum_{i,j}\bar c_iK_\Xi(z_i,z_j)c_j
   =
   \sum_{\alpha,\beta\ge0}C_{\alpha\beta}(z_0)\,v_\alpha\overline{v_\beta}
   =
   \lim_{N\to\infty}\sum_{\alpha,\beta<N}C_{\alpha\beta}(z_0)\,v_\alpha\overline{v_\beta}
   \ge0,
\]
exactly as in the proof of Theorem~\ref{thm:one-point-jets}.  Thus every finite Pick matrix of $M_\Xi$
with nodes in $D_r(z_0)$ is positive semidefinite.  The disk $D_r(z_0)$ lies in $\Omega_-$ and avoids
$Z(\Xi)$, so $M_\Xi$ is holomorphic there.  Choose a countable, pairwise distinct node sequence in
$D_r(z_0)$ with an accumulation point in $D_r(z_0)$; Lemma~\ref{lem:prefix-extension} (scalar case,
$G=M_\Xi$) removes every pole of $M_\Xi$ in $\Omega_-$.  Every zero $z_\rho$ of $\Xi$ is a pole of
$M_\Xi$ with principal part $-m_\rho/(z-z_\rho)\ne0$ --- the detection that required the dense channel
family in Theorem~\ref{thm:one-point-jets} is automatic for the scalar coordinate.  Hence no zero of
$\Xi$ lies in $\Omega_-$; the symmetries of Lemma~\ref{lem:xi-basic}(ii)--(iii) exclude the upper
half-plane, and Lemma~\ref{lem:rh-translation} gives RH.
\end{proof}

\begin{remark}[what the scalar theorem repairs]\label{rmk:scalar-jets-role}
Theorem~\ref{thm:scalar-jets} is the single-channel companion of
Theorem~\ref{thm:one-point-jets}: the detection quantifier ``for every $\kappa$'' is traded for the
automatic nonvanishing of the scalar residues.  It is the statement actually consumed by the
whitened-terminal coordinates (Theorem~\ref{thm:whitened-terminal}), by the reference-whitened defect
(Definition~\ref{def:terminal-defect-fixed}, Corollary~\ref{cor:arp-terminal-defect}), and --- through
case (b) --- by the Li--Keiper dictionary (Theorem~\ref{thm:li-dictionary}) directly at the Li gauge
$z_0=-\tfrac i2$, without any passage of positivity from interior gauges to the boundary.
\end{remark}

\begin{lemma}[prime-scale window content of jets]\label{lem:prime-window}
For $s=\sigma+it$ with $\sigma\ge1+\epsilon_0$,
\[
   \frac1{k!}\left|\left(\frac{\zeta'}{\zeta}\right)^{(k)}(s)\right|
   \le
   \frac1{k!}\sum_{n\ge2}\Lambda(n)(\log n)^k n^{-\sigma}
   \le
   C(\epsilon_0)\left(\frac{2}{\epsilon_0}\right)^k.
\]
Moreover, for every $A\ge2$, the contribution of terms with
\[
   \left|\log n-\frac{k}{\epsilon_0}\right|>\frac{A\sqrt k}{\epsilon_0}
\]
is at most $e^{-cA^2}$ times the same majorant, with an absolute $c>0$.
\end{lemma}

\begin{proof}
The derivative Dirichlet series is absolutely convergent in $\sigma>1$, and differentiation gives
coefficients bounded in modulus by $\Lambda(n)(\log n)^k n^{-\sigma}$.  With $u=\log n$, the elementary
bound
\[
   u^ke^{-\epsilon_0u/2}\le \left(\frac{2k}{e\epsilon_0}\right)^k
\]
and the convergence of $\sum_{n\ge2}\Lambda(n)n^{-1-\epsilon_0/2}$ give the displayed exponential
bound after Stirling.  The localization statement is the Chernoff bound for the gamma-type weight
$u^ke^{-\epsilon_0u}/k!$, applied as a dominating envelope to the positive coefficients.
\end{proof}

\begin{lemma}[Szeg\H{o} jet Gram identity]\label{lem:szego-jet}
For $z_0=t-iy$ with $y\ge\half$ (i.e.\ $z_0\in \overline{U_-}$, including the boundary line $y=\half$),
the jets of the Szeg\H{o} kernel
\[
   S(z,w)=\frac{-i}{z-\bar w}
\]
satisfy
\[
   \frac1{j!\,k!}\partial_z^j\partial_{\bar w}^kS(z,w)\big|_{z=w=z_0}
   =
   \int_0^\infty
      \frac{(-i\tau)^j}{j!}e^{-iz_0\tau}\,
      \overline{\frac{(-i\tau)^k}{k!}e^{-iz_0\tau}}\,d\tau
   =
   \varepsilon_j\overline{\varepsilon_k}
   \binom{j+k}{j}\frac1{(2y)^{j+k+1}},
\]
where $\varepsilon_j=(-i)^j$.  Hence the jet matrix is unitarily diagonal-equivalent to the Gram
matrix of $\{\tau^ke^{-y\tau}/k!\}_{k<N}$ in $L^2(0,\infty)$, and is positive definite.  Write
\[
   G_N(z_0):=\left[\frac1{j!\,k!}\partial_z^j\partial_{\bar w}^kS(z,w)\big|_{z=w=z_0}\right]_{j,k=0}^{N-1}
   \succ0
\]
for this jet matrix; $G_N(z_0)$ is the reference Gram used in
Definition~\ref{def:terminal-defect-fixed}, defined and strictly positive definite for every
$z_0=t-iy\in U_-$ (equivalently every $y>\half$), with entries as displayed above.  Denote its
least eigenvalue by $\lambda_N^{\mathrm{jet}}(y)>0$.
\end{lemma}

\begin{proof}
For $z,w$ with $\Im z,\Im w<0$,
\[
   S(z,w)=\int_0^\infty e^{-iz\tau}e^{i\bar w\tau}\,d\tau,
\]
the integral converging absolutely since $|e^{-iz\tau}|=e^{\tau\Im z}$ decays for $\tau\to\infty$; this
covers every $z_0=t-iy$ with $y>0$, in particular every $y\ge\half$.
Differentiating under the integral gives the first identity.  At $z=w=z_0=t-iy$, the exponential
weight is $e^{-2y\tau}$, so
\[
   \frac{(-i)^ji^k}{j!\,k!}\int_0^\infty \tau^{j+k}e^{-2y\tau}\,d\tau
   =
   (-i)^ji^k\frac{(j+k)!}{j!\,k!}\frac1{(2y)^{j+k+1}},
\]
which is the displayed formula.  The functions $\tau^ke^{-y\tau}/k!$ are linearly independent, so
their Gram matrix is positive definite.
\end{proof}

\begin{lemma}[archimedean jet lower bound]\label{lem:arch-jet}
Fix $N$, $\epsilon_0>0$, and $Y<\infty$.  For $z_0=t-iy$ with
$y\in[\half+\epsilon_0,Y]$ and $|t|$ large, the archimedean part $J_N^\infty(z_0)$ of the scalar jet
matrix satisfies, in the unitary diagonal normalization of Lemma~\ref{lem:szego-jet},
\[
   J_N^\infty(z_0)
   \succeq
   \left(
      \lambda_N^{\mathrm{jet}}(y)\log\frac{|t|}{2\pi}
      -
      \frac{C_{N,Y,\epsilon_0}}{|t|}
   \right)I.
\]
\end{lemma}

\begin{proof}
The proof of Lemma~\ref{lem:arch-cluster} is an analytic estimate for the two-variable kernel:
on a fixed polydisk around $(z_0,z_0)$ contained in $U_-\times U_-$,
\[
   K_\infty(z,w)-\log\frac{|t|}{2\pi}S(z,w)=O_{Y,\epsilon_0}(|t|^{-1})
\]
as a holomorphic-antiholomorphic function.  Cauchy's estimates on a slightly smaller polydisk bound
all mixed derivatives of total order at most $2N$ by $C_{N,Y,\epsilon_0}/|t|$.  Taking the jet matrix
therefore gives
\[
   J_N^\infty(z_0)
   =
   \log\frac{|t|}{2\pi}J_N^S(z_0)+O_{N,Y,\epsilon_0}(|t|^{-1}),
\]
where $J_N^S$ is the positive definite Szeg\H{o} jet Gram matrix of Lemma~\ref{lem:szego-jet}.  The
eigenvalue bound follows.
\end{proof}

\begin{theorem}[fixed-$N$ coincident jet positivity at large height]\label{thm:jet-fixedN}
Fix $N$, $\epsilon_0>0$, and $Y<\infty$.  There is
$T(N,\epsilon_0,Y)$ such that for every $z_0=t-iy$ with
$y\in[\half+\epsilon_0,Y]$ and $|t|\ge T(N,\epsilon_0,Y)$, the scalar jet matrix
$J_N(z_0)$ is positive definite.  Its margin is proportional to
\[
   \lambda_N^{\mathrm{jet}}(y)\log\frac{|t|}{2\pi}.
\]
\end{theorem}

\begin{proof}
Decompose $J_N=J_N^\infty+J_N^{\mathrm{prime}}$.  The entries of $J_N^{\mathrm{prime}}$ involve
derivatives of $\zeta'/\zeta$ of order at most $2N$ at points with
$\Re s\ge1+\epsilon_0$.  Lemma~\ref{lem:prime-window} bounds each such normalized derivative by a
constant depending only on $N$ and $\epsilon_0$; the denominators in the Pick kernel are bounded away
from zero because $|z_0-\bar z_0|=2y\ge1+2\epsilon_0$.  Hence
\[
   \|J_N^{\mathrm{prime}}\|\le C_{N,\epsilon_0}.
\]
Lemma~\ref{lem:arch-jet} gives an archimedean lower bound growing like
$\lambda_N^{\mathrm{jet}}(y)\log(|t|/(2\pi))$.  Choosing $T(N,\epsilon_0,Y)$ so this dominates
$C_{N,\epsilon_0}$ and the $O(|t|^{-1})$ error proves positivity.
\end{proof}

\begin{corollary}[below-threshold coalescence for fixed $N$]\label{cor:below-threshold}
Fix $N$, $\epsilon_0>0$, and $Y<\infty$.  There is
$h_N=h_N(\epsilon_0,Y)>0$, independent of $t$, such that every scalar $N$-node cluster of diameter at
most $h_N$, centered at $z_0=t-iy$ with $y\in[\half+\epsilon_0,Y]$ and
$|t|\ge T(N,\epsilon_0,Y)$, has positive definite Pick matrix.  Together with
Theorem~\ref{thm:cluster-quantitative}, this covers single scalar clusters from full coalescence
$h=0$ through the Cauchy-determinant scale, with overlap for large $|t|$.
\end{corollary}

\begin{proof}
Pass from node values to divided-difference coordinates by the standard invertible confluent
Vandermonde matrix $E_h$ associated with the cluster.  In those coordinates,
\[
   E_h^*\mathsf P E_h=J_N(z_0)+O_N(h\log|t|)
\]
entrywise: this is Taylor's formula for the holomorphic-antiholomorphic kernel, with the archimedean
derivatives bounded by $C_N\log|t|$ and the prime derivatives bounded by Lemma~\ref{lem:prime-window}.
By Theorem~\ref{thm:jet-fixedN}, $J_N(z_0)\succeq c_N\log|t|\,I$ for large $|t|$.  Taking
$h_N<c_N/(2C_N')$ makes the perturbation smaller than half of this margin.  Congruence by the
invertible matrix $E_h$ preserves positive definiteness.
\end{proof}

\begin{lemma}[scale selection]\label{lem:scale-selection}
Let $w_1,\ldots,w_N$ lie in a box of diameter $D$.  Consider the disjoint rings
\[
   R_r=[DN^{-r},DN^{-r+1}),\qquad r=1,\ldots,N^2.
\]
Since there are fewer than $N^2$ pairwise distances, some ring $R_{r_*}$ contains none of them.  Set
$h_*=DN^{-r_*}$ and define clusters as the connected components of the graph with edges
$|w_i-w_j|<h_*$.  Then
\[
   \operatorname{diam}(\textnormal{each cluster})<Nh_*,
   \qquad
   \operatorname{dist}(\textnormal{distinct clusters})\ge Nh_*,
   \qquad
   h_*\ge DN^{-N^2}.
\]
\end{lemma}

\begin{proof}
Two nodes in the same component are joined by a chain of fewer than $N$ edges, each of length $<h_*$,
giving the diameter bound.  If two nodes lie in distinct components, their distance is at least
$h_*$.  Since the ring $[h_*,Nh_*)$ is distance-free, the distance is in fact at least $Nh_*$.  The
lower bound for $h_*$ follows from $r_*\le N^2$.
\end{proof}

\begin{definition}[confluent Newton basis]\label{def:newton-basis}
For a multiset $\zeta_0,\ldots,\zeta_k$ in $U_-$, the divided difference of
$e_\tau(z)=e^{-iz\tau}$ has the Hermite--Genocchi representation
\[
   e[\zeta_0,\ldots,\zeta_k](\tau)
   =
   (-i\tau)^k
   \int_{\Delta_k}e^{-i\zeta(s)\tau}\,ds,
   \qquad
   \zeta(s)=\zeta_0+\sum_{i=1}^k s_i(\zeta_i-\zeta_0).
\]
The formula is continuous through confluent limits and satisfies
\[
   |e[\zeta_0,\ldots,\zeta_k](\tau)|
   \le
   \frac{\tau^k}{k!}e^{-(\half+\epsilon_0)\tau}
\]
whenever all $\Im\zeta_i\le-(\half+\epsilon_0)$.  Given a clustered configuration, order each cluster
and let $E$ be the block-diagonal Newton change of basis.  For pairwise distinct nodes, $E$ is
invertible.  For every kernel $K$ holomorphic in $z$ and antiholomorphic in $w$, the entries of
$E^*\mathsf K E$ are the corresponding two-variable divided differences of $K$.
\end{definition}

\begin{lemma}[uniform Newton--Szeg\H{o} positivity]\label{lem:newton-szego}
Fix $N$, $\epsilon_0>0$, $Y$, and $A$, and set
\[
   d_0=(2A+2Y)N^{-N^2}.
\]
On the compact space of clustered configurations in the box
\[
   |a|\le A,\qquad y\in[\half+\epsilon_0,Y],
\]
with cluster types fixed, confluent intra-cluster multisets allowed, and distinct cluster centers
separated by at least $d_0$, there is a constant
$c_N^{\mathcal S}=c_N^{\mathcal S}(\epsilon_0,Y,A)>0$ such that
\[
   \widetilde S:=E^*SE\succeq c_N^{\mathcal S}I
\]
for every configuration, where $S$ is the Szeg\H{o} matrix of Lemma~\ref{lem:szego-pattern}.
\end{lemma}

\begin{proof}
The matrix $S$ is the Gram matrix of the functions $e^{-iz_j\tau}$ in $L^2(0,\infty)$.  After the
Newton congruence, $\widetilde S$ is the Gram matrix of the confluent Newton basis functions of
Definition~\ref{def:newton-basis}.  Within each cluster these functions form a triangular basis of
the Hermite span $\operatorname{span}\{\tau^r e^{-iz_c\tau}\}$.  Across distinct clusters, the
generalized frequencies are distinct because the cluster centers are separated by at least $d_0$.
Thus the Gram matrix is positive definite at every point of the parameter space.

The entries depend continuously on the configuration by dominated convergence, using the majorant
$C_N\tau^{2N}e^{-(1+2\epsilon_0)\tau}$.  The parameter space is a finite union of compact sets, one
for each partition type and cluster ordering.  Hence the continuous function
$\lambda_{\min}(\widetilde S)$ attains a strictly positive minimum.
\end{proof}

\begin{lemma}[divided differences of bounded kernels]\label{lem:dd-bounded}
Let $K(z,w)$ be holomorphic in $z$, antiholomorphic in $w$, and bounded by $B$ on the
$\epsilon_0/2$-enlarged box-domain.  Then, for every clustered configuration as above, the entries of
$E^*\mathsf K E$ are bounded by $B(2/\epsilon_0)^{2N}$, and therefore
\[
   \|E^*\mathsf K E\|\le NB(2/\epsilon_0)^{2N}.
\]
\end{lemma}

\begin{proof}
Each entry is a double Hermite--Genocchi average of
\[
   \frac{1}{j!\,k!}\partial_z^j\partial_{\bar w}^kK
\]
with $j,k<N$, over hulls contained in the enlarged box-domain.  Cauchy's estimates on polydisks of
radius $\epsilon_0/2$ give
\[
   \frac{|\partial_z^j\partial_{\bar w}^kK|}{j!\,k!}
   \le B(2/\epsilon_0)^{j+k}
   \le B(2/\epsilon_0)^{2N}.
\]
The simplex averages have mass one, and the operator-norm bound follows from the entrywise bound.
\end{proof}

\begin{theorem}[band closure of scalar ARP-P$(N)$]\label{thm:band-closure}
Fix $N$, $\epsilon_0\in(0,\half]$, $Y$, and $A$.  There is
$T^*=T^*(N,\epsilon_0,Y,A)$ such that, for every $|t|\ge T^*$ and every configuration of $N$ pairwise
distinct scalar nodes
\[
   z_j=t+a_j-iy_j,\qquad |a_j|\le A,\qquad y_j\in[\half+\epsilon_0,Y],
\]
the scalar Pick matrix satisfies $\mathsf P[M_\Xi;z_1,\ldots,z_N]\succ0$.  Coalescent limits are
positive semidefinite by continuity, with coincident blocks governed by Theorem~\ref{thm:jet-fixedN}.
\end{theorem}

\begin{proof}
Apply Lemma~\ref{lem:scale-selection} to $w_j=y_j+ia_j$ inside a box of diameter at most $2A+2Y$ and
use the resulting clustered structure.  On the enlarged box-domain, decompose the scalar Pick kernel as
\[
   K=\log\frac{|t|}{2\pi}S+R+K_{\mathrm{prime}},
\]
where $R=K_\infty-\log(|t|/(2\pi))S$.  The proof of Lemma~\ref{lem:arch-cluster} is uniform in
coalescing configurations: its numerator error is $O((1+A+Y)/|t|)$ and the denominators satisfy
$|z-\bar w|\ge1+2\epsilon_0$ on the enlarged domain.  Hence $|R|\le C_{\mathrm{box}}/|t|$.  Also
$|K_{\mathrm{prime}}|\le C(\epsilon_0,A,Y)$ there, because $\zeta'/\zeta$ is bounded on
$\Re s\ge1+\epsilon_0/2$.

Apply the Newton congruence $E$:
\[
   E^*\mathsf P E
   =
   \log\frac{|t|}{2\pi}\,\widetilde S
   +
   \widetilde R
   +
   \widetilde K_{\mathrm{prime}}.
\]
Lemma~\ref{lem:newton-szego} gives $\widetilde S\succeq c_N^{\mathcal S}I$.  Lemma~\ref{lem:dd-bounded}
gives
\[
   \|\widetilde K_{\mathrm{prime}}\|\le C_N(\epsilon_0,A,Y),
   \qquad
   \|\widetilde R\|\le C_N'(\epsilon_0,A,Y)/|t|.
\]
Therefore
\[
   E^*\mathsf P E
   \succeq
   \left(
      c_N^{\mathcal S}\log\frac{|t|}{2\pi}
      -
      C_N
      -
      \frac{C_N'}{|t|}
   \right)I.
\]
This is positive for $|t|\ge T^*$.  Since $E$ is invertible for pairwise distinct nodes, congruence
preserves positive definiteness of the original Pick matrix.
\end{proof}

\begin{corollary}[H7 band closure]\label{cor:H7-band}
For $N=2$, Theorem~\ref{thm:band-closure} closes scalar level-two positivity for every two-node
configuration in a fixed band at all heights $|t|\ge T^*(2,\epsilon_0,Y,A)$, including degenerating
and coincident limits.  The Andreief/SOS program remains useful as an identity-level and effectivity
program, but is no longer essential for scalar fixed-band positivity.
\end{corollary}

\begin{remark}[subsumption, ineffectivity, and the guardrail]\label{rmk:band-honest}
Theorem~\ref{thm:band-closure} subsumes Theorems~\ref{thm:cluster},
\ref{thm:cluster-quantitative}, \ref{thm:jet-fixedN}, and Corollary~\ref{cor:below-threshold} as
fixed-$N$ scalar band-positivity statements.  Those earlier results are retained because their
constants are effective; the compactness constant $c_N^{\mathcal S}$ in
Lemma~\ref{lem:newton-szego} is not.  There is no conflict with Lemma~\ref{lem:live-H0}: the
Nevanlinna--Pick bridge needs all finite prefixes of a single accumulating sequence, while here
$T^*(N)$ may grow without control as $N\to\infty$.  The band floor $\epsilon_0>0$ is also essential
for this theorem, but the minimal scalar bridge may be placed inside a fixed interior band.  After
this closure, the essential Step~5 content is the terminal scalar condition of
Theorem~\ref{thm:minimal-input}; PDH at the edge and effectivity of the general compactness constant
$c_N^{\mathcal S}$ are important but non-minimal subproblems.  In the one-point terminal gauge this
effectivity is replaced by the explicit Pascal--Laguerre floor of Lemma~\ref{lem:pascal-effective}.
\end{remark}

\begin{theorem}[minimal scalar open input]\label{thm:minimal-input}
Fix a box
\[
   \mathcal B=\{a-iy:\ |a|\le A,\ y\in[\half+\epsilon_0,Y]\}\subset U_-,
\]
a height shift $t_0\in\R$, and a pairwise distinct node sequence
\[
   z_j=t_0+a_j-iy_j\in t_0+\mathcal B
\]
with an accumulation point in $t_0+\mathcal B$.  Then
\[
   \mathrm{RH}
   \quad\Longleftrightarrow\quad
   \mathsf P[M_\Xi;z_1,\ldots,z_N]\succeq0
   \quad\text{for every }N\ge1 .
\]
\end{theorem}

\begin{proof}
Assume RH.  Lemma~\ref{lem:mxi-kernel}(iv) gives the finite Pick positivity of $M_\Xi$ at every finite
node set in $\Omega_-\setminus Z(\Xi)$, in particular at the given nodes.  (The raw zero measure has
infinite mass and $M_\Xi$ is not its Cauchy transform; the positivity comes from the paired Hadamard
kernel identity of Lemma~\ref{lem:mxi-kernel}(ii), not from Lemma~\ref{lem:cauchy-pick}.)

Conversely, assume the displayed positivity for every prefix of the fixed sequence.  The nodes are
pairwise distinct, accumulate in $t_0+\mathcal B\subset U_-$, and $M_\Xi$ is meromorphic on $\C$ and
holomorphic at every node and near the accumulation point, since its poles lie in the zero strip
$|\Im z|<\half$.  Lemma~\ref{lem:prefix-extension} (scalar case, $G=M_\Xi$) therefore extends $M_\Xi$
holomorphically to $\Omega_-$.

Every zero $z_\rho$ of $\Xi$ is a pole of
\[
   M_\Xi(z)=-\frac{\Xi'(z)}{\Xi(z)}
\]
with nonzero residue $-m_\rho$.  Hence no zero of $\Xi$ lies in $\Omega_-$.  The symmetries
Lemma~\ref{lem:xi-basic}(ii)--(iii) gives $\Xi(-z)=\Xi(z)$ and
$\Xi(\bar z)=\overline{\Xi(z)}$, excluding the upper half-plane.  Thus all zeros of $\Xi$ are real,
which is RH by Lemma~\ref{lem:rh-translation}.
\end{proof}

\begin{definition}[terminal Pick inequality]\label{def:terminal-inequality}
Fix a box, height $t_0$, and accumulating scalar node sequence as in
Theorem~\ref{thm:minimal-input}.  For each prefix of length $N$, apply the scale selection and Newton
congruence of Lemmas~\ref{lem:scale-selection}--\ref{lem:newton-szego}.  The exact terminal condition is
\[
   \lambda_{\min}\!\left(
      \log\frac{|t_0|}{2\pi}\,\widetilde S_N
      +\widetilde R_N
      +\widetilde K_{\mathrm{prime},N}
   \right)\ge0
   \qquad(N\ge1),
\]
with the obvious interpretation for bounded $|t_0|$ by replacing the logarithmic asymptotic with the
exact archimedean kernel.  The sufficient norm form is
\[
   c_N^{\mathcal S}\log\frac{|t_0|}{2\pi}
   \ge
   \|\widetilde K_{\mathrm{prime},N}\|
   +
   \|\widetilde R_N\|.
\]
\end{definition}

\begin{remark}[terminal inequality status]\label{rmk:terminal-status}
By Theorem~\ref{thm:minimal-input}, the exact terminal condition in
Definition~\ref{def:terminal-inequality} for one accumulating scalar sequence is equivalent to RH.
The displayed norm inequality is a useful sufficient quantitative form, not an equivalent reformulation
unless the norm loss is removed.  The left side is governed by the Newton--Szeg\H{o} lower bound
$c_N^{\mathcal S}$; the right side is governed by divided differences of the prime kernel.  The crude
Cauchy estimate grows exponentially in $N$, while the confluent Szeg\H{o}/Laguerre constants may decay
exponentially.  The remaining essential problem is therefore cancellation in the prime divided
differences strong enough to control this race uniformly in $N$ for a fixed height.
\end{remark}

\begin{theorem}[sequence independence of the terminal input]\label{thm:sequence-gauge}
Let $\{z_j\}$ and $\{z_j'\}$ be two pairwise distinct node sequences contained in interior boxes of
$U_-$, each with an accumulation point in its box.  Then
\[
   \bigl[\mathsf P[M_\Xi;z_1,\ldots,z_N]\succeq0\ \textnormal{for every }N\bigr]
   \Longleftrightarrow
   \bigl[\mathsf P[M_\Xi;z_1',\ldots,z_N']\succeq0\ \textnormal{for every }N\bigr]
   \Longleftrightarrow
   \mathrm{RH}.
\]
Thus the choice of an accumulating sequence is a coordinate gauge for attacking the terminal input, not
a weakening of its mathematical content.
\end{theorem}

\begin{proof}
Prefix positivity for either sequence implies RH by Theorem~\ref{thm:minimal-input}, applied to the
corresponding box and accumulation point.  Conversely, RH implies scalar Pick positivity for every finite
node set by Lemma~\ref{lem:mxi-kernel}(iv).  Applying this to all prefixes of
either sequence proves the equivalence.
\end{proof}

\begin{lemma}[effective Pascal--Laguerre floor]\label{lem:pascal-effective}
Let
\[
   P_N=\left[\binom{j+k}{j}\right]_{j,k=0}^{N-1},
   \qquad
   L_N=\left[\binom{j}{k}\right]_{0\le k\le j<N}.
\]
Then $P_N=L_NL_N^*$ and
\[
   \|L_N\|^2\le4^N,\qquad \|L_N^{-1}\|^2\le4^N,
\]
hence
\[
   4^{-N}\le\lambda_{\min}(P_N)\le\lambda_{\max}(P_N)\le4^N .
\]
Consequently the Szeg\H{o} jet Gram matrix of Lemma~\ref{lem:szego-jet} satisfies
\[
   \lambda_N^{\mathrm{jet}}(y)\ge4^{-N}(2y)^{-(2N-1)},\qquad y\ge\half .
\]
\end{lemma}

\begin{proof}
The Cholesky identity is Vandermonde convolution:
\[
   \binom{j+k}{j}=\sum_{r\ge0}\binom{j}{r}\binom{k}{r}.
\]
The inverse of the lower binomial triangle is the signed lower binomial triangle,
$(L_N^{-1})_{jk}=(-1)^{j-k}\binom{j}{k}$.  The row and column sums of $|L_N|$ and $|L_N^{-1}|$ are
bounded by powers of $2$; in particular
\[
   \|L_N\|^2\le\|L_N\|_1\|L_N\|_\infty\le4^N,
   \qquad
   \|L_N^{-1}\|^2\le4^N .
\]
Therefore
$\lambda_{\max}(P_N)=\|L_N\|^2\le4^N$ and
$\lambda_{\min}(P_N)=\|L_N^{-1}\|^{-2}\ge4^{-N}$.

By Lemma~\ref{lem:szego-jet}, after removing unimodular phase factors the $N$-jet Szeg\H{o} Gram is
$D_yP_ND_y$, where $D_y=\diag((2y)^{-j-1/2})_{j=0}^{N-1}$.  For $y\ge\half$,
$\min_j(D_y)_{jj}=(2y)^{-N+1/2}$, so
\[
   \lambda_{\min}(D_yP_ND_y)
   \ge (2y)^{-2N+1}\lambda_{\min}(P_N)
   \ge4^{-N}(2y)^{-(2N-1)}.
\]
\end{proof}

\begin{theorem}[whitened terminal operator]\label{thm:whitened-terminal}
Fix $z_0=t_0-iy\in U_-$ and an integer $N$ such that the archimedean jet matrix
$J_N^\infty(z_0)$ is positive definite.  Define
\[
   T_N(z_0):=
   \bigl(J_N^\infty(z_0)\bigr)^{-1/2}
   J_N^{\mathrm{prime}}(z_0)
   \bigl(J_N^\infty(z_0)\bigr)^{-1/2}.
\]
Then
\[
   J_N(z_0)=J_N^\infty(z_0)-J_N^{\mathrm{prime}}(z_0)\succeq0
   \quad\Longleftrightarrow\quad
   \lambda_{\max}(T_N(z_0))\le1 .
\]
Consequently, for every $N\le N_*(t_0)$ in the range where $J_N^\infty(z_0)\succ0$
(Lemma~\ref{lem:pole-pair}), $J_N(z_0)\succeq0$ is equivalent to $\lambda_{\max}(T_N(z_0))\le1$; the
reference-whitened observable of Definition~\ref{def:terminal-defect-fixed} carries the equivalence with
RH to \emph{all} $N$ without restriction.
\end{theorem}

\begin{proof}
The first assertion is congruence by the invertible matrix
$(J_N^\infty)^{-1/2}$:
\[
   J_N\succeq0
   \Longleftrightarrow
   I-(J_N^\infty)^{-1/2}J_N^{\mathrm{prime}}(J_N^\infty)^{-1/2}\succeq0.
\]
Since the displayed matrix is Hermitian, this is equivalent to
$\lambda_{\max}(T_N)\le1$.  If the archimedean matrices are positive for every $N$, the collection of
conditions $J_N(z_0)\succeq0$ for all $N$ is the scalar one-point jet serialization of RH,
Theorem~\ref{thm:scalar-jets} (its distance hypothesis holds by case (a), since $\Im z_0<-\half$; the
detection is automatic because the scalar residues $-m_\rho$ are nonzero at every zero).
\end{proof}

\begin{lemma}[pole-pair obstruction in the archimedean gauge]\label{lem:pole-pair}
Write the archimedean scalar as $g_\infty=g_{\rm arch}+g_{\rm pole}$, where
$g_{\rm pole}(z)=\frac1{i/2-z}+\frac1{-i/2-z}$ is the contribution of the trivial factor $s(s-1)$ (each
pole $\pm i/2$ carrying weight $+1$, since $i(-\tfrac1s-\tfrac1{s-1})=\tfrac1{i/2-z}+\tfrac1{-i/2-z}$).
The pole pair has zero boundary density on $\R$ --- so the archimedean weight $w_A$ of
Lemma~\ref{lem:signed-arch} does not see it --- but it is an interior conjugate pair for every gauge
$z_0=t_0-iy$, with exterior Blaschke modulus
\[
   |b(-i/2)|=\Bigl(\tfrac{t_0^2+(y+1/2)^2}{t_0^2+(y-1/2)^2}\Bigr)^{1/2}
   =1+\tfrac{y}{t_0^2}+O(t_0^{-4}).
\]
Consequently the archimedean jet matrix $J_N^\infty(z_0)$ is \emph{not} positive definite for all $N$:
it acquires a negative direction at scale $N_*(t_0)\asymp y^{-1}t_0^2\log t_0$, by the mechanism of
Proposition~\ref{prop:blaschke-rate}.  For $N\le N_*(t_0)$ positivity holds, with the threshold constant
made explicit in Proposition~\ref{prop:blaschke-rate}.  The effect is invisible numerically
($N_*\sim10^5$ at $t_0=300$).
\end{lemma}

\begin{definition}[terminal defect, reference-whitened]\label{def:terminal-defect-fixed}
For $z_0=t_0-iy$ with $y\ge\half$, let $J_N(z_0)=[C_{jk}(z_0)]_{j,k<N}$ be the scalar jet matrix of
Lemma~\ref{lem:mxi-kernel}(iii) and $G_N(z_0)$ the reference Gram of Lemma~\ref{lem:szego-jet}, defined
and strictly positive definite on the same domain $y\ge\half$.  Set
\[
   \delta_N(z_0):=\min_{c\ne0}\,\frac{c^*J_N(z_0) c}{c^*G_N(z_0) c}.
\]
Since $G_N(z_0)\succ0$ for every $N$ and every such $z_0$, immediately
$\delta_N(z_0)\ge0\iff J_N(z_0)\succeq0$, at \emph{every} gauge $z_0=t_0-iy$, $y\ge\half$, without
restriction on $t_0$ --- in particular at the Li gauge $z_0=-\tfrac i2$ ($t_0=0$, $y=\half$).  Hence
$\delta_N\ge0\ \forall N\iff J_N\succeq0\ \forall N\iff$ RH (Theorem~\ref{thm:scalar-jets}, at any
$z_0$ satisfying its distance hypothesis; see Remark~\ref{rmk:defect-scope}), for
\emph{all} $N$ without restriction.

The sign-defect is deliberately kept separate from asymptotic scaling: where a height-normalized
magnitude is wanted (the large-height asymptotics of Theorem~\ref{thm:cluster} and the sampling
theorems), use the rescaled variant
\[
   \widetilde\delta_N(z_0):=\frac{\delta_N(z_0)}{L},\qquad L=\log\tfrac{|t_0|}{2\pi},
\]
defined only in the regime $L>0$ (i.e.\ $|t_0|>2\pi$), where it has the same sign as $\delta_N$.  All
numerical tables in this paper are computed at gauges with $|t_0|\gg2\pi$, where the two normalizations
differ by the positive factor $L$ and every stated sign and every stated order of magnitude in $N$ is
unaffected.  The $J_N^\infty$-whitened coordinate $T_N$ of
Theorem~\ref{thm:whitened-terminal} is equivalent to this one exactly in the range $N\le N_*(t_0)$ of
Lemma~\ref{lem:pole-pair} and is used only there; the numerical harness already computes the
reference-whitened form (it whitens with the reference Gram $G_N$, not with $J_N^\infty$), so its
outputs are unaffected.
\end{definition}

\begin{remark}[exact scope of the admissible gauge]\label{rmk:defect-scope}
Definition~\ref{def:terminal-defect-fixed} defines $\delta_N(z_0)$ for every $z_0=t_0-iy$ with
$y\ge\half$; this is the full domain on which the reference Gram $G_N(z_0)$ of
Lemma~\ref{lem:szego-jet} is positive definite.  The equivalence
$\delta_N(z_0)\ge0\ \forall N\iff\mathrm{RH}$ used in
Corollary~\ref{cor:arp-terminal-defect} additionally requires the distance hypothesis of
Theorem~\ref{thm:scalar-jets}, namely $d_0(z_0):=\operatorname{dist}(z_0,Z(\Xi))>0$.  By
Theorem~\ref{thm:scalar-jets}(a)--(b) this holds at \emph{every} such gauge: for $y>\half$ it holds
because $Z(\Xi)\subset\{|\Im z|<\half\}$ forces $d_0(z_0)\ge y-\half>0$, and at $y=\half$ (the Li line)
it holds unconditionally by the classical zero-free region, uniformly in $t_0$.  Hence the admissible
scope of $z_0$ in Corollary~\ref{cor:arp-terminal-defect} is exactly
\[
   z_0=t_0-iy,\qquad y\ge\half,\qquad t_0\in\R,
\]
with no further restriction: the corollary holds at every gauge on or above the Li line, uniformly in
$t_0$, including $t_0=0$.
\end{remark}

\begin{corollary}[ARP-P--terminal-defect equivalence]\label{cor:arp-terminal-defect}
Fix $z_0=t_0-iy$ with $y\ge\half$ (Remark~\ref{rmk:defect-scope}).  For the fixed scalar channel and the
reference-whitened defect of Definition~\ref{def:terminal-defect-fixed} at this $z_0$,
\[
   \textnormal{ARP-P}
   \quad\Longleftrightarrow\quad
   \delta_N(z_0)\ge0\quad\text{for every }N.
\]
\end{corollary}

\begin{proof}
By Definition~\ref{def:terminal-defect-fixed}, $\delta_N(z_0)\ge0$ for every $N$ is equivalent to
$J_N(z_0)\succeq0$ for every $N$, since the reference Gram $G_N(z_0)$ is positive definite at every
$y\ge\half$; by Remark~\ref{rmk:defect-scope} the distance hypothesis of
Theorem~\ref{thm:scalar-jets} holds at this $z_0$ (case (a) if $y>\half$, case (b) if $y=\half$), so
$J_N(z_0)\succeq0\ \forall N$ is equivalent to RH; and by
Corollary~\ref{cor:live-equivalence}, RH is equivalent to ARP-P.
\end{proof}

\begin{theorem}[prime shift-operator representation]\label{thm:prime-shift}
On $L^2(0,\infty)$ let
\[
   (S_\lambda f)(\tau)=f(\tau-\lambda)\mathbf 1_{\tau\ge\lambda}
\]
and define the von Mangoldt shift form
\[
   \mathcal P_\Lambda
   :=
   \sum_{n\ge2}\Lambda(n)n^{-1/2}\bigl(S_{\log n}+S_{\log n}^*\bigr)
\]
on finite spans of exponentials $e^{-iz\tau}$ with $\Im z\le-(\half+\epsilon_0)$.  If
\[
   F(\tau)=\sum_j\bar c_j e^{-iz_j\tau},
\]
then the scalar prime part of the Pick form satisfies
\[
   c^*\mathsf P_{\mathrm{prime}}c
   =
   -\,\langle \mathcal P_\Lambda F,F\rangle_{L^2(0,\infty)}.
\]
\end{theorem}

\begin{proof}
On $\Re s>1$,
\[
   -\frac{\zeta'}{\zeta}(s)=\sum_{n\ge2}\Lambda(n)n^{-s},
   \qquad
   n^{-s}=n^{-1/2}e^{-i(\log n)z}.
\]
For a single frequency $\lambda>0$, using
\[
   \frac1{z-\bar w}=i\int_0^\infty e^{-i(z-\bar w)\tau}\,d\tau
   \qquad(\Im(z-\bar w)<0),
\]
one obtains
\[
   \frac{i\bigl(e^{-i\lambda z}+e^{i\lambda\bar w}\bigr)}{z-\bar w}
   =
   -\int_0^\infty
      \left[
        e^{-iz(\tau+\lambda)}e^{i\bar w\tau}
        +
        e^{-iz\tau}e^{i\bar w(\tau+\lambda)}
      \right]d\tau .
\]
After summing against $\bar c_ic_j$, the right side is
\[
   -2\Re\int_0^\infty F(\tau+\lambda)\overline{F(\tau)}\,d\tau
   =
   -\langle(S_\lambda+S_\lambda^*)F,F\rangle .
\]
The sum over $n$ is absolutely convergent on the indicated spans because
$|F(\tau)|\le Ce^{-(\half+\epsilon_0)\tau}$ and hence
\[
   |\langle(S_{\log n}+S_{\log n}^*)F,F\rangle|
   \le C'n^{-\half-\epsilon_0},
\]
so it is dominated by $\sum_n\Lambda(n)n^{-1-\epsilon_0}$.
\end{proof}

\begin{theorem}[terminal spectral form]\label{thm:terminal-spectral}
Fix an interior-band accumulating scalar sequence and let $\mathcal E$ be the finite linear span of
the corresponding exponentials $e^{-iz_j\tau}$.  The terminal scalar input is equivalent to
\[
   \langle\mathcal P_\Lambda F,F\rangle\le \mathcal A(F,F)
   \qquad(F\in\mathcal E),
\]
where $\mathcal A$ is the exact archimedean/pole Pick form.  In a large-height fixed band,
\[
   \mathcal A(F,F)=\log\frac{|t_0|}{2\pi}\|F\|_{L^2(0,\infty)}^2+\mathcal R(F,F),
\]
with $\mathcal R$ the explicitly bounded Stirling remainder from Lemma~\ref{lem:arch-cluster}.
\end{theorem}

\begin{proof}
By Theorem~\ref{thm:prime-shift}, the prime contribution to the scalar Pick form is
$-\langle\mathcal P_\Lambda F,F\rangle$.  The remaining archimedean and pole terms define
$\mathcal A(F,F)$.  Therefore positivity of every finite scalar Pick prefix is exactly
\[
   \mathcal A(F,F)-\langle\mathcal P_\Lambda F,F\rangle\ge0
\]
on the finite exponential span.  The equivalence with RH is Theorem~\ref{thm:minimal-input}.  The
large-height expansion of $\mathcal A$ is the same Stirling/Szeg\H{o} expansion used in
Lemma~\ref{lem:arch-cluster}.
\end{proof}

\begin{proposition}[Laguerre matrix elements of the terminal operator]\label{prop:laguerre-entries}
At $z_0=t_0-iy$, the leading Szeg\H{o} part of the archimedean jet Gram is diagonalized by the
orthonormal Laguerre basis
\[
   L_k^{(y,t_0)}(\tau)=e^{-it_0\tau}\sqrt{2y}\,\ell_k(2y\tau)e^{-y\tau},
\]
where $\ell_k$ are the standard orthonormal Laguerre polynomials on $[0,\infty)$ with weight
$e^{-u}du$.  In this basis, the whitened prime operator has entries
\[
   (T_N)_{jk}
   =
   -\frac1{\log(|t_0|/2\pi)}
   \sum_{n\ge2}\frac{\Lambda(n)}{\sqrt n}
   \left[
      n^{-it_0}m_{jk}^{(y)}(\log n)
      +n^{it_0}\overline{m_{kj}^{(y)}(\log n)}
   \right]
   +O\!\left(\frac{E_N(y)}{\log |t_0|}\right),
\]
where
\[
   m_{jk}^{(y)}(\lambda)
   =
   \int_0^\infty
      \sqrt{2y}\,\ell_j(2y(u+\lambda))e^{-y(u+\lambda)}
      \sqrt{2y}\,\ell_k(2yu)e^{-yu}\,du .
\]
Each $m_{jk}^{(y)}(\lambda)$ is a polynomial in $\lambda$ times $e^{-y\lambda}$ and is computable in
closed form.  The error $E_N(y)$ is effective and comes only from the archimedean Stirling remainder.
\end{proposition}

\begin{proof}
Lemma~\ref{lem:szego-jet} identifies the leading archimedean jet Gram with the Gram of
$\tau^ke^{-iz_0\tau}/k!$ in $L^2(0,\infty)$.  Orthonormalizing the polynomial factor against the
weight $e^{-2y\tau}d\tau$ gives the displayed Laguerre functions.  The shift representation
Theorem~\ref{thm:prime-shift} gives the matrix element of $S_{\log n}+S_{\log n}^*$ as the two lag
integrals above, with the phases $n^{\mp it_0}$ supplied by the carrier $e^{-it_0\tau}$.  The
Stirling remainder is holomorphic on a fixed polydisk in the interior band, so Cauchy's estimates make
the corresponding finite jet correction effective.
\end{proof}

\begin{corollary}[effective $\log\log$ contractivity range]\label{cor:loglog-effective}
Fix $y\ge\half+\epsilon_0$.  There are effective constants
$\beta=\beta(\epsilon_0,y)$, $c_0$, and $T_2$ such that, for $|t_0|\ge T_2$,
\[
   \lambda_{\max}(T_N(t_0-iy))
   \le
   \frac{c_0N^2e^{\beta N}}{\log(|t_0|/2\pi)}
\]
for all $N\le N_*(t_0)$, the range where the archimedean jet matrix is positive
(Lemma~\ref{lem:pole-pair}).  In particular $\lambda_{\max}(T_N)\le1$ is proved effectively for
\[
   N\le
   \frac{\log\log |t_0|-\log(c_0N^2)}{\beta}.
\]
\end{corollary}

\begin{proof}
By Theorem~\ref{thm:whitened-terminal},
\[
   \lambda_{\max}(T_N)\le\frac{\|J_N^{\mathrm{prime}}\|}{\lambda_{\min}(J_N^\infty)}.
\]
Lemma~\ref{lem:prime-window} bounds the prime jet entries by an exponential
$C(\epsilon_0)(2/\epsilon_0)^{2N}$, hence the matrix norm by $N^2$ times the same exponential.
Lemma~\ref{lem:arch-jet}, with the effective floor from Lemma~\ref{lem:pascal-effective}, gives
\[
   \lambda_{\min}(J_N^\infty)
   \ge
   4^{-N}(2y)^{-(2N-1)}\log(|t_0|/2\pi)-C_N'/|t_0|,
\]
with $C_N'$ effective by Cauchy estimates on the fixed polydisk.  For $|t_0|\ge T_2$ the remainder is
absorbed, yielding the stated exponential bound with, for example, a crude admissible
$\beta\le\log(64y^2/\epsilon_0^2)$.
\end{proof}

\begin{definition}[terminal defect observable]\label{def:terminal-defect}
For a fixed one-point terminal gauge $z_0=t_0-iy$ and a scalar jet section, define
\[
   \delta_N(z_0):=1-\lambda_{\max}(T_N(z_0)).
\]
The terminal input is $\delta_N(z_0)\ge0$ for every $N$.  Numerical evidence should therefore be read
through the sign and decay law of $\delta_N$, not through the existence of a fixed positive margin.  This
$J_N^\infty$-whitened defect is well defined only for $N\le N_*(t_0)$ (Lemma~\ref{lem:pole-pair}); the
reference-whitened observable of Definition~\ref{def:terminal-defect-fixed} is the all-$N$ form
equivalent to RH, and is the one the harness computes.
\end{definition}

\begin{lemma}[signed archimedean decomposition]\label{lem:signed-arch}
Let $h_A$ denote the archimedean part of
\[
   h(z)=-\frac{\Xi'(z)}{\Xi(z)}=h_A(z)+h_P(z),
\]
and set
\[
   w_A(u):=-\frac1\pi\Im h_A(u)
   =
   \frac1{2\pi}\log\frac{|u|}{2\pi}+O((1+u^2)^{-1})
   \qquad (u\in\R,\ |u|\to\infty).
\]
For $v_j(u)=(u-z_0)^{-(j+1)}$, let $G_N$ be the Lebesgue Gram matrix of
$v_0,\ldots,v_{N-1}$.  The archimedean jet matrix admits
\[
   A_N=\int_\R vv^*\,w_A(u)\,du+H_N,
\]
and the remainder satisfies the uniform form bound
\[
   \|G_N^{-1/2}H_NG_N^{-1/2}\|
   \le C(y)\frac{\log |t_0|}{|t_0|}
   \qquad(N\ge1)
\]
in a fixed interior band.
\end{lemma}

\begin{proof}[Proof sketch]
The cancellation of the elementary real parts on the critical line gives the displayed $w_A$ after
Stirling.  Subtracting the Cauchy transform of $w_A$ leaves a function with zero boundary imaginary
part on $\R$; by the Schwarz reflection principle and Morera's theorem, this continuation step (given
here as a sketch) extends the function analytically to a disk of radius comparable to
$|t_0|$ around the working point, with derivative $O(\log |t_0|/|t_0|)$.  Cauchy's estimates applied
to the two-variable divided difference kernel give entries bounded by
$M(t_0/4)^{-(j+k)}$, and Schur comparison with the exact Szeg\H{o}/Lebesgue Gram gives the displayed
uniform form estimate.

\emph{Herglotz normalization of the weight.}  We fix the normalization of $w_A$, which is the only
point previously flagged.  Write the archimedean part as
$h_A(z)=i\bigl[-(\tfrac1s+\tfrac1{s-1})+\tfrac12\log\pi-\tfrac12\psi(s/2)\bigr]$, $s=\half+iz$, with
$\psi=\Gamma'/\Gamma$.  The digamma function is Nevanlinna: $\psi(s)=-\gamma+\sum_{n\ge0}
(\tfrac1{n+1}-\tfrac1{n+s})$ has $\Im\psi(s)>0$ for $\Im s>0$.  In $w_A(u)=-\tfrac1\pi\Im h_A(u)$ the
term $-\tfrac12\psi(s/2)$ contributes $\tfrac1{2\pi}\Re\psi(\tfrac14+\tfrac{iu}{2})=\tfrac1{2\pi}
\log\tfrac{|u|}{2}+O(u^{-2})$ (Stirling), the constant $\tfrac12\log\pi$ contributes the shift
$-\tfrac1{2\pi}\log\pi$, and the two rational terms $-(\tfrac1s+\tfrac1{s-1})$ are anti-Nevanlinna and
contribute only $O(u^{-2})$; together these give the nonnegative boundary density
$\tfrac1{2\pi}\log\tfrac{|u|}{2\pi}$.  Hence
\[
   w_A(u)=\frac1{2\pi}\log\frac{|u|}{2\pi}+O(u^{-2})\ \ge\ 0\qquad(|u|\ge 2\pi),
\]
so the representing weight is a genuine nonnegative density on the support of every gauge with
$|t_0|\ge 2\pi+W$, and $A_N=\int vv^*w_A+H_N$ is a positive form there.  (Numerically, at $y=1$ the
identity holds to full working precision --- e.g.\ $w_A(300)=0.615278=\tfrac1{2\pi}\log\frac{300}{2\pi}$
--- and the archimedean jet $A_{12}$ is positive definite for $t_0=50,300,1000$; see
Section~\ref{sec:h8-harness}.)  No entrywise jet estimate is used.
\end{proof}

\begin{theorem}[leading sampling form of the terminal defect]\label{thm:defect-sampling}
Assume RH and fix $z_0=t_0-iy$ in an interior band.  Put
\[
   L=\log\frac{|t_0|}{2\pi}
\]
and let
\[
   \mathcal R_N(z_0)=
   \left\{
      R(u)=\sum_{j=0}^{N-1}\bar c_j(u-z_0)^{-(j+1)}:\ c\in\C^N
   \right\}.
\]
Then, in the scalar one-point terminal gauge,
\[
   \delta_N(z_0)
   =
   \min_{0\ne R\in\mathcal R_N(z_0)}
   \frac{\sum_\rho m_\rho |R(z_\rho)|^2}
        {\dfrac{L}{2\pi}\displaystyle\int_\R |R(u)|^2\,du}
   \left(1+O_y(L^{-1})+O_y(|t_0|^{-1})\right),
\]
where under RH the points $z_\rho$ are real, in the constant-weight regimes specified precisely in
Theorem~\ref{thm:defect-sampling-v2}.  Outside those regimes this display is a leading model formula;
the weighted statement of Theorem~\ref{thm:defect-sampling-v2} is the governing audited version.
\end{theorem}

\begin{proof}
Under RH, the scalar zero-side jet matrix is the Gram matrix
\[
   c^*J_Nc=\sum_\rho m_\rho |R_c(z_\rho)|^2,
   \qquad
   R_c(u)=\sum_{j=0}^{N-1}\bar c_j(u-z_0)^{-(j+1)}.
\]
The leading archimedean part is $L$ times the Szeg\H{o} jet Gram.  The latter is exactly the Lebesgue
Gram of the same rational functions:
\[
   \frac1{2\pi}\int_\R
   \frac{du}{(u-z_0)^{j+1}(u-\bar z_0)^{k+1}}
   =
   i(-1)^k\binom{j+k}{j}(2iy)^{-(j+k+1)}.
\]
This residue computation is the real-line form of Lemma~\ref{lem:szego-jet}.  Lemma~\ref{lem:arch-jet}
adds the uniform weighted-form Stirling remainder.  Replacing the exact weight $w_A$ by
$L/(2\pi)$ is justified in the ranges of Theorem~\ref{thm:defect-sampling-v2}; there it contributes
the indicated constant-weight relative error.  Finally,
Theorem~\ref{thm:whitened-terminal} identifies
$\delta_N=1-\lambda_{\max}(T_N)$ with the minimum of the quotient $c^*J_Nc/c^*A_Nc$, which yields the
displayed sampling form.
\end{proof}

\begin{theorem}[audited weighted sampling form]\label{thm:defect-sampling-v2}
Assume RH.  Let $\mathsf d_N^w$ be the minimum sampling quotient with denominator weighted by
$w_A$:
\[
   \mathsf d_N^w
   :=
   \min_{0\ne R\in\mathcal R_N(z_0)}
   \frac{\sum_\rho m_\rho |R(z_\rho)|^2}
        {\displaystyle\int_\R |R(u)|^2w_A(u)\,du}.
\]
Then, uniformly in $N$,
\[
   \delta_N=\mathsf d_N^w
   \left(1+O_y\!\left(\frac{\log |t_0|}{|t_0|}\right)\right).
\]
Replacing $w_A$ by the constant density $L/(2\pi)$ gives the constant-weight quotient
$\mathsf d_N$ in the following regimes:
\[
   \delta_N=\mathsf d_N\left(1+O\!\left(\frac N{|t_0|}\right)+O(e^{-N})\right)
   \qquad (N\ge\tfrac32L),
\]
and
\[
   \delta_N=\mathsf d_N\left(1+O\!\left(\frac{(16y^2)^N}{|t_0|}\right)\right)
   \qquad
   \left(N\le(1-\varepsilon)\frac{L}{\log(16y^2)}\right).
\]
The intermediate window between these two ranges is left as a declared audit gap for the
constant-weight replacement, not for the weighted sampling identity.
\end{theorem}

\begin{proof}
The first assertion is Theorem~\ref{thm:defect-sampling} with the denominator kept in its exact
weighted form and Lemma~\ref{lem:signed-arch} used as a form estimate.  The constant-weight replacement
requires concentration of near-minimizers in the window where
$w_A(u)=L/(2\pi)+O(N/t_0)$; for large $N$ this follows from the lower sampling contribution of distant
zeros together with the factorial upper bound supplied by the explicit interpolant.  In the small-$N$
range the Pascal--Laguerre lower floor controls the error after whitening.  The stated middle range is
kept out of the theorem because the present argument does not give a uniform comparison there.
\end{proof}

\begin{theorem}[conditional saturation of the terminal operator]\label{thm:saturation}
Assume RH and fix $z_0=t_0-iy\in U_-$.  In the scalar one-point jet gauge,
\[
   \limsup_{N\to\infty}\lambda_{\max}(T_N(z_0))=1
\]
along the terminal sections.  More concretely, if $R_c(u)=\sum_{j=0}^{N-1}\bar c_j(u-z_0)^{-(j+1)}$
is chosen to vanish at the $N-1$ zero ordinates nearest $t_0$, then
\[
   0<1-\lambda_{\max}(T_N)
   \le
   \frac{\sum_{\rho\ {\rm far}}m_\rho |R_c(z_\rho)|^2}
        {c^*A_Nc}
   \qquad(N\le N_*(t_0)),
\]
where $A_N=J_N^\infty$ is the archimedean jet form (positive definite in this range,
Lemma~\ref{lem:pole-pair}); the saturation $\limsup=1$ is read for all $N$ through the reference-whitened
defect of Definition~\ref{def:terminal-defect-fixed}.  The right side is the zero-tail energy after interpolation
of the nearest atoms; the Pascal--Laguerre floor of Lemma~\ref{lem:pascal-effective} gives the
effective denominator needed to make this tail estimate quantitative.
\end{theorem}

\begin{proof}
Under RH, the scalar Pick jet matrix is the Gram matrix of the functions
$(u-z_0)^{-(j+1)}$ against the positive zero measure
\[
   d\Sigma_\Xi(u)=\sum_\rho m_\rho\,\delta_{z_\rho}(u),
\]
while
$J_N=A_N-J_N^{\mathrm{prime}}$.  For any nonzero $c$,
\[
   1-\lambda_{\max}(T_N)
   \le
   \frac{c^*J_Nc}{c^*A_Nc}
   =
   \frac{\sum_\rho m_\rho |R_c(z_\rho)|^2}{c^*A_Nc}.
\]
Choose $c$ so that the numerator vanishes at the $N-1$ zero ordinates closest to $t_0$; this is
ordinary rational interpolation in the basis $(u-z_0)^{-(j+1)}$.  Only the far zero tail remains,
which gives the displayed inequality.  Positivity under RH gives
$\lambda_{\max}(T_N)\le1$.  Since the zero measure has infinite support and the interpolation degree
increases, the far-tail quotient can be made arbitrarily small (Proposition~\ref{thm:sat-quant} and
Theorem~\ref{thm:two-sided}, closed modulo the Stirling/RvM bookkeeping flagged there), so the largest
eigenvalues are arbitrarily close to $1$; equivalently, the terminal boundary has no fixed positive
margin.
\end{proof}

\begin{corollary}[no fixed margin at the terminal boundary]\label{cor:no-margin}
The statement
\[
   \lambda_{\max}(T_N)\le1-\delta\qquad(N\ge1)
\]
is false for every fixed $\delta>0$ (equivalently, the reference-whitened defect $\delta_N$ of
Definition~\ref{def:terminal-defect-fixed} has no fixed positive lower bound).  Under RH it contradicts
Theorem~\ref{thm:saturation}; under $\neg$RH, Theorem~\ref{thm:minimal-input} implies that
$\lambda_{\max}(T_N)>1$ for some finite section in every terminal gauge detecting the off-line pole.
Thus every proof strategy for the terminal input that loses a fixed multiplicative factor at leading
order is structurally incapable of proving RH through this route.
\end{corollary}

\begin{proposition}[Blaschke detection rate for an off-line zero]\label{prop:blaschke-rate}
Suppose $\Xi$ has a non-real pair $\{z_w,\bar z_w\}$ with $\Im z_w<0$.  For the terminal base point
$z_0=t_0-iy$, set the exterior Blaschke factor
\[
   b(\zeta)=\frac{\zeta-\bar z_0}{\zeta-z_0}.
\]
Then
\[
   \liminf_{N\to\infty}\lambda_{\max}(T_N)^{1/N}\ge |b(z_w)|>1.
\]
If this pair is the unique dominant off-line pair in Blaschke distance and the cost of avoiding the
nearby real zeros is subexponential, then
\[
   \lim_{N\to\infty}\lambda_{\max}(T_N)^{1/N}=|b(z_w)|.
\]
For several dominant pairs the right side is replaced by $\max_w |b(z_w)|$.
\end{proposition}

\begin{proof}[Proof sketch]
The pair contribution
\[
   g_{\rm pair}(z)=\frac1{z_w-z}+\frac1{\bar z_w-z}
\]
has real boundary values on $\R$, and its Pick kernel decomposes into the rank-two hyperbolic block
\[
   e_{z_w}\otimes\overline{e_{\bar z_w}}
   +
   e_{\bar z_w}\otimes\overline{e_{z_w}},
\]
so the form contains $2\Re[X\overline Y]$ with $X$ and $Y$ the evaluations of the rational test at
$z_w$ and $\bar z_w$.  Evaluation at $\bar z_w\in\C^+$ is bounded in the model space, while evaluation
at the exterior point $z_w\in\C^-$ has reproducing-kernel norm growing like
\[
   \frac{|b(z_w)|^{2N}-1}{4\pi|\Im z_w|}.
\]
This gives the lower exponential rate.  The matching upper rate follows when no other off-line pair has
larger Blaschke modulus and when the finite Blaschke factors used to avoid nearby real zeros carry only
subexponential cost.  In the Davenport--Heilbronn test, $|b(z_w)|=2.25545$ in exact coordinates,
matching the observed $2.2556$.
\end{proof}

\begin{corollary}[finite detection scale]\label{cor:blaschke-detection}
An off-line zero at Blaschke radius $|b(z_w)|=e^\beta>1$ is detected in terminal sections on the scale
$N\asymp \beta^{-1}\log C$, where $C$ is the finite threshold separating $\lambda_{\max}(T_N)$ from
$1$.  Two distinct terminal base points give two Blaschke level curves and hence, in principle,
triangulate the off-line zero.
\end{corollary}

\begin{conjecture}[terminal defect law]\label{conj:defect-law}
Fix $z_0=t_0-iy$ in an interior band and set
\[
   L=\log\frac{|t_0|}{2\pi}.
\]
For the scalar one-point terminal gauge,
\[
   \delta_N(z_0)=1-\lambda_{\max}(T_N(z_0))
   =
   \left(\frac{\kappa_{\rm sin}(y)L}{N}\right)^{2N(1+o(1))}.
\]
The symmetric sine-model numerical data at $y=1$ currently give the estimate
\[
   \kappa_{\rm sin}(1)=0.686\pm0.004.
\]
Asymmetric offsets display discrete left/right assignment transitions and should not be extrapolated by
the same smooth tail fit.  This is an empirical conjecture, not an input to the proof path.
\end{conjecture}

\begin{definition}[lattice defect constant]\label{def:kappa-lattice}
Let
\[
   \Gamma_L=\left\{t_0+\frac{2\pi k}{L}:k\in\Z\right\}
\]
be the exact lattice of density $L/(2\pi)$.  Define the model constant
\[
   \kappa_{\mathrm{lat}}(y):=
   \lim_{N/L\to\infty}
   \frac NL
   \left[
      \min_{0\ne R\in\mathcal R_N(t_0-iy)}
      \frac{\sum_{\gamma\in\Gamma_L}|R(\gamma)|^2}
           {\dfrac L{2\pi}\|R\|_{L^2(\R)}^2}
   \right]^{1/(2N)},
\]
when the limit exists.  This is a Christoffel-type extremal problem for a lattice sampled against
Lebesgue measure in a rational Hardy model space.  It is independent of RH and is the clean model
problem for the constant in Conjecture~\ref{conj:defect-law}.
\end{definition}

\begin{proposition}[sine-model form of the lattice constant]\label{prop:sine-model}
The lattice defect of Definition~\ref{def:kappa-lattice} is the terminal defect of the sine model.  Let
\[
   \Xi_{\rm lat}(z)=\sin\!\left(\frac{L(z-u_0)}2\right),
   \qquad
   M_{\rm lat}(z)=-\frac{\Xi_{\rm lat}'(z)}{\Xi_{\rm lat}(z)}
   =-\frac L2\cot\!\left(\frac{L(z-u_0)}2\right).
\]
For $\Im z<0$,
\[
   M_{\rm lat}(z)
   =
   -\frac{iL}{2}
   -iL\sum_{m\ge1}e^{-imL(z-u_0)}.
\]
The constant term produces the leading archimedean/Szeg\H{o} form $L\,S$, while the oscillatory part
is a one-frequency shift model with lags $mL$.  Thus $\kappa_{\rm lat}(y)$ is a constant of a classical
approximation problem for $\sin$, with no arithmetic input.
\end{proposition}

\begin{proof}
The zeros of $\Xi_{\rm lat}$ are precisely the lattice
$u_0+2\pi k/L$.  The logarithmic derivative gives the displayed cotangent expression.  Since
$\Im z<0$, the standard expansion
\[
   \cot w=i+2i\sum_{m\ge1}e^{-2imw}
   \qquad(\Im w<0)
\]
with $w=L(z-u_0)/2$ gives the series for $M_{\rm lat}$.  The constant $-iL/2$ has one-node Pick
quotient $L/(2y)$ at $z=-iy$, matching the leading $L$ times Szeg\H{o} kernel.  The remaining terms
are exactly shifts at the lags $mL$, as in Theorem~\ref{thm:prime-shift}.  Running the same whitening
pipeline therefore computes the lattice defect.
\end{proof}

\begin{definition}[sine concentration constant]\label{def:kappa-sine}
Let $\delta_N^{\rm lat}(y,L,u_0)$ be the terminal defect of the sine model of
Proposition~\ref{prop:sine-model}.  When the limit exists, define
\[
   \kappa_{\rm sin}(y)
   :=
   \lim_{N\to\infty}
   \frac NL
   \left(\delta_N^{\rm lat}(y,L,u_0)\right)^{1/(2N)}.
\]
The current numerical evidence supports the factorial law and the smooth symmetric-offset extrapolation
\[
   \kappa_{\rm sin}(1)=0.686\pm0.004.
\]
Asymmetric offsets may show discrete jumps in finite $N$ caused by left/right assignment changes of the
extremal zeros; these are recorded as finite-size transitions rather than a change in the limiting
constant.
Analytically, this is a constrained-equilibrium/prolate concentration problem: minimize
lattice-sampled energy against continuous $L^2$ energy over the rational model space
$\mathcal R_N\subset H^2(\C^+)$.
\end{definition}

\begin{problem}[constrained equilibrium form of $\kappa_{\rm sin}$]\label{prob:kappa-final}
Determine the constant $\kappa_{\rm sin}(y)$ by the following scaled constrained-equilibrium problem.
Among measures $\sigma\ge0$ with
\[
   \sigma(\R)=1,\qquad \sigma\le\tfrac12\,{\rm Leb},
\]
and obstacle condition
\[
   U^\sigma(t)\le\log|t|\qquad (t\in\R\setminus S(\sigma)),
\]
maximize the center field $U^\sigma(0)$.  If
\[
   G^*:=\sup_\sigma U^\sigma(0),
\]
the predicted relation is
\[
   \kappa_{\rm sin}(y)=\frac y\pi e^{-G^*}.
\]
The uniform interpolant gives $G=-1$ and hence the rigorous candidate upper scale
$\kappa_{\rm sin}(y)\le ey/\pi$.  Numerics at $y=1$ ask for
$G^*=-0.768\pm0.007$.  The expected extremal has a saturated core and free boundary/contact regions,
as in Dragnev--Saff constrained equilibrium for discrete orthogonal polynomial asymptotics.
\end{problem}

\begin{proposition}[quantitative saturation upper-bound reduction]\label{thm:sat-quant}
Assume RH.  Let $R_N$ be the rational interpolant in Theorem~\ref{thm:saturation}, vanishing at the
$N-1$ zero ordinates nearest $t_0$.  If the far-zero product estimate
\[
   \sum_{\rho\ {\rm far}}m_\rho |R_N(z_\rho)|^2
   \le
   c^*A_Nc\left(\frac{C(y)L}{N}\right)^{2N}
\]
holds with $L=\log(|t_0|/2\pi)$, then
\[
   \delta_N(z_0)\le\left(\frac{C(y)L}{N}\right)^{2N}.
\]
The estimate is the exact analytic form of the remaining quantitative saturation calculation: it asks
for the Riemann--von Mangoldt zero density, viewed through the rational interpolator and the
Pascal--Laguerre floor, to supply a factorial tail bound.
\end{proposition}

\begin{proof}
This is Theorem~\ref{thm:saturation} with the displayed product estimate substituted into the
zero-tail quotient.  All normalization loss is carried by the denominator $c^*A_Nc$, whose effective
lower control is Lemma~\ref{lem:pascal-effective}.  Thus the problem of proving the factorial upper
bound is reduced to the stated far-zero product estimate.
\end{proof}

\begin{theorem}[two-sided factorial saturation bounds]\label{thm:two-sided}
Assume RH.  For $L\le N\le |t_0|^{1/4}$,
\[
   \left(\frac{c_1(y)L}{N}\right)^{2N}
   \le
   \delta_N(z_0)
   \le
   \left(
      \frac{eyL}{\pi N}
      \left(1+O\!\left(\frac1{\log\log |t_0|}\right)\right)
   \right)^{2N},
\]
where $c_1(y)>0$ is effective.  The upper bound is closed modulo effective Stirling/Riemann--von
Mangoldt bookkeeping; the lower bound uses a Turan--Remez separation estimate and is retained as the
principal audit point for this theorem.
\end{theorem}

\begin{proof}[Proof sketch]
For the upper bound, choose the rational interpolant vanishing at the $N-1$ zero ordinates nearest
$t_0$.  On the scale $W=\pi N/L$, the logarithmic potential of the uniform interpolation pattern has
far-field maximum at infinity, giving the factor $(eyL/\pi N)^{2N}$; the Riemann--von Mangoldt
fluctuation in the window contributes the stated $O(1/\log\log |t_0|)$ under RH.  For the lower bound,
among the $4N$ nearest zeros, a Turan--Remez selection gives at least $N$ points whose product distance
against any degree-$N$ numerator is bounded below by $(c\,{\rm gap})^N$.  Combining this selection with
the pole field at depth $y$ yields the lower factorial scale.  The selection-to-constant step is the
remaining place where a full proof must record the effective $c_1(y)$.
\end{proof}

\begin{remark}[why the factorial scale matters]\label{rmk:factorial-wall}
Conjecture~\ref{conj:defect-law} sharpens Corollary~\ref{cor:no-margin}: the terminal positivity is
not merely close to the boundary, but close at a factorial scale.  Therefore any argument losing a
factor of size $\exp(o(N\log N))$ at the terminal boundary is still too coarse for the final input.
The only known mechanisms compatible with this scale are exact Gram identities or equivalently exact
positivity identities for the arithmetic data.
\end{remark}

\begin{proposition}[root recovery from the extremal vector]\label{prop:root-recovery}
Assume RH and let $c_N$ be a near-extremal vector for $\delta_N$.  In the main interpolation window,
the numerator of
\[
   R_{c_N}(u)=\sum_{j=0}^{N-1}\overline{(c_N)_j}(u-z_0)^{-(j+1)}
\]
has zeros exponentially close to the zero ordinates that carry the dominant sampling constraints.  More
precisely, the distance in each stable gap is governed by the local potential drop between the center
field and that gap.  The statement is quantitative after a Rouche estimate on each gap.
\end{proposition}

\begin{proof}[Proof sketch]
The quotient in Theorem~\ref{thm:defect-sampling-v2} forces
$|R_{c_N}(\gamma)|^2$ to be small at every sampled zero ordinate in the active window.  The potential
profile gives a lower bound on the maximum of $R_{c_N}$ in each gap, while the small endpoint values
give a Rouche comparison with the local linear factor.  The numerical root-forensics block matches the
predicted gradient of accuracy: central zeros are recovered to more digits than edge zeros.
\end{proof}

\begin{remark}[asymmetric lattice offsets]\label{rmk:asymmetric-offsets}
The lattice offsets $0$ and half-gap are symmetric relative to the projection of $z_0$ and give smooth
tail extrapolations.  Asymmetric offsets can undergo discrete left/right assignment transitions in the
extremal zero pattern.  Such transitions alter the finite-$N$ prefactor by
\[
   \Delta\kappa_N\sim \frac{\Delta\log A}{2N},
\]
which explains observed finite jumps such as the repaired $a/4$, $N=56$ value without changing the
expected limiting sine constant.
\end{remark}

\begin{theorem}[terminal wall as Weil positivity]\label{thm:wall-weil}
For the Laguerre unilateral test class
\[
   \widehat g(\tau)\sim e^{it_0\tau}q(\tau)e^{-y\tau}\mathbf 1_{\tau>0}
\]
convolved with its reflected partner, the terminal quadratic form satisfies
\[
   c^*J_Nc=W(|R_c|^2),
\]
where $W$ is the Weil explicit-formula functional.  Consequently, after completion of this Laguerre
test class in the Weil topology,
\[
   \mathsf d_N(z_0)\ge0\ \text{for every }N
   \quad\Longleftrightarrow\quad
   \text{Weil positivity on this class}
   \quad\Longleftrightarrow\quad
   \mathrm{RH}.
\]
\end{theorem}

\begin{proof}[Proof sketch]
The equality of forms is the explicit formula applied to the test $|R_c|^2$ in the unilateral
Laguerre representation: the zero side is $c^*J_Nc$, while the archimedean and prime sides are exactly
the terminal decomposition used above.  Completeness of Laguerre functions gives density in the
unilateral window class.  The remaining analytic point is continuity of the Weil functional in this
completed topology, of the same type as the standard Bombieri--Weil extension.  This is
Lemma~\ref{lem:weil-continuity} below; with it recorded, the equivalence with RH is the classical Weil
criterion.
\end{proof}

\begin{lemma}[continuity of the Weil functional on the completed Laguerre class]\label{lem:weil-continuity}
Equip the test class with the norm
\[
   \|g\|:=\sup_u(1+u^2)|g(u)|+\sup_\tau e^{(\half+\epsilon_0/2)|\tau|}|\widehat g(\tau)|.
\]
Then the three sides of $W$ are $\|\cdot\|$-continuous: the zero side by
$\sum_\rho(1+|\Re z_\rho|^2)^{-1}<\infty$ and $|\Im z_\rho|<\half$; the prime side dominated by
$\|g\|\sum_n\Lambda(n)n^{-1-\epsilon_0/2}$; the archimedean side by the weight
$(1+u^2)^{-1}\log(2+|u|)$.  Laguerre tests at depth $y\ge\half+\epsilon_0$ satisfy
$|\widehat g(\tau)|\lesssim e^{-y|\tau|}$ and are dense in the completion.  This closes the completion
step of Theorem~\ref{thm:wall-weil}.
\end{lemma}



\subsubsection{The interpolation certificate tower}

\begin{proposition}[Nevanlinna--Pick certificate tower]\label{prop:np-tower}
Assume ARP-P, fix a finite channel $F$, and fix the node sequence used in
Theorem~\ref{thm:pick-rh}.  For every $N$ there exists a matrix Pick-class interpolant $F_N$ satisfying
$F_N(z_j)=g^F(z_j)$ for $j\le N$.  Any admissible choice of such interpolants has the following
properties:
\begin{enumerate}[label=\textnormal{(\roman*)}]
\item $\{F_N\}$ has pinned-normal subsequential limits on $\Omega_-$ in the sense of
Lemma~\ref{lem:pinned-normality};
\item $F_N\to g^F$ locally uniformly on $U_-$;
\item the analytic bridge applies to the limiting target exactly as in Theorem~\ref{thm:pick-rh}.
\end{enumerate}
Conversely, the existence of any tower of Cauchy transforms of positive measures converging to $g^F$ on
$U_-$ implies ARP-P.
\end{proposition}

\begin{proof}
Existence for each $N$ is the matrix Nevanlinna--Pick theorem applied to Definition~\ref{def:ARP-P}.
Pinned normality and the identification of every subsequential limit with $g^F$ on $U_-$ are proved by
Lemma~\ref{lem:pinned-normality} as used in Theorem~\ref{thm:pick-rh}; since all subsequential limits
agree, the full sequence converges locally uniformly on $U_-$.  The bridge conclusion is the last
paragraph of that theorem.  The converse is Lemma~\ref{lem:cauchy-pick} plus closure of the positive
semidefinite cone under pointwise limits at finite node sets.
\end{proof}

\begin{remark}[role of the von Mangoldt tower]\label{rmk:vm-tower-role}
The von Mangoldt tower is reclassified as a candidate certificate for ARP-P.  It is valuable because its
finite spectral measures are positive, ultimately reflecting $\Lambda\ge0$.  It is not the minimal
essential object: any proof of tower-ARP proves ARP-P by Theorem~\ref{thm:tower-implies-pick}, while
any proof of ARP-P supplies Pick-class interpolation certificates by Proposition~\ref{prop:np-tower}.
\end{remark}

\begin{proposition}[direct-sum tower exposes the placement problem]\label{prop:placement-problem}
In a fully decoupled direct-sum realization with one-dimensional free cells,
\[
   M_P^F(z)_{\alpha\beta}
   =
   \sum_{v\in\mathcal V_P}
   \frac{c_v\,\widehat\phi_\beta(\tau_v)\overline{\widehat\phi_\alpha(\tau_v)}}{\lambda_v-z}
   +M_{\infty,P}^F(z)_{\alpha\beta},
   \qquad
   \tau_v=k\log p,\quad c_v=(\log p)p^{-k/2}.
\]
Therefore tower-ARP for such a decoupled model is a spectral placement problem for the real numbers
$\lambda_v$ and the archimedean component.  Without an independent arithmetic rule forcing the weighted
empirical spectral measures to converge to the zero measure detected by $g^F$, the direct-sum tower does
not close ARP.
\end{proposition}

\begin{proof}
For a one-dimensional self-adjoint free cell, $A_v=\lambda_v I$ with $\lambda_v\in\R$, and the cell
resolvent is $(\lambda_v-z)^{-1}$.  Compressing against the source vector
$\sqrt{c_v}\widehat\phi_\alpha(\tau_v)e_v$ gives the displayed contribution.  Summing over the finite
cutoff and adding the archimedean compression gives the formula.  The limit is therefore controlled by
the weak convergence of the corresponding positive spectral measures.  Choosing the $\lambda_v$ to make
that limit equal to $g^F$ is precisely the missing arithmetic placement mechanism.
\end{proof}

\subsubsection{Target inequality and correlation form}

\begin{proposition}[Weil splitting of ARP-P]\label{prop:target-ineq}
Theorem~\ref{thm:explicit-channel} decomposes $g^F$ on $U_-$ as a Weil explicit-formula functional:
\[
   g^F=g_\infty^F-g_{\mathrm{prime}}^F,
\]
where $g_\infty^F$ denotes the archimedean and pole/degree contribution and
$g_{\mathrm{prime}}^F$ denotes the von-Mangoldt prime contribution, both interpreted in the same
regularized Weil sense as Theorem~\ref{thm:explicit-channel}.  Since the Pick matrix is linear in the
function, ARP-P is equivalent to the finite-node dominance
\[
   \mathsf P[g_\infty^F;z_1,\ldots,z_N]
   \succeq
   \mathsf P[g_{\mathrm{prime}}^F;z_1,\ldots,z_N]
\]
for every finite source plane $F$ and every finite node set in $U_-$.
\end{proposition}

\begin{proof}
Apply the linear map
\[
   h\longmapsto
   \left[
      \frac{h(z_i)-h(z_j)^*}{z_i-\bar z_j}
   \right]_{i,j=1}^{N}
\]
to the Weil splitting of Theorem~\ref{thm:explicit-channel}.  The inequality is exactly the assertion
that the Pick matrix of $g^F=g_\infty^F-g_{\mathrm{prime}}^F$ is positive semidefinite.
\end{proof}

\begin{remark}[Bochner correlation viewpoint]\label{rmk:bochner}
For $z\in U_-$,
\[
   \frac1{z_\rho-z}=-i\int_0^\infty e^{-izt}e^{iz_\rho t}\,dt.
\]
Thus the zero-side target can be written formally as
\[
   g^F(z)=-i\int_0^\infty e^{-izt}\psi^F(t)\,dt,
   \qquad
   \psi^F(t)_{\alpha\beta}
   =
   \sum_\rho
   m_\rho\,\ell_{z_\rho}(\phi_\beta)
   \overline{\ell_{\bar z_\rho}(\phi_\alpha)}
   e^{iz_\rho t}.
\]
Through the explicit formula, $\psi^F$ is represented by arithmetic distributions.  ARP-P is the
finite-node Pick form of the same positivity problem; in Fourier language it asks for the corresponding
extended correlation distribution to be positive definite.  This is a global condition, not a
prime-cell-by-prime-cell assertion.
\end{remark}

\begin{remark}[Davenport--Heilbronn falsifier]\label{rmk:dh-pick}
The analogue of ARP-P must fail for Davenport--Heilbronn data.  Indeed, the Davenport--Heilbronn function
has off-critical zeros; if its arithmetic channel data satisfied ARP-P, the proof of
Theorem~\ref{thm:pick-rh} would force those zeros onto the critical line.  Therefore some finite Pick
matrix at nodes in the absolute-convergence region must have a negative eigenvalue.  This gives a
finite-dimensional numerical falsifier for the criterion.
\end{remark}




\subsection{Step 5 substructure (continued): H8 falsification harness and the $\Omega$-chain to $\Omega_7$}\label{sec:h8-harness}

This appendix records a self-contained, independent reimplementation of the terminal observable,
built for a single purpose: to make every candidate step for the open input Step~5 (H8) cheap to
falsify before any analytic effort is spent on it.  The code is provided with the paper as
\texttt{h8lab.py} (kernels), \texttt{h8\_validate.py} (three-law validation), and
\texttt{h8\_forensic.py} (inverse forensics), together with a cache \texttt{zeros\_100\_175.json}.
It is numerical evidence, not proof; the governing statements remain
Theorems~\ref{thm:whitened-terminal}, \ref{thm:prime-shift}, \ref{thm:saturation}.

\subsubsection*{Two independent coordinates}

The harness computes the terminal defect in two ways that use disjoint inputs, so that agreement is a
genuine cross-check of the explicit formula rather than a tautology.

\begin{enumerate}[label=\textnormal{(\alph*)}]
\item \emph{Sampling coordinate} (the audited governing form,
Theorem~\ref{thm:defect-sampling}): $\widetilde\delta_N=\delta_N/L=\min_c\ c^*J_{\rm zero}c\,/\,(L\,c^*Gc)$
--- the rescaled variant of Definition~\ref{def:terminal-defect-fixed}, of the same sign as $\delta_N$
at these gauges since $L>0$ --- where
$G$ is the exact Szeg\H{o}/Lebesgue reference Gram
$G_{jk}=i(-1)^k\binom{j+k}{j}(2iy)^{-(j+k+1)}$, and $J_{\rm zero}$ is assembled from rank-one zero
atoms $v(\gamma)v(\gamma)^*$, $v_j(u)=(u-z_0)^{-(j+1)}$, plus off-line pair blocks.  This uses only the
\emph{zero} configuration.
\item \emph{Arithmetic coordinate} (the object a certificate must factor): the Pick jets
$A_N$ (archimedean, closed form via $\psi^{(p)}$) and $\mathcal P_\Lambda$ (prime), assembled from the
exact log-derivative $-\zeta'/\zeta$ through the coefficient recursion
$\delta\,J_{jk}+J_{j-1,k}-J_{j,k-1}=h_j\mathbf 1_{k=0}-\overline{h_k}\mathbf 1_{j=0}$,
$\delta=z_0-\bar z_0$.  Crucially this uses \emph{no zeros}: the prime part is obtained from
$\zeta^{(k)}(s_0)$ by series division, avoiding the impossible $n^{-(1/2+y)}$ Dirichlet truncation.
The terminal input is $J_{\rm arith}:=A_N-\mathcal P_\Lambda\succeq0$.
\end{enumerate}

By the explicit formula $J_{\rm arith}=J_{\rm zero}$; the harness confirms this at jet level to the
finite-zero truncation error.

\subsubsection*{Three-law validation}

Run at $y=1$ on the gauge $z_0=t_0-i$ with $t_0$ centered on the $\zeta$-zero window
$\#100$--$175$ ($t_0\approx297.238$, $L\approx3.88$):

\begin{longtable}{@{}p{0.30\linewidth}p{0.62\linewidth}@{}}
\toprule
Law & Measured behaviour of the harness \\
\midrule
\endhead
$\zeta$: arithmetic PSD & $\delta_N(\text{arith})>0$ with factorial decay
$1.6\!\cdot\!10^{-2},\,8.0\!\cdot\!10^{-5},\,7.7\!\cdot\!10^{-8},\,1.8\!\cdot\!10^{-11}$
($N=4,6,8,10$); $J_{\rm arith}$ matches the $76$-zero Gram to the truncation tail. \\
DH: falsifier fires & Sampling $\delta_N$ crosses negative for the known off-line pair
$v=0.308517$; the negative branch ratios converge to the Blaschke rate
$2.673\to2.256\to\mathbf{2.2558}$, matching Proposition~\ref{prop:blaschke-rate}
($|b(z_w)|=2.2554$). \\
Sine/lattice: trivial & $\delta_N>0$ with factorial decay and empty cascade, as in
Remark~\ref{rmk:cascade-validation}. \\
\bottomrule
\end{longtable}

\subsubsection*{Inverse-forensic certificate}

Let $c$ be the generalized minimizer of the arithmetic pencil $(J_{\rm arith},A_N)$ realizing
$\delta_N$, and let $R_c(u)=\sum_j\overline{c_j}(u-z_0)^{-(j+1)}$ be the associated rational function.
The following table lists, for the roots of the numerator of $R_c$ lying in the active window
$|u-t_0|\le W=\pi N/L$, the maximal distance off the real axis and the maximal distance to the actual
$\zeta$-zero set.

\begin{longtable}{@{}p{0.06\linewidth}p{0.14\linewidth}p{0.18\linewidth}p{0.26\linewidth}@{}}
\toprule
$N$ & $W=\pi N/L$ & $\max|\Im\,\text{root}|$ & $\max$ recovery error in $W$ \\
\midrule
\endhead
$6$ & $4.86$ & $8.3\cdot10^{-59}$ & $0.43$ \\
$8$ & $6.48$ & $3.4\cdot10^{-55}$ & $0.64$ \\
$10$ & $8.10$ & $5.1\cdot10^{-53}$ & $0.23$ \\
$12$ & $9.71$ & $4.4\cdot10^{-49}$ & $0.45$ \\
$14$ & $11.3$ & $1.1\cdot10^{-44}$ & $0.65$ \\
\bottomrule
\end{longtable}

Two facts are visible.  First, the minimizer roots are \emph{real to $10^{-44}$ or better}, computed
from the arithmetic side alone: a per-gauge numerical certificate that the arithmetic terminal operator
places no mass off the critical line.  Second, the innermost roots coincide with actual $\zeta$-zeros to
full working precision (for $N=12$ the four central roots reproduce
$294.96537,\,295.57325,\,297.97928,\,299.84033$ exactly to eight figures); the larger window-error entries
are the expected edge roots that interpolate between consecutive zeros.  This is
Proposition~\ref{prop:root-recovery} made concrete: the arithmetic data, with no zero input, reconstructs
the Riemann zeros nearest the gauge and puts them on $\R$.

\begin{remark}[status of the certificate]\label{rmk:harness-scope}
The forensic certificate is unconditional numerical evidence at each finite $N$ and each gauge, and it
falsifies instantly for Davenport--Heilbronn.  It is the falsification instrument that any proposed exact Gram factorization
$A_N-\mathcal P_\Lambda=B^*B$ must pass on $\zeta$ and fail on DH, and the harness decides this in hours.
\end{remark}

\begin{theorem}[scalar level-one positivity]\label{thm:level-one-scalar}
For the scalar coordinate
\[
   M_\Xi(z)=i\left(-\frac{\xi'}{\xi}(s)\right),\qquad s=\half+iz,
\]
the one-node Pick inequality on $U_-$ is equivalent to
\[
   \Re\frac{\xi'}{\xi}(s)\ge0,\qquad \Re s>1.
\]
This inequality is classical.
\end{theorem}

\begin{proof}
Let $z=x+iy\in U_-$.  Then $\Re s=\half-y>1$, so $y=-(\Re s-\half)$.  The diagonal Pick quotient is
\[
   \frac{M_\Xi(z)-M_\Xi(z)^*}{z-\bar z}
   =
   \frac{\Im M_\Xi(z)}{\Im z}
   =
   \frac{-\Re(\xi'/\xi)(s)}{-(\Re s-\half)}
   =
   \frac{\Re(\xi'/\xi)(s)}{\Re s-\half}.
\]
Since $\Re s-\half>0$, positivity is equivalent to the displayed inequality.

By the Hadamard product for $\xi$, with the constant term fixed as in
Davenport~\cite[Sec.~12]{DavenportMNT},
\[
   \frac{\xi'}{\xi}(s)
   =
   \sum_\rho\left(\frac1{s-\rho}+\frac1\rho\right)+B,
\]
with the usual real normalization in which the real parts of the constant terms cancel after pairing
zeros.  Taking real parts gives
\[
   \Re\frac{\xi'}{\xi}(s)
   =
   \sum_\rho
   \frac{\sigma-\beta}{|s-\rho|^2},
   \qquad s=\sigma+it,
\]
in the standard symmetric limiting sense.  The classical zero-free theorem for $\zeta$ on
$\Re s=1$ and the functional equation give $0<\beta<1$ for nontrivial zeros.  Hence for $\sigma>1$ every
summand has nonnegative real part, and the inequality follows.
\end{proof}

\begin{remark}[scope of the level-one theorem]\label{rmk:level-one-scope}
Theorem~\ref{thm:level-one-scalar} closes the scalar one-node seed.  It is not the full matrix-channel
ARP-P(1), because the channel target contains source weights
$\ell_{z_\rho}(\phi_\beta)\overline{\ell_{\bar z_\rho}(\phi_\alpha)}$.  The full channel hierarchy is the
remaining Pick-positivity problem.
\end{remark}

\subsubsection*{H8: the reduced contractivity, a candidate route, and its gap}

We record the strongest current attempt on the terminal input $\delta_N\ge0$, in the exact form used at the
workbench, so that the point where it stalls is explicit and independently attackable.

\emph{The exact reduction.}  Pass to the additive model $L^2(0,\infty)$ of Theorem~\ref{thm:prime-shift}.
A rational $R\in\mathcal R_N$ corresponds to $F(\tau)=\sum_jc_j\tau^je^{-iz_0\tau}$, and the scalar
terminal form is
\[
   c^*(A_N-\mathcal P_\Lambda)c
   =\langle(\mathcal W_A-\mathcal T_\Lambda)F,F\rangle,
   \qquad
   \mathcal T_\Lambda=\sum_{n\ge2}\frac{\Lambda(n)}{\sqrt n}\,(S_{\log n}+S_{\log n}^*),
\]
with $\mathcal W_A$ the (positive) archimedean weight operator of Lemma~\ref{lem:signed-arch}.  Thus
\[
   \boxed{\ \delta_N\ge0\ \forall N\ \Longleftrightarrow\
   \lambda_{\max}\bigl(\mathcal W_A^{-1/2}\mathcal T_\Lambda\mathcal W_A^{-1/2}\bigr)\le1\ \text{on the Laguerre tower}\ }
\]
i.e.\ the whitened prime shift operator $\mathcal T:=\mathcal W_A^{-1/2}\mathcal T_\Lambda
\mathcal W_A^{-1/2}$ is dominated from above by the identity.  This is
Theorem~\ref{thm:whitened-terminal} in operator form.  \emph{Warning (measured):} the domination is
strictly one-sided --- the \emph{two-sided} norm genuinely exceeds $1$ at interior gauges (at
$t_0=30$, $y=1$ the bottom eigenvalue of the whitened jet compression is $-3.34,-7.62,-11.90$ at
$N=4,8,12$, driver \texttt{omega7\_operator\_interference.py}); an earlier phrasing of this box as
$\|\mathcal T\|\le1$ was incorrect and is hereby corrected.

\emph{Candidate route (Fej\'er--Riesz over-domination).}  The whitened operator $\mathcal T$ is a
one-sided (Toeplitz) operator with symbol
$\sigma(\xi)=2\sum_n\Lambda(n)n^{-1/2}\cos(\xi\log n)/w_A$-normalization, and the scalar level-one
theorem (Theorem~\ref{thm:level-one-scalar}) gives the pointwise positivity
$\Re(\xi'/\xi)(s)\ge0$ for $\Re s>1$, i.e.\ positivity of the whitened symbol at every interior gauge.
For a Toeplitz operator a \emph{nonnegative symbol on the boundary} would give a Fej\'er--Riesz
spectral factor $\mathcal W_A-\mathcal T_\Lambda=\mathcal D^*\mathcal D$, closing $\delta_N\ge0$.

\emph{Where it stalls (the isolated gap).}  The interior positivity of
Theorem~\ref{thm:level-one-scalar} is on $\Re s>1$; the contractivity requires the whitened symbol to
stay nonnegative as the gauge is pushed to the boundary $\Re s\to\half^+$, i.e.\ on the critical line.
There $\Re(\xi'/\xi)(\tfrac12+it)=\sum_\rho(\tfrac12-\beta_\rho)|s-\rho|^{-2}$ is a signed sum whose
sign is governed by the very off-line zeros in question.  So the candidate route reduces, cleanly and
without loss, to a single named inequality:
\[
   \textbf{(C)}\qquad
   \liminf_{\sigma\downarrow1/2}\ \bigl(\mathcal W_A-\mathcal T_\Lambda\bigr)\ \succeq\ 0
   \quad\text{as forms on the Laguerre tower.}
\]
Inequality (C) is equivalent to RH; we state it as the reduced target rather than assert it.

\emph{Experimental evidence bearing on (C).}  The harness supplies three data points.
\begin{enumerate}[label=\textnormal{(E\arabic*)}]
\item \emph{Reality certificate.}  The arithmetic minimizer roots are real to $10^{-44}$ and land on
the zeros (table above); the near-null direction of $\mathcal W_A-\mathcal T_\Lambda$ is thus a genuine
interpolant of on-line atoms.
\item \emph{Analytic factor.}  The Cholesky factor $B_N(t_0)$ of $A_N-\mathcal P_\Lambda$ varies
smoothly across a zero gap (e.g.\ $B_{11}$ decreasing monotonically $1.72\!\cdot\!10^{-4}\to
6.2\!\cdot\!10^{-5}$ over $[295.57,297.98]$ at $N=8$), so the certificate is an analytic matrix
function of the gauge while $\delta_N>0$; its terminal diagonal is $\asymp\sqrt{\delta_N}$ and vanishes
exactly where an off-line zero would appear.
\item \emph{No-margin threshold.}  Displacing one zero off the line by $\varepsilon$ flips $\delta_N<0$
at a threshold that collapses super-exponentially in $N$:
$\varepsilon^*(N)\approx1.3\!\cdot\!10^{-2},\,5.2\!\cdot\!10^{-4},\,1.1\!\cdot\!10^{-5}$ for
$N=6,8,10$.  Positivity is therefore a genuine phase property, sharpening with $N$, exactly as (C)
predicts.
\end{enumerate}
The reproducible drivers are \texttt{h8\_forensic.py}, \texttt{exp\_cholesky.py}, and
\texttt{exp\_threshold.py}.  The attack currently stands at inequality (C): a boundary/Fej\'er--Riesz argument
for the whitened symbol on the critical line would close $\delta_N\ge0$ for all $N$, and the harness
would immediately confirm it on $\zeta$ and reject it on Davenport--Heilbronn.

\subsubsection*{The terminal jets are the Li--Keiper coefficients pulled into the convergent region}

Inequality (C) is not new in the landscape: it is the boundary form of the Li--Keiper criterion, and
identifying it as such tells us exactly which classical technique the attack must supply.

\begin{lemma}[boundary regularity of the fixed-$N$ jets]\label{lem:boundary-continuity}
Fix $t_0$ and $N$.  Since the zero set of $\Xi$ is discrete and contains no point with
$|\Im z_\rho|=\half$ (Lemma~\ref{lem:xi-basic} gives $0<\beta<1$, hence $|\Im z_\rho|=|\beta-\half|<\half$
strictly), the compact segment $\{t_0-iy:\,y\in[\half,Y]\}$ has positive distance to it.  Hence the jets
$C_{jk}(t_0-iy)$, $j,k<N$, are continuous on $y\in[\half,Y]$ (dominated convergence with the
$O(|z_\rho|^{-2})$ majorant of Lemma~\ref{lem:mxi-kernel}(iii), uniform on the compact segment), and
$\delta_N(t_0-\tfrac i2)=\lim_{y\downarrow1/2}\delta_N(t_0-iy)$.
In particular interior positivity at fixed $N$ passes to the Li boundary, unconditionally.  This lemma
transfers positivity in one direction only (interior $\to$ boundary); the equivalence \emph{at} the
boundary gauge is supplied independently by Theorem~\ref{thm:scalar-jets}(b), and this lemma is never
used to argue from boundary positivity back to interior positivity.
\end{lemma}

\begin{convention}[symmetric zero summation]\label{conv:symmetric-sum}
Every sum $\sum_\rho$ over nontrivial zeros in this paper, when the individual terms are not summable
in absolute value, denotes the symmetric limit
\[
   \sum_\rho(\cdots)
   :=
   \lim_{T\to\infty}\sum_{|\Im\rho|\le T}(\cdots),
\]
zeros counted with multiplicity, and every pairing invoked (e.g.\ $\rho\leftrightarrow1-\rho$,
$\rho\leftrightarrow\bar\rho$) is applied within a common symmetric truncation before the limit is
taken.  This is the classical convention under which the Li generating function and the Li--Keiper
criterion are stated \cite{Li1997,Keiper1992,BombieriLagarias1999}, and it is the convention used for
$\lambda_n$, $S_m=\sum_\rho\rho^{-m}$, and every partial-fraction manipulation of these sums below.  The
absolutely convergent sums $\sum_\rho m_\rho|z_\rho|^{-2}$, $C_{jk}(z_0)$, and $K_\Xi(z,w)$ of
Lemma~\ref{lem:mxi-kernel} require no such convention; the symmetric limit is needed only where $S_1$
(equivalently $\lambda_1$) itself is involved, since $\sum_\rho|\rho|^{-1}$ diverges.
\end{convention}

\begin{theorem}[Li--Keiper dictionary for the terminal gauge]\label{thm:li-dictionary}
Fix the Cayley coordinate $w_{z_0}(z)=\dfrac{z-z_0}{z-\bar z_0}$ at a gauge $z_0=t_0-iy$, $y\ge\half$.
For $y=\half$, $t_0=0$ (the boundary point $z_0=-\tfrac{i}{2}$, i.e.\ $s_0=1$), and a nontrivial zero
$\rho$ with $z$-coordinate $z_\rho=(\rho-\half)/i$,
\[
   w_{-i/2}(z_\rho)\ =\ 1-\frac1\rho .
\]
Consequently, writing $a_\rho=(z_\rho-z_0)^{-1}$, the boundary terminal jet is
$C_{jk}(-\tfrac{i}{2})=\sum_\rho m_\rho\,a_\rho^{\,j+1}(z_\rho-\bar z_0)^{-(k+1)}$
(Lemma~\ref{lem:mxi-kernel}(iii)); under RH this reads
$\sum_\rho m_\rho\,a_\rho^{\,j+1}\overline{a_\rho}^{\,k+1}$ and is the Gram (moment)
matrix of the pushforward of the positive zero measure under the Cayley map $z\mapsto w_{-i/2}(z)=
1-1/\rho$; and
\[
   \delta_N(-\tfrac{i}{2})\ge0\ \ \forall N
   \quad\Longleftrightarrow\quad
   \lambda_n:=\sum_\rho\Bigl(1-(1-\tfrac1\rho)^n\Bigr)\ge0\ \ \forall n,
\]
the zero sum taken in the symmetric sense of Convention~\ref{conv:symmetric-sum},
both being equivalent to RH --- the left side by the scalar one-point jet serialization,
Theorem~\ref{thm:scalar-jets}, applied \emph{directly at the Li gauge} $z_0=-\tfrac i2$ (case (b): the
gauge lies on the boundary of the Pick domain $U_-$ but at positive, unconditional distance from the
zero set of $\Xi$, so it is a regular base point for the kernel), the right side by the
Li~\cite{Li1997}--Keiper~\cite{Keiper1992} criterion (see also Bombieri--Lagarias
\cite{BombieriLagarias1999}).  For $y>\half$ the terminal jet at $z_0=t_0-iy$ is the analytic
continuation of the same Cayley/Li data into the half-plane of absolute convergence
$\Re s_0=\half+y>1$, and inequality \textnormal{(C)} is precisely the boundary limit $y\downarrow\half$.
\end{theorem}

\begin{proof}
Write $z_\rho=(\rho-\half)/i=-i(\rho-\half)$.  Then $z_\rho+\tfrac{i}{2}=-i(\rho-\half)+\tfrac{i}{2}
=-i\rho+i=-i(\rho-1)$ and $z_\rho-\tfrac{i}{2}=-i(\rho-\half)-\tfrac{i}{2}=-i\rho$, so
\[
   w_{-i/2}(z_\rho)=\frac{z_\rho-(-\tfrac{i}{2})}{z_\rho-(+\tfrac{i}{2})}
   =\frac{-i(\rho-1)}{-i\rho}=\frac{\rho-1}{\rho}=1-\frac1\rho,
\]
the Li variable, of modulus $1$ under RH.  The scalar jet identity
$C_{jk}(z_0)=\sum_\rho m_\rho\,(z_\rho-z_0)^{-(j+1)}(z_\rho-\bar z_0)^{-(k+1)}$ is
Lemma~\ref{lem:mxi-kernel}(iii); under RH, $z_\rho-\bar z_0=\overline{z_\rho-z_0}$, so with
$a_\rho=(z_\rho-z_0)^{-1}$ it reads $C_{jk}=\sum_\rho m_\rho\,a_\rho^{\,j+1}\overline{a_\rho}^{\,k+1}$
and exhibits $C(z_0)$ as the Gram matrix of the vectors
$(a_\rho^{\,j+1})_j$ against the positive weights $m_\rho$, hence positive semidefinite whenever the
$z_\rho$ are real.  Its positivity for all $N$ is equivalent to RH by
Theorem~\ref{thm:scalar-jets}, applied directly at $z_0=-\tfrac i2$: the distance hypothesis is case (b)
of that theorem --- a zero of $\Xi$ within distance $\delta$ of $-\tfrac i2$ would be a zero of $\zeta$
within $\delta$ of $s=1$, excluded unconditionally by the classical zero-free region --- so no passage
of positivity from the interior gauges to the boundary is used in either direction.
The Li--Keiper positivity $\lambda_n\ge0$ for all $n$ is equivalent to
RH by \cite{Li1997,Keiper1992}.  The two positivity families are therefore equivalent to each other, and
the boundary jet is built from the very Li variable $1-1/\rho$ by the first identity.  For $y>\half$,
$s_0=\half+iz_0$ has $\Re s_0=\half+y>1$, and $C_{jk}(z_0)$ is holomorphic in $z_0$ there
(Lemma~\ref{lem:mxi-kernel}(iii)), agreeing with the boundary data as $y\downarrow\half$; thus the interior
jets are the analytic continuation and (C) is their boundary value.
\end{proof}

\begin{remark}[the dictionary is an all-orders equivalence, not a levelwise transform]\label{rmk:li-levelwise}
The equivalence of Theorem~\ref{thm:li-dictionary} holds between the \emph{complete} families
($\forall N$ versus $\forall n$), through RH; it is not a finite-order positivity-preserving triangular
transform, and the two truncated hierarchies are genuinely different at each finite level.  This is
visible already at $N=2$.  Partial fractions against the pairing $\rho\leftrightarrow1-\rho$
(so $\sum_\rho1/(\rho-1)=-\lambda_1$ and $\sum_\rho1/\rho^2=2\lambda_1-\lambda_2$) give, for the
boundary jets $C_{jk}=i^{\,j+k+2}\sum_\rho m_\rho(\rho-1)^{-(j+1)}\rho^{-(k+1)}$,
\[
   C_{00}=2\lambda_1,\qquad
   C_{01}=\overline{C_{10}}=i\,(4\lambda_1-\lambda_2),\qquad
   C_{11}=8\lambda_1-2\lambda_2,
\]
hence the exact identity
\[
   \det J_2(-\tfrac i2)
   =
   C_{00}C_{11}-|C_{01}|^2
   =
   \lambda_2\,(4\lambda_1-\lambda_2).
\]
Thus $J_1\succeq0$ is exactly $\lambda_1\ge0$ (with the harmless factor $2$), while $J_2\succeq0$ is
equivalent to $\lambda_1\ge0$, $\lambda_2\ge0$, \emph{and} $\lambda_2\le4\lambda_1$ --- strictly
stronger than finite Li positivity at that level, and not implied by it: every jet entry lies in the
integer span of $\lambda_1,\ldots,\lambda_{j+k+1}$, yet the two truncated positivity hierarchies
coincide only in the $\forall N/\forall n$ limit, where both are RH.
Numerically $4\lambda_1-\lambda_2=3.71\cdot10^{-5}>0$ and
$\det J_2=\lambda_2(4\lambda_1-\lambda_2)\approx3.4\cdot10^{-6}>0$, consistent with RH.  This is one
more manifestation of conservation of difficulty: the dictionary changes coordinates on the same
terminal positivity, it does not decompose it into weaker levelwise statements.
\end{remark}

\begin{remark}[what the dictionary buys]\label{rmk:li-buys}
The Li--Keiper coefficients live exactly at the edge $\Re s=1$, where the prime series is only
conditionally summable and the positivity $\lambda_n\ge0$ has resisted proof since 1997.  The terminal
gauges $y>\half$ move the identical positivity into $\Re s_0>1$, where the prime series converges
absolutely and the level-one positivity $\Re(\xi'/\xi)>0$ (Theorem~\ref{thm:level-one-scalar}) is
unconditional and manifest.  The open content is therefore isolated as a single passage to the boundary:
prove that the interior positivity, uniform in $N$, survives $y\downarrow\half$ --- a Fatou/boundary-value
problem for the whitened Weil symbol rather than a positivity to be conjured from scratch.  Numerically
the boundary values are reproduced exactly: the closed-form $\tfrac1{(n-1)!}\frac{d^n}{ds^n}
[s^{n-1}\log\xi]_{s=1}$ gives $\lambda_{1..5}=0.023096,\,0.092346,\,0.207639,\,0.368790,\,0.575543$,
matching the zero-side sum $\sum_\rho(1-(1-1/\rho)^n)$, and $|w_{-i/2}(u)|=1$ on $\R$ to full precision
(driver \texttt{li\_check.py}).
\end{remark}

\begin{remark}[the boundary passage is regular; the difficulty is the $N\to\infty$ race]
\label{rmk:boundary-regular}
The Fatou passage of Remark~\ref{rmk:li-buys} splits into two limits, $y\downarrow\half$ and
$N\to\infty$, and the harness locates the difficulty entirely in the second.  Fixing $N$ and lowering
$y$ toward the Li boundary, the sampling margin $\delta_N(t_0-iy)$ stays strictly positive and varies
smoothly, with no degeneration at $y=\half$; e.g.\ at $t_0\approx297.24$,
\[
   \begin{array}{c|ccc}
      y & N=6 & N=8 & N=10\\\hline
      1.00 & 7.9\!\cdot\!10^{-5} & 7.6\!\cdot\!10^{-8} & 1.8\!\cdot\!10^{-11}\\
      0.60 & 3.9\!\cdot\!10^{-7} & 4.0\!\cdot\!10^{-11} & 1.2\!\cdot\!10^{-15}\\
      0.505& 6.1\!\cdot\!10^{-8} & 3.1\!\cdot\!10^{-12} & 4.6\!\cdot\!10^{-17}
   \end{array}
\]
(driver \texttt{exp\_boundary.py}).  Thus for each fixed $N$ the boundary value is a positive limit,
and (C) is \emph{not} obstructed by the passage $y\downarrow\half$; the whole content is the uniformity
of positivity as $N\to\infty$ \emph{at} the boundary --- equivalently the growth race in
$\lambda_n=\lambda_n^{\rm arch}+\lambda_n^{\rm prime}$, where the archimedean ceiling
$\lambda_n^{\rm arch}\sim\tfrac{n}{2}\log n$ (Bombieri--Lagarias \cite{BombieriLagarias1999}) must
dominate the oscillatory prime part $\lambda_n^{\rm prime}$ for every $n$.  This is the sharpest form of
the remaining target: an unconditional lower bound on the archimedean Li coefficient that beats the
prime oscillation uniformly in $n$, with the interior gauges $y>\half$ available as an absolutely
convergent regularization in which $\Re(\xi'/\xi)>0$ holds outright.
\end{remark}

\subsubsection*{Chain $\Omega$: the end-to-end path and the shape of its single residue}

We assemble the whole route as one linear chain, marking each $\Omega$-link Closed or, for the one
terminal residue, stating it exactly.  This is the object we offer for end-to-end audit.  The
$\Omega$-chain is the third and deepest level of the logical hierarchy of \S\ref{sec:status-map}: it
lives below the H-level analysis of Step~5, and its entries are typed --- $\Omega_1$--$\Omega_6$ are
closed reductions or equivalences ($\Omega$-links), while $\Omega_7$ is not a link but the terminal
statement the chain isolates.

\begin{center}
\small
\renewcommand{\arraystretch}{1.3}
\begin{tabular}{@{}p{0.05\linewidth}p{0.68\linewidth}p{0.21\linewidth}@{}}
\toprule
& Statement & Status\\\midrule
$\Omega_1$ & RH $\Longleftrightarrow$ all $\Xi$-zeros real (Lemma~\ref{lem:rh-translation}) & Closed\\
$\Omega_2$ & real zeros $\Longleftrightarrow$ ARP-P (Steps 1--15, Corollary~\ref{cor:live-equivalence}) & Closed\\
$\Omega_3$ & ARP-P $\Longleftrightarrow$ terminal positivity $\delta_N(z_0)\ge0\ \forall N$ in the
reference-whitened defect (Corollary~\ref{cor:arp-terminal-defect}; the observable is
Definition~\ref{def:terminal-defect-fixed}) & Closed, in the reference-whitened form\\
$\Omega_4$ & $\delta_N\ge0\ \forall N\Longleftrightarrow$ one-sided whitened domination
$\lambda_{\max}(\mathcal W_A^{-1/2}\mathcal T_\Lambda\mathcal W_A^{-1/2})\le1$ (the two-sided norm
exceeds $1$ at interior gauges; see the warning below) & Closed; $J_N^\infty$-whitening valid for
$N\le N_*(t_0)$ (Lemma~\ref{lem:pole-pair})\\
$\Omega_5$ & interior positivity ($y>\half$) has a positive, regular boundary limit $y\downarrow\half$
(Lemma~\ref{lem:boundary-continuity}) & Closed, per $N$\\
$\Omega_6$ & boundary positivity $\Longleftrightarrow$ Li--Keiper $\lambda_n\ge0\ \forall n$, via
$w_{-i/2}=1-1/\rho$ (Theorem~\ref{thm:li-dictionary}) & Closed\\
$\Omega_7$ & $\lambda_n\ge0$ for all $n$: archimedean ceiling $\tfrac n2\log n$ dominates the prime
oscillation & \textbf{terminal residue}: open; $\equiv$ RH (Li)\\
\bottomrule
\end{tabular}
\end{center}

Every $\Omega$-link $\Omega_1$--$\Omega_6$ is proved above.  The terminal residue $\Omega_7$ is the classical Li--Keiper
positivity in the sharp uniform-in-$n$ form isolated by the chain.  Two proved facts about $\Omega_7$
delimit what a closing argument may look like.

\begin{proposition}[no positive prime truncation (computational) --- the terminal positivity is a global cancellation]
\label{thm:no-truncation}
Fix an interior gauge and $N\ge6$.  For the whitened truncated form
$m_N(X):=\lambda_{\min}\bigl(A_N^{-1/2}(A_N-\mathcal T_\Lambda^{\le X})A_N^{-1/2}\bigr)$, where
$\mathcal T_\Lambda^{\le X}=\sum_{n\le X}\Lambda(n)n^{-1/2}(S_{\log n}+S_{\log n}^*)$, one has (to the
stated working precision) $m_N(X)<0$ for a range of finite $X$ even though the full defect
$m_N(\infty)=\delta_N>0$.  Consequently $A_N-\mathcal P_\Lambda$ is \emph{not} of the form
$(\text{positive truncated part})-(\text{small tail})$: no finite prime cutoff yields a positive
operator to perturb from.
\end{proposition}

\begin{proof}[Proof (computational, with the harness)]
Driver \texttt{omega\_g1g2\_truncation.py} evaluates $m_N(X)$ and the whitened tail norm at
$t_0\approx297.24$, $y=1$.  Already at $X=50$ one finds $m_6=-0.53$, $m_8=-0.66$, $m_{10}=-0.73$, while
$\delta_6=8.0\!\cdot\!10^{-5}$, $\delta_8=7.7\!\cdot\!10^{-8}$, $\delta_{10}=1.8\!\cdot\!10^{-11}$; the
tail norm exceeds $|m_N(X)|$ at every listed $X$.  Thus the truncated operator is indefinite while the
full one is positive: the positivity is restored only by the entire prime tail.
\end{proof}

\begin{remark}[why this fixes the method, not just a method]\label{rmk:omega-shape}
Theorem~\ref{thm:no-truncation} says the archimedean part does \emph{not} dominate the prime part termwise
or up to a controllable tail; positivity is an exact cancellation across all primes, equivalently the
exact evaluation of the explicit formula at the (real) zeros.  The same obstruction defeats the two
perturbative backups we tested: mollifying the symbol (a deeper gauge) is the $X\to\infty$ end of the
same truncation and inherits the indefiniteness at finite scale (driver
\texttt{omega\_g1g2\_truncation.py}), and the backward-heat-flow route (driver
\texttt{omega\_g3\_backflow.py}, from $t>\Lambda_{\rm dBN}$ where Polymath15 \cite{Polymath15} gives
real zeros) propagates the factorially small margin $\delta_N\asymp(\kappa L/N)^{2N}$ against a
polynomial Gronwall drift, hence closes only a range $N\le N_*(L)$.  Both confirm, from independent
directions, that $\Omega_7$ carries the full arithmetic content: a proof must supply the exact
archimedean-over-prime domination of the Li coefficients uniformly in $n$, not a dominated inequality.
The chain reduces RH to precisely that classical statement and to nothing softer.
\end{remark}

\subsubsection*{$\Omega_7$ up close: exact structure, the archimedean trend, and the detection mechanism}

We record what a direct assault on $\Omega_7$ meets, with exact numbers (drivers
\texttt{omega7\_li\_decomp.py}, \texttt{omega7\_li\_inject.py}).  Two structural facts organize
everything.

\begin{proposition}[the pairing and the on-line sum-of-squares]\label{prop:li-sos}
Nontrivial zeros pair as $\rho\leftrightarrow1-\rho$, and the Li variables are reciprocal:
$w_{1-\rho}=1/w_\rho$, $w_\rho=1-1/\rho$.  Moreover
\[
   |w_\rho|=\frac{|\rho-1|}{|\rho|}
   \begin{cases} =1,&\Re\rho=\tfrac12,\\ >1,&\Re\rho<\tfrac12,\\<1,&\Re\rho>\tfrac12.\end{cases}
\]
Consequently, if RH holds then, pairing $\rho=\tfrac12+i\gamma$ with $\bar\rho=1-\rho$ and writing
$w_\rho=e^{i\theta_\gamma}$,
\[
   \lambda_n=\sum_{\gamma>0}\bigl(2-2\cos n\theta_\gamma\bigr)
   =4\sum_{\gamma>0}\sin^2\!\tfrac{n\theta_\gamma}{2}\ \ge\ 0,
\]
a manifest sum of squares.  Thus $\Omega_7$ is trivial \emph{given} RH; its content is establishing the
same nonnegativity from the arithmetic definition of $\lambda_n$, in which the zeros do not appear.
\end{proposition}

\begin{proof}
$w_\rho w_{1-\rho}=\frac{\rho-1}{\rho}\cdot\frac{-\rho}{1-\rho}=1$; and $|w_\rho|=|\rho-1|/|\rho|$ with
$|\rho-1|^2-|\rho|^2=1-2\Re\rho$ gives the trichotomy.  For $\Re\rho=\tfrac12$, $1-\rho=\bar\rho$, so
the pair contributes $[1-w^n]+[1-\overline{w}^{\,n}]=2-2\Re w^n=2-2\cos n\theta_\gamma$.  Summing over
$\gamma>0$ and using $\theta_\gamma=O(1/\gamma)$ (whence the terms are $O(n^2/\gamma^2)$ for
$\gamma\gg n$) gives an absolutely convergent nonnegative series.
\end{proof}

\begin{remark}[the exact trend and the measured margin]\label{rmk:li-trend}
Splitting $\log\xi=\bigl[\log s-\tfrac s2\log\pi+\log\Gamma(\tfrac s2)\bigr]+\log\bigl((s-1)\zeta(s)
\bigr)$ gives, through $\lambda_n=n\sum_{k=1}^n\binom{n-1}{k-1}\sigma_k$ with
$\sigma_k=[(s-1)^k]\log\xi(s)$ (closed form: polygamma for the archimedean part, Stieltjes constants
$\gamma_j$ for the $\zeta$ part), an \emph{unconditional} exact decomposition
$\lambda_n=\lambda_n^{\rm arch}+\lambda_n^{\rm prime}$.  Computed to $n=120$: every $\lambda_n>0$; the
archimedean part follows the Voros trend $\lambda_n\sim\tfrac n2\log\tfrac n{2\pi}$
(ratio $0.823,0.845,0.873$ at $n=40,80,120$, increasing to $1$); and the prime oscillation is a small
relative fluctuation, $|\lambda_n^{\rm prime}|/\lambda_n^{\rm arch}$ falling to
$0.002$--$0.04$ for $n\ge30$ (still $0.083$ at $n=20$).  Under RH this smallness is forced by $|w_\rho|=1$; unconditionally it is
exactly what is at stake.
\end{remark}

\begin{remark}[the detection mechanism, exhibited]\label{rmk:li-detection}
The obstruction is a single geometric term.  A putative off-line zero $\rho_0$ ($\Re\rho_0\ne\tfrac12$)
has a functional-equation partner with $|w|>1$, and its quartet contributes
$-\,2\Re w^n+\cdots$ with $|w|^n$ growing geometrically; for large $n$ this overturns the
$\tfrac n2\log n$ trend and drives $\lambda_n<0$.  Injecting a low off-line quartet at
$\rho_0=0.9+3i$ (with $\max|w|=1.0434$ over the quartet) into the true $\lambda_n$ turns it negative
first at $n=97$ ($\lambda_{97}^{\rm inj}=-5.95$ against the true $\lambda_{97}=113.4$), while an on-line
control $\rho_0=\tfrac12+3i$ keeps every $\lambda_n>0$ (it only adds $4\sin^2$).  This is the
Li-coordinate form of the terminal forensic: positivity is destroyed precisely, and only, by the
geometric growth $|w|>1$ of an off-line zero.  Closing $\Omega_7$ therefore means bounding
$|\lambda_n^{\rm prime}|<\lambda_n^{\rm arch}$ for all $n$ \emph{unconditionally}, i.e.\ forbidding
$|w_\rho|>1$ --- which is RH itself, now displayed as a single explicit inequality between an
archimedean trend of size $\tfrac n2\log n$ and an arithmetic oscillation whose only possible unbounded
source is an off-line zero.
\end{remark}

\subsubsection*{Which inputs cannot close $\Omega_7$: an insufficiency lemma and the height--range correspondence}

We explored bounding $\lambda_n^{\rm prime}$ against $\lambda_n^{\rm arch}$ for all $n$ using inputs
weaker than RH.  The outcome is an \emph{input-class restriction}, not an impossibility: it rules out
one family of hypotheses (zero-location statements weaker than RH) and leaves the direct arithmetic
attack fully open.

\begin{proposition}[insufficiency of zero-location inputs]\label{prop:unboundable}
If $\xi$ has a zero $\rho_0$ with $\Re\rho_0\ne\tfrac12$, then $\lambda_n<0$ for infinitely many $n$
(this is the known converse direction of the Li criterion
\cite{Li1997,BombieriLagarias1999}).  Consequently no hypothesis \emph{consistent with} the existence
of a single off-line zero can imply $\lambda_n\ge0$ for all $n$.  In particular neither the classical
zero-free region $\beta<1-c/\log(|\gamma|+2)$ nor any zero-density estimate
$N(\sigma,T)\ll T^{a(1-\sigma)}$ suffices as the sole input for $\Omega_7$: both permit off-line
zeros.
\end{proposition}

\begin{proof}[Proof sketch, with attribution]
The full statement is Bombieri--Lagarias \cite{BombieriLagarias1999} applied to the zero multiset of
$\xi$; we record the mechanism for a single dominant quartet.  By the functional equation and
conjugation, $\rho_0$ belongs to a quartet containing a zero $\rho$ with $\Re\rho<\tfrac12$, for which
$r:=|w_\rho|=|\rho-1|/|\rho|>1$ (Proposition~\ref{prop:li-sos}).  Its quartet contributes
$-\,(r^n+r^{-n})\cos(n\varphi)+O(1)$, $\varphi=\arg w_\rho$, of amplitude $r^n$ exceeding, for large
$n$ along a positive-density set where $\cos(n\varphi)\ge\tfrac12$ (Weyl equidistribution if
$\varphi/2\pi$ is irrational, periodicity if rational), both the trend
$\lambda_n^{\rm arch}\sim\tfrac n2\log\tfrac n{2\pi}$ and the contribution of the remaining zeros
\emph{provided} $\rho$ realizes the maximal $|w|$ alone.  When several off-line zeros share the maximal
modulus, cancellations among their phases must be excluded; this is exactly what the
Bombieri--Lagarias argument handles in general, and we rely on their theorem rather than the sketch for
the full claim.
\end{proof}

\begin{remark}[scope: what this does and does not say]\label{rmk:unboundable-scope}
Proposition~\ref{prop:unboundable} restricts the \emph{input class}: $\Omega_7$ cannot be closed by
citing a zero-location statement weaker than RH, because any such statement tolerates the zero that
destroys positivity.  It does \emph{not} say that $\Omega_7$ is unprovable, and the slogan ``the only
sufficient input is RH itself'' would be circular: since $\lambda_n\ge0\ \forall n$ is
\emph{equivalent} to RH \cite{Li1997}, any proof of $\Omega_7$ tautologically yields RH --- that is the
goal of the chain, not an obstruction to it.  The open route is the one the decomposition itself
exhibits: bound $|\lambda_n^{\rm prime}|<\lambda_n^{\rm arch}$ directly on the \emph{prime side}
(explicit-formula / Weil-positivity inputs, test-function constructions, or the operator-theoretic
mechanisms M1--M3 discussed below), none of which passes through a zero-location hypothesis and
none of which is excluded here.  $\Omega_7$ therefore remains \textbf{open}, with this lemma serving
only to prune dead-end input families.
\end{remark}

\begin{remark}[height--range correspondence: the achievable unconditional traction]\label{rmk:height-range}
The geometric growth is slow for high zeros, and this is quantitative.  A putative off-line zero at
height $\gamma$ and displacement $\delta=\tfrac12-\beta$ has
$\log|w_\rho|=\dfrac{\delta}{\beta^2+\gamma^2}\,(1+o(1))$, so it can force $\lambda_n<0$ only once
$r^n\gtrsim\lambda_n^{\rm arch}$, i.e.\ for
\[
   n\ \gtrsim\ n_*(\gamma,\delta)\ \approx\ \frac{\beta^2+\gamma^2}{\delta}\,
   \log\!\Bigl(\tfrac n2\log\tfrac n{2\pi}\Bigr).
\]
The measured thresholds obey $n_*\cdot\log|w|$ nearly constant (values
$4.85,4.92,5.73,5.98,7.50,6.71,8.62$ across $\gamma=3\!-\!10$, driver
\texttt{omega7\_li\_threshold.py}), confirming $n_*\asymp(\beta^2+\gamma^2)\delta^{-1}\log(\cdots)$.
Contrapositively, if RH is verified up to height $H$ --- no off-line zeros with $|\gamma|\le H$ ---
then the lowest possible off-line zero sits at $\gamma>H$ and cannot overturn the trend until
$n\gtrsim H^2\log H$: the single-zero analysis gives $\lambda_n>0$ unconditionally for
$n\lesssim H^2\log H$.  With RH verified to $H\approx3\times10^{12}$ this is $\lambda_n>0$ for
$n$ up to about $10^{25}$.  The step from the single-zero bound to a fully rigorous combined bound over
all $\gamma>H$ requires a zero-density input to control the phases of many high zeros; that is the exact,
and only, place where unconditional work on $\Omega_7$ via \emph{zero-location} inputs can still gain
ground: such inputs yield a finite range of $n$, never the all-$n$ statement
(Proposition~\ref{prop:unboundable}).  The all-$n$ statement must come from the prime side
(Remark~\ref{rmk:unboundable-scope}), where no such barrier is proved.
\end{remark}

\subsubsection*{The prime side made exact: $\Omega_7$ as an inequality about an explicit prime sum}

Remark~\ref{rmk:unboundable-scope} leaves one open front: bound $|\lambda_n^{\rm prime}|$ against
$\lambda_n^{\rm arch}$ directly on the prime side.  The first step is to make
$\lambda_n^{\rm prime}$ an \emph{explicit} prime sum.  Both identities below were numerically
gated before being written (driver \texttt{omega7\_prime\_kernel.py}).

\begin{proposition}[Laguerre kernel of the prime side]\label{prop:prime-kernel}
For every $n\ge1$ and $x>0$,
\[
   n\sum_{k=1}^{n}\binom{n-1}{k-1}\frac{(-x)^k}{k!}\;=\;-\,x\,L^{(1)}_{n-1}(x),
\]
with $L^{(1)}_{n-1}$ the generalized Laguerre polynomial.  Consequently, for $\Re s_0>1$,
applying the Li binomial transform to the Taylor coefficients
$\sigma_k^{\zeta}(s_0)=[(s-s_0)^k]\log\zeta(s)$ gives the absolutely convergent representation
\[
   T_n(s_0)\;:=\;n\sum_{k=1}^{n}\binom{n-1}{k-1}\,\sigma_k^{\zeta}(s_0)
   \;=\;-\sum_{m\ge2}\frac{\Lambda(m)}{m^{s_0}}\,L^{(1)}_{n-1}(\log m),
\]
and the prime part of the Li coefficient is the boundary value
\[
   \lambda_n^{\rm prime}
   \;=\;\lim_{\varepsilon\downarrow0}\Bigl[\;
   n\sum_{k=1}^{n}\binom{n-1}{k-1}\frac{(-1)^{k-1}}{k\,\varepsilon^{k}}
   \;-\;\sum_{m\ge2}\frac{\Lambda(m)}{m^{1+\varepsilon}}\,L^{(1)}_{n-1}(\log m)\Bigr],
\]
the divergences of the pole and prime pieces cancelling in the limit.
\end{proposition}

\begin{proof}
The kernel identity follows from the generating function
$\sum_{n\ge0}L^{(1)}_{n}(x)z^{n}=(1-z)^{-2}e^{-xz/(1-z)}$ together with
$\sum_n\lambda_n^{\bullet}z^{n-1}$ applied to $m^{-s}=m^{-s_0}e^{-(s-s_0)\log m}$; it was verified to
$25$ digits at mixed $(n,x)$.  The representation for $T_n(s_0)$ is then term-by-term rearrangement of
the absolutely convergent double sum ($\Re s_0>1$).  The boundary statement is the same computation for
$\log\bigl((s-1)\zeta(s)\bigr)=\log(s-1)+\log\zeta(s)$ at $s_0=1+\varepsilon$, whose left side is
regular at $\varepsilon=0$ with Taylor data the Stieltjes constants.  Numerical gates: $T_n(s_0)$
agrees with the exact log-$\zeta$ derivative route to $10^{-9}$ at $s_0=2$ and $10^{-6}$ at $s_0=1.5$
(prime powers to $2\times10^6$ plus PNT tail; the residual has exactly the size of the
$\psi(y)-y$ fluctuation, i.e.\ the zero term); the $\varepsilon$-limit converges linearly to the exact
Stieltjes values of $\lambda_n^{\rm prime}$.
\end{proof}

\begin{remark}[what the representation buys, and what it costs]\label{rmk:prime-kernel-shape}
$\Omega_7$ is now the single explicit inequality
\[
   \Bigl|\lim_{\varepsilon\downarrow0}\bigl[\text{pole}_n(\varepsilon)-
   \textstyle\sum_m\Lambda(m)m^{-1-\varepsilon}L^{(1)}_{n-1}(\log m)\bigr]\Bigr|
   \;<\;\lambda_n^{\rm arch}\sim\tfrac n2\log\tfrac n{2\pi}\qquad(\forall n),
\]
with two structural assets not available to zero-location inputs: the weights $\Lambda(m)\ge0$ are
nonnegative, and the kernel $L^{(1)}_{n-1}$ is classical (Plancherel--Rotach: oscillatory for
$\log m\lesssim4n$, monotone beyond).  The cost is equally visible: the kernel's effective support
reaches $m\approx e^{4n}$, so the inequality demands cancellation in $\psi(y)-y$ over exponentially
long ranges --- the measured cancellation ratios
$|{\sum}|/{\sum}|\cdot|$ at $s_0=1.5$ are $0.28$--$0.49$ for $n\le12$ and must, at
$\varepsilon=0$, come from the oscillation of the prime counting error against the Laguerre
oscillation.  This is the precise prime-side face of $\Omega_7$.
\end{remark}

\subsubsection*{Step 2: pointwise prime-error bounds are insufficient --- the $\Theta(n)$ cancellation deficit}

With Proposition~\ref{prop:prime-kernel} in hand, the natural first attack is partial summation:
integrate by parts (exactly, as Stieltjes integrals; the pole term is the average-$\Lambda$
compensation, $\mathrm{pole}_n(\varepsilon)=\int_0^\infty L^{(1)}_{n-1}(x)e^{-\varepsilon x}dx$,
verified to $10^{-20}$) to get
\[
   \lambda_n^{\rm prime}\;=\;-\lim_{\varepsilon\downarrow0}\int_1^\infty
   \bigl(\psi(y)-y\bigr)\,\frac{d}{dy}\Bigl[\frac{L^{(1)}_{n-1}(\log y)}{y^{1+\varepsilon}}\Bigr]\,dy ,
\]
and then insert a pointwise bound $|\psi(y)-y|\le B(y)$.  We measured the resulting ceiling
$A_n[B]:=\int_1^\infty B(y)\,\bigl|\tfrac{d}{dy}[L^{(1)}_{n-1}(\log y)/y]\bigr|\,dy$ against
$\lambda_n^{\rm arch}$ (driver \texttt{omega7\_partial\_summation.py}; true $|\psi(y)-y|$ from a sieve
to $10^7$, tails by model):

\begin{center}\small
\begin{tabular}{r|rrrr}
$n$ & $\lambda_n^{\rm arch}\!\sim\!\tfrac n2\log\tfrac n{2\pi}$ & $A_n$ (true$+$RH tail) & $A_n$ (RH model $\tfrac{\sqrt y\log^2y}{8\pi}$) & $A_n$ (VK shape) \\\hline
 8 & $0.97$ & $45.3$ & $37.8$ & $2.8\times10^{11}$\\
20 & $11.6$ & $334$  & $311$  & $4.6\times10^{23}$\\
35 & $30.1$ & $1221$ & $1178$ & $1.0\times10^{33}$\\
\end{tabular}
\end{center}

\begin{remark}[second insufficiency: pointwise bounds on $\psi(y)-y$]\label{rmk:pointwise-insufficient}
The ceiling exceeds the archimedean budget at \emph{every} $n\ge2$ --- even for the RH-optimal
pointwise bound $B(y)=\sqrt y\,\log^2y/(8\pi)$.  Unconditional shapes are hopeless outright: the
Vinogradov--Korobov saving $e^{-c(\log y)^{3/5}}$ cannot suppress the $y$-size of the trivial bound
against the polynomial Laguerre tail, and $A_n$ explodes exponentially.  But the RH-model deficit is
only \emph{polynomial}: across $2\le n\le35$ the measured law is
$A_n\approx0.27\,n^2\log n$ (the ratio $A_n/(n^2\log n)$ is $0.26$--$0.27$ throughout), against a
budget $\lambda_n^{\rm arch}\approx\tfrac n2\log n$.  So partial summation with absolute values
overshoots by a factor $\Theta(n)$, no more: closing $\Omega_7$ requires extracting exactly one
factor $n$ of sign-cancellation from the oscillation of $\psi(y)-y$ against the oscillation of the
Laguerre kernel (both oscillate in $\log m\lesssim4n$; their relative phase is where the remaining
content of RH sits).  Together with Proposition~\ref{prop:unboundable} this sharpens the target of
the prime-side programme: not pointwise magnitude, not zero location --- \emph{phase}.  The
structural asset $\Lambda(m)\ge0$ and the operator mechanisms M1--M3, which see signs, are the
remaining instruments (Step 3 of the attack plan).
\end{remark}

\subsubsection*{Step 3: the $\Lambda\ge0$ asset in whitened coordinates --- factorial per-term loss}

The last instrument of the prime-side plan is the nonnegativity of the von Mangoldt weights.  At an
interior gauge $z_0=t_0-iy$ with $y>\tfrac12$ (so $\Re s_0=\tfrac12+y>1$) the prime jet decomposes
absolutely per prime power, $\mathcal P_\Lambda=\sum_m P_m$, each $P_m$ carrying the weight
$\Lambda(m)\ge0$ (decomposition cross-checked against the exact $\zeta$-derivative jet:
error $3\times10^{-4}$ at $N=4$, $2.5\times10^{-2}$ at $N=8$; $N=12$ is truncation-limited at
$X=2\times10^6$).  Whitening by $S=A_N^{-1/2}$, we measured the per-term ceiling
$S_{\rm abs}(N)=\sum_m\Lambda(m)\,\|S P_m^{(1)}S\|$ (the best any term-by-term absolute argument can
achieve using only $\Lambda\ge0$) against the true one-sided extremes of $S\mathcal P_\Lambda S$
(driver \texttt{omega7\_operator\_interference.py}):

\begin{center}\small
\begin{tabular}{rr|rrrr}
$t_0$ & $N$ & $\lambda_{\max}(S\mathcal P S)$ & $\lambda_{\min}$ & $S_{\rm abs}$ & interference $I$ \\\hline
10 & 4 & $0.99999993$ & $-0.097$ & $219$ & $2.2\times10^2$\\
10 & 8 & $1-10^{-12}$ & $-0.369$ & $1.9\times10^4$ & $1.9\times10^4$\\
30 & 4 & $0.9982$ & $-3.34$ & $59.5$ & $17.8$\\
30 & 8 & $1-2\times10^{-12}$ & $-7.62$ & $4.5\times10^3$ & $5.8\times10^2$\\
30 & 12 & $1^{-}$ & $-11.90$ & $\gtrsim2.9\times10^5$ & $\gtrsim2.5\times10^4$\\
\end{tabular}
\end{center}

\begin{remark}[third insufficiency, and the sharp shape of H8]\label{rmk:whitened-interference}
Three structural facts, each measured above.
(i)~\emph{The one-sided domination is razor-sharp:} $\lambda_{\max}(S\mathcal P_\Lambda S)$ saturates
at $1$ to twelve digits already at $N=8$ --- the H8 inequality holds with margin
$\to0$ exponentially fast in $N$, as difficulty-equivalence with RH demands.
(ii)~\emph{Per-term positivity is factorially insufficient:} the ceiling $S_{\rm abs}$ grows
factorially in $N$ (e.g.\ $219\to1.9\times10^4$ from $N=4$ to $8$ at $t_0=10$), so the
$\Lambda(m)\ge0$ asset used absolutely --- any argument that bounds prime contributions one at a
time, in any whitened jet coordinates --- loses \emph{factorially}, even worse than the $\Theta(n)$
deficit of the Li diagonal (Remark~\ref{rmk:pointwise-insufficient}).  The surviving positivity is
carried entirely by cross-prime phase interference, the decoherence mechanism of
Remark~\ref{rmk:H1-matrix}.
(iii)~\emph{The reduction is one-sided only:} the bottom of the whitened prime jet grows like $-N$,
so no proof route through two-sided norm bounds (contraction arguments, unitary dilations of
$\mathcal T$) can be correct as stated; the target is exactly
$\lambda_{\max}\le1$.  Taken together with the first two insufficiencies
(Proposition~\ref{prop:unboundable}: zero-location inputs; Remark~\ref{rmk:pointwise-insufficient}:
pointwise error bounds), the prime-side programme has now excluded every input class that ignores
phase: what remains of RH in $\Omega_7$/H8 is precisely a phase-coherence statement about
$\{m^{-it_0}\}$ against the Laguerre oscillation, sharp to factor $1$, with zero margin.  This is the
honest terminal shape of the exercise.
\end{remark}



\subsection{Step 5 substructure (continued): closure attempts on \texorpdfstring{$\Omega_7$}{Omega7} and a closure specification}\label{subsec:omega7-attack}

The terminal residue $\Omega_7$ --- equivalently $\delta_N\ge0$ for all $N$ --- is the single open mathematical
input of the whole route.  We record here, for the audit trail, the strongest attacks executed against
it, each run to its exact failure point, together with the unconditional by-products those failures
leave behind.  Difficulty conservation (H0) is not an editorial convention but modus ponens: any
purported proof contains exactly one step that is either RH-strength new mathematics or an error.  The
productive form of the exercise is therefore to choose the best route by elimination, execute it, and
perform the autopsy of the death point --- which is itself a theorem.

\subsubsection*{Route elimination}

The proved walls exclude fixed margins (Corollary~\ref{cor:no-margin}), losses $e^{o(N\log N)}$ (the
factorial wall), zero-location inputs (Proposition~\ref{prop:unboundable}), pointwise bounds on
$\psi(y)-y$ (the $\Theta(n)$ deficit of Remark~\ref{rmk:pointwise-insufficient}), and term-by-term
$\Lambda\ge0$ positivity (the measured factorial loss of Remark~\ref{rmk:whitened-interference}).  Three
further mechanisms are dismissed quickly.  \emph{Krein extension}: short-window positivity
($\operatorname{supp}<\log2$, no primes) is provable, and Krein extends any interval-positive-definite
function to a positive-definite function on $\R$; but the route needs \emph{our} kernel to be the
positive extension, not the mere existence of one.  \emph{Fej\'er--Riesz}: factoring
$\mathcal A-\mathcal T_\Lambda=B^*B$ requires the boundary symbol to be nonnegative $\iff$ RH (circular,
already recorded).  \emph{Reflection positivity}: the per-prime symbol $2\Re[x/(1-x)]$, $|x|=p^{-1/2}$,
changes sign, so positivity lives in cross-prime interference.  What remains unexplored is the
\emph{gauge structure}: unconditional positivity at large height for each fixed $N$ (band closure)
against geometric negativity near an off-line zero ($\neg$RH).  Between them lies only a quantifier
exchange, whose natural tool is almost-periodicity of the prime side.

\subsubsection*{The attack: almost-periodic gauge transport}

\begin{lemma}[gauge transport of the arithmetic form]\label{lem:transport}
Fix $c\in\C^N$ and gauges $z_0=t_0-iy$, $z_0''=t_0''-iy$ with $N\le N_*(\min(t_0,t_0''))$.  Then
\[
   c^*J_N(t_0'')c
   =
   c^*J_N(t_0)c+(L''-L)\,c^*S_Nc
   +\mathcal E_{\rm arch}+\mathcal E_{\rm prime},
\]
where $|\mathcal E_{\rm arch}|\le C(y)\bigl(\tfrac{\log t_0}{t_0}+\tfrac{\log t_0''}{t_0''}\bigr)
c^*(LS_N)c$ \textnormal{[modulo Lemma~\ref{lem:signed-arch}]}, and
\[
   |\mathcal E_{\rm prime}|
   \le
   \Bigl(\sup_{n\le X}\bigl|n^{-it_0''}-n^{-it_0}\bigr|\Bigr)\,C_1(N)\,\|c\|^2
   +C_2\,X^{-(y-\frac12)}{\rm poly}(N)\,\|c\|^2 .
\]
\end{lemma}

\begin{lemma}[geometric detection]\label{lem:geom-detection}
Under $\neg$RH with an off-line zero $z_w=\gamma_w-i\delta_w$, at the gauge $t_0=\gamma_w$ there is a
witness $c$ with $c^*J_Nc\le-m_N$, $m_N\asymp|b|^{2N}$, $|b|=\frac{y+\delta_w}{y-\delta_w}>1$ ---
geometric growth in $N$ (Proposition~\ref{prop:blaschke-rate}, validated to five figures on
Davenport--Heilbronn).
\end{lemma}

The intended contradiction runs as follows.  By Kronecker, the recurrence times of the phases
$\{n^{-it}\}_{n\le X}$ to precision $\varepsilon$ are syndetic with gap $\ell(X,\varepsilon)$.  Transport
the negative witness of Lemma~\ref{lem:geom-detection} to a gauge $t_0''$ with paired phases, away from
the quartet: there the zero side gives $c^*J_N(t_0'')c\ge-C$ (the local real zeros give a Gram form
$\succeq0$; the quartet from distance $\Delta\ge\sqrt N$ contributes $e^{2Ny\delta_w/\Delta^2}=O(1)$),
while Lemma~\ref{lem:transport} gives $c^*J_N(t_0'')c\le-m_N+(L''-L)c^*S_Nc+\mathcal E$.  If the
geometric $m_N$ dominates the polynomial errors and the controlled term $(L''-L)c^*S_Nc$ (chosen via
$t_0''\in[t_0,2t_0]$), the two estimates collide and RH follows.

\subsubsection*{Autopsy: the recurrence--detection imbalance}

The attack dies at a quantified point.  The $N$-jets see primes up to $X=e^{cN}$, and the margin
requires $\varepsilon\le m_Ne^{-CN}$.  The Kronecker recurrence gap is
\[
   \ell(X,\varepsilon)\;\asymp\;\varepsilon^{-\pi(X)}\;=\;\exp\!\bigl(e^{cN}\cdot C'N\bigr)\;=\;
   \operatorname{Tower}(N),
\]
while the detection gain $m_N=e^{c\delta_wN}$ is only geometric.  The admissible window $[t_0,2t_0]$
contains a recurrence time only if $t_0\ge\ell(X,\varepsilon)$, i.e.\ $\gamma_w\ge
\operatorname{Tower}(\log\gamma_w/\delta_w)$, unsatisfiable for every $\gamma_w$: full phase matching
costs a tower, detection pays geometric, and the two never meet.  One need not match all phases, only
force a single scalar $F(t):=c^*P_N(t)c$ to return near its value at $t_0$ --- a level-return of an
almost-periodic function, cheap (positive density) if the value is not extremal.  Whether $F(t_0)$ is
extremal is decided by the following rigidity.

\begin{proposition}[ceiling rigidity]\label{prop:ceiling}
Fix $c$ and let $F(t):=c^*P_N(t)c$, an almost-periodic--structured function of the gauge with frequency
module $\{\log n\}$.  Whenever the zeros within the detection window of $t$ are real,
\[
   F(t)\;\le\;c^*A_N(t)c\;+\;O(\textnormal{far-pair tails}),
\]
because $c^*J_N(t)c=\sum_\gamma|R_c^{(t)}(\gamma)|^2\ge0$ is a Gram form.  The event
$F(t)>c^*A_N(t)c$ --- the prime form exceeding the archimedean ceiling --- occurs \emph{exactly} at
gauges detecting an off-line zero.  Hence the super-ceiling level set of $F$ coincides with the
(off-line) exceptional set itself.
\end{proposition}

Thus $F(t_0)=c^*A_Nc+m_N$ sits \emph{above} the ceiling: it is not a typical fluctuation of $F$ but the
detection event itself, and the returns of the required super-ceiling deviation are, one for one, other
off-line zeros.  The recurrence argument is exactly equivalent to producing the object whose
nonexistence it seeks to prove.  This is difficulty conservation with an explicit mechanism.

\begin{theorem}[W4: recurrence--detection imbalance]\label{thm:W4}
No proof of $\Omega_7$ can proceed by gauge transport of the arithmetic form.  Full phase-matching
transport requires recurrence times of size $\exp(\pi(e^{cN})\log(1/\varepsilon))$, super-exponentially
larger than the geometric detection gain $|b|^{2N}$; and scalar-level transport requires returns of a
super-ceiling deviation of $c^*P_N(t)c$, which by Proposition~\ref{prop:ceiling} occur exactly at
off-line zeros.  Gauge-recurrence arguments are therefore excluded alongside
\textnormal{W1--W3} (margins, truncations, zero-location and pointwise inputs).  In the classification of
\S\ref{sec:walls} this is $I(\mathcal I;C)$ with $\mathcal I=$ gauge-transport/recurrence arguments and
$C=\Omega_7$.
\end{theorem}

The classical landing is the most useful by-product: the question ``how often can
$F(t)=c^*P_N(t)c$ be large?'' is, in classical coordinates, the theory of \emph{large values of
Dirichlet polynomials} (Montgomery, Huxley, and the current front Guth--Maynard 2024).  Large-value
estimates bound the \emph{measure} of the bad-gauge set; by this route the apparatus re-derives
zero-density--type theorems, not RH --- exactly what difficulty conservation predicts, and now with the
precise classical problem attached to each closing move (scalar transport $\Longleftarrow$ optimal
large values $\Longleftarrow$ a density hypothesis $\Longleftarrow\cdots\Longleftarrow$ RH).  This
yields a concrete achievable target, sharper than the deterministic $\log\log$ range of
Corollary~\ref{cor:loglog-effective}:
\[
   \delta_N(t_0)\ge0\ \text{ for all }t_0\in[T,2T]\text{ except measure }O(T^{1-c}),\quad
   \text{uniformly for }N\le(\log T)^\theta,
\]
provable with existing large-value technology and offered here as the most immediately achievable
incremental result of this line of attack.

\subsubsection*{Round two: transplanting the finite-field proof}

RH for curves over finite fields is proved (Weil) through positivity of the Hodge index on
$C\times C$ --- the Castelnuovo inequality --- and Weil positivity for $\zeta$ is its numerical shadow.
Executing the dictionary in the present coordinates, every piece transfers except one, and it is always
the same piece.

\begin{center}\small
\begin{tabular}{@{}p{0.44\linewidth}p{0.5\linewidth}@{}}
\toprule
Function field (proved) & Present framework \\
\midrule
Correspondence algebra, $\mathrm{Fr}_n\circ\mathrm{Fr}_m=\mathrm{Fr}_{nm}$ & shift semigroup
$S_{\log n}S_{\log m}=S_{\log nm}$ \; (exists) \\
Lefschetz fixed-point formula & explicit formula (Theorem~\ref{thm:explicit-channel}) \; (exists) \\
Trivial classes $H^0,H^2$ (fibre/plane) & the two pole atoms at $\pm i/2$ (Lemma~\ref{lem:pole-pair})
\; (exists) \\
Projecting off the trivial plane & Schur shorting of the pole vector \; (exists) \\
Intersection form on $H^1$ & the form $J_N$ \\
Ample class $\theta$ & the archimedean form $\mathcal A$ \\
\textbf{Riemann--Roch with effectivity} (Castelnuovo engine) & $\varnothing$ --- does not exist \\
\bottomrule
\end{tabular}
\end{center}

\begin{definition}[closure specification]\label{def:closure-spec}
A \emph{positivity-bearing realization} is a graded object $H=H^0\oplus H^1\oplus H^2$ with a pairing
$\langle\cdot,\cdot\rangle$ and an algebra action of the shift semigroup, such that:
\begin{enumerate}[label=\textnormal{(X\arabic*)}]
\item the semigroup $\{S_{\log n}\}$ acts as correspondences with
$S_{\log n}\circ S_{\log m}=S_{\log nm}$;
\item the pairing of correspondences is computed by the Weil explicit formula
(Theorem~\ref{thm:explicit-channel}), with $\dim H^0=\dim H^2=1$ realized by the pole atoms at
$\pm i/2$;
\item \textnormal{(effectivity/Riemann--Roch)} every class $\alpha\perp H^0\oplus H^2$ with
$\langle\alpha,\alpha\rangle>0$ admits an effective representative, and effective classes pair
nonnegatively against the archimedean ample class $\theta=\mathcal A$;
\item the finite jets $J_N(z_0)$ are compressions of the $H^1$-pairing.
\end{enumerate}
\end{definition}

\begin{theorem}[conditional closure]\label{thm:conditional-closure}
If a realization satisfying \textnormal{(X1)--(X4)} exists, then $\delta_N(z_0)\ge0$ for every $N$,
hence $\Omega_7$, hence \textnormal{RH}.
\end{theorem}
\begin{proof}[Proof (Castelnuovo transplant)]
Suppose $c^*J_Nc<0$ for some $N,c$.  By \textnormal{(X4)} this lifts to a class $\alpha\in H^1$ with
$\langle\alpha,\alpha\rangle<0$ after shorting the trivial plane \textnormal{(X2)}; the Hodge-index
argument of Weil, powered by \textnormal{(X3)} exactly as effectivity powers the
Castelnuovo--Severi inequality on $C\times C$, forces
$\langle\alpha,\alpha\rangle\ge-\langle\alpha,\theta\rangle^2/\langle\theta,\theta\rangle$ with the
trivial-plane contribution already removed, yielding $\langle\alpha,\alpha\rangle\ge0$; contradiction.
\end{proof}

The consistency tests pass non-trivially: Davenport--Heilbronn fails \emph{exactly} at (X1) (no Euler
product, no correspondence semigroup), so the falsifier discriminates at the correct axiom rather than
by accident; the sine model satisfies (X1)--(X2) trivially with $H^1$ of controlled rank (empty cascade);
and any (X3) candidate is falsifiable in hours against the measured Cholesky flow and factorial defect
law.  Two cross-coordinate attempts confirm the same gap.  The \emph{Nyman--Beurling} pincer
($\mathrm{RH}\iff d_N\to0$, a twin arithmetic Gram of Vasyunin sums) collapses because every
inter-frame inequality factors through the common spectral measure, i.e.\ through the zeros; it does,
however, supply the M\"obius mollifier as an explicit falsification test (a certificate must
reproduce Burnol's rate $d_N^2\sim C/\log N$).  The \emph{multiplicative-independence/large-values}
attempt reproduces the achievable measure-theoretic target above, not RH.

\subsubsection*{Creating (X3): the collapse obstruction and the \texorpdfstring{$\theta$}{theta}-square}

\begin{lemma}[the collapse obstruction]\label{lem:collapse}
In commutative rings, $\Z\otimes_\Z\Z=\Z$, hence
$\operatorname{Spec}\Z\times_{\operatorname{Spec}\Z}\operatorname{Spec}\Z=\operatorname{Spec}\Z$: the
diagonal is an isomorphism and the square degenerates to the curve itself.  Any realization of
\textnormal{(X3)} must therefore first change the base category (a base ``below'' $\Z$).  The gap is
not theorem-shaped but \emph{definition-shaped}: before proving Riemann--Roch on the square, the square
must exist.
\end{lemma}

Weil had the surface $C\times_{\mathbb F_q}C$ for free; here it must be manufactured.  Three known
substrates de-collapse the product, and the falsification battery discriminates among them.

\begin{center}\small
\begin{tabular}{@{}p{0.2\linewidth}p{0.08\linewidth}p{0.1\linewidth}p{0.28\linewidth}p{0.24\linewidth}@{}}
\toprule
Substrate & X1 & X2 & X3 (effectivity) & Test vs.\ our data \\
\midrule
Connes--Consani (scaling site, char.\ $1$) & $\surd$ & $\surd$ (trace formula) & partial: tropical
RR exists, but idempotent semirings have no subtraction $\Rightarrow$ no \emph{signature} (Hodge index
needs an abelian category) & restricted-support positivity $=$ our finite-$N$ regime \\
Borger ($\lambda$-rings / Witt: the square as $W(\Z)$ with two Frobenius families) & $\surd\surd$
(the $\psi_p$ commute: our shift semigroup) & unproved & no finite cohomology theory & right structure,
no engine \\
Deninger (foliated dynamical system) & $\surd$ (periodic orbits of length $\log p=$ our lags) &
conjectural & needs leafwise Hodge theory & his flow $=$ our gauge translation; ceiling rigidity (W4) $=$
his recurrence structure \\
\bottomrule
\end{tabular}
\end{center}

The construction starts from the one thing that already exists exactly in dimension one: arithmetic
Riemann--Roch by theta invariants (van der Geer--Schoof, Bost).  For a metrized line bundle $\bar L$
over $\overline{\operatorname{Spec}\Z}$,
\[
   h^0_\theta(\bar L):=\log\sum_{x\in L}e^{-\pi\|x\|^2},\qquad
   h^0_\theta(\bar L)-h^0_\theta(\bar\omega-\bar L)=\deg\bar L\quad\text{(exact)},
\]
where Serre duality is literally Poisson summation and the effectivity engine is Minkowski ($\deg$ large
$\Rightarrow$ short vector $\Rightarrow$ effective section).  In dimension one the full Weil dictionary
is thus already proved, with Poisson as duality and geometry of numbers as effectivity; (X3) is its lift
to the square.  The divisors of the square have a natural candidate here: correspondences are the
Dirichlet polynomials (finite shift combinations), the ample class is $\theta=\mathcal A$, and the (X2)
pairing is Poisson along the diagonal.  Three definitions of ``effective'' were tested.  \emph{(i) The
cone of real-point evaluation atoms} is circular (it presupposes real zeros).  \emph{(ii) The measured
Cholesky factor $B_N(t_0)$} exists $\iff$ positivity (presupposes the conclusion), though it records the
shape an effective section would take.  \emph{(iii) The prime-shift cone}
$E:=\operatorname{cone}\{\Lambda(n)n^{-1/2}S_{\log n}\}\cup(\text{positive archimedean cone})$ survives
the circularity tests: it is defined from $\Lambda(n)\ge0$ alone, with no zeros, and discriminates
Davenport--Heilbronn at the correct axiom.  With this cone, (X3) becomes a single statement of
geometry-of-numbers type.

\begin{conjecture}[M2: Minkowski for correspondences]\label{conj:M2}
Let $D$ be a correspondence class orthogonal to the trivial plane $H^0\oplus H^2$ (the pole atoms at
$\pm i/2$) with $\langle D,D\rangle_W>0$ under the Weil pairing.  Then the class of $D$, modulo the
trivial plane, contains an effective representative in the shift cone $E$; and every element of $E$
pairs nonnegatively against the ample class $\theta=\mathcal A$.
\end{conjecture}

\begin{theorem}[conditional closure, sharpened]\label{thm:M2-closure}
\textnormal{(X1)}, \textnormal{(X2)}, \textnormal{(X4)} \textnormal{(all existing)} together with
\textnormal{M2} imply \textnormal{(X3)}, hence $\delta_N\ge0$ for all $N$, hence $\Omega_7$, hence
\textnormal{RH} --- by the Castelnuovo transplant of Theorem~\ref{thm:conditional-closure}.
\end{theorem}

\begin{problem}[M2$_N$: finite falsifiable shadow]\label{prob:M2N}
In the $N$-jet compression at gauge $z_0$: does every $c\perp(\textnormal{pole plane})$ with
$c^*J_Nc>0$ admit a decomposition $c=e_+-e_-$ with $e_\pm$ in the compressed shift cone $E_N$ and
$\langle e_+,e_+\rangle\ge\langle c,c\rangle-\varepsilon_N$?  This is a conic-feasibility problem
(SDP), decidable at $N\le12$ in hours.  Protocol: \textnormal{(i)} verify the ampleness pairings
$\langle\Lambda(n)n^{-1/2}S_{\log n}\,v,\ \mathcal A\,v\rangle\ge0$ numerically on the Laguerre frame;
\textnormal{(ii)} run feasibility on $\zeta$ across gauges; \textnormal{(iii)} confirm structural
failure on Davenport--Heilbronn at the cone-existence step; \textnormal{(iv)} check that the extremal
$e_+$ reproduces the measured Cholesky flow.
\end{problem}

\subsubsection*{A candidate construction of (X3): thermal Riemann--Roch (Gibbs effectivity)}

We record one concrete proposal for the missing effectivity engine (X3).  Everything in this
subsubsection has \emph{conjectural} status; its value is that it is precise enough to be attacked, and
it comes with the map of where to attack it.

\emph{The idea.}  Every attempt to build the square
$\operatorname{Spec}\Z\times_{\mathbb F_1}\operatorname{Spec}\Z$ as a \emph{geometric} object dies in the
collapse $\Z\otimes\Z=\Z$ (Lemma~\ref{lem:collapse}).  The proposal here is that the square is not a
space but a \emph{thermodynamic system}: its points are KMS states, its Serre duality is temperature
inversion $\beta\mapsto1/\beta$, its trivial plane $H^0\oplus H^2$ is a phase transition, and --- the
core --- effectivity is not a count of sections but \emph{positivity of free energy}.  The role played by
Minkowski in dimension one (and by Riemann--Roch in Weil's proof) is taken by the unconditional Gibbs
variational principle; the positivity engine ceases to be spectral (where every argument is circular) and
becomes entropic (where positivity is structural --- a variance).

\emph{The objects.}  \textnormal{(D1)} The Frobenius semiring $\{[n]\}_{n\ge1}$, $[m][n]=[mn]$, is the
shift semigroup; its analytic carrier already exists as the Bost--Connes system, whose
$\mathrm{KMS}_\beta$ states, twisted by the archimedean weight $\kappa$ (below), are the proposed family
$\omega_\beta$.  Bost--Connes has a phase transition at $\beta=1$ with partition function $\zeta(\beta)$.

\begin{definition}[thermal sections and effectivity]\label{def:thermal-effectivity}
For a finitely supported real class $D=\sum_n c_n[n]$, define the free energy
\[
   F_\beta(D):=\log\,\omega_\beta\!\left(e^{D}\right),
   \qquad e^{D}:=\sum_{k\ge0}\frac{D^k}{k!}\ \ \text{(well-defined by multiplicativity of }[n]\text{)}.
\]
$D$ is \emph{effective} iff $F_\beta(D)\ge0$ for all $\beta$ in the self-dual window.  The effective cone
is multiplicatively closed by construction (H\"older on states).
\end{definition}

\textnormal{(D3)} Riemann--Roch as a theta flip: the duality $\beta\mapsto1/\beta$ is implemented not as a
symmetry of Bost--Connes (it is not one) but by the adelic Fourier transform inside the GNS space ---
Tate's local functional equations as local RR, assembled into
\[
   F_\beta(D)-F_{1/\beta}(\kappa-D)=\deg D,\qquad \kappa=\mathcal A\ \text{(the ample archimedean class)},
\]
with $\deg$ read off the residue of the Bost--Connes transition.

\begin{conjecture}[Gibbs realization of the Weil pairing]\label{conj:X3G}
There is a normalization $c(\beta)>0$ and a heat-regularization $D\mapsto\widehat D$ such that, for every
real class $D$ orthogonal to the trivial plane,
\[
   \langle D,D\rangle_W \;=\; -\,\lim_{\beta\to1^+} c(\beta)\,\mathrm{Var}_{\omega_\beta}\!\big(\widehat
   D\big)\;\le\;0 .
\]
That is: on primitive classes, the negative of the Weil pairing \emph{is a variance} of the twisted KMS
state.  Combined with the two-dimensional positive trivial plane, this is the Hodge index statement, and
\textnormal{RH} follows by the Castelnuovo transplant (Theorem~\ref{thm:conditional-closure}).
\end{conjecture}

The structural non-circularity, which distinguishes this from the buried mechanisms, is that a variance
is $\ge0$ by Cauchy--Schwarz, not by spectral information, and $\omega_\beta$ is defined from purely
multiplicative and thermal data (the Bost--Connes algebra plus the KMS condition) with no zeros in any
input: it passes W2 by design.

\emph{Three structural coincidences suggesting novelty.}
\textnormal{(i)} \emph{Temperature dictionary.}  Bost--Connes KMS weights scale as $n^{-\beta}$, our
effective weights as $n^{-1/2-y}$, giving $\beta=\tfrac12+y$; the boundary limit $y\to\tfrac12^+$ where
all the programme's difficulty lives is \emph{exactly} the Bost--Connes phase transition
$\beta\to1^+$, and the pole of $\zeta$ as a divergent partition function coincides with the pole pair
$\pm i/2=H^0\oplus H^2$ appearing at symmetry breaking.
\textnormal{(ii)} \emph{Native multiplicativity.}  The measured $\sim1/\delta_N$ cancellation in the
\emph{linear} cone (sums do not buy effectivity) is repaired by $F_\beta=\log\omega(e^D)$, which is
log-convex and multiplicatively closed --- the closure the data demanded.
\textnormal{(iii)} \emph{Torus lift.}  The characters of $\Q_{>0}^*$ form the infinite torus
$\mathbb T^\infty$ (one angle per prime), and the explicit-formula line is the Kronecker flow inside it;
the conjecture is equivalent to the statement that the Weil distribution is the restriction to the
Kronecker line of a positive kernel on $\mathbb T^\infty$.  Since the vanishing of $\Lambda$ at mixed
integers $p^aq^b$ forces the mixed cumulants to vanish, the kernel is quasi-product (a Gaussian field
over independent coordinates), and the whole content compresses into a single lemma: the archimedean
term (the only non-product coupling) lifts to the torus diagonal dominating the per-prime negatives.

\emph{The load map.}
\textnormal{(L1, pencil and paper, start here)} Compute the cumulants of the regularized $[p]$ in
$\omega_\beta$ from the known Bost--Connes KMS formulae and compare with the explicit-formula
coefficients; kill-test: if the candidate's mixed cumulants do not vanish at non-prime-powers, the state
family is wrong and is discarded in an afternoon.
\textnormal{(L2, the core, where circularity would hide)} The twisted state
$\omega_\beta^\kappa(\cdot):=\omega_\beta(e^{\kappa/2}\cdot e^{\kappa/2})/\omega_\beta(e^\kappa)$ must
produce the archimedean density as a diagonal conditional variance; the exact place to look for the trap
is the positivity of the $\kappa$-twist --- if proving that $\omega_\beta^\kappa$ is a state turns out to
be equivalent to $W\le0$, the proposal dies the standard death, and that is the first
theorem-or-counterexample to attempt.
\textnormal{(L3, the weakest joint)} Bost--Connes lacks the flip symmetry; the proposal derives it from
adelic Poisson à la Tate inside the GNS space, to be proved first in the semilocal $\{p,\infty\}$ case
(two places, fully explicit, with the Connes--Consani semilocal trace formulae as control) --- this baby
case is the first complete attainable theorem of X3-G.  Suggested order: L1 $\to$ one-prime twist
positivity (L2) $\to$ full semilocal (L3), each with a kill-test of days, none requiring code.


\subsubsection*{X3-G, round two: three free lemmas, the corrected flip, and the archimedean partition problem}

The three design constraints of the recalibration are adopted as such: (C1) the hard core is the identity
$\langle D,D\rangle_W=-c\,\mathrm{Var}(\widehat D)$, not the positivity of the state; (C2) the flip
requires resolution across the phase transition; (C3) the regularization $\widehat D$ and the effective
window must be built explicitly.  A technical correction refines (C2): in Bost--Connes the
$\mathrm{KMS}_\beta$ state is unique for $\beta\in(0,1]$ --- unique at $\beta=1$ exactly --- and the
Galois simplex opens only for $\beta>1$.  So $\omega_1$ is well defined; the flip difficulty is instead
that $\beta\leftrightarrow1/\beta$ carries the unique-state region to the simplex region, hence cannot be
a system isomorphism.  This is used below as data, not obstacle.

\emph{Realization.}  In the type-I representation ($\beta>1$): $\mathcal H=\ell^2(\N)$, $He_n=(\log
n)e_n$, $\mu_me_n=e_{mn}$, and $[n]:=\mu_n+\mu_n^*$ is the shift $S_{\log n}$ verbatim; the Gibbs state is
$\omega_\beta(\cdot)=\mathrm{Tr}(e^{-\beta H}\cdot)/\zeta(\beta)$.

\begin{lemma}[raw covariance is exactly product]\label{lem:bc-covariance}
For $m,n\ge2$, $\ \omega_\beta([n])=0$ and
$\mathrm{Cov}_{\omega_\beta}([m],[n])=\delta_{mn}\,(1+n^{-\beta})$.
\end{lemma}
\begin{proof}
$\mathrm{Tr}(e^{-\beta H}\mu_m\mu_n^*)=\sum_k k^{-\beta}\langle e_k,\mu_m e_{k/n}\rangle[n\mid k]
=\delta_{mn}\sum_{n\mid k}k^{-\beta}=\delta_{mn}n^{-\beta}\zeta(\beta)$, so
$\omega_\beta(\mu_m\mu_n^*)=\delta_{mn}n^{-\beta}$; likewise $\omega_\beta(\mu_m^*\mu_n)=\delta_{mn}$; and
$\omega_\beta(\mu_{mn})=\sum_kk^{-\beta}\delta_{mnk,k}=0$.  Expanding
$[m][n]=(\mu_m+\mu_m^*)(\mu_n+\mu_n^*)$ into its four terms and using $\omega_\beta([n])=0$ gives
$\mathrm{Cov}=\delta_{mn}(1+n^{-\beta})$.
\end{proof}

\begin{lemma}[no-go for diagonal twists]\label{lem:bc-nogo-diagonal}
If $\kappa=f(H)$, the twisted state $\omega_\beta^\kappa$ still has
$\mathrm{Cov}([m],[n])\propto\delta_{mn}$.  Hence the archimedean class $\kappa$ cannot be a function of
the number operator: it must be affiliated to the transverse (flow) variable, i.e.\ the twist lives on
the adelic GNS space, not on $\ell^2(\N)$.
\end{lemma}
\begin{proof}
With $\kappa=f(H)$ diagonal, $\mathrm{Tr}(e^{-\beta H}e^{\kappa/2}\mu_m\mu_n^*e^{\kappa/2})=\sum_k
k^{-\beta}e^{f(\log k)}\langle e_k,\mu_m\mu_n^*e_k\rangle\propto\delta_{mn}$: no diagonal weight creates
off-diagonal correlation.
\end{proof}

\begin{lemma}[twist positivity is free]\label{lem:bc-twist-positivity}
For any admissible $\kappa$, $\omega_\beta(e^{\kappa/2}\,\cdot\,e^{\kappa/2})$ is positive (it is
$\omega(B^*\!\cdot B)$), modulo domain questions only.
\end{lemma}
\begin{proof}
$A\ge0\Rightarrow B^*AB\ge0$.
\end{proof}

Three readings.  Lemma~\ref{lem:bc-covariance}: the raw state is \emph{exactly} the product kernel ---
zero arithmetic content, so all of the Weil information must be created by the twist.
Lemma~\ref{lem:bc-nogo-diagonal}: the archimedean twist is \emph{forced} onto the variable conjugate to
the flow, which is exactly where $\mathcal A$ was built (on the boundary flow, not the number operator) ---
the structure pushes towards the correct geometry by itself.  Lemma~\ref{lem:bc-twist-positivity}
discharges (C1): the difficulty lands entirely on the moment identity, not on the positivity of the state.

\emph{The corrected duality (flip with Galois averaging).}  Since the flip crosses
unique~$\leftrightarrow$~simplex, define the symmetrized free energy
\[
   \widehat F_\beta(D):=
   \begin{cases}
      F_\beta(D), & \beta\le1,\\[3pt]
      \displaystyle\int_{\hat\Z^\times}F_\beta^{\sigma}(D)\,d\sigma\ \ (\text{Haar over the extremal
      KMS states}), & \beta>1,
   \end{cases}
\]
and thermal Riemann--Roch v2: $\widehat F_\beta(D)-\widehat F_{1/\beta}(\kappa-D)=\deg D$.  The average
is not a patch: it is the assertion that Serre duality of the arithmetic square passes through averaging
over $\mathrm{Gal}(\Q^{\rm ab}/\Q)$ --- the $H^2$ side carries the Galois torsor.  It remains the weakest
joint, but both sides are now canonical.  Explicit regularization:
$\widehat D_\varepsilon:=e^{-\varepsilon H}De^{-\varepsilon H}$, window $\beta\in(1,1+\eta)$ open to the
right of $1$ --- criticality is reached as a limit, never touched; everything converges in the window
because $\zeta(\beta)<\infty$, discharging (C3) with no hidden collapse.

\emph{The two-prime audit, executed.}  \textnormal{(Rigid layer, a theorem.)}  For purely almost-periodic
kernels (the atomic part), Bochner on compacts together with uniqueness of Bohr coefficients along the
dense winding give: positive-definiteness on $\mathbb T^\infty$ $\iff$ positive-definiteness on the
Kronecker line.  The suspected conclusion is thus confirmed exactly: for the atomic part the torus is a
change of coordinates, no new engine.  \textnormal{(The remainder.)}  The archimedean density is
\emph{not} almost-periodic (its Bohr coefficients vanish), so it does not live on $\mathbb T^\infty$; the
correct home of the lift is $\R\times\mathbb T^\infty$ --- the adele class space --- and there the
positive-definite completion freedom is exactly one: the transverse profile with which the archimedean
mass is distributed among the prime coordinates.  This is the one new dial.

\begin{problem}[archimedean partition problem]\label{prob:arch-partition}
Find nonnegative densities $\{\sigma_p\}$ with $\sum_p\sigma_p\le\mathcal A$ (the archimedean density)
such that for \emph{every} prime $p$ the single-frequency kernel
\[
   W_p(t)\;:=\;\sigma_p(t)\;-\;\sum_{k\ge1}\Lambda(p^k)\,p^{-k/2}\cos(kt\log p)\ (\mathrm{reg})
\]
is positive definite on $\R$.  Existence $\Rightarrow$ $W=\sum_pW_p+(\text{leftover}\ge0)$ is positive
definite $\Rightarrow$ Weil positivity $\Rightarrow$ \textnormal{RH}.
\end{problem}

Each subproblem is classical: a single-frequency train against a density is one-dimensional
Herglotz/Toeplitz, completely characterized.  The global problem is a resource allocation with LP/SDP
duality structure, and the dual certificate is a new falsification format.  The H0 ledger is unadorned:
by the prime number theorem the total prime demand exhausts the archimedean density at leading order, so
the margin is of RH size --- $\Omega_7$ is conserved, as the audit required.  What changed is the format:
from a monolithic spectral positivity to (i) a moment identity with a free variance, plus (ii) an
allocation with an explicit dual and classical subproblems.  One measured datum already speaks: W3
showed that the \emph{uniform} partition fails (the per-prime symbol changes sign), so the partition, if
it exists, is necessarily \emph{adaptive} --- $\sigma_p$ loaded near the resonances of $p$ and emptied
far from them.  This is the first empirical constraint on the ansatz.

\emph{Revised work order.}  (1) L1 is half done: Lemma~\ref{lem:bc-covariance} computed the raw moments;
the twisted layer remains --- guided by Lemma~\ref{lem:bc-nogo-diagonal} (twist on the flow variable),
compute $\mathrm{Cov}^\kappa([m],[n])$ for the simplest Gaussian twist in the adelic flow generator and
compare with the explicit-formula coefficients.  (2) Single-prime partition \emph{[executed below:
Proposition~\ref{prop:W5}, Lemma~\ref{lem:perprime-feasible}]}: the density formulation fails at large
resonances for every prime (fifth wall), while every fixed-cutoff subproblem is feasible with explicit
converging mass.  (3) Global accounting \emph{[executed below]}: the stacked per-prime demand diverges,
but the true frequency-local accounting reduces exactly to $\Omega_7$.  (4) The flip (L3) only after (1)--(3) survive, and only
then the semilocal $\{p,\infty\}$ case with the Galois-averaged v2 duality.  The entire new content of
X3-G is now compressed into two named objects: the Weil--variance identity (Conjecture~\ref{conj:X3G})
and the archimedean partition problem (Problem~\ref{prob:arch-partition}), the second with a property no
earlier candidate had --- classical subproblems, a dual certificate, and an ansatz already constrained by
our own data.


\subsubsection*{The partition problem executed: a fifth wall, a per-prime feasibility theorem, and the reduction back to $\Omega_7$}

The revised work order was carried out numerically (scripts \texttt{omega7\_partition\_p2.py},
\texttt{omega7\_partition\_global.py}).  The outcome is sharp and, as $H0$ demands, closes no gap ---
but it changes the format decisively and removes one candidate route.

\emph{Exact single-prime condition.}  Testing the per-prime Weil piece against positive-definite
autocorrelations $g=h\star\tilde h$ (so $\hat g\ge0$) and dualizing, the single-prime domination
$\int\sigma_p\,g\ge\sum_k(\log p)p^{-k/2}g(k\log p)$ for all admissible $g$ is equivalent, by Bochner,
to the pointwise frequency inequality
\begin{equation}\label{eq:sigma-domination}
   \hat\sigma_p(\xi)\;\ge\;D_p(\xi):=2\sum_{k\ge1}(\log p)\,p^{-k/2}\cos\!\big(k\,\xi\log p\big)
   \qquad(\forall\xi\in\R),
\end{equation}
whose right-hand side is \emph{almost-periodic} in $\xi$ with period $2\pi/\log p$ and peak value
$D_p(0)=2\log p/(\sqrt p-1)$ at every resonance $\xi\in(2\pi/\log p)\Z$.

\begin{proposition}[fifth wall: the density formulation is per-prime infeasible]\label{prop:W5}
If $\sigma_p$ is an $L^1$ density then $\hat\sigma_p$ is a positive-definite continuous function with
$\hat\sigma_p(\xi)\to0$ as $|\xi|\to\infty$ (Riemann--Lebesgue), while the demand $D_p$ does not decay.
Hence \eqref{eq:sigma-domination} fails at all sufficiently large resonances, for \emph{every} prime and
\emph{every} finite mass budget.  The archimedean class $\sigma_p$ cannot be a density: it must be a
measure with an atomic (almost-periodic) part, or the test class must be cut off.  In the classification
of \S\ref{sec:walls} this is $I(\mathcal I;C)$ with $\mathcal I=L^1$-density archimedean allocations and
$C=$ the per-prime domination \eqref{eq:sigma-domination}.
\end{proposition}
\noindent This was verified directly (gate GA1): with ten times the peak budget and a near-atomic
Gaussian density (std $10^{-2}$), the deficit still reaches the full peak at a finite resonance
($n=86$ for $p=2$, $n=199$ for $p=5$).  W5 is the frequency-side twin of the pole-pair obstruction
(Lemma~\ref{lem:pole-pair}): a smooth archimedean object cannot, unaided, dominate an atomic prime
demand across all scales.

\emph{The cutoff regime is feasible per prime.}  Restricting test functions to support $[-R,R]$ (phase-space
cutoff, exactly the compression of Connes--Consani~\cite{ConnesConsani2021}) truncates the demand to the
finite frequency set $\{k\log p:\ k\log p\le 2R\}$.  The finite trigonometric-moment problem
\eqref{eq:sigma-domination} is then solvable, and the linear program minimizing $\int\sigma_p$ returns a
bounded density whose mass \emph{converges} in $R$:

\begin{center}\small
\begin{tabular}{lrrrrr}
\toprule
 & $R=2$ & $R=4$ & $R=8$ & $R=16$ & $R=32$\\
\midrule
$p=2$ (peak $3.35$) & $2.76$ & $3.32$ & $3.60$ & $3.72$ & $3.79$\\
$p=3$ (peak $3.00$) & $2.43$ & $2.94$ & $3.00$ & $3.24$ & $3.48$\\
$p=5$ (peak $2.60$) & $2.08$ & $2.51$ & $2.64$ & $2.60$ & $2.86$\\
\bottomrule
\end{tabular}
\end{center}
\noindent The increments halve as $R$ doubles: per prime, at any fixed cutoff, a bounded density of mass
$\lesssim1.15\cdot 2\log p/(\sqrt p-1)$ dominates (gate GA2b confirms a genuine \emph{smooth} density of
bounded height suffices; atoms are not essential at fixed $R$).  This is a clean unconditional lemma:
\emph{every single-prime Toeplitz subproblem is feasible with explicit converging mass}.

\begin{lemma}[per-prime cutoff feasibility]\label{lem:perprime-feasible}
For each prime $p$ and each cutoff $R<\infty$ there is a nonnegative density $\sigma_p^{(R)}$ with
$\int\sigma_p^{(R)}\le c_p(R)\cdot\tfrac{2\log p}{\sqrt p-1}$, $c_p(R)\uparrow c_p\le\tfrac{6}{5}$, satisfying
\eqref{eq:sigma-domination} on all frequencies $|\xi|\le2R$.  Consequently the obstruction to
$\Omega_7$ is neither per-prime nor at any finite cutoff; it lies entirely in the simultaneous limit
$R\to\infty$ over all primes.
\end{lemma}

\emph{The obstruction is global, and equals $\Omega_7$.}  Summing the per-prime peaks,
$\sum_{p\le P}2\log p/(\sqrt p-1)$ \emph{diverges} ($\sim2\sum_p\log p/\sqrt p$); the measured
cumulative demand is $13.4,\,23.6,\,41.5,\,129.6,\,405.3,\,1272.1$ for $P=11,31,101,10^3,10^4,10^5$
(gate GA2a).  A single fixed archimedean budget therefore cannot cover all primes \emph{if the demands
stacked}.  They do not: the resonances $k\log p$ sit at disjoint frequencies, so the additive sum
overcounts, and the true accounting is \emph{frequency-local} --- at each $\xi$, compare
$\hat{\mathcal A}(\xi)$ against $\sum_{p,k}(\log p)p^{-k/2}\delta$-mass near $\xi$.  That local comparison
is precisely the pointwise symbol positivity of the terminal form at the boundary, i.e.\ $\Omega_7$
itself.  So the partition reformulation \emph{is equivalent} to $\Omega_7$ (no reduction of difficulty,
as $H0$ predicts), but in a strictly sharper format: all finite subproblems are feasible
(Lemma~\ref{lem:perprime-feasible}), the failure mode is isolated (the uniform-in-$R$ constant $c_p$
must stay bounded \emph{simultaneously} across all primes as the cutoff opens), and infeasibility, when
probed at finite $R$, is certified by an explicit LP dual.

\emph{Alignment with the constructed engine.}  This frequency-side picture is dual to an operator-side
engine built elsewhere in the programme, pursuing the same obstruction through Connes's colligation
route: the regularized colligation margin that survives the critical continuation $\omega\downarrow0$
and is positive for $\zeta$ on every tested configuration, the Kre\u\i n--Langer negative-square index
$\mathrm{ind}_-=0$ for $\zeta$ that counts off-line zeros exactly, and the finite-prime colligation whose
\emph{scalar} truncation diverges at $\omega=0$ because the Euler product is a boundary object --- the
exact operator-side shadow of W5.  The two sides agree: the difficulty is uniformity across the full family
(all configurations / all primes / the open cutoff), never any single member.  The unconditional
$\sigma_\infty$ constructed by Connes--Consani at the archimedean place~\cite{ConnesConsani2021} is the
prototype of $\sigma_p^{(R)}$; and their Riemann--Roch for $\overline{\operatorname{Spec}\Z}$
\cite{ConnesConsani2023} realizes (X3) in dimension one, leaving the surface (self-intersection / Hodge
index) version as the precise residue of the closure specification (Definition~\ref{def:closure-spec}).


\subsubsection*{The ``to infinity'' formulation: a single scale-uniform envelope inequality is $\Omega_7$}

The partition analysis localizes the entire difficulty in the simultaneous limit over all scales
(Lemma~\ref{lem:perprime-feasible}).  On the Li side this limit has a strikingly clean shape.  Write,
via $\Omega_6$, $\lambda_n=\lambda_n^{\mathrm{arch}}+\lambda_n^{\mathrm{prime}}$ with the exact
Stieltjes-built split of $\log\xi$ (script \texttt{omega7\_uniform\_margin.py}, dps~$120$, exact to
$n=170$).  The archimedean part is the \emph{ceiling}
\[
   \lambda_n^{\mathrm{arch}}=\tfrac{n}{2}\log\tfrac{n}{2\pi}+O(n),
\]
growing with scaling exponent $1$; the prime part is a \emph{bounded-exponent oscillation}.  Measured
envelope $E(n):=\max_{m\le n}|\lambda_m^{\mathrm{prime}}|$:

\begin{center}\small
\begin{tabular}{lrrrrrr}
\toprule
$n$ & $10$ & $40$ & $80$ & $120$ & $150$ & $170$\\
\midrule
$\lambda_n^{\mathrm{prime}}$ & $1.32$ & $1.16$ & $0.36$ & $2.09$ & $0.68$ & $-1.01$\\
ceiling $\tfrac{n}{2}\log\tfrac{n}{2\pi}$ & $2.32$ & $37.0$ & $101.8$ & $177.0$ & $238.0$ & $280.3$\\
ratio ceiling$/|\lambda^{\mathrm{prime}}|$ & $1.8$ & $31.8$ & $284$ & $84.8$ & $350$ & $276$\\
\bottomrule
\end{tabular}
\end{center}
\noindent A least-squares fit gives $E(n)\sim n^{\alpha}$ with $\alpha=0.264$ (below $\tfrac12$), and
$E(n)/(\sqrt n\log n)$ is bounded ($\max 0.111$ over $n\le170$).  The prime oscillation is therefore
$O(\sqrt n\log n)$, dominated by the ceiling with a margin whose exponent is $1-\alpha\approx0.74>0$:
the ceiling wins at every scale, and $\lambda_n>0$ for all $n\le170$ with a margin that grows without
bound.  This is the honest ``to infinity'' picture the programme was missing: our (X3) does not merely
match finitely many primes (as in the primes-$<13$ reconstruction of~\cite{Connes2026}); the
archimedean growth outpaces the prime oscillation at every scale, provided one uniform bound holds.

\begin{conjecture}[the scale-uniform envelope bound $=$ (X3) to infinity]\label{conj:envelope}
There is $c>0$ with $|\lambda_n^{\mathrm{prime}}|\le c\,\sqrt n\,\log n$ for all $n$.  Equivalently
(Bombieri--Lagarias), $\lambda_n=\tfrac n2\log\tfrac{n}{2\pi}+O(\sqrt n\log n)$; and since the ceiling
dominates any $O(\sqrt n\log n)$ term for $n$ beyond an explicit $n_0$, together with the direct
verification $\lambda_n>0$ for $n\le n_0$ this yields $\lambda_n\ge0$ for all $n$, i.e.\ \textnormal{RH}.
\end{conjecture}

\begin{remark}[the finite window is closed unconditionally]\label{rmk:finite-window}
The implication is fully explicit and its finite part is RH-free.  With the measured constant
$c=\max_{n\le120}|\lambda_n^{\mathrm{prime}}|/(\sqrt n\log n)=0.986$, the rigorous ceiling minorant
$\tfrac n2\log\tfrac{n}{2\pi}-\tfrac n2$ exceeds $c\sqrt n\log n$ for all $n\ge n_0=51$
(script \texttt{omega7\_finite\_window.py}), and $\lambda_n>0$ is verified \emph{exactly} (from the
Stieltjes split, no zeros used) for $n=1,\dots,51$.  Hence the entire content of ``envelope $\Rightarrow$
RH'' is the single asymptotic inequality of Conjecture~\ref{conj:envelope}; the small-$n$ regime carries
none of the difficulty.
\end{remark}

The prize of this reformulation: $\Omega_7$ is reduced to a \emph{single classical, scale-uniform
inequality} on the explicit prime-built sequence $\lambda_n^{\mathrm{prime}}=-\!\sum_m\Lambda(m)m^{-1}
L_{n-1}^{(1)}(\log m)$ (continued) --- no surface, no $F_1$, no Hodge index in the statement.  The
attack is analytic: in the oscillatory (Plancherel--Rotach) regime
$L_{n-1}^{(1)}(\log m)\sim (\text{amplitude})\cos\!\big(2\sqrt{n\log m}-\theta\big)$, so
$\lambda_n^{\mathrm{prime}}$ is a stationary-phase sum of $\Lambda(m)m^{-1}$ against a chirp; the
$\sqrt n\log n$ envelope is exactly the square-root cancellation of that sum, and proving it is the
prime-number-theorem error term at RH strength.  The falsifier discriminates: for a function with an
off-line zero $\rho=\beta+i\gamma$ the prime side acquires a term $\sim n^{\beta'}$
($\beta'=\max(\beta,1-\beta)>\tfrac12$) that eventually overtakes the ceiling and drives $\lambda_n<0$
--- exactly the Davenport--Heilbronn behaviour already isolated in the harness.  So
Conjecture~\ref{conj:envelope} is faithful: it holds iff RH, and its failure mode is the detected
off-line pole.

\emph{The stationary-phase reduction is exact --- and terminates at RH with no slack.}  Carrying the
classical attack to its end pins the irreducible step precisely.  Write the oscillation on the zero
side,
\[
   \lambda_n=(\text{arch.\ ceiling})-O(n),\qquad O(n):=\sum_\rho\Big(1-\tfrac1\rho\Big)^{\!n}.
\]
For a zero on the critical line $\rho=\tfrac12+i\gamma$ one has $1-\tfrac1\rho=e^{i\theta(\gamma)}$ with
$\theta(\gamma)=2\arctan\tfrac{1}{2\gamma}$ \emph{real} and $|1-\tfrac1\rho|=1$ \emph{exactly}, so
$(1-\tfrac1\rho)^n=e^{in\theta(\gamma)}$ is pure oscillation; pairing $\rho$ with $\bar\rho$ and using
$\theta(\gamma)\sim1/\gamma$ together with the density
$dN(\gamma)=\tfrac{1}{2\pi}\log\tfrac{\gamma}{2\pi}\,d\gamma$,
\[
   O(n)\approx2\,\mathrm{Re}\!\int e^{in/\gamma}\,\tfrac{1}{2\pi}\log\tfrac{\gamma}{2\pi}\,d\gamma .
\]
The phase $\varphi(\gamma)=n/\gamma$ has $\varphi'(\gamma)=-n/\gamma^2\neq0$ (no stationary point); the
substitution $u=n/\gamma$ exhibits $\asymp\sqrt n$ effective oscillations against a $\log$ weight, and
the square-root cancellation gives $|O(n)|\asymp\sqrt n\log n$.  The mechanism \emph{works}.

\begin{proposition}[exact reduction of the envelope bound]\label{prop:envelope-reduction}
The stationary-phase estimate above is lossless: its sole load-bearing hypothesis is $\theta(\gamma)
\in\R$ for every zero, i.e.\ $|1-\tfrac1\rho|=1$ for every $\rho$.  If some zero has $\rho=\beta+i\gamma$
with $\beta\neq\tfrac12$, then $(1-\tfrac1\rho)^n=r^n e^{in\theta}$ with
$r=|1-\tfrac1\rho|=\big(1+\tfrac{1-2\beta}{|\rho|^2}\big)^{1/2}$; the functional-equation partner with
$\beta<\tfrac12$ has $r>1$, so its term grows geometrically and no cancellation can hold, whence
$|O(n)|\gtrsim r^n$ eventually exceeds $\sqrt n\log n$.  Therefore
\[
   |O(n)|\le c\,\sqrt n\,\log n\ (\forall n)
   \iff |1-\tfrac1\rho|=1\ (\forall\rho)
   \iff \textnormal{RH},
\]
with no slack at any step.  Equivalently: the square-root cancellation is \emph{contingent exactly} on
the reality of $\theta(\gamma)$; every phase-blind method (stationary phase, saddle point,
Plancherel--Rotach) is by construction insensitive to zero locations and hence structurally unable to
supply this hypothesis.
\end{proposition}
\noindent This is $H0$ in its sharpest form: the classical analysis of Conjecture~\ref{conj:envelope}
is complete and correct, and it reduces the bound to the statement that no zero has $\beta\neq\tfrac12$
--- which classical zero-free-region technology narrows (to $\beta>1-c/\log|\gamma|$) but has never
pushed to the critical line.  The gap between the classical zero-free region and $\beta=\tfrac12$ is
exactly the residue, and it is not closed by any finite numerical verification (which certifies finitely
many $n$ over zeros already known to be on the line, never the $n\to\infty$ tail of a putative off-line
zero at arbitrary height).

\emph{Three classical routes, one nucleus.}  The classical attack on
Conjecture~\ref{conj:envelope} can be mounted from three independent directions; each is carried here
to the exact line where $\sqrt n\log n$ is produced, and each terminates at a \emph{distinct} classical
equivalent of RH.

\begin{proposition}[the envelope bound is a classical RH nucleus]\label{prop:three-routes}
Write $w_\rho:=1-\tfrac1\rho$ and $O(n)=\sum_\rho w_\rho^{\,n}$.

\smallskip\noindent\textbf{(R1) Direct bound.} The functional equation gives the exact identity
$w_{1-\rho}=1/w_\rho$ and $|w_\rho|^2-1=(1-2\beta)/|\rho|^2$; hence the pair $\{\rho,1-\rho\}$ contributes
$w_\rho^{\,n}+w_\rho^{-n}$, which is bounded \emph{iff} $|w_\rho|=1$.  The triangle bound diverges
($\sum_\rho|w_\rho|^n=\infty$ on the line), so cancellation in $\sum_\gamma\cos(n\theta_\gamma)$ is
mandatory, and it requires $\theta_\gamma\in\R$.  \emph{Load-bearing step:} $|w_\rho|=1$ for all $\rho$.

\smallskip\noindent\textbf{(R2) Closing the strip.} Off-line zeros contribute $\sum 2r_\rho^{\,n}\cos(n\theta_\rho)$
with $r_\rho>1$ for the $\beta<\tfrac12$ partner; a single such zero forces $|O(n)|>\sqrt n\log n$ beyond
a finite $n_*$.  Unconditional zero-density estimates ($N(\sigma,T)\ll T^{a(1-\sigma)}\log T$; Ingham
$a=3$, improved through Huxley to Guth--Maynard~\cite{GuthMaynard}) lower the exponent but never give
$0$ in the open strip.  \emph{Load-bearing step:} $N(\sigma,T)=0$ on $\tfrac12<\sigma<1$.

\smallskip\noindent\textbf{(R3) Prime-side rewrite.} With $L_{n-1}^{(1)}(\log m)\sim A\cos(2\sqrt{n\log m}-\varphi)$,
$\lambda_n^{\mathrm{prime}}\sim-\int x^{-1}A\cos(2\sqrt{n\log x})\,d\psi(x)$.  The main term ($d\psi\approx
dx$, unconditional) is $O(\sqrt n\log n)$; the error $\psi(x)-x=-\sum_\rho x^\rho/\rho$ re-injects the
zeros, and its contribution is $O(\sqrt n\log n)$ \emph{iff} $\psi(x)-x=O(\sqrt x\log^2 x)$.
\emph{Load-bearing step:} the RH-strength prime-counting error term.
\smallskip
The three load-bearing steps are the three canonical classical equivalents of RH (Riemann; Landau;
von~Koch).  In every route the main term is unconditional and yields $\sqrt n\log n$; the entire
RH-strength weight sits in the tail/error, and it is the same weight.  Thus
Conjecture~\ref{conj:envelope} is a genuine classical \emph{nucleus}: no classical route closes it
without closing RH, because all three routes converge to RH itself.
\end{proposition}
\noindent This is the sharpest honest statement the programme can make about the terminal step: a
lossless, three-front classical reduction with the irreducible line exposed in each front --- a closed
\emph{reduction} path (auditable line by line), not a closed \emph{proof}.

\emph{The archimedean perturbative route is quantitatively closed.}  The one remaining candidate source
of positivity --- the archimedean place dominating the prime term (Connes--Consani~\cite{ConnesConsani2021},
the partition Problem~\ref{prob:arch-partition}, the entropic X3-G) --- can now be diagnosed sharply,
and the diagnosis rules out the \emph{perturbative} form of the route.  Two unconditional facts collide.
\emph{(i)} The archimedean positivity margin is only logarithmic in the support: for the Weil operator on
an interval of length $a$, the lowest eigenvalue satisfies (screw-function representation,
\cite{ScrewWeil}) $\lambda_a=\log(1/a)+\mu_1-\log(2\pi)+\psi(2)-1+O(a)$ as $a\to0^+$, with $\mu_1>0$, and
this margin \emph{shrinks} as the support grows.  \emph{(ii)} The prime perturbation grows without bound
as the support opens (each prime enters once $k\log p\le$ support, with negative oscillatory weight, and
their $\ell^1$ mass is $\asymp e^{a/2}$).  A logarithmically-growing margin cannot dominate an
unboundedly-growing perturbation, so \emph{no uniform operator-norm comparison
$\|\text{prime}\|<(\text{archimedean gap})$ can hold} without deciding RH; the eigenvalue reformulation
$\lambda_a\ge0\ \forall a\iff\textnormal{RH}$ (ibid.) makes the equivalence exact.  The only
unconditional positivity in this family is the small-support regime, where finitely many primes
contribute and the $\log(1/a)$ term dominates a \emph{finite} sum --- precisely
Lemma~\ref{lem:perprime-feasible}, which never reaches the arithmetic scale that decides RH.  The
precedent is exact: Conrey--Li~\cite{ConreyLi} falsified the de~Branges positivity conditions that would
have delivered RH structurally.  Conclusion: any archimedean route to $\Omega_7$ must be
\emph{non-perturbative} --- it must exploit the fine cancellation of the prime comb against the
archimedean symbol, not a norm bound --- and that cancellation is, once again, exactly RH.


\subsubsection*{Honest status of the \texorpdfstring{$\Omega_7$}{Omega7} attack}

Seven mechanisms were executed, seven autopsies performed, four walls (W1--W4) established, and one
closure specification produced.  The gap changed shape --- from ``a missing cohomology'' (undefined) to
a single cone statement of Minkowski type (M2), the first formulation of the programme in which the
missing input is conic rather than spectral, hence probeable by convex conic programming at finite $N$
(Problem~\ref{prob:M2N}).  The unconditional by-products, incorporable to the programme after
verification, are: Lemma~\ref{lem:transport} (gauge transport of the arithmetic form);
Proposition~\ref{prop:ceiling} (ceiling rigidity: the RH exceptional set is the super-level set of an
explicit prime-built almost-periodic function); Theorem~\ref{thm:W4} (the fourth wall); the closure
specification (X1)--(X4) with (X1),(X2),(X4) already realized inside the programme; and the dictionary
to Dirichlet large values connecting the residue to the Guth--Maynard front.  What is \emph{not}
created is the engine of M2: Minkowski in dimension one runs on compactness and the covolume of
$\R^n/L$; the analogue for correspondence classes needs ``linear equivalence on the square'' --- what it
means to move $D$ within its class --- which remains definition-shaped, the residue of the residue, now
named as a cone.  If M2$_N$ passes numerically on $\zeta$ and fails on Davenport--Heilbronn, the shift
cone is evidenced as the correct effective cone and the problem reduces to endowing it with a volume
theorem; if M2$_N$ fails on $\zeta$, this (X3) is falsified in hours rather than years.  Either outcome
advances the programme, and $\Omega_7$ remains, as stated throughout, the sole open input, now specified
to the precision of a finite convex-feasibility experiment.

\subsection{Step 5 substructure (continued): de Bruijn--Newman cascade characterization}\label{subsec:cascade}

This subsection is self-contained: it characterizes the de Bruijn--Newman constant by finite terminal
sections, but it does not evaluate that constant and therefore does not prove RH.

\begin{definition}[de Bruijn--Newman flow and terminal cascade]\label{def:flow-defect}
Let
\[
   \Xi_t(z)=\int_0^\infty \Phi(u)e^{tu^2}\cos(zu)\,du,
   \qquad \partial_t\Xi_t=-\partial_z^2\Xi_t,\qquad \Xi_0=\Xi,
\]
and let $\Lambda_{\rm dBN}$ be the de Bruijn--Newman constant.  Fix $z_0=t_0-iy$ and the reference
form
\[
   c^*A_Nc=\frac L{2\pi}\|R_c\|_{L^2(\R)}^2.
\]
For the Pick jet matrix $J_N(t)$ of $-\Xi_t'/\Xi_t$ at $z_0$, define
\[
   \delta_N(t):=\min_{c\ne0}\frac{c^*J_N(t)c}{c^*A_Nc},
   \qquad
   t_N^*:=\sup\{t:\delta_N(t)<0\}.
\]
\end{definition}

\begin{lemma}[$\alpha1$: positivity in the real-zero region]\label{lem:alpha1}
For $t\ge\Lambda_{\rm dBN}$ all zeros of $\Xi_t$ are real, hence $J_N(t)$ is a Gram matrix and
$\delta_N(t)\ge0$.  Consequently
\[
   t_N^*\le\Lambda_{\rm dBN}\qquad(N\ge1).
\]
\end{lemma}

\begin{proof}
This is the defining property of $\Lambda_{\rm dBN}$ and the Cauchy-transform Pick-kernel Gram
identity of Lemma~\ref{lem:cauchy-pick}, applied to the zero measure of $\Xi_t$.
\end{proof}

\begin{lemma}[$\alpha2$: regular flow derivative]\label{lem:alpha2}
In the real-zero region, write the zeros as $\gamma_j(t)$.  With principal-value pairing,
\[
   \dot\gamma_j=2\sum_{k\ne j}\frac1{\gamma_j-\gamma_k}.
\]
For fixed $c$ and $f=|R_c|^2$,
\[
   \frac{d}{dt}\,c^*J_N(t)c
   =
   2\sum_{j<k}
   \frac{f'(\gamma_j)-f'(\gamma_k)}{\gamma_j-\gamma_k},
\]
and each summand is regular as $\gamma_j\to\gamma_k$.
\end{lemma}

\begin{proof}
Differentiate the zero-side Gram sum $c^*J_N(t)c=\sum_j f(\gamma_j(t))$ and substitute the Calogero
equation, giving $\frac{d}{dt}c^*J_Nc=\sum_j f'(\gamma_j)\,\dot\gamma_j
=2\sum_j f'(\gamma_j)\sum_{k\ne j}\frac1{\gamma_j-\gamma_k}$.  Antisymmetrizing $j\leftrightarrow k$ in
$f'(\gamma_j)/(\gamma_j-\gamma_k)$ turns this into the displayed paired divided difference; the numerator
$f'(\gamma_j)-f'(\gamma_k)$ vanishes linearly as $\gamma_j\to\gamma_k$, so every summand is regular.
It remains to justify the two convergences, which we do unconditionally.

\emph{Step 1: the inner field is a convergent principal value with an unavoidable local collision term.}
Fix $j$ and set $S_j:=2\,\mathrm{PV}\sum_{k\ne j}(\gamma_j-\gamma_k)^{-1}$, the pairing being the
symmetric one $|\gamma_k-\gamma_j|\le R$, $R\to\infty$.  Then $S_j$ exists as a symmetric principal value
at every simple zero, with
\[
   |S_j|\le C\bigl(\log(2+|\gamma_j|)+g_j^{-1}\bigr),\qquad
   g_j:=\operatorname{dist}(\gamma_j,\textnormal{other zeros}),
\]
the local $g_j^{-1}$ term being unavoidable near collisions --- the velocity of a zero diverges as the
gap closes, which is exactly the physics the cascade exploits.  For fixed $t$ in the \emph{open}
real-zero region all gaps are positive, only finitely many $g_j^{-1}$ terms meet the window where $f'$
has not yet decayed, so the block $|\gamma_k-\gamma_j|\le1$ (with $O(\log(2+|\gamma_j|))$ zeros by
Lemma~\ref{lem:alpha-L2}) is a finite sum; for consecutive tail zeros the paired sum
$\frac1{\gamma_j-\gamma_k}+\frac1{\gamma_j-\gamma_{k'}}$ straddling a point is
$O(|\gamma_j-\gamma_k|^{-2})$ times the gap, hence summable.  For the tail, Abel
summation against the counting function $N(T)=\frac{T}{2\pi}\log\frac{T}{2\pi}+O(\log T)$ gives
\[
   \sum_{1<|\gamma_k-\gamma_j|\le R}\frac1{\gamma_j-\gamma_k}
   =\int_{1}^{R}\frac{dN_j(u)}{\mp u}
   =O(\log(2+|\gamma_j|)),
\]
uniformly in $R$, the odd part of the smooth density cancelling in the principal value and the
$O(\log)$ remainder of $N$ producing the stated bound.  Thus $S_j$ exists as a principal value and
$|S_j|\le C\log(2+|\gamma_j|)$ with an absolute $C$.

\emph{Step 2: the outer sum converges absolutely.}  Since $R_c(u)=\sum_{j=0}^{N-1}\overline{c_j}
(u-z_0)^{-(j+1)}$ has leading decay $(u-z_0)^{-1}$, the sampling weight $f=|R_c|^2$ satisfies
$f(u)=O(|u|^{-2})$ and $f'(u)=O(|u|^{-3})$ as $|u|\to\infty$, with constants depending only on $c$ and
$z_0$.  Hence
\[
   \Bigl|\sum_j f'(\gamma_j)S_j\Bigr|
   \le C\sum_j\frac{\log(2+|\gamma_j|)}{(1+|\gamma_j|)^{3}}<\infty,
\]
the series converging by $N(T)=O(T\log T)$.  Both the inner principal value and the outer sum therefore
converge, and the antisymmetrized paired form equals $\frac{d}{dt}c^*J_Nc$; this is exactly the paired
divided-difference expression of the statement.
\end{proof}

\begin{theorem}[$\alpha3$: Gronwall rigidity in the real-zero region]\label{thm:alpha3}
There exists $C_N={\rm poly}(N,L)$ such that, while all zeros are real,
\[
   \left|\frac{d}{dt}\delta_N(t)\right|\le C_N\sqrt{\delta_N(t)}.
\]
Thus positivity at level $\delta_N=m$ persists backward in flow time for at least
$2\sqrt m/C_N$.
\end{theorem}

\begin{proof}[Proof sketch and correction note]
Use Feynman--Hellmann at a minimizing vector for $\delta_N(t)$.  The exact minimizer need not satisfy
$f'(\gamma_j)=0$ at sampled zeros; instead one uses
$|f'(\gamma_j)|^2\le2f(\gamma_j)F_j''$ and Cauchy--Schwarz in the paired derivative formula of
Lemma~\ref{lem:alpha2}.  Since the sampling sum of $f(\gamma_j)$ is $\delta_N$ times the reference
denominator and $F_j''$ is bounded by a polynomial in $N,L$ on the active window, the derivative is
$O(C_N\sqrt{\delta_N})$.  This corrects the stronger but false log-Lipschitz exponent $1$ that would
hold only for an idealized interpolator.
\end{proof}

\begin{lemma}[$\alpha4$-L1: de Bruijn confinement]\label{lem:alpha-L1}\label{lem:alpha-L1-full}
Let
\[
   B(t):=\sup\{|\Im\rho|:\Xi_t(\rho)=0\}.
\]
Then:
\begin{enumerate}[label=\textnormal{(\roman*)}]
\item the non-real zeros of $\Xi_t$ occur in conjugate pairs;
\item \textnormal{(de Bruijn)} for $t_2\ge t_1$,
\[
   B(t_2)^2\le \max\{B(t_1)^2-2(t_2-t_1),0\};
\]
\item for every $t\in\R$,
\[
   B(t)\le\sqrt{\frac14+2\max(-t,0)}.
\]
\end{enumerate}
\end{lemma}

\begin{proof}
Since $\Phi$ is real and even, $\Xi_t$ is real on $\R$ and even; the zero set is therefore invariant
under conjugation, proving (i).  Assertion (ii) is de Bruijn's strip-narrowing theorem for the forward
heat flow \cite[Thm.~13]{deBruijn1950}, applied to the evolution from $t_1$ to $t_2$.  For $t\le0$,
use the complementary backward (strip-widening) direction of de Bruijn's theorem
\cite[Thm.~13]{deBruijn1950}: evolving backward from time $0$ to time $t$, the strip half-width can
grow at most at the matching quadratic rate, $B(t)^2\le B(0)^2+2|t|$; with the critical strip
$B(0)\le1/2$ this gives the bound.  For $t\ge0$, apply (ii) from $0$ to $t$.
\end{proof}

\begin{lemma}[$\alpha4$-L2: uniform zero counting with honest exponent]\label{lem:alpha-L2}\label{lem:alpha-L2-full}
Fix a compact flow interval $t\in[-T_0,T_1]$.  There are constants $C,X_0$ such that, for
$T\ge X_0$,
\[
   N_{\Xi_t}(T+1)-N_{\Xi_t}(T)\le C(\log T)^2,
   \qquad
   N_{\Xi_t}(T+H)-N_{\Xi_t}(T)
   =
   \frac{H}{2\pi}\log\frac{T}{2\pi}+O((\log T)^2),
\]
uniformly for $0<H\le T$ and uniformly in $t$ on the compact interval.  Equivalently, the flowed
argument-error term satisfies $S_t(T)=O((\log T)^2)$ uniformly on compact $t$-ranges.
\end{lemma}

\begin{proof}
\emph{Step 1: upper bound by contour shift.}  The kernel $\Phi$ extends holomorphically to
$|\Im u|<\pi/4$ and satisfies
\[
   |\Phi(u+i\alpha)|\le C e^{9u/2}e^{-\pi\cos(2\alpha)e^{2u}}.
\]
For $z=x+iy$, $|y|\le5$, and large $x$, shift $u\mapsto u+i\alpha$ with
$\delta:=\pi/4-\alpha$.  Then
\[
   |\Xi_t(z)|
   \le
   e^{-(\pi/4-\delta)x}
   \int_\R |\Phi(u+i\alpha)|e^{|t|u^2+5|u|}\,du
   \le
   e^{-(\pi/4-\delta)x}\exp\{C+|t|(\tfrac12\log(1/\delta))^2\}\delta^{-C}.
\]
Taking $\delta=1/x$ gives
\[
   |\Xi_t(z)|\le e^{-\pi x/4}\exp\{C_t(\log x)^2\},
\]
uniformly for $t$ in the chosen compact interval.

\emph{Step 2: saddle lower bound at height $4$.}  The $n=1$ term in $\Phi$ dominates with relative
error $O(e^{-2u})$.  For
\[
   h(u)=\frac92u+tu^2-\pi e^{2u}+izu
\]
the saddle is
\[
   u^*(z)=\frac12\log\frac{x}{2\pi}+\frac{i\pi}{4}+O_t\!\left(\frac{\log x}{x}\right),
   \qquad
   h''(u^*)=-2iz+O(1).
\]
At height $y=4$, Laplace's method gives
\[
   \Xi_t(x+4i)=A_t(x+4i)(1+O(x^{-1/2})),
\]
where
\[
   \log|A_t(x+4i)|
   =
   -\frac{\pi x}{4}
   +\frac{t}{4}\left(\log\frac{x}{2\pi}\right)^2
   +2\log\frac{x}{2\pi}
   +O_t(1).
\]
The terms $n\ge2$ are $O(n^{-1/2-8})$ relative at this height, so the series is absolutely
convergent and the estimate is zero-free on the horizontal line.

\emph{Step 3: Jensen.}  Apply Jensen to disks $D(x_0+4i,5)$, which cover the relevant strip over
$[x_0-3,x_0+3]$.  The boundary upper bound from Step 1 and the center lower bound from Step 2 give
\[
   n(D(x_0+4i,4.6))
   \le
   \log\frac{\max_{\partial D}|\Xi_t|}{|\Xi_t(x_0+4i)|}
   \le
   O_t((\log x_0)^2).
\]

\emph{Step 4: main term.}  The argument principle on
$[T,T+H]\times[-4,4]$, using the saddle approximation on the horizontal sides and Step 3 on the
vertical sides, gives the Riemann--von Mangoldt main term with error $O((\log T)^2)$.  The expected
error is $O(\log T)$; the stronger bound would require a Phragmen--Lindelof refinement with the
quadratic logarithmic compensator.  The $(\log T)^2$ exponent is the one used below.
\end{proof}

\begin{lemma}[$\alpha4$-L3: interpolator tail]\label{lem:alpha-L3}\label{lem:alpha-L3-full}
Let $R=P/(u-z_0)^N$, with all roots of $P$ in the active window
$|u-t_0|\le W=\pi N/L$.  Then
\[
   \sum_{|\gamma-t_0|>2W}|R(\gamma)|^2
   \le
   e^{2N[\widehat\varphi_{\rm out}-\widehat\varphi(t_c)](1+o(1))}\,c^*A_Nc
   =
   \left(\frac{CyL}{N}\right)^{2N(1+o(1))}c^*A_Nc.
\]
\end{lemma}

\begin{proof}
\emph{Step 1: sampling to $L^2$ without exponential loss.}  In the far region $R$ has no zeros or
poles at distance less than a fixed fraction of the distance to the active window.  On each unit
interval $I_d$ at distance $d$, Nikolskii's inequality for $P$ gives
\[
   \sup_{I_d}|R|^2\le C N^2\int_{2I_d}|R|^2.
\]
By Lemma~\ref{lem:alpha-L2}, the number of zeros in $I_d$ is $O((\log d)^2)$.  Summing over far
intervals yields
\[
   \sum_{\rm far}|R(\gamma)|^2
   \le
   C N^2\sum_d(\log d)^2\int_{2I_d}|R|^2
   \le
   C N^2(\log X_R)^2\|R\|^2_{L^2({\rm far})},
\]
where $X_R$ is polynomial in $N,t_0$ on the scales used here.  This is subfactorial.

\emph{Step 2: far $L^2$ mass by potential.}  For $|u-t_0|>2W$,
\[
   \log|R(u)|=N\bigl[U^{\mu_N}(u)-\log|u-z_0|\bigr],
\]
where $\mu_N$ is the empirical root measure.  Comparing $\mu_N$ with the uniform model by partial
summation against $N_t-\mathrm{mean}$ and using $S_t=O(L^2)$ from Lemma~\ref{lem:alpha-L2}, the exponent
error is $O(L^2\log N)=o(N\log N)$ in the regime of the theorem.  The potential calculation is the
upper side of Theorem~\ref{thm:two-sided}, and gives the displayed bound.
\end{proof}

\begin{lemma}[$\alpha4$-L4: central derivative, honest loss]\label{lem:alpha-L4}\label{lem:alpha-L4-full}
Let the explicit interpolator have a real root at $\gamma_c$, with
$|\gamma_c-t_0|\le t_1W$, and set
\[
   \eta:=\min\left\{1,\frac{L}{2\pi}\operatorname{dist}(\gamma_c,\textnormal{another active object})\right\}.
\]
Then
\[
   \frac{2|R'(\gamma_c)|^2}{c^*A_Nc}
   \ge
   \frac{\pi L}{y}\,\eta^2 e^{-C L^2\log N}
   \exp\{2N[\widehat\varphi(t_d)-\widehat\varphi(t_c)](1+o(1))\}.
\]
\end{lemma}

\begin{proof}
We have
\[
   R'(\gamma_c)=\frac{\prod_j(\gamma_c-x_j)}{(\gamma_c-z_0)^N},
\]
where the product omits the root at $\gamma_c$.  The two nearest-neighbor factors contribute at least
$(\eta\,2\pi/L)^2$.  For the remaining factors, compare
\[
   \sum_j\log|\gamma_c-x_j|
\]
with the corresponding potential integral by partial summation.  Lemma~\ref{lem:alpha-L2} gives the
error bound $C L^2\log N$.  The pole contributes $(d^2+y^2)^{-N}$ exactly.  The denominator
$c^*A_Nc\le (L/2\pi)\|R\|^2$ is bounded from above by the same potential calculation, with the same
loss.  Thus the displayed lower bound follows.  The measured normalization is closer to an
$N^{-3/2}$ loss, but the proof only uses the subfactorial factor $e^{-C L^2\log N}$.
\end{proof}

\begin{theorem}[$\alpha4$: final quadratic detection of a complex pair]\label{thm:alpha4}
Let $t<\Lambda_{\rm dBN}$ and let $\{z_w,\bar z_w\}$, $z_w=\gamma_c-iv$, be a non-real pair of
$\Xi_t$ with $|\gamma_c-t_0|\le t_1\pi N/L$.  Then
\[
   \delta_N(t)
   \le
   \left(\frac{C yL}{N}\right)^{2N(1+o(1))}
   -
   v^2\,\frac{\pi L}{y}\,\eta^2e^{-C L^2\log N}
   \exp\{2N[\widehat\varphi(t_d)-\widehat\varphi(t_c)](1+o(1))\}
   +O(v^3\|\cdot\|),
\]
where
\[
   \widehat\varphi(t)=\psi(t)-\log|t|-1,\qquad
   \psi(t)=\tfrac12[(1-t)\log|1-t|+(1+t)\log(1+t)].
\]
In particular, $\delta_N(t)<0$ once
\[
   v^2\ge
   e^{C L^2\log N}\,\frac{y}{\pi L}\,m_N\,e^{-2N\widehat\varphi(t_d)} ,
\]
with only subfactorial loss.
\end{theorem}

\begin{proof}
Use the explicit test function
\[
   R(u)=\frac{P(u)}{(u-z_0)^N}.
\]
The numerator $P$ is chosen with real roots at the real zeros in the active window, one real root at
$\gamma_c$, and one complex root at one member of every additional non-real pair in the window.
Those additional pairs contribute
$2\Re\{R(z')\overline{R(\bar z')}\}=0$ exactly.  For the target pair, expanding around the real
root at $\gamma_c$ gives, without any real-coefficient assumption,
\[
   R(z_w)\overline{R(\bar z_w)}
   =
   -v^2|R'(\gamma_c)|^2+O(v^3),
\]
so the rank-two hyperbolic contribution is
$-2v^2|R'(\gamma_c)|^2+O(v^3)$.  The real far-zero contribution is bounded by
Lemma~\ref{lem:alpha-L3}, while Lemma~\ref{lem:alpha-L4} supplies the displayed lower bound for the
central derivative.  Lemma~\ref{lem:alpha-L2} gives the root budget in the active window.
\end{proof}

\begin{theorem}[$\alpha5$: terminal cascade formula for de Bruijn--Newman]\label{thm:alpha5}
\label{thm:cascade-main}
\[
   \lim_{N\to\infty}t_N^*=\Lambda_{\rm dBN}.
\]
\end{theorem}

\begin{proof}
Lemma~\ref{lem:alpha1} gives $t_N^*\le\Lambda_{\rm dBN}$ for every $N$.  Conversely, fix
$t<\Lambda_{\rm dBN}$.  By Lemma~\ref{lem:alpha-L1}, some conjugate pair is non-real, with finite
height $v>0$ and center $\gamma_c$.  For all
\[
   N\ge \max\{C L|\gamma_c-t_0|,\ N_1(v)\},
\]
the center lies inside the detection window of Theorem~\ref{thm:alpha4}, and the factorial threshold
there is smaller than $v^2$.  Hence $\delta_N(t)<0$ and $t_N^*\ge t$.  Since this holds for every
$t<\Lambda_{\rm dBN}$, $\liminf_N t_N^*\ge\Lambda_{\rm dBN}$, while the opposite inequality was already
proved.
\end{proof}

\begin{corollary}[cascade form of RH]\label{cor:alpha-rh}
Using Rodgers--Tao, $\Lambda_{\rm dBN}\ge0$, one has
\[
   \mathrm{RH}\quad\Longleftrightarrow\quad \sup_N t_N^*\le0
   \quad\Longleftrightarrow\quad
   \lim_N t_N^*=0 .
\]
\end{corollary}

In summary, the cascade formula (Theorem~\ref{thm:alpha5}) does not prove RH; it gives a finite-section
characterization of the de Bruijn--Newman constant.  The open input is the same as before, in flow
coordinates: for $\zeta$, exclude a cascade accumulating at some $\Lambda_{\rm dBN}>0$. What is proved,
precisely, is isolated in the following components: principal-value bookkeeping in Lemma~\ref{lem:alpha2};
the honest $(\log T)^2$ flow-counting bound in Lemma~\ref{lem:alpha-L2}; the product-tail estimate of
Lemma~\ref{lem:alpha-L3}; and the derivative lower bound of Lemma~\ref{lem:alpha-L4}, whose proved
loss is $e^{O(L^2\log N)}$ rather than the numerically observed $N^{3/2}$.  These losses are
subfactorial and do not affect Theorem~\ref{thm:alpha5}.  The numerical experiments described below
validate the quadratic law, the Blaschke rate, and the cascade formula in the tested regimes.

\begin{remark}[computability and validation]\label{rmk:cascade-validation}
Each $\delta_N(t)$ is a finite-section computation.  Numerical verification confirms: the quadratic
pair mechanism with shape constant $0.75$ to $10^{-5}$; the off-line Blaschke rate $2.25545$
predicted versus $2.2556$ measured for Davenport--Heilbronn; and the cascade formula with ratio
$1.000$ through the tested range, with the finite-threshold correction visible only at the smallest
tested order.  The sine model, a stationary Calogero lattice, has an empty cascade.
\end{remark}

\begin{remark}[rate heuristic]\label{rmk:cascade-rate}
Under GUE statistics for $\zeta$,
\[
   |t_N^*|
   \approx
   \frac{\pi^2}{2L^2}
   \left(\frac{3\pi^2}{12.8\,N}\right)^{2/3}.
\]
The cascade exponent $-2/3$ measures the repulsion exponent $\beta=2$.  Thus the cascade observable
is a finite-section bridge between the de Bruijn--Newman flow and local zero statistics.
\end{remark}

\begin{remark}[formal critical-line symbol]\label{rmk:critical-symbol}
If the half-line shifts in $\mathcal P_\Lambda$ are replaced formally by full-line shifts, the Fourier
symbol is
\[
   2\sum_{n\ge2}\Lambda(n)n^{-1/2}\cos(\theta\log n)
   =
   -2\Re\frac{\zeta'}{\zeta}\!\left(\half+i\theta\right)
\]
in the sense of formal Dirichlet series.  The series is not pointwise convergent on the critical line;
the half-line compression and the exponential depth $y>\half$ are precisely the regularization.  This
explains why the terminal operator sees the critical-line obstruction while remaining computable in
finite sections.
\end{remark}

\begin{remark}[terminal numerical observable]\label{rmk:terminal-numerics}
The finite-section numerical experiment is now read as follows.  Compute
\[
   C_N(t_0):=\lambda_{\max}(T_N(t_0-iy))\log(|t_0|/2\pi),
   \qquad
   \delta_N(t_0):=1-\lambda_{\max}(T_N(t_0-iy)).
\]
For $\zeta$, the expected signature is $\delta_N>0$ with $\delta_N\to0$; for a Davenport--Heilbronn
falsifier, some finite section must cross $\lambda_{\max}>1$.  The observed saturation therefore
supports Theorem~\ref{thm:saturation} and rules out margin-based proof attempts, but it does not close
the terminal inequality.
\end{remark}

\begin{remark}[Scale Correspondence]\label{rmk:scale-correspondence}
Lemma~\ref{lem:prime-window} identifies jet order with prime scale: the order-$k$ jet at depth
$\epsilon_0$ probes prime powers with $\log n$ concentrated near $k/\epsilon_0$, i.e. the H5 window of
length $L\asymp k/\epsilon_0$.  Together with Theorems~\ref{thm:one-point-jets},
\ref{thm:det-serialization}, and~\ref{thm:live-H5}, this gives four equivalent serializations of the
same global quantifier: node Pick positivity, determinant positivity, window positivity, and one-point
jet positivity.  The fixed-$N$ band regime is closed by Theorem~\ref{thm:band-closure}; the terminal
one-point gauge has the effective Pascal--Laguerre floor of Lemma~\ref{lem:pascal-effective} and the
effective $\log\log$ range of Corollary~\ref{cor:loglog-effective}.  The open uniform H8 problem is now
exactly the no-margin condition $\delta_N\ge0$ for all $N$, equivalently the prime-scale quantifier as
$N\to\infty$.
\end{remark}



\subsection{Realization layer: source core, finite channels, and the measure bridge}\label{sec:shorting}

Throughout this subsection, $A_P^\circ$ denotes a self-adjoint operator on a Hilbert space
$\cH_P^\circ$ containing the shorted prime-cell direct sum, with the source maps $\Phi_P^F$ of
Definition~\ref{def:source-embedding}. The construction of such operators --- the von Mangoldt
canonical-system realization, built from the positive Hamiltonian $H_P\succeq0$ of
Theorem~\ref{thm:live-H1} by the standard de Branges correspondence (\S\ref{sec:relwork}) --- is carried
out elsewhere in this programme and is not reproved here; the present layer uses only
\textnormal{(i)} self-adjointness, hence $\|(A_P^\circ-z)^{-1}\|\le|\Im z|^{-1}$, and
\textnormal{(ii)} the $P$-uniform source bound of Lemma~\ref{lem:sourcebound}. All statements below are
conditional on these two properties.

\begin{lemma}[Mertens pole growth]\label{lem:mertens-pole}
For $x\ge2$,
\[
   \sum_{n\le x}\frac{\Lam(n)}{n}=\log x+O(1),\qquad
   \sum_{n\le x}\frac{\Lam(n)\log n}{n}=\half(\log x)^2+O(\log x).
\]
\end{lemma}

\begin{proof}
The identity $\log m=\sum_{d\mid m}\Lam(d)$ gives
\[
   \log\lfloor x\rfloor!=\sum_{n\le x}\Lam(n)\left\lfloor\frac{x}{n}\right\rfloor.
\]
Writing $\lfloor x/n\rfloor=x/n-\{x/n\}$ and using Stirling's formula together with Chebyshev's bound
$\sum_{n\le x}\Lam(n)=O(x)$ gives
\[
   x\sum_{n\le x}\frac{\Lam(n)}{n}=x\log x+O(x),
\]
hence the first estimate.  Abel summation applied to
$A(t)=\sum_{n\le t}\Lam(n)/n=\log t+O(1)$ gives
\[
   \sum_{n\le x}\frac{\Lam(n)\log n}{n}
   =A(x)\log x-\int_2^x\frac{A(t)}{t}\,dt
   =\half(\log x)^2+O(\log x).
\]
\end{proof}

\begin{definition}[pole vector and primitive complement]\label{def:pole-vector}
Let $e_0$ be the normalized pole/degree vector in the finite canonical Hilbert space $\cH_P$, and write
\[
   \cH_P=\C e_0\oplus e_0^\perp.
\]
The pole cell is the positive rank-one form
\[
   H_{\mathrm{pole},P}=a_P\,|e_0\rangle\langle e_0|,\qquad
   a_P=\sum_{n\le P^2}\frac{\Lam(n)\log n}{n}+O(\log P).
\]
The primitive complement is $\cH_P^\circ=e_0^\perp$ equipped with the Schur-shorted form.
\end{definition}

\begin{lemma}[rank-one pole isolation]\label{lem:pole-isolation}
The only divergent direction in the finite von Mangoldt Gram form is the pole vector $e_0$.  Its
self-pairing satisfies
\[
   [e_0,e_0]_P=a_P=\half(\log P^2)^2+O(\log P),
\]
and every vector $\phi\in\cS_a^\circ$ has uniformly bounded unshorted primitive mass away from the
$e_0$ direction.
\end{lemma}

\begin{proof}
By Definition~\ref{def:pole-vector}, the pole contribution is a positive rank-one form supported on
$e_0$.  Lemma~\ref{lem:mertens-pole} gives
\[
   a_P=\sum_{n\le P^2}\frac{\Lam(n)\log n}{n}+O(\log P)
      =\half(\log P^2)^2+O(\log P).
\]
All non-pole prime cells are evaluated on the pole-shorted source profile.  For
$\phi\in\cS_a^\circ$,
\[
   \sum_{p,k}(\log p)p^{-k/2}|\widehat\phi(k\log p)|^2
   \le C\sum_{p,k}(\log p)p^{-k(1/2+2a_\phi)}<\infty
\]
because $a_\phi>\half>\tfrac14$.  Thus no source direction in the primitive core carries the Mertens
divergence.
\end{proof}

\begin{lemma}[Schur shorting closes A3]\label{lem:short}
The divergent pole contribution is a positive rank-one block, and Schur shorting it preserves the
negative-square index.  More precisely, if
\[
   K_P=\begin{pmatrix}\alpha_P&\beta_P^*\\ \beta_P&C_P\end{pmatrix},\qquad \alpha_P>0,
\]
then
\[
   K_P^\circ=C_P-\beta_P\alpha_P^{-1}\beta_P^*
\]
satisfies
\[
   \sqm(K_P)=\sqm(\alpha_P)+\sqm(K_P^\circ)=\sqm(K_P^\circ).
\]
\end{lemma}

\begin{proof}
Lemma~\ref{lem:pole-isolation} identifies the pole block as the positive scalar
$\alpha_P=[e_0,e_0]_P>0$ on the rank-one space $\C e_0$.  Haynsworth inertia additivity follows from the
congruence
\[
   \begin{pmatrix}I&0\\-\beta_P\alpha_P^{-1}&I\end{pmatrix}^*
   K_P
   \begin{pmatrix}I&0\\-\beta_P\alpha_P^{-1}&I\end{pmatrix}
   =
   \begin{pmatrix}\alpha_P&0\\0&K_P^\circ\end{pmatrix}.
\]
Since $\alpha_P>0$, it contributes no negative square.
\end{proof}

\begin{definition}[finite source embedding]\label{def:source-embedding}
For a finite cutoff $P$ and $\phi\in\cS_a^\circ$, define the unshorted primitive source vector
$\widetilde\iota_P\phi$ in the finite prime-cell Hilbert space by its cell coordinates
\[
   (\widetilde\iota_P\phi)_{p,k}
   :=
   \bigl((\log p)p^{-k/2}\bigr)^{1/2}\widehat\phi(k\log p),
   \qquad p^k\le P.
\]
Equivalently, if the primitive cell carriers $u_{p,k}$ are normalized, then
\[
   \widetilde\iota_P\phi
   =
   \sum_{p^k\le P}
   \bigl((\log p)p^{-k/2}\bigr)^{1/2}\widehat\phi(k\log p)\,u_{p,k}.
\]
The primitive source vector used in the resolvent compression is the Schur-shorted projection
\[
   \iota_P\phi:=\Pi_P^\circ \widetilde\iota_P\phi\in\cH_P^\circ.
\]
For $F=\operatorname{span}\{\phi_1,\ldots,\phi_m\}$, the source map is
\[
   \Phi_P^F(c_1,\ldots,c_m)=\sum_{\alpha=1}^m c_\alpha\,\iota_P\phi_\alpha.
\]
\end{definition}

\begin{lemma}[source embedding norm identity]\label{lem:source-embedding-norm}
Before Schur shorting,
\[
   \|\widetilde\iota_P\phi\|^2
   =
   \sum_{p^k\le P}(\log p)p^{-k/2}|\widehat\phi(k\log p)|^2.
\]
After shorting,
\[
   \|\iota_P\phi\|_{\cH_P^\circ}^2
   \le
   \|\widetilde\iota_P\phi\|^2.
\]
\end{lemma}

\begin{proof}
The first identity is the norm in the orthogonal finite prime-cell direct sum of
Definition~\ref{def:source-embedding}.  The second inequality is the monotonicity of Schur shorting,
Lemma~\ref{lem:shorting-monotone}, applied to the positive pole block.
\end{proof}

\begin{lemma}[compact logarithmic tests belong to the source core]\label{lem:compact-core}
Every $\phi\in C_c^\infty(0,\infty)$ belongs to $\cS_a^\circ$.
\end{lemma}

\begin{proof}
Let $\widehat\phi$ be the logarithmic Mellin profile and assume
$\operatorname{supp}\widehat\phi\subset[\alpha,\beta]\subset(0,\infty)$.  Fix any
$a_\phi>\half$.  For each $N$,
\[
   e^{a_\phi t}(1+t)^N|\widehat\phi(t)|
\]
is continuous on the compact interval $[\alpha,\beta]$ and is zero outside it, hence has a finite
supremum $C_N$.  Therefore
\[
   |\widehat\phi(t)|\le C_Ne^{-a_\phi t}(1+t)^{-N}.
\]
The pole-shorted condition is obtained by projecting away the single pole vector of
Definition~\ref{def:pole-vector}; this changes only a finite-dimensional component and leaves the compact
logarithmic test in the primitive source space.

For the strip estimate, write $\operatorname{supp}\phi\subset[0,T]$ and fix
$\eps_\phi\in(0,a_\phi-\half)$.  If $|\Im w|\le\half+\eps_\phi/2$, integration by parts gives, for
every $N$,
\[
   |\ell_w(\phi)|
   =
   \left|\int_0^T\phi(t)e^{-iwt}\,dt\right|
   \le
   C_{N,\phi}(1+|\Re w|)^{-N},
\]
because all boundary terms vanish for a compactly supported smooth test and
${d^N}(e^{-iwt})/{dt^N}=(-iw)^Ne^{-iwt}$, while
$|e^{-iwt}|\le e^{(\half+\eps_\phi/2)T}$ on the enlarged closed substrip.  Thus
$\phi\in\cS_a^\circ$.
\end{proof}

\begin{lemma}[shorting does not increase primitive source mass]\label{lem:shorting-monotone}
Let $K_P\ge0$ be written in the pole/primitive splitting as in Lemma~\ref{lem:short}.  For every primitive
vector $x\in e_0^\perp$,
\[
   \langle K_P^\circ x,x\rangle
   =
   \langle C_Px,x\rangle-\langle \alpha_P^{-1}\beta_P^*x,\beta_P^*x\rangle
   \le \langle C_Px,x\rangle.
\]
Consequently passing from the unshorted primitive block to the Schur-shorted primitive form cannot
increase the norm of a source vector.
\end{lemma}

\begin{proof}
Since $\alpha_P>0$, the second term
$\langle \alpha_P^{-1}\beta_P^*x,\beta_P^*x\rangle$ is nonnegative.  The displayed inequality is
immediate.  This is the order-theoretic form of Schur shorting: $0\le K_P^\circ\le C_P$ as quadratic
forms on the primitive complement.
\end{proof}

\begin{lemma}[uniform finite-channel source bound]\label{lem:sourcebound}
For every $\phi\in\cS_a^\circ$,
\[
   \sup_P\|\iota_P\phi\|_{\cH_P^\circ}<\infty.
\]
\end{lemma}

\begin{proof}
By Lemma~\ref{lem:source-embedding-norm}, before shorting the primitive diagonal mass is
\[
   \sum_{p,k}(\log p)p^{-k/2}|\widehat\phi(k\log p)|^2
   \le
   C\sum_{p,k}(\log p)p^{-k(1/2+2a_\phi)}.
\]
The series converges because $1/2+2a_\phi>1$.  The bound is independent of $P$.  By
Lemma~\ref{lem:source-embedding-norm}, Schur shorting the pole mode can only decrease this primitive
quadratic form.  Therefore the same $P$-independent bound holds in $\cH_P^\circ$.
\end{proof}

\begin{corollary}[finite-channel normal family]\label{cor:normal}
For every finite source plane $F\subset\cS_a^\circ$, the functions
\[
   M_P^F(z)=(\Phi_P^F)^*(A_P^\circ-z)^{-1}\Phi_P^F
\]
satisfy the uniform bound required in Theorem~\ref{thm:vitali}.
\end{corollary}

\begin{proof}
Combine self-adjointness of $A_P^\circ$ with Lemma~\ref{lem:sourcebound}:
\[
   \|M_P^F(z)\|\le \|\Phi_P^F\|^2|\Im z|^{-1}\le C_F|\Im z|^{-1}.
\]
\end{proof}

\begin{lemma}[subsequential spectral-measure compactness]\label{lem:helly-cauchy}
Fix a finite source plane $F\subset\cS_a^\circ$.  For every sequence of cutoffs $P_j$ there is a
subsequence $P_{j_\nu}$ and a positive matrix measure $\Sigma^F$ on $\R$, with finite total mass, such
that
\[
   M_{P_{j_\nu}}^F(z)
   \longrightarrow
   \int_\R\frac{d\Sigma^F(t)}{t-z}
\]
locally uniformly on $\C\setminus\R$.
\end{lemma}

\begin{proof}
Let $E_P$ be the spectral measure of the self-adjoint operator $A_P^\circ$.  Define the positive
matrix-valued measure
\[
   \Sigma_P^F(B):=(\Phi_P^F)^*E_P(B)\Phi_P^F,\qquad B\subset\R\ \text{Borel}.
\]
The spectral theorem gives
\[
   M_P^F(z)=\int_\R\frac{d\Sigma_P^F(t)}{t-z}.
\]
Corollary~\ref{cor:normal} and Lemma~\ref{lem:sourcebound} give a $P$-independent bound on the total
mass:
\[
   \Sigma_P^F(\R)=(\Phi_P^F)^*\Phi_P^F\le C_F I_F .
\]
For each pair of basis vectors $e_\alpha,e_\beta$ in $F$, the scalar measures
$\langle\Sigma_P^F e_\beta,e_\alpha\rangle$ therefore have uniformly bounded total variation.  By
Banach--Alaoglu, equivalently the Helly selection theorem for finite measures on the locally compact space
$\R$, a diagonal subsequence converges weak-* against $C_0(\R)$ for all finitely many matrix entries.  The
limit is a positive matrix measure $\Sigma^F$ because positivity is preserved under weak-* limits:
for every vector $c$, the scalar measures
$\langle\Sigma_{P_{j_\nu}}^F c,c\rangle$ are positive and converge to
$\langle\Sigma^F c,c\rangle$.

For each $z\in\C\setminus\R$, the function $t\mapsto(t-z)^{-1}$ belongs to $C_0(\R)$, hence the displayed
Cauchy transforms converge pointwise.  On a compact $K\subset\C\setminus\R$, the kernels
$(t-z)^{-1}$ are uniformly bounded and uniformly Lipschitz in $z$ by constants depending only on
$\operatorname{dist}(K,\R)$.  A finite-net argument upgrades pointwise convergence to uniform convergence
on $K$.  Thus convergence is locally uniform on $\C\setminus\R$.
\end{proof}

\begin{theorem}[moment criterion for operator-tower convergence]\label{thm:moment-arp}
Fix a point $z_\bullet\in U_-$.  For a finite source plane $F\subset\cS_a^\circ$, the standing tower
convergence
\[
   M_P^F(z)\longrightarrow \GXi^F(z)
\]
locally uniformly on $U_-$ is equivalent to the convergence of all resolvent moments at $z_\bullet$:
\[
   (\Phi_P^F)^*(A_P^\circ-z_\bullet)^{-k-1}\Phi_P^F
   \longrightarrow
   \frac{1}{k!}\,\partial_z^k\GXi^F(z_\bullet),
   \qquad k=0,1,2,\ldots .
\]
\end{theorem}

\begin{proof}
If $M_P^F\to\GXi^F$ locally uniformly on $U_-$, Cauchy's integral formula gives convergence of all
derivatives at $z_\bullet$.  Since
\[
   \partial_z^k(A_P^\circ-z)^{-1}=k!(A_P^\circ-z)^{-k-1},
\]
the displayed moment convergence follows.

Conversely, assume all displayed moments converge.  Corollary~\ref{cor:normal} makes
$\{M_P^F\}_P$ a normal family on $\Omega_-$.  Let $H$ be any subsequential locally uniform limit on
$\Omega_-$.  Cauchy's formula again gives
\[
   \frac{1}{k!}\,\partial_z^kH(z_\bullet)
   =
   \lim_P(\Phi_P^F)^*(A_P^\circ-z_\bullet)^{-k-1}\Phi_P^F
   =
   \frac{1}{k!}\,\partial_z^k\GXi^F(z_\bullet).
\]
The point $z_\bullet$ lies in $U_-$, where $\Xi$ has no zeros because $\Re s>1$ and the Euler product for
$\zeta$ is nonzero there; hence $\GXi^F$ is holomorphic near $z_\bullet$.  The Taylor series of $H$ and
$\GXi^F$ therefore agree in a disk around $z_\bullet$.  By the identity theorem they agree on the
connected component of $U_-$, and then on $\Omega_-$ after removing any apparent poles of the meromorphic
target, because $H$ is holomorphic on all of $\Omega_-$.  Every subsequential limit is the same; normality
then implies the full sequence converges locally uniformly to $\GXi^F$ on $U_-$.
\end{proof}

\begin{theorem}[measure-form tower bridge]\label{thm:measure-bridge}
Assume the standing operator-theoretic tower realizes tower-ARP, i.e.\
$M_P^F\to\GXi^F$ on $U_-$ for every finite source plane $F$.  Then every finite-channel subsequential
spectral limit supplied by
Lemma~\ref{lem:helly-cauchy} has Cauchy transform equal to $\GXi^F$ on $U_-$ and holomorphically removes
all possible poles of $\GXi^F$ in $\Omega_-$.  Consequently RH follows.
\end{theorem}

\begin{proof}
Fix $F$ and a cutoff sequence.  Lemma~\ref{lem:helly-cauchy} gives a subsequence whose transforms converge
locally uniformly on $\C\setminus\R$ to a Cauchy transform
\[
   H^F(z)=\int_\R\frac{d\Sigma^F(t)}{t-z}.
\]
The function $H^F$ is holomorphic on $\Omega_-$.  The tower convergence identifies the same
subsequential limit with $\GXi^F$ on $U_-$.  Hence $H^F=\GXi^F$ on $U_-$ and, by the identity theorem,
on all of $\Omega_-$ after
removing the poles of the meromorphic target.  Lemma~\ref{lem:residue} detects every non-real zero by
some finite channel, so no such pole can exist in $\Omega_-$.  The evenness and real symmetry of $\Xi$
reflect the conclusion to the upper half-plane by Lemma~\ref{lem:xi-basic}(ii)--(iii):
$\Xi(-z)=\Xi(z)$ gives $M_\Xi(-z)=-M_\Xi(z)$ and $\Xi(\bar z)=\overline{\Xi(z)}$ gives
$M_\Xi(\bar z)=\overline{M_\Xi(z)}$ away from poles.  Therefore all zeros of $\Xi$ lie on the real
$z$-axis, which is RH.
\end{proof}


\subsection{Step 6: positive Cauchy transforms have positive Pick kernels}

Steps 6--15 are carried out formally in Lemma~\ref{lem:cauchy-pick}, Definition~\ref{def:ARP},
Theorem~\ref{thm:tower-implies-pick}, Lemma~\ref{lem:pinned-normality}, Theorem~\ref{thm:pick-rh},
Lemma~\ref{lem:residue}, Theorem~\ref{thm:rh-pick}, and Corollary~\ref{cor:live-equivalence}; the
surrounding discussion explains their content but the formal statements are authoritative.

Let $\Sigma$ be a positive finite matrix measure on $\R$ and
\[
   g(z)=\int_\R\frac{d\Sigma(t)}{t-z}.
\]
Then, for real $t$,
\[
   \frac{1}{t-z}-\frac{1}{t-\bar w}
   =
   \frac{z-\bar w}{(t-z)(t-\bar w)}.
\]
Hence
\[
   \frac{g(z)-g(w)^*}{z-\bar w}
   =
   \int_\R\frac{d\Sigma(t)}{(t-z)(t-\bar w)}.
\]
For finite nodes $z_i$ and vectors $c_i$,
\[
   \sum_{i,j}c_i^*
   \left[
      \int_\R\frac{d\Sigma(t)}{(t-z_i)(t-\bar z_j)}
   \right]c_j
   =
   \int_\R
      \left(\sum_i\frac{c_i}{t-\bar z_i}\right)^*
      d\Sigma(t)
      \left(\sum_j\frac{c_j}{t-\bar z_j}\right)\ge0.
\]
This proves the Pick positivity of all positive spectral Cauchy transforms.


\begin{lemma}[Cauchy transforms have positive Pick kernels]\label{lem:cauchy-pick}
Let $\Sigma$ be a positive finite matrix measure on $\R$ and
\[
   g(z)=\int_\R\frac{d\Sigma(t)}{t-z}.
\]
Then, for all $z,w\in\C\setminus\R$ with $z\ne\bar w$ \textnormal{(in particular for all $z,w$ in the
same half-plane)}, the integral defining $g$ converges absolutely since
$|t-z|\ge|\Im z|$ and $\Sigma(\R)<\infty$, and
\[
   \frac{g(z)-g(w)^*}{z-\bar w}
   =
   \int_\R
      \frac{d\Sigma(t)}{(t-z)(t-\bar w)}
\]
as a positive matrix kernel.
\end{lemma}

\begin{proof}
Since $t$ is real,
\[
   \frac{1}{t-z}-\frac{1}{t-\bar w}
   =
   \frac{z-\bar w}{(t-z)(t-\bar w)}.
\]
Dividing by $z-\bar w$ gives the displayed identity.  If $z_1,\ldots,z_N$ are
nodes and $c_1,\ldots,c_N$ are vectors in the channel space, then
\[
   \sum_{i,j}
   c_i^*
   \left[
      \int_\R\frac{d\Sigma(t)}{(t-z_i)(t-\bar z_j)}
   \right]
   c_j
   =
   \int_\R
      \left(\sum_i\frac{c_i}{t-\bar z_i}\right)^*
      d\Sigma(t)
      \left(\sum_j\frac{c_j}{t-\bar z_j}\right)
   \ge0.
\]
Thus the block Pick matrices are positive semidefinite.
\end{proof}

\subsection{Step 7: tower-ARP implies ARP-P}

If a positive tower satisfies tower-ARP, then each finite cutoff produces a positive spectral Cauchy
transform $M_P^F$.  By Step~6 every finite Pick matrix of $M_P^F$ is positive semidefinite.
Tower-ARP gives pointwise convergence
\[
   M_P^F(z_i)\longrightarrow g^F(z_i)
\]
for each finite node set.  The cone of positive semidefinite matrices is closed, so
\[
   \mathsf P[g^F;z_1,\ldots,z_N]\succeq0.
\]
Thus tower-ARP is a stronger possible certificate for ARP-P, but it is not required by the minimal
route.


\begin{definition}[positive tower; tower-ARP]\label{def:ARP}
A \emph{positive tower} for the finite channel $F$ is a sequence $\{\Sigma_P^F\}_P$ of positive
finite matrix measures on $\R$ with $\sup_P\Sigma_P^F(\R)<\infty$, whose Cauchy transforms
\[
   M_P^F(z)=\int_\R\frac{d\Sigma_P^F(t)}{t-z}
\]
converge pointwise on $U_-$ to $g^F(z)$.  The channel data satisfy \emph{tower-ARP} if every finite
source plane $F\subset\cS_a^\circ$ admits a positive tower.  Tower-ARP is a convergence property, not a
per-cutoff identity.
\end{definition}

\begin{remark}[operator-theoretic towers]\label{rmk:operator-towers}
Any family of self-adjoint operators $A_P$ on Hilbert spaces $\cH_P$ with $P$-uniformly bounded source
maps $\Phi_P^F:F\to\cH_P$ produces the measures of Definition~\ref{def:ARP} through the spectral
compressions $\Sigma_P^F(B)=(\Phi_P^F)^*E_P(B)\Phi_P^F$, for which
$M_P^F(z)=(\Phi_P^F)^*(A_P-z)^{-1}\Phi_P^F$.  The von Mangoldt canonical-system realization described
above is one candidate family; the realization layer of Section~\ref{sec:shorting} records its source
bounds.
The chain uses nothing beyond Definition~\ref{def:ARP}.
\end{remark}

\begin{theorem}[tower-ARP implies ARP-P]\label{thm:tower-implies-pick}
If the channel data satisfy tower-ARP, Definition~\ref{def:ARP}, then ARP-P holds.
\end{theorem}

\begin{proof}
Each $M_P^F$ is, by definition, the Cauchy transform of a positive finite matrix measure.  Hence
Lemma~\ref{lem:cauchy-pick} gives
\[
   \mathsf P[M_P^F;z_1,\ldots,z_N]\succeq0
\]
for every finite node set in $U_-$.  Tower-ARP gives $M_P^F(z_i)\to g^F(z_i)$ at each node.  Since the
cone of positive semidefinite matrices is closed, the limit matrix
$\mathsf P[g^F;z_1,\ldots,z_N]$ is positive semidefinite.
\end{proof}

\subsection{Step 8: ARP-P gives Nevanlinna--Pick interpolants}

Fix a finite channel $F$ and a finite node set $z_1,\ldots,z_N\subset U_-$.  ARP-P says exactly that the
matrix Pick condition for the lower half-plane is satisfied.  The classical matrix-valued
Nevanlinna--Pick theorem therefore supplies a holomorphic matrix function $F_N$ on $\Omega_-$ such that
\[
   F_N(z_j)=g^F(z_j),\qquad j=1,\ldots,N,
\]
and
\[
   \frac{F_N(z)-F_N(w)^*}{z-\bar w}\succeq0.
\]
The hypotheses of the classical theorem are precisely the finite node condition, the target matrices
$g^F(z_j)$, and the positive semidefiniteness of the displayed block Pick matrix; these are exactly
provided by Step~5.

The classical theorem is used in its degenerate form: positive \emph{semi}definiteness of the block
Pick matrix, singular cases included, is necessary and sufficient for solvability in the matrix Pick
class.  The exact statement used, in the half-plane normalization of this paper and with the Cayley
congruence written out, is displayed as Proposition~\ref{prop:np-degenerate} below; its disk core is the
commutant-lifting Schur--Pick theorem \cite{FoiasFrazho}.  ARP-P supplies
exactly this semidefiniteness, so no strict positivity is assumed anywhere in the bridge.

\begin{lemma}[pinned normality of the matrix Pick class]\label{lem:pinned-normality}
Let $\{F_n\}$ be holomorphic $M_m(\C)$-valued functions on $\Omega_-$ with $\Im F_n\preceq0$, and
suppose $F_n(z_1)=W_1$ for all $n$, at some fixed $z_1\in\Omega_-$.  Then a subsequence converges
locally uniformly on $\Omega_-\setminus Z$, for a discrete set $Z$, to a function which extends
holomorphically to all of $\Omega_-$ as a matrix Pick-class function $F$; moreover
$F(z)=\lim_kF_{n_k}(z)$ whenever the right side eventually stabilizes or converges.
\end{lemma}

\begin{proof}
\emph{Cayley step.}  For $W$ with $\Im W\preceq0$ one computes, for every vector $x$,
\[
   \|(W-iI)x\|^2=\|Wx\|^2+\|x\|^2-2\Im\langle Wx,x\rangle\ge\|x\|^2,
   \qquad
   \|(W+iI)x\|^2\le\|(W-iI)x\|^2 .
\]
Thus $W-iI$ is invertible with $\|(W-iI)^{-1}\|\le1$ and
$S:=(W+iI)(W-iI)^{-1}$ is a contraction, with the exact identities
\[
   S-I=2i(W-iI)^{-1},\qquad W=i(S+I)(S-I)^{-1}.
\]
Thus $S_n(z):=(F_n(z)+iI)(F_n(z)-iI)^{-1}$ is a holomorphic contraction-valued family.

\emph{Montel and pinning.}  The entries of $S_n$ are bounded by $1$, so a diagonal subsequence
converges locally uniformly to a holomorphic contraction-valued $S_\infty$.  At the pinned node,
$S_\infty(z_1)-I=2i(W_1-iI)^{-1}$ is invertible.  Hence
$\varphi:=\det(S_\infty-I)$ is holomorphic and not identically zero, and
$Z:=\{\varphi=0\}$ is discrete in $\Omega_-$.

\emph{Back-transform off $Z$.}  On compacts of $\Omega_-\setminus Z$, $|\varphi|$ is bounded below,
so $(S_{n_k}-I)^{-1}\to(S_\infty-I)^{-1}$ uniformly and
$F_{n_k}\to F:=i(S_\infty+I)(S_\infty-I)^{-1}$ locally uniformly; $F$ is holomorphic with
$\Im F\preceq0$ on $\Omega_-\setminus Z$ as a limit of such.

\emph{Removability across $Z$.}  Fix $\zeta_0\in Z$ and a unit vector $c$, and set
$f:=c^*Fc$ on a punctured disk around $\zeta_0$ inside $\Omega_-\setminus Z$.  Since
$\Im f\le0$, $|f-i|\ge1$, so $h:=(f-i)^{-1}$ is holomorphic and bounded by $1$, hence extends across
$\zeta_0$ by Riemann's theorem.  If $h(\zeta_0)=0$, then $f=i+1/h$ would have a pole at $\zeta_0$;
near a pole the values of $f$ cover all sufficiently large moduli in every direction of the plane, in
particular values with positive imaginary part, contradicting $\Im f\le0$.  Hence
$h(\zeta_0)\ne0$ and $f$ extends holomorphically.  Applying this to
$c\in\{e_j,\,e_k,\,e_j+e_k,\,e_j+ie_k\}$ and polarizing recovers every entry of $F$, which therefore
extends holomorphically across $Z$; the extension is Pick-class by continuity.

\emph{Node values.}  If $F_{n_k}(w)=W$ for all large $k$ at some $w\in\Omega_-$, then when
$w\notin Z$ this passes to the limit directly.  When $w\in Z$, choose a circle around $w$ inside
$\Omega_-\setminus Z$ and use the Cauchy integral:
\[
   F(w)=\frac1{2\pi i}\oint\frac{F(\zeta)}{\zeta-w}\,d\zeta
   =\lim_k\frac1{2\pi i}\oint\frac{F_{n_k}(\zeta)}{\zeta-w}\,d\zeta=W,
\]
since each $F_{n_k}$ is holomorphic at $w$.
\end{proof}

\begin{proposition}[degenerate matrix Nevanlinna--Pick theorem, half-plane form]\label{prop:np-degenerate}
Let $z_1,\ldots,z_N\in\Omega_-$ be pairwise distinct and let $W_1,\ldots,W_N\in M_m(\C)$.  There exists
a holomorphic $F:\Omega_-\to M_m(\C)$ with $\Im F\preceq0$ and $F(z_j)=W_j$ for all $j$ if and only if
the block Pick matrix
\[
   \left[\frac{W_i-W_j^*}{z_i-\bar z_j}\right]_{i,j=1}^N
\]
is positive \emph{semi}definite; singular (degenerate) Pick matrices are included.
\end{proposition}

\begin{proof}
\emph{Reduction to the disk.}  The M\"obius map $\zeta(z)=(z+i)/(z-i)$ takes $\Omega_-$ biholomorphically
onto the unit disk $\mathbb D$, and the node images $\zeta_j:=\zeta(z_j)$ are pairwise distinct.  If the Pick matrix is
positive semidefinite, its diagonal blocks give $(W_j-W_j^*)/(2i\Im z_j)\succeq0$ with $\Im z_j<0$, hence
$\Im W_j\preceq0$; by the Cayley computation in the proof of Lemma~\ref{lem:pinned-normality},
$W_j-iI$ is then invertible and $S_j:=(W_j+iI)(W_j-iI)^{-1}$ is a contraction, with
$S_j-I=2i(W_j-iI)^{-1}$ invertible.  The algebraic identities
\[
   I-S_iS_j^*
   =
   2i\,(W_i-iI)^{-1}(W_i-W_j^*)\,(W_j^*+iI)^{-1},
   \qquad
   1-\zeta_i\bar\zeta_j
   =
   \frac{2i\,(z_i-\bar z_j)}{(z_i-i)(\bar z_j+i)},
\]
verified by direct expansion (the factors $(W\pm iI)$ and their inverses commute), give the exact block
congruence
\[
   \left[\frac{I-S_iS_j^*}{1-\zeta_i\bar\zeta_j}\right]
   =
   D\left[\frac{W_i-W_j^*}{z_i-\bar z_j}\right]D^*,
   \qquad
   D=\diag\bigl((z_i-i)(W_i-iI)^{-1}\bigr)_{i=1}^N,
\]
with $D$ invertible.  Hence the half-plane Pick matrix is positive semidefinite exactly when the disk
Schur--Pick matrix is.

\emph{Sufficiency.}  Assume the Pick matrix is positive semidefinite.  By the classical matrix
Schur--Pick theorem via commutant lifting \cite{FoiasFrazho} --- valid verbatim for positive
\emph{semi}definite, possibly singular, Pick matrices --- there is a holomorphic contraction-valued
$S:\mathbb D\to M_m(\C)$ with $S(\zeta_j)=S_j$.  Since $S(\zeta_1)-I=S_1-I$ is invertible,
$\varphi:=\det(S\circ\zeta-I)\not\equiv0$ and its zero set $Z\subset\Omega_-$ is discrete and avoids
every node.  On $\Omega_-\setminus Z$ define $F:=i(S\circ\zeta+I)(S\circ\zeta-I)^{-1}$.  For a
contraction $S_0$ with $S_0-I$ invertible and $u:=(S_0-I)^{-1}x$,
\[
   \Im\langle i(S_0+I)(S_0-I)^{-1}x,\,x\rangle
   =
   \|S_0u\|^2-\|u\|^2\le0,
\]
so $\Im F\preceq0$ off $Z$.  The removability step in the proof of
Lemma~\ref{lem:pinned-normality} (scalar compressions $c^*Fc$, boundedness of $(c^*Fc-i)^{-1}$, and
polarization) extends $F$ holomorphically across $Z$ with $\Im F\preceq0$ on all of $\Omega_-$.  At
the nodes, the exact inverse-Cayley identity $W=i(S+I)(S-I)^{-1}$ of Lemma~\ref{lem:pinned-normality}
gives $F(z_j)=W_j$.

\emph{Necessity.}  If such an $F$ exists, then $S:=(F+iI)(F-iI)^{-1}\circ\zeta^{-1}$ is a holomorphic
contraction on $\mathbb D$ by the same Cayley computation, the disk Schur--Pick matrix at
$(\zeta_j,S(\zeta_j))$ is positive semidefinite by the classical theorem
\cite{FoiasFrazho}, and the displayed congruence transfers the positivity back to the half-plane Pick
matrix.
\end{proof}

\begin{lemma}[prefix serialization: single-channel Pick positivity removes all poles]\label{lem:prefix-extension}
Let $G$ be an $M_m(\C)$-valued meromorphic function on $\Omega_-$.  Let $\{z_j\}_{j\ge1}\subset\Omega_-$
be pairwise distinct with an accumulation point $z_*\in\Omega_-$, and suppose $G$ is holomorphic at every
$z_j$ and in a neighborhood of $z_*$.  If
\[
   \mathsf P[G;z_1,\ldots,z_N]
   =
   \left[\frac{G(z_i)-G(z_j)^*}{z_i-\bar z_j}\right]_{i,j=1}^N
   \succeq0
   \qquad\text{for every }N\ge1,
\]
then every pole of $G$ in $\Omega_-$ is removable, and the resulting holomorphic extension is a matrix
Pick-class function ($\Im G\preceq0$) on $\Omega_-$.
\end{lemma}

\begin{proof}
For each $N$, Proposition~\ref{prop:np-degenerate} supplies a holomorphic $F_N$ on $\Omega_-$ with
$\Im F_N\preceq0$ and $F_N(z_j)=G(z_j)$ for $j\le N$.  Every $F_N$ interpolates the pinned node $z_1$
with the common value $G(z_1)$, so Lemma~\ref{lem:pinned-normality} produces a subsequence converging to
a holomorphic matrix Pick-class function $H$ on $\Omega_-$ with $H(z_j)=G(z_j)$ for every $j$, each node
being eventually interpolated.  Both $H$ and $G$ are holomorphic near $z_*$, and $\{z_j\}$ accumulates at
$z_*$; the identity theorem gives $H=G$ on a neighborhood of $z_*$.  Since $\Omega_-$ is connected and
$G$ is meromorphic there, $G=H$ on $\Omega_-$ minus the polar set of $G$; hence every pole of $G$ in
$\Omega_-$ is removable and the extension is $H$.
\end{proof}

\begin{theorem}[ARP-P implies RH]\label{thm:pick-rh}
Assume ARP-P.  Then every nontrivial zero of $\zeta$ has real part $\half$.
\end{theorem}

\begin{proof}
Fix a finite channel $F$ and choose a countable, pairwise distinct node set
$\{z_j\}_{j\ge1}\subset U_-$ with an accumulation point $z_*\in U_-$.  ARP-P gives
$\mathsf P[g^F;z_1,\ldots,z_N]\succeq0$ for every $N$.  By Lemma~\ref{lem:GXi-meromorphic}, $\GXi^F$ is
meromorphic on $\C$ and equals $g^F$ on $U_-$; its poles lie in the zero strip $|\Im z|<\half$
(Lemma~\ref{lem:xi-basic}(i)), so it is holomorphic at every node and near $z_*$.
Lemma~\ref{lem:prefix-extension} therefore removes every pole of $\GXi^F$ in $\Omega_-$.
Lemma~\ref{lem:residue} detects every non-real zero in some source
channel, so no zero of $\Xi$ can correspond to a point in $\Omega_-$.  The symmetries
$M_\Xi(-z)=-M_\Xi(z)$ and $M_\Xi(\bar z)=\overline{M_\Xi(z)}$ reflect the conclusion to the upper
half-plane.  Therefore all zeros of $\Xi$ lie on the real $z$-axis, which is RH.
\end{proof}

\subsection{Step 9: NP normality and Montel extension}

Choose a countable node set $\{z_j\}_{j\ge1}\subset U_-$ with an accumulation point in $U_-$.  For each
$N$, Step~8 gives a Pick-class interpolant $F_N$ through the first $N$ nodes.  By
Lemma~\ref{lem:pinned-normality}, applied with the pinned node $z_1$ interpolated by every $F_N$, a
subsequence converges to a holomorphic matrix Pick-class function $H$ on $\Omega_-$, and
$H(z_j)=g^F(z_j)$ for every $j$, since each node is eventually interpolated.  Since the node set
accumulates in $U_-$ and both $H$ and $g^F$ are holomorphic there, the identity theorem gives
\[
   H=g^F\qquad\text{on }U_-.
\]
Thus $g^F$ has a holomorphic extension to all of $\Omega_-$.


\begin{theorem}[finite-channel Vitali bridge]\label{thm:vitali}
Let $M_P^F:\Omega_-\to M_m(\C)$ be holomorphic functions satisfying
\[
   \|M_P^F(z)\|\le \frac{C_F}{|\Im z|}
\]
locally uniformly in $P$.  If $M_P^F(z)\to\GXi^F(z)$ pointwise on the nonempty open set $U_-$, where
$\GXi^F$ is meromorphic on $\Omega_-$, then $\GXi^F$ has no pole in $\Omega_-$.
\end{theorem}

\begin{proof}
The estimate makes $\{M_P^F\}$ a normal family on $\Omega_-$ by Montel.  Let $H$ be any subsequential
locally uniform limit.  Then $H$ is holomorphic on $\Omega_-$.  On $U_-$ the full sequence converges to
$\GXi^F$, hence $H=\GXi^F$ on $U_-$.  The identity theorem forces every subsequential limit to coincide on
all of $\Omega_-$.  The full sequence therefore converges locally uniformly to a holomorphic extension of
$\GXi^F|_{U_-}$, and all poles of $\GXi^F$ in $\Omega_-$ are removable.
\end{proof}

\subsection{Step 10: residue detection by channels}

Let $z_\rho$ be a non-real zero of $\Xi$.  The principal part of $\GXi^F$ in a one-dimensional channel
$F=\operatorname{span}\{\phi\}$ is proportional to
\[
   \ell_{z_\rho}(\phi)\overline{\ell_{\bar z_\rho}(\phi)}.
\]
Because $C_c^\infty(0,\infty)\subset\cS_a^\circ$, neither functional $\ell_{z_\rho}$ nor
$\ell_{\bar z_\rho}$ can vanish on the whole source core: if $\ell_a$ vanished on all compactly supported
smooth tests, then the distribution $e^{-iat}$ would vanish on $(0,\infty)$, impossible.  The kernels of
two nonzero linear functionals cannot cover the whole vector space.  Therefore one can choose
$\phi\in C_c^\infty(0,\infty)$ with
\[
   \ell_{z_\rho}(\phi)\overline{\ell_{\bar z_\rho}(\phi)}\ne0.
\]
That channel detects the pole at $z_\rho$.


\begin{lemma}[residue detection by the source core]\label{lem:residue}
For the source core $\cS_a^\circ$ of Definition~\ref{def:core}, and for every non-real pole $z_\rho$ of
$M_\Xi$, there is a one-dimensional source plane
$F=\operatorname{span}\{\phi\}\subset\cS_a^\circ$ such that $\GXi^F$ has nonzero residue at $z_\rho$.
\end{lemma}

\begin{proof}
At a zero $z_\rho$ of multiplicity $m_\rho$, $M_\Xi$ has principal part
$-m_\rho/(z-z_\rho)$.  The residue in a source channel is a scalar multiple of
$\ell_{z_\rho}(\phi)\overline{\ell_{\bar z_\rho}(\psi)}$, where in the logarithmic model
\[
   \ell_a(\phi)=\int_0^\infty \phi(t)e^{-iat}\,dt.
\]
By Definition~\ref{def:core}, $C_c^\infty(0,\infty)\subset\cS_a^\circ$.  If $\ell_{z_\rho}$ vanished on
all $C_c^\infty(0,\infty)$, the locally integrable distribution $e^{-iz_\rho t}$ would be zero on
$(0,\infty)$.  This is impossible: on any compact interval it is a continuous function, and a continuous
function whose distribution is zero is identically zero, whereas $e^{-iz_\rho t}\ne0$ for every real
$t$.  The same argument applies to $\ell_{\bar z_\rho}$.  Hence the kernels of $\ell_{z_\rho}$ and
$\ell_{\bar z_\rho}$ are two proper linear subspaces of $C_c^\infty(0,\infty)$.  Their union cannot be the
whole vector space: if $u$ is outside the first kernel and $v$ is outside the second, then either $u$ is
outside both kernels or $v$ is outside both kernels or $u+v$ is outside both.  Choose
$\phi\in C_c^\infty(0,\infty)\subset\cS_a^\circ$ outside both kernels.  Then
$\ell_{z_\rho}(\phi)\overline{\ell_{\bar z_\rho}(\phi)}\ne0$, and for
$F=\operatorname{span}\{\phi\}$ the residue of $\GXi^F$ is therefore nonzero.
\end{proof}

\subsection{Step 11: ARP-P removes off-line poles in \texorpdfstring{$\Omega_-$}{Omega-}}

Assume ARP-P.  Step~9 gives a holomorphic extension of each $\GXi^F|_{U_-}$ to $\Omega_-$.  But Step~3
says $\GXi^F$ is meromorphic with explicitly known principal parts.  Hence every pole of $\GXi^F$ in
$\Omega_-$ is removable.  Step~10 says every non-real zero of $\Xi$ lying in $\Omega_-$ is visible as a
non-removable pole in some finite channel.  Therefore no such zero exists.

\subsection{Step 12: symmetries cover all of \texorpdfstring{$\C\setminus\R$}{C minus R}}

By Lemma~\ref{lem:xi-basic}(ii)--(iii), the completed function satisfies $\Xi(-z)=\Xi(z)$ and is real on
the real axis.  Consequently
\[
   M_\Xi(-z)=-M_\Xi(z),\qquad
   M_\Xi(\bar z)=\overline{M_\Xi(z)}
\]
away from poles.  A non-real zero in the upper half-plane reflects by these symmetries to a non-real zero
in $\Omega_-$.  Step~11 rules out the latter under ARP-P.  Hence ARP-P rules out all non-real zeros of
$\Xi$.

\subsection{Step 13: real zeros of \texorpdfstring{$\Xi$}{Xi} are RH}

By Step~1, a zero $\rho=\beta+i\gamma$ corresponds to $z_\rho=\gamma-i(\beta-\half)$.  This number is
real if and only if $\beta=\half$.  Thus all zeros of $\Xi$ being real is exactly the Riemann Hypothesis.


\begin{theorem}[analytic bridge implies RH]\label{thm:bridge-rh}
Assume that for every finite $F\subset\cS_a^\circ$ there are functions $M_P^F$ satisfying the hypotheses
of Theorem~\ref{thm:vitali} and converging to $\GXi^F$ on $U_-$.  Then RH holds.
\end{theorem}

\begin{proof}
If $M_\Xi$ had a pole in $\Omega_-$, Lemma~\ref{lem:residue} would choose a finite channel where that pole
is visible, contradicting Theorem~\ref{thm:vitali}.  By Lemma~\ref{lem:xi-basic}(ii)--(iii), poles in
$\C^+$ are reflected from poles in $\Omega_-$.  Explicitly, $\Xi(-z)=\Xi(z)$ gives
\[
   M_\Xi(-z)=-M_\Xi(z),
\]
and the reality of $\Xi$ on $\R$ gives $M_\Xi(\bar z)=\overline{M_\Xi(z)}$ away from poles.  Hence a
non-real pole in any quadrant reflects to a pole in $\Omega_-$.  Therefore every pole of
$M_\Xi=-\Xi'/\Xi$ is real, so every zero of $\Xi$ is real.  By Lemma~\ref{lem:rh-translation}, this is
RH.
\end{proof}

\subsection{Step 14: RH implies ARP-P}

Assume RH.  Then every $z_\rho$ is real, the target becomes the Cauchy transform of a positive finite
matrix measure supported on the real zero coordinates, and Step~6 gives the Pick positivity of every
finite node matrix.  The formal statement and proof follow.

\begin{theorem}[RH implies ARP-P]\label{thm:rh-pick}
Assume RH.  Then ARP-P holds.
\end{theorem}

\begin{proof}
Under RH every zero coordinate $z_\rho$ is real.  For each finite source plane $F$, Lemma~\ref{lem:GXi-meromorphic}
and the rapid decay of the source functionals give a positive finite matrix measure
\[
   \Sigma_\Xi^F
   =
   \sum_\rho
   m_\rho
   \left[
      \ell_{z_\rho}(\phi_\beta)\overline{\ell_{z_\rho}(\phi_\alpha)}
   \right]_{\alpha,\beta}
   \delta_{z_\rho}.
\]
The total mass is finite because the entries decay faster than any power along the real zero set and
$N(T)=O(T\log T)$.  The target is its Cauchy transform:
under RH $\bar z_\rho=z_\rho$, so the general channel weight
$\ell_{z_\rho}(\phi_\beta)\overline{\ell_{\bar z_\rho}(\phi_\alpha)}$ in \eqref{eq:GXiF} is exactly the
displayed rank-one positive atom.  Thus
\[
   g^F(z)=\int_\R\frac{d\Sigma_\Xi^F(t)}{t-z}.
\]
Lemma~\ref{lem:cauchy-pick} gives the positive semidefiniteness of every finite Pick matrix.  Hence
ARP-P holds.
\end{proof}

\subsection{Step 15: ARP-P \texorpdfstring{$\Longleftrightarrow$}{<=>} RH}

Steps 8--13 prove
\[
   \textnormal{ARP-P}\Longrightarrow\textnormal{RH}.
\]
Step~14 proves
\[
   \textnormal{RH}\Longrightarrow\textnormal{ARP-P}.
\]
Therefore the route has the exact structural status
\[
   \boxed{\textnormal{ARP-P}\Longleftrightarrow\textnormal{RH}}.
\]
This is not a claim that RH is proved; it is a proof that the remaining mathematical work has been
localized to the arithmetic Pick positivity of Step~5.


\begin{corollary}[the ARP-P equivalence]\label{cor:live-equivalence}
\[
   \textnormal{ARP-P}\Longleftrightarrow\textnormal{RH}.
\]
\end{corollary}

\begin{proof}
The forward implication is Theorem~\ref{thm:pick-rh}.  The reverse implication is
Theorem~\ref{thm:rh-pick}.
\end{proof}


\begin{corollary}[exact status of the remaining gap]\label{cor:gap-status}
ARP-P is equivalent to RH and is implied by tower-ARP:
\[
   \textnormal{tower-ARP}\Longrightarrow\textnormal{ARP-P}\Longleftrightarrow\textnormal{RH}.
\]
Thus the remaining essential input of the route is not FCC, determinant normality, or a per-cutoff
trace-log identity.  It is the arithmetic Pick positivity of Definition~\ref{def:ARP-P}.
\end{corollary}


\section{Epilogue}\label{sec:epilogue}

It is worth saying plainly what made this hard, and where, honestly, it still is hard.

The programme's real product is not a finished proof but a \emph{map}: a precise account of where the
difficulty of the Riemann Hypothesis lives. Part~I is the negative geography --- the no-gos and the seven
structural walls, each carrying its explicit logical classification (equivalence, framework obstruction,
or insufficiency theorem; \S\ref{sec:walls}), all unified by the single master quantifier (the passage
from average, almost-all, or interior to specific, uniform, or boundary). Part~II is the proof path those
walls shaped: the ARP-P chain, which does not try to cross any wall but instead \emph{reduces} the
Hypothesis, step by step, to one inequality --- the terminal-defect positivity $\Omega_7$, $\delta_N\ge0$
for all $N$, equivalently $\lambda_n\ge0$ for all $n$. Fourteen of the fifteen Steps of the implication
chain are closed in the precise sense recorded in the status map (\S\ref{sec:status-map}); the sole
unresolved input occurs within Step~5. The internal reduction of Step~5 terminates at $\Omega_7$, which
is equivalent to Li--Keiper positivity and hence to RH; we do not close it, and we say so.

The value of the reduction is that it names the residue exactly. When the roads explored in this
programme repeatedly ended at the same wall,
the task was not to find another road but to understand the wall so precisely that one can state it as a
single inequality, and then to attack \emph{that}. We attacked it, from four independent directions ---
an archimedean partition, a stationary-phase saddle, a positive de Branges kernel, and a perturbative
archimedean domination --- and each returned, provably, to the walls of Part~I: the partition to the
uniform-norm wall, the perturbative bound to the wrong-sign wall (a logarithmic margin against an
unbounded perturbation), the positive kernel to the Weil-positivity wall (the de Branges positivity
conditions being falsified by Conrey--Li), and the saddle to the same magnitude-blind cancellation the
programme met at the start. That $\Omega_7$ resists all four in the \emph{same} way is not a failure of
effort; it is the sharpest statement the programme can make about why the last step is the whole problem.

The road to this map was long and, for most of its length, negative. The author recalls printing a
particular computation three times because each printing revealed a sign error introduced in correcting
the previous one; recalls confirming, once again, that positivity is orthogonal to magnitude, and later
that it is a \emph{sign} property no size functional can see; recalls a great many letters and emails,
some of them stern, every one of them useful, in which an objection one had hoped was pedantic turned out
to be fatal and, in being fatal, pointed at the next door. The chain in Part~II is built almost entirely
out of the wreckage of those objections.

A last word, in the same register of honesty in which the programme was conducted. We are aware of how
many careful arguments for this problem have failed at a step their authors believed secure, and we do not
exempt ourselves --- two of our own results are recorded above as withdrawn. We have written every closed
step as a proved statement and named the classical facts it rests on; we have isolated the one place where
all the weight is carried, $\Omega_7$, and asked, explicitly, that it be attacked. What we offer is a
complete \emph{reduction} and an honest open problem --- not a settled theorem. A result of this kind is
not finished when its authors are persuaded; it is finished when the community is. The standard we ask
the reader to hold this paper to is the one recorded in \S\ref{sec:withdrawn}: read each claim at the
exact strength of its classification, and where a claim fails that reading, treat it as we treated our
own withdrawn results. We submit the programme and the path in that spirit.

\appendix


\begin{thebibliography}{99}

%% --- Foundational and classical ---
\bibitem{Riemann} B.~Riemann, \emph{\"Uber die Anzahl der Primzahlen unter einer gegebenen Gr\"osse},
Monatsber.\ Berliner Akad.\ (1859), 671--680.
\bibitem{Weil1952} A.~Weil, \emph{Sur les ``formules explicites'' de la th\'eorie des nombres premiers},
Comm.\ S\'em.\ Math.\ Univ.\ Lund (1952), 252--265.
\bibitem{Weil1948} A.~Weil, \emph{Sur les courbes alg\'ebriques et les vari\'et\'es qui s'en d\'eduisent},
Actualit\'es Sci.\ Ind.\ \textbf{1041}, Hermann, Paris, 1948.
\bibitem{Titchmarsh} E.~C.~Titchmarsh, \emph{The Theory of the Riemann Zeta-Function}, 2nd ed.\ (rev.\
D.~R.~Heath-Brown), Oxford Univ.\ Press, 1986.
\bibitem{Edwards} H.~M.~Edwards, \emph{Riemann's Zeta Function}, Academic Press, 1974; Dover, 2001.
\bibitem{Davenport} H.~Davenport, \emph{Multiplicative Number Theory}, 3rd ed., Grad.\ Texts in Math.\
\textbf{74}, Springer, 2000.
\bibitem{Hardy1914} G.~H.~Hardy, \emph{Sur les z\'eros de la fonction $\zeta(s)$ de Riemann}, C.~R.\ Acad.\
Sci.\ Paris \textbf{158} (1914), 1012--1014.

%% --- de Branges / Krein-Langer / functional analysis ---
\bibitem{deBranges} L.~de~Branges, \emph{Hilbert Spaces of Entire Functions}, Prentice--Hall, Englewood
Cliffs NJ, 1968.
\bibitem{KreinLanger} M.~G.~Kre\u{\i}n and H.~Langer, \emph{\"Uber die verallgemeinerten Resolventen und
die charakteristische Funktion eines isometrischen Operators im Raume $\Pi_\kappa$}, in: Hilbert Space
Operators and Operator Algebras, Colloq.\ Math.\ Soc.\ J\'anos Bolyai \textbf{5} (1972), 353--399.
\bibitem{Aronszajn} N.~Aronszajn, \emph{Theory of reproducing kernels}, Trans.\ Amer.\ Math.\ Soc.\
\textbf{68} (1950), 337--404.
\bibitem{Bognar} J.~Bogn\'ar, \emph{Indefinite Inner Product Spaces}, Springer, 1974.
\bibitem{Azizov} T.~Ya.~Azizov and I.~S.~Iokhvidov, \emph{Linear Operators in Spaces with an Indefinite
Metric}, Wiley, 1989.
\bibitem{ReedSimon} M.~Reed and B.~Simon, \emph{Methods of Modern Mathematical Physics}, vols.\ I, IV,
Academic Press, 1972--1978.
\bibitem{Simon} B.~Simon, \emph{Trace Ideals and Their Applications}, 2nd ed., Amer.\ Math.\ Soc., 2005.
\bibitem{Akhiezer} N.~I.~Akhiezer, \emph{The Classical Moment Problem and Some Related Questions in
Analysis}, Oliver \& Boyd, 1965.
\bibitem{Remmert} R.~Remmert, \emph{Classical Topics in Complex Function Theory}, Grad.\ Texts in Math.\
\textbf{172}, Springer, 1998. \emph{(Montel/Vitali normal families.)}
\bibitem{HornJohnson} R.~A.~Horn and C.~R.~Johnson, \emph{Matrix Analysis}, 2nd ed., Cambridge Univ.\
Press, 2013. \emph{(Sylvester's law of inertia; Haynsworth additivity.)}

%% --- Adelic--spectral (trace formula / arithmetic site / spectral triples) ---
\bibitem{connes1999} A.~Connes, \emph{Trace formula in noncommutative geometry and the zeros of the
Riemann zeta function}, Selecta Math.\ (N.S.) \textbf{5} (1999), no.~1, 29--106. arXiv:math/9811068.
\bibitem{connesconsani2021} A.~Connes and C.~Consani, \emph{Weil positivity and trace formula: the
archimedean place}, Selecta Math.\ (N.S.) \textbf{27} (2021), no.~77.
\bibitem{connesconsani-site} A.~Connes and C.~Consani, \emph{The scaling site}, J.\ Reine Angew.\ Math.\
\textbf{755} (2019), 1--43; \emph{The arithmetic site}, C.~R.\ Acad.\ Sci.\ Paris (2014). arXiv:1507.05818.
\bibitem{ccm-prolate} A.~Connes and C.~Consani, \emph{Spectral triples and $\zeta$-cycles}, and related
work on the prolate/Sonin space, arXiv:2006.13771.
\bibitem{ccm-zeta} A.~Connes, C.~Consani, and H.~Moscovici, \emph{Zeta spectral triples} (2025),
arXiv:2511.22755; and \emph{Zeta zeros and prolate wave operators}, Ann.\ Funct.\ Anal.\ \textbf{15}
(2024), no.~4, Paper No.~71.
\bibitem{bostconnes} J.-B.~Bost and A.~Connes, \emph{Hecke algebras, type III factors and phase
transitions with spontaneous symmetry breaking in number theory}, Selecta Math.\ (N.S.) \textbf{1}
(1995), no.~3, 411--457.
\bibitem{Deligne} P.~Deligne, \emph{La conjecture de Weil. I}, Publ.\ Math.\ IH\'ES \textbf{43} (1974),
273--307.
\bibitem{Deninger} C.~Deninger, \emph{Some analogies between number theory and dynamical systems on foliated
spaces}, in: Proceedings of the International Congress of Mathematicians, Vol.\ I (Berlin, 1998), Doc.\ Math.\
Extra Vol.\ I (1998), 163--186.

%% --- Analytic number theory and the zeros ---
\bibitem{Montgomery} H.~L.~Montgomery, \emph{The pair correlation of zeros of the zeta function}, in:
Analytic Number Theory, Proc.\ Sympos.\ Pure Math.\ \textbf{24}, Amer.\ Math.\ Soc., 1973, 181--193.
\bibitem{Selberg} A.~Selberg, \emph{On the zeros of Riemann's zeta-function}, Skr.\ Norske Vid.-Akad.\
Oslo I \textbf{10} (1942), 1--59.
\bibitem{Levinson} N.~Levinson, \emph{More than one third of zeros of Riemann's zeta-function are on
$\sigma=\tfrac12$}, Adv.\ Math.\ \textbf{13} (1974), 383--436.
\bibitem{Conrey1989} J.~B.~Conrey, \emph{More than two fifths of the zeros of the Riemann zeta function
are on the critical line}, J.\ Reine Angew.\ Math.\ \textbf{399} (1989), 1--26.
\bibitem{ConreyNotices} J.~B.~Conrey, \emph{The Riemann Hypothesis}, Notices Amer.\ Math.\ Soc.\
\textbf{50} (2003), no.~3, 341--353.
\bibitem{IwaniecKowalski} H.~Iwaniec and E.~Kowalski, \emph{Analytic Number Theory}, Amer.\ Math.\ Soc.\
Colloq.\ Publ.\ \textbf{53}, 2004.
\bibitem{Bombieri} E.~Bombieri, \emph{Problems of the Millennium: The Riemann Hypothesis}, Clay
Mathematics Institute, 2000.
\bibitem{Sarnak} P.~Sarnak, \emph{Problems of the Millennium: The Riemann Hypothesis} (2004), Clay
Mathematics Institute report.
\bibitem{Tate} J.~Tate, \emph{Fourier analysis in number fields and Hecke's zeta-functions}, thesis,
Princeton, 1950; in: Cassels--Fr\"ohlich (eds.), \emph{Algebraic Number Theory}, 1967.

%% --- Spectral / random-matrix / heat-flow ---
\bibitem{BerryKeating} M.~V.~Berry and J.~P.~Keating, \emph{The Riemann zeros and eigenvalue
asymptotics}, SIAM Rev.\ \textbf{41} (1999), no.~2, 236--266.
\bibitem{KeatingSnaith} J.~P.~Keating and N.~C.~Snaith, \emph{Random matrix theory and
$\zeta(\tfrac12+it)$}, Comm.\ Math.\ Phys.\ \textbf{214} (2000), no.~1, 57--89.
\bibitem{RudnickSarnak} Z.~Rudnick and P.~Sarnak, \emph{Zeros of principal $L$-functions and random
matrix theory}, Duke Math.\ J.\ \textbf{81} (1996), no.~2, 269--322.
\bibitem{RodgersTao} B.~Rodgers and T.~Tao, \emph{The de Bruijn--Newman constant is non-negative}, Forum
Math.\ Pi \textbf{8} (2020), e6.
\bibitem{Newman} C.~M.~Newman, \emph{Fourier transforms with only real zeros}, Proc.\ Amer.\ Math.\ Soc.\
\textbf{61} (1976), no.~2, 245--251.

%% --- Beurling--Nyman / Li ---
\bibitem{Beurling} A.~Beurling, \emph{A closure problem related to the Riemann zeta-function}, Proc.\
Nat.\ Acad.\ Sci.\ U.S.A.\ \textbf{41} (1955), 312--314.
\bibitem{BaezDuarte} L.~B\'aez-Duarte, \emph{A strengthening of the Nyman--Beurling criterion for the
Riemann hypothesis}, Atti Accad.\ Naz.\ Lincei \textbf{14} (2003), no.~1, 5--11.
\bibitem{Li} X.-J.~Li, \emph{The positivity of a sequence of numbers and the Riemann hypothesis}, J.\
Number Theory \textbf{65} (1997), no.~2, 325--333.
\bibitem{ConreyLi} J.~B.~Conrey and X.-J.~Li, \emph{A note on some positivity conditions related to zeta
and $L$-functions}, Int.\ Math.\ Res.\ Not.\ \textbf{2000}, no.~18, 929--940.
\bibitem{ScrewWeil} \emph{Weil's quadratic form via the screw function}, preprint (2026), arXiv:2606.09096.

%% --- The falsifier ---
\bibitem{DavenportHeilbronn} H.~Davenport and H.~Heilbronn, \emph{On the zeros of certain Dirichlet
series} I, II, J.\ London Math.\ Soc.\ \textbf{11} (1936), 181--185 and 307--312.

%% --- Large values, zero-density, zero-free regions ---
\bibitem{Ingham} A.~E.~Ingham, \emph{On the estimation of $N(\sigma,T)$}, Quart.\ J.\ Math.\ Oxford
\textbf{11} (1940), 291--292.
\bibitem{Korobov} N.~M.~Korobov, \emph{Estimates of trigonometric sums and their applications}, Uspekhi
Mat.\ Nauk \textbf{13} (1958), no.~4, 185--192.
\bibitem{Huxley} M.~N.~Huxley, \emph{On the difference between consecutive primes}, Invent.\ Math.\
\textbf{15} (1972), 164--170.
\bibitem{BourgainDemeterGuth} J.~Bourgain, C.~Demeter, and L.~Guth, \emph{Proof of the main conjecture in
Vinogradov's mean value theorem for degrees higher than three}, Ann.\ of Math.\ (2) \textbf{184} (2016),
633--682.
\bibitem{Wooley} T.~D.~Wooley, \emph{The cubic case of the main conjecture in Vinogradov's mean value
theorem}, Adv.\ Math.\ \textbf{294} (2016), 532--561.
\bibitem{GuthMaynard} L.~Guth and J.~Maynard, \emph{New large value estimates for Dirichlet polynomials},
preprint (2024), arXiv:2405.20552.
\bibitem{MTY} M.~J.~Mossinghoff, T.~S.~Trudgian, and A.~Yang, \emph{Explicit zero-free regions for the
Riemann zeta-function}, Res.\ Number Theory \textbf{10} (2024), Paper No.~11.
\bibitem{TTY} T.~Tao, T.~S.~Trudgian, and A.~Yang, \emph{New exponent pairs, zero density estimates, and
zero additive energy estimates}, Math.\ Comp.\ (2025).
\bibitem{PlattTrudgian} D.~J.~Platt and T.~S.~Trudgian, \emph{The Riemann hypothesis is true up to
$3\times10^{12}$}, Bull.\ Lond.\ Math.\ Soc.\ \textbf{53} (2021), no.~3, 792--797.
\bibitem{Odlyzko} A.~M.~Odlyzko, \emph{On the distribution of spacings between zeros of the zeta
function}, Math.\ Comp.\ \textbf{48} (1987), no.~177, 273--308.

%% --- de Bruijn--Newman flow ---
\bibitem{deBruijn} N.~G.~de~Bruijn, \emph{The roots of trigonometric integrals}, Duke Math.\ J.\
\textbf{17} (1950), 197--226.
\bibitem{PolymathDBN} D.~H.~J.~Polymath, \emph{Effective approximation of heat flow evolution of the
Riemann $\xi$ function, and a new upper bound for the de Bruijn--Newman constant}, Res.\ Math.\ Sci.\
\textbf{6} (2019), no.~3, Paper No.~31.

%% --- Hyperbolicity / Jensen polynomials / interlacing ---
\bibitem{GORZ} M.~Griffin, K.~Ono, L.~Rolen, and D.~Zagier, \emph{Jensen polynomials for the Riemann zeta
function and other sequences}, Proc.\ Natl.\ Acad.\ Sci.\ USA \textbf{116} (2019), no.~22, 11103--11110.
\bibitem{MSS} A.~Marcus, D.~A.~Spielman, and N.~Srivastava, \emph{Interlacing families II: mixed
characteristic polynomials and the Kadison--Singer problem}, Ann.\ of Math.\ (2) \textbf{182} (2015),
327--350.
\bibitem{BorceaBranden} J.~Borcea and P.~Br\"and\'en, \emph{The Lee--Yang and P\'olya--Schur programs.
I}, Comm.\ Pure Appl.\ Math.\ \textbf{62} (2009), no.~12, 1595--1631.

%% --- Random matrices, extreme values, multiplicative chaos ---
\bibitem{Soundararajan} K.~Soundararajan, \emph{Extreme values of zeta and $L$-functions}, Math.\ Ann.\
\textbf{342} (2008), no.~2, 467--486.
\bibitem{BondarenkoSeip} A.~Bondarenko and K.~Seip, \emph{Large greatest common divisor sums and extreme
values of the Riemann zeta function}, Duke Math.\ J.\ \textbf{166} (2017), no.~9, 1685--1701.
\bibitem{FHK} Y.~V.~Fyodorov, G.~A.~Hiary, and J.~P.~Keating, \emph{Freezing transition, characteristic
polynomials of random matrices, and the Riemann zeta function}, Phys.\ Rev.\ Lett.\ \textbf{108} (2012),
170601.
\bibitem{SaksmanWebb} E.~Saksman and C.~Webb, \emph{The Riemann zeta function and Gaussian multiplicative
chaos: statistics on the critical line}, Ann.\ Probab.\ \textbf{48} (2020), no.~6, 2680--2754.
\bibitem{RhodesVargas} R.~Rhodes and V.~Vargas, \emph{Gaussian multiplicative chaos and applications: a
review}, Probab.\ Surv.\ \textbf{11} (2014), 315--392.

%% --- Universality ---
\bibitem{Voronin} S.~M.~Voronin, \emph{Theorem on the ``universality'' of the Riemann zeta-function},
Izv.\ Akad.\ Nauk SSSR Ser.\ Mat.\ \textbf{39} (1975), no.~3, 475--486.
\bibitem{Bagchi} B.~Bagchi, \emph{Recurrence in topological dynamics and the Riemann hypothesis}, Acta
Math.\ Hungar.\ \textbf{50} (1987), 331--340.

%% --- Further adelic / noncommutative-geometry programme ---
\bibitem{ConnesConsaniBC} A.~Connes and C.~Consani, \emph{$\mathrm{BC}$-system, absolute cyclotomy and the
quantized calculus}, EMS Surv.\ Math.\ Sci.\ \textbf{9} (2023), no.~2, 447--475.
\bibitem{ConnesConsaniRR} A.~Connes and C.~Consani, \emph{Riemann--Roch for $\Spec\mathbb Z$}, C.~R.\
Math.\ Acad.\ Sci.\ Paris \textbf{362} (2024), 1727--1732, arXiv:2306.00456.

%% --- Other functional-analytic tools ---
\bibitem{Carleson} L.~Carleson, \emph{Interpolations by bounded analytic functions and the corona
problem}, Ann.\ of Math.\ (2) \textbf{76} (1962), 547--559.
\bibitem{Zagier} D.~Zagier, \emph{The Rankin--Selberg method for automorphic functions which are not of
rapid decay}, J.\ Fac.\ Sci.\ Univ.\ Tokyo \textbf{28} (1981), 415--437.
\bibitem{IwaniecSpectral} H.~Iwaniec, \emph{Spectral Methods of Automorphic Forms}, 2nd ed., Grad.\ Stud.\
Math.\ \textbf{53}, Amer.\ Math.\ Soc., 2002.

\bibitem{deBruijn1950} N.~G. de~Bruijn, \emph{The roots of trigonometric integrals}, Duke Math.\
J.\ \textbf{17} (1950), 197--226.

\bibitem{ConnesConsani2021} A.~Connes and C.~Consani, \emph{Weil positivity and trace formula, the
archimedean place}, Selecta Math.\ \textbf{27} (2021), no.~4, art.~77. arXiv:2006.13771.

\bibitem{ConnesConsani2023} A.~Connes and C.~Consani, \emph{Riemann--Roch for
$\overline{\operatorname{Spec}\Z}$}, Bull.\ Sci.\ Math.\ \textbf{187} (2023), 103293. arXiv:2205.01391.

\bibitem{Connes2026} A.~Connes, \emph{The Riemann Hypothesis: Past, Present and a Letter Through Time},
preprint 2026. arXiv:2602.04022.

\bibitem{DavenportMNT} H.~Davenport, \emph{Multiplicative Number Theory}, 3rd ed., revised by
H.~L.~Montgomery, Springer, 2000.

\bibitem{FoiasFrazho} C.~Foias and A.~E.~Frazho, \emph{The Commutant Lifting Approach to Interpolation
Problems}, Birkh\"auser, 1990.

\bibitem{MontgomeryVaughanMVT} H.~L.~Montgomery and R.~C.~Vaughan, \emph{Hilbert's inequality},
J.\ London Math.\ Soc.\ (2) \textbf{8} (1974), 73--82.

\bibitem{SimonTrace} B.~Simon, \emph{Trace Ideals and Their Applications}, 2nd ed., AMS, 2005.

\bibitem{Li1997} X.-J.~Li, \emph{The positivity of a sequence of numbers and the Riemann hypothesis},
J.\ Number Theory \textbf{65} (1997), 325--333.

\bibitem{BombieriLagarias1999} E.~Bombieri and J.~C.~Lagarias, \emph{Complements to Li's criterion for
the Riemann hypothesis}, J.\ Number Theory \textbf{77} (1999), 274--287.

\bibitem{Keiper1992} J.~B.~Keiper, \emph{Power series expansions of Riemann's $\xi$ function}, Math.\
Comp.\ \textbf{58} (1992), 765--773.

\bibitem{Polymath15} D.~H.~J.~Polymath, \emph{Effective approximation of heat flow evolution of the
Riemann $\xi$ function, and a new upper bound for the de Bruijn--Newman constant}, Res.\ Math.\ Sci.\
\textbf{6} (2019), no.~3, art.~31.

\end{thebibliography}

\end{document}
